% Generated by build_chapter.py from the GY1972 files. Do not edit by hand; edit the files and rebuild.
\chapter{The academic year 1972--73}\label{ch:1972}
This chapter covers the 28 requested catalogue records: 27 talks and the 1973 year heading. Its 272 occurrences are ordered by the manuscript, with status assessed on 9 November 2026. The dominant feature is sustained exposition across several meetings: Hodge structures and local fields; formal groups and their homological algebra; higher algebraic K-theory; quivers and singularities; surgery and stable homotopy. Berezin's quantization, Lomonosov's invariant-subspace argument, Henkin's analytic announcements, and the reports on spectral asymmetry and p-adic L-functions connect these longer sequences to current research.

\begin{center}\small
\begin{tabular}{@{}clrr@{}}\toprule
Class & Meaning & Items & Percent\\\midrule
A & Speaker/Gelfand publication soon afterward & 25 & 9.2\\
B & Developed by other authors & 14 & 5.1\\
C & Partly developed, exact residual stated & 1 & 0.4\\
D & No development identified & 0 & 0.0\\
E & Known before the talk & 219 & 80.5\\
F & Not identifiable from the notes & 13 & 4.8\\
\midrule & Total & 272 & 100.0\\\bottomrule
\end{tabular}
\end{center}

The classes describe historical fate, not truth. Several A items contain literal errors, while an E theorem may be substantially harder than a new announcement. Lomonosov's note was received before his seminar talk; Berezin's symbol paper was already published, whereas the quantization programme was announced in 1973 and expanded in 1974--75. Henkin's bounded extension theorem was old by March 1973, but the convex embedding announcement has a later precise publication. The APS chronology is especially instructive: a 1973 unrestricted regularity announcement and later geometric proofs must be distinguished from Gilkey's 1981 completion of the general eta theorem \citep{lomonosov1973invariant,berezin1972symbols,henkin1972continuation,atiyahpatodisinger1976spectral3,gilkey1981residue}.

Eleven problem entries track definite questions. Most were already resolved or led promptly to established results. The sole C occurrence, 1972-287.11, is explicitly a reconstructed all-degree programme for the numerical homotopy groups of \(G/O\), with the stable-stem and extension data still required by its exact sequence. No particular finite-degree group is asserted to be unknown. The notebook supplies no bounded target for a new contribution, so no C/D item is selected for research development. Berezin's nonhomogeneous analytic-unit phrase and the proposed complex-character spectral series remain F because their missing categories prevent a uniquely formulated problem; nearby modern results do not fill those omissions.

The development appendix instead gives a complete, independently checked correction to Vasserstein's exact published flag assertion. Three elementary matrices are pairwise adjacent but do not lie in one common triangular coset, over every nonzero unital ring and after stabilization. A common acyclic-support criterion replaces the false test. This remains A, with literal correctness recorded separately; the appendix makes no novelty or priority claim \citep[\S11, Remark~(c)]{vaserstein1976foundations}. Shorter marked checks elsewhere repair the cubic cocycle's two-torsion, the reverse Adams formula, the modular-period parity and the raw p-adic zeta continuity claim. Exact computations and separate derivations accompany these corrections.

The cross-year links are specific. Gabrielov continues the Hodge material of 1969-182.2 and 1971-248/251; Buchstaber continues 1971-268/270 on formal laws and 1968-171/173 on normal maps. Vasserstein connects to 1971-253--257 and the earlier K-theory comparison problem, without turning its exposition into a repetition of every earlier conjecture. Quivers continue the four-subspace thread of 1970-198. Lomonosov revisits the compact-operator theme of 1964-50, while eta invariants connect to GP-1971-04 on spectral flow. Berezin's symbol calculus links back to quantization in 1969-178 and 1971-229.

Two ownership issues require explicit treatment. P000449 is a preceding-year tail, already represented by 1971-270.10. P000480 is the Dobrushin discussion already indexed as 1972-288.1--5 in the supplied previous index; record 288 is absent from the requested list, so those occurrences are linked, not duplicated. Within this chapter, 5 March runs 290,292,291 and 26 March runs 296,298,297. Broad catalogue scan windows are not treated as exact talk boundaries. The correction file preserves uncertain readings separately from mathematical emendations, and the bibliography's verification annotations distinguish reading a result from opening metadata.

\GStalk{272}{1972-09-25}{A. M. Gabrielov}{P000450-P000451}{Variations of Hodge structure and their classifying domains}
Gabrielov defines an integral variation of polarized Hodge structure, its flat holomorphic bundle and Griffiths transversality, then describes the classifying flag domain and its homogeneous-space presentation. This is an expository continuation of 1969-182.2 and 1971-251.1--5, not an announcement of a new theorem. The actual heading is on P000450; P000449 is the preceding talk's tail. Two statements require mathematical emendation rather than transcription repair: the orthogonality index on P000451, and the unrestricted claim that the associated Riemannian symmetric space has a complex structure.

\begin{GSitem}{1972-272.1}{E}{Integral local system and its scalar extensions}
\GSfield{Kind; source}{definition; \GSpage{P000450}}
\GSfield{The notes}{Over a nonsingular complex analytic space \(S\), the notes take a locally free finite-rank sheaf \(H_{\mathbb Z}\), form \(H_{\mathbb C}=H_{\mathbb Z}\otimes\mathbb C\), \(H_{\mathbb R}\), and \(H_{\mathcal O}=H_{\mathbb Z}\otimes\mathcal O_S\), and call a local system a bundle with connection.}
\GSfield{Statement}{An integral local system is a sheaf locally isomorphic to \(\underline{\mathbb Z}^{,r}\). Tensoring over \(\mathbb Z\) with the constant real or complex sheaf gives local systems; \(\mathcal H=H_{\mathbb Z}\otimes_{\mathbb Z}\mathcal O_S\) is a holomorphic vector bundle. Locally constant complex sheaves correspond to holomorphic bundles with integrable holomorphic connection; an integral structure is additional data.}
\GSfield{Reconstruction}{The word integrable (equivalently flat) is necessary in the connection equivalence. The complex local system and its associated holomorphic bundle must be distinguished, although the notebook sometimes writes \(H_{\mathbb C}\) for both.}
\GSfield{Fate}{The exposition predates this talk: \citet[\S2, pp.~232--235]{griffiths1970periods} summarizes the period-domain formalism in an invited lecture delivered on 29 August 1968, received on 18 September 1969, and published in March 1970.}
\GSfield{Present status}{As of 9 November 2026: standard local-system/VHS terminology; no research problem is posed.}
\GSfield{Assessment}{literal statement: correct with missing hypotheses; publication match: exact; verification: read the result; development: none; confidence: high}
\end{GSitem}
\begin{GSitem}{1972-272.2}{E}{Canonical flat connection and horizontal sections}
\GSfield{Kind; source}{construction; \GSpage{P000450}}
\GSfield{The notes}{The displayed connection is \(\nabla:H_{\mathcal O}\to\Omega_S^1\otimes_{\mathcal O_S}H_{\mathcal O}\), called Gauss--Manin; the kernel is the sheaf of locally constant complex sections.}
\GSfield{Statement}{For a complex local system \(L\), the associated bundle \(\mathcal H=L\otimes_{\mathbb C}\mathcal O_S\) has the flat connection \(\nabla(v\otimes f)=v\otimes df\), with \(\ker\nabla=L\). For a smooth proper family, the construction applied to a cohomology local system is the Gauss--Manin connection.}
\GSfield{Reconstruction}{The first construction needs no family. The name Gauss--Manin is geometrically justified only after identifying the local system with the cohomology of such a family; the notes do not supply a family here. This repeats the setting of 1969-182.2 and 1971-251.5.}
\GSfield{Fate}{The exposition predates this talk: \citet[\S2, pp.~232--235]{griffiths1970periods} summarizes the period-domain formalism in an invited lecture delivered on 29 August 1968, received on 18 September 1969, and published in March 1970. Its formula (2.1) records the connection's compatibility with the Hodge filtration.}
\GSfield{Present status}{As of 9 November 2026: the local-system construction and its geometric Gauss--Manin specialization remain standard.}
\GSfield{Assessment}{literal statement: correct; publication match: exact; verification: read the result; development: none; confidence: high}
\end{GSitem}
\begin{GSitem}{1972-272.3}{E}{Hodge filtration and Griffiths transversality}
\GSfield{Kind; source}{definition; \GSpage{P000450}}
\GSfield{The notes}{A decreasing filtration \(F^0\supset F^1\supset\cdots\supset F^n\) by holomorphic subbundles is required to satisfy \(\nabla\mathcal F^p\subset\Omega_S^1\otimes\mathcal F^{p-1}\). A locally constant bilinear form is to give a polarized Hodge structure in each fibre.}
\GSfield{Statement}{A polarized integral variation of pure Hodge structure of weight \(n\) consists of an integral local system, holomorphic subbundles \(F^p\mathcal H\), a flat bilinear form \(Q\), fibrewise Hodge opposition and positivity, and \(\nabla(F^p\mathcal H)\subset\Omega_S^1\otimes F^{p-1}\mathcal H\). The displayed finite range presupposes effective Hodge types \(0\leq p,q\leq n\); the general definition uses all integer indices, with eventually constant endpoints.}
\GSfield{Reconstruction}{The bundle \(\mathcal H\), not a locally constant sheaf, carries the holomorphic filtration. Flatness and the effective-weight convention make explicit the notebook's abbreviated notation.}
\GSfield{Fate}{The exposition predates this talk: \citet[\S2, pp.~232--235]{griffiths1970periods} summarizes the period-domain formalism in an invited lecture delivered on 29 August 1968, received on 18 September 1969, and published in March 1970. See equation (2.1) for transversality and the preceding definition of the classifying space.}
\GSfield{Present status}{As of 9 November 2026: established definition; the notes do not pose a transversality conjecture.}
\GSfield{Assessment}{literal statement: correct with missing hypotheses; publication match: exact; verification: read the result; development: none; confidence: high}
\end{GSitem}
\begin{GSitem}{1972-272.4}{E}{Compact dual and the erroneous orthogonality index}
\GSfield{Kind; source}{construction; \GSpage{P000450-P000451}}
\GSfield{The notes}{The notes impose fixed flag dimensions and the displayed condition \(\psi(F^p,F^{n-p})=0\). The original P000451 scan has this same index, without \(+1\).}
\GSfield{Statement}{Fix a real vector space \(H_{\mathbb R}\), Hodge numbers of weight \(n\), and a nondegenerate \((-1)^n\)-symmetric bilinear form \(Q\). The compact dual is the projective variety of flags with the prescribed dimensions and \(Q(F^p,F^{n-p+1})=0\) for every \(p\), taking the appropriate connected flag component if needed.}
\GSfield{Reconstruction}{The index is emended to \(n-p+1\), also written correctly in Gabrielov's continuation on P000458. This is a mathematical error faithfully transcribed: the stronger printed condition would generally force the intended pairing of complementary Hodge types to vanish.}
\GSfield{Fate}{The exposition predates this talk: \citet[\S2, pp.~232--235]{griffiths1970periods} summarizes the period-domain formalism in an invited lecture delivered on 29 August 1968, received on 18 September 1969, and published in March 1970. The bilinear relations and flag-domain construction in \S2 support the corrected formulation, not the literal index.}
\GSfield{Present status}{As of 9 November 2026: the standard compact-dual construction is established; the notebook's literal orthogonality formula is false.}
\GSfield{Assessment}{literal statement: false as written; publication match: partial; verification: read the result; development: annotated; confidence: high}
\end{GSitem}
\begin{GSitem}{1972-272.5}{E}{Opposed filtrations, Hodge decomposition, and Weil operator}
\GSfield{Kind; source}{definition; \GSpage{P000451}}
\GSfield{The notes}{The page imposes \(F^p\oplus\overline{F^{n-p+1}}=H_{\mathbb C}\), defines \(H^{p,n-p}=F^p\cap\overline{F^{n-p}}\), and lets \(C\) act by \(i^{2p-n}\) on this summand.}
\GSfield{Statement}{For a decreasing exhaustive separated filtration satisfying the displayed opposition, \(H_{\mathbb C}=\bigoplus_{p+q=n}H^{p,q}\), \(\overline{H^{p,q}}=H^{q,p}\), and \(F^p=\bigoplus_{r\geq p}H^{r,n-r}\). The real-linear Weil operator complexifies to \(C|_{H^{p,q}}=i^{p-q}\).}
\GSfield{Reconstruction}{The direct-sum decomposition and recovery of the filtration are stated explicitly to unpack the notebook's concise definition; the notation and exponent themselves are faithful.}
\GSfield{Fate}{The exposition predates this talk: \citet[\S2, pp.~232--235]{griffiths1970periods} summarizes the period-domain formalism in an invited lecture delivered on 29 August 1968, received on 18 September 1969, and published in March 1970.}
\GSfield{Present status}{As of 9 November 2026: the standard equivalence between opposed filtrations and pure real Hodge decompositions.}
\GSfield{Assessment}{literal statement: correct; publication match: exact; verification: read the result; development: none; confidence: high}
\end{GSitem}
\begin{GSitem}{1972-272.6}{E}{Polarization positivity and the open period domain}
\GSfield{Kind; source}{definition; \GSpage{P000451}}
\GSfield{The notes}{The form \(\psi(x,C\bar x)\) is positive, and imposing this positivity and opposition gives an open subset \(D\) of the compact flag variety \(\overline D\).}
\GSfield{Statement}{Normalize the real polarization \(Q\) so that \(h(x,y)=Q(x,C\overline y)\) is a positive-definite Hermitian form. For fixed polarized Hodge numbers and fixed \(Q\), the flags satisfying Hodge opposition and this positivity form an open complex submanifold \(D\) of the compact dual. Its connected components may be selected separately.}
\GSfield{Reconstruction}{The harmless overall sign convention for \(Q\) is fixed explicitly; positivity means for every nonzero complex vector. Integrality of a lattice does not make the period domain discrete.}
\GSfield{Fate}{The exposition predates this talk: \citet[\S2, pp.~232--235]{griffiths1970periods} summarizes the period-domain formalism in an invited lecture delivered on 29 August 1968, received on 18 September 1969, and published in March 1970.}
\GSfield{Present status}{As of 9 November 2026: standard period-domain construction with the stated polarization convention.}
\GSfield{Assessment}{literal statement: correct; publication match: exact; verification: read the result; development: none; confidence: high}
\end{GSitem}
\begin{GSitem}{1972-272.7}{E}{Homogeneous period domain and compact stabilizer}
\GSfield{Kind; source}{theorem; \GSpage{P000451}}
\GSfield{The notes}{The real group preserving \(\psi\) acts transitively on \(D\); the stabilizer \(H\) is compact.}
\GSfield{Statement}{Let \(G_{\mathbb R}\) be the full real isometry group of the polarization with the relevant flag-component action. A period-domain orbit is \(G_{\mathbb R}/V\), where the stabilizer \(V\) of a polarized Hodge decomposition is compact. The identity component \(G_{\mathbb R}^{0}\) acts transitively on each connected component of the period domain.}
\GSfield{Reconstruction}{The notebook's stabilizer \(H\) is renamed \(V\) to distinguish it from the Hodge vector space. Connected-component qualifications are necessary if one replaces the full real isometry group by its identity component; the continuation on P000458 does so.}
\GSfield{Fate}{The exposition predates this talk: \citet[\S2, pp.~232--235]{griffiths1970periods} summarizes the period-domain formalism in an invited lecture delivered on 29 August 1968, received on 18 September 1969, and published in March 1970. See equation (2.3) and the surrounding homogeneous-space description.}
\GSfield{Present status}{As of 9 November 2026: established homogeneous-space description, component by component.}
\GSfield{Assessment}{literal statement: correct; publication match: exact; verification: read the result; development: none; confidence: high}
\end{GSitem}
\begin{GSitem}{1972-272.8}{E}{Projection to a Riemannian symmetric space}
\GSfield{Kind; source}{claim; \GSpage{P000451}}
\GSfield{The notes}{For a maximal compact \(K\supset H\), the map \(G/H\to G/K\) is said not to be holomorphic; the following sentence says \(G/K\) has a complex structure.}
\GSfield{Statement}{A maximal compact subgroup \(K\) containing the compact period stabilizer \(V\) gives a smooth fibration \(K/V\to G_{\mathbb R}/V\to G_{\mathbb R}/K\). The base is a Riemannian symmetric space and is not, in general, Hermitian symmetric. In the Hermitian cases the projection need not be holomorphic; when \(V=K\), it is the identity and is holomorphic for the compatible structure.}
\GSfield{Reconstruction}{Both universal assertions need qualification. No complex structure on the symmetric space is supplied in the notes. The correct general statement is smooth homogeneity, with additional hypotheses for complex geometry.}
\GSfield{Fate}{The exposition predates this talk: \citet[\S2, pp.~232--235]{griffiths1970periods} summarizes the period-domain formalism in an invited lecture delivered on 29 August 1968, received on 18 September 1969, and published in March 1970. Equation (2.3) explicitly describes a differentiable fibration and qualifies holomorphicity; it does not assert that every symmetric quotient is complex.}
\GSfield{Present status}{As of 9 November 2026: the literal blanket claim is false; the qualified period-domain fibration is classical.}
\GSfield{Assessment}{literal statement: false as written; publication match: partial; verification: read the result; development: annotated; confidence: high}
\end{GSitem}

\GStalk{273}{1972-09-25}{Yu. I. Manin}{P000452-P000453}{Retelling of Manin's first lecture on p-adic numbers: the lattice tree}
The heading explicitly calls this a \ru{пересказ}, a retelling of Manin's first lecture. The two pages construct the regular tree of homothety classes of rank-two lattices over a nonarchimedean local field, with its \(\mathrm{PGL}_2\)-action and invariant distance. The analogy with the Lobachevsky plane is motivational. A crossed-out equal-characteristic alternative is not promoted to an assertion by the speaker, and the vague reference to A. Weil is not assigned to an invented publication.

\begin{GSitem}{1972-273.1}{E}{Local field, valuation ring, and projective linear action}
\GSfield{Kind; source}{definition; \GSpage{P000452}}
\GSfield{The notes}{The uncancelled text takes \(K\) to be a finite extension of \(\mathbb Q_p\), with integer ring \(A\), maximal ideal \(\mathfrak m\), and fractional linear transformations of \(\mathbb P^1(K)\).}
\GSfield{Statement}{For a finite extension \(K/\mathbb Q_p\), put \(A=\mathcal O_K\), \(\mathfrak m=(\pi)\), and \(k=A/\mathfrak m\cong\mathbb F_q\). Then \(G=\mathrm{PGL}_2(K)=\mathrm{GL}_2(K)/K^\times\) acts on \(\mathbb P^1(K)\) by \(z\mapsto(az+b)/(cz+d)\), with the usual interpretation at poles and infinity.}
\GSfield{Reconstruction}{Invertibility of the representing matrix and the projective convention at infinity are made explicit. The construction also works for equal-characteristic local fields, but that extension is editorial: the relevant phrase is cancelled in the original.}
\GSfield{Fate}{The local-field tree was already part of the established theory: \citet[Introduction, pp.~130--131, and \S1]{mumford1972degenerating} uses the lattice tree and its \(q+1\) neighbours in a paper received on 16 March 1971 and published in 1972. \citet[Chapter~II]{serre1977arbres} is a later systematic exposition based on his 1968--1969 Coll\`ege de France course; its 1977 publication is not the construction's first circulation.}
\GSfield{Present status}{As of 9 November 2026: standard local-field setting and projective action; no open question.}
\GSfield{Assessment}{literal statement: correct; publication match: exact; verification: read the result; development: none; confidence: high}
\end{GSitem}
\begin{GSitem}{1972-273.2}{E}{Lattices, homothety classes, and the vertex homogeneous space}
\GSfield{Kind; source}{definition; \GSpage{P000452}}
\GSfield{The notes}{A lattice is written \(L=Ae_1\oplus Ae_2\); lattices are equivalent when they differ by a scalar in \(K^\times\). Vertices are these equivalence classes, also written \(G/U\) for \(U=\mathrm{GL}_2(A)/A^\times\).}
\GSfield{Statement}{A full lattice in \(K^2\) is a finitely generated \(A\)-submodule spanning \(K^2\), necessarily free of rank two. The homothety classes form a transitive \(G\)-set. The stabilizer of \([A^2]\) is the image of \(K^\times\mathrm{GL}_2(A)\) in \(G\), naturally \(\mathrm{GL}_2(A)/A^\times\); hence the vertex set is \(G/U\).}
\GSfield{Reconstruction}{Fullness and spanning are essential: a rank-two module embedded degenerately would not define the intended vertex. The stabilizer identification unpacks the concise quotient notation.}
\GSfield{Fate}{The local-field tree was already part of the established theory: \citet[Introduction, pp.~130--131, and \S1]{mumford1972degenerating} uses the lattice tree and its \(q+1\) neighbours in a paper received on 16 March 1971 and published in 1972. \citet[Chapter~II]{serre1977arbres} is a later systematic exposition based on his 1968--1969 Coll\`ege de France course; its 1977 publication is not the construction's first circulation.}
\GSfield{Present status}{As of 9 November 2026: standard vertex model of the Bruhat--Tits tree of \(\mathrm{PGL}_2(K)\).}
\GSfield{Assessment}{literal statement: correct; publication match: exact; verification: read the result; development: none; confidence: high}
\end{GSitem}
\begin{GSitem}{1972-273.3}{E}{Simultaneous elementary divisors for two lattices}
\GSfield{Kind; source}{theorem; \GSpage{P000452}}
\GSfield{The notes}{For two lattices the notes assert a basis with \(L_1=Ae_1\oplus Ae_2\) and \(L_2=A\pi^{\nu_1}e_1\oplus A\pi^{\nu_2}e_2\), with integers \(\nu_i\).}
\GSfield{Statement}{For full lattices \(L_1,L_2\subset K^2\), there is an \(A\)-basis \(e_1,e_2\) of \(L_1\) and integers \(\nu_1\leq\nu_2\) such that \(L_2=\pi^{\nu_1}Ae_1\oplus\pi^{\nu_2}Ae_2\). The unordered exponents are uniquely determined; simultaneous scaling of \(L_2\) shifts both by the same integer.}
\GSfield{Reconstruction}{This is Smith normal form over the discrete valuation ring after multiplying \(L_2\) by a sufficiently large power of \(\pi\). The ordering and uniqueness, needed for distance, are supplied explicitly.}
\GSfield{Fate}{The local-field tree was already part of the established theory: \citet[Introduction, pp.~130--131, and \S1]{mumford1972degenerating} uses the lattice tree and its \(q+1\) neighbours in a paper received on 16 March 1971 and published in 1972. \citet[Chapter~II]{serre1977arbres} is a later systematic exposition based on his 1968--1969 Coll\`ege de France course; its 1977 publication is not the construction's first circulation.}
\GSfield{Present status}{As of 9 November 2026: classical elementary-divisor theorem; its use here is expository.}
\GSfield{Assessment}{literal statement: correct; publication match: exact; verification: read the result; development: none; confidence: high}
\end{GSitem}
\begin{GSitem}{1972-273.4}{E}{Invariant distance between lattice classes}
\GSfield{Kind; source}{construction; \GSpage{P000453}}
\GSfield{The notes}{Representatives are normalized to \(L_2=Ae_1\oplus A\pi^\nu e_2\subset L_1\), \(\nu\geq0\), with \(L_1/L_2\cong A/\mathfrak m^\nu\). The distance is \(\nu\), equivalently \(|\nu_1-\nu_2|\).}
\GSfield{Statement}{For homothety classes \([L_1],[L_2]\), define \(d([L_1],[L_2])=|\nu_1-\nu_2|\), where \(\nu_i\) are the simultaneous elementary divisors. This is independent of bases and representatives, invariant under \(\mathrm{PGL}_2(K)\), and is a metric. The stated cyclic quotient holds after the indicated normalized inclusion; arbitrary representatives need not admit a quotient \(L_1/L_2\).}
\GSfield{Reconstruction}{The scan explicitly boxes \(L_2\subset L_1\), so the initial normalized quotient is not erroneous. The later variant with two quotient summands must also be read after scaling to an inclusion; negative exponents cannot be used in \(A/\pi^{\nu_i}A\).}
\GSfield{Fate}{The local-field tree was already part of the established theory: \citet[Introduction, pp.~130--131, and \S1]{mumford1972degenerating} uses the lattice tree and its \(q+1\) neighbours in a paper received on 16 March 1971 and published in 1972. \citet[Chapter~II]{serre1977arbres} is a later systematic exposition based on his 1968--1969 Coll\`ege de France course; its 1977 publication is not the construction's first circulation.}
\GSfield{Present status}{As of 9 November 2026: the standard graph distance, with the inclusion convention restored for quotient formulas.}
\GSfield{Assessment}{literal statement: correct with missing hypotheses; publication match: exact; verification: read the result; development: none; confidence: high}
\end{GSitem}
\begin{GSitem}{1972-273.5}{E}{Adjacency and the q plus one neighbours}
\GSfield{Kind; source}{construction; \GSpage{P000453}}
\GSfield{The notes}{Two vertices are joined when their distance is one, and exactly \(q+1\) edges meet at each vertex.}
\GSfield{Statement}{The neighbours of \([L]\) are represented uniquely by sublattices \(L'\) with \(\pi L\subsetneq L'\subsetneq L\). They correspond to one-dimensional subspaces of \(L/\pi L\cong k^2\), so there are \(|\mathbb P^1(k)|=q+1\) neighbours.}
\GSfield{Reconstruction}{The residue-line correspondence is the modern explanation of the stated valency; the notebook states the result without this proof.}
\GSfield{Fate}{The local-field tree was already part of the established theory: \citet[Introduction, pp.~130--131, and \S1]{mumford1972degenerating} uses the lattice tree and its \(q+1\) neighbours in a paper received on 16 March 1971 and published in 1972. \citet[Chapter~II]{serre1977arbres} is a later systematic exposition based on his 1968--1969 Coll\`ege de France course; its 1977 publication is not the construction's first circulation.}
\GSfield{Present status}{As of 9 November 2026: a classical local combinatorial description of the tree.}
\GSfield{Assessment}{literal statement: correct; publication match: exact; verification: read the result; development: none; confidence: high}
\end{GSitem}
\begin{GSitem}{1972-273.6}{E}{The lattice graph is a tree}
\GSfield{Kind; source}{theorem; \GSpage{P000453}}
\GSfield{The notes}{The graph \(\Delta\) is called a connected graph without cycles, with its natural \(G\)-action.}
\GSfield{Statement}{The graph of homothety classes of full rank-two \(A\)-lattices and distance-one edges is connected and acyclic. Its graph metric is the elementary-divisor metric, and the action of \(\mathrm{PGL}_2(K)\) preserves adjacency. Thus its geometric realization is the Bruhat--Tits tree in rank one.}
\GSfield{Reconstruction}{The name Bruhat--Tits is modern identification, not an extra theorem attributed to this retelling. No assertion about higher-rank buildings is inferred.}
\GSfield{Fate}{The local-field tree was already part of the established theory: \citet[Introduction, pp.~130--131, and \S1]{mumford1972degenerating} uses the lattice tree and its \(q+1\) neighbours in a paper received on 16 March 1971 and published in 1972. \citet[Chapter~II]{serre1977arbres} is a later systematic exposition based on his 1968--1969 Coll\`ege de France course; its 1977 publication is not the construction's first circulation.}
\GSfield{Present status}{As of 9 November 2026: established rank-one building; the notebook contains no unresolved tree theorem.}
\GSfield{Assessment}{literal statement: correct; publication match: exact; verification: read the result; development: none; confidence: high}
\end{GSitem}
\begin{GSitem}{1972-273.7}{E}{Analogy with the hyperbolic plane and the unspecified Weil reference}
\GSfield{Kind; source}{remark; \GSpage{P000452-P000453}}
\GSfield{The notes}{The graph is described as an analogue of the Lobachevsky plane, with a disk drawing and a brief \ru{см. А. Вейль} near the distance discussion.}
\GSfield{Statement}{The precise content available is the analogy between \(\mathrm{PGL}_2(K)/\mathrm{PGL}_2(A)\) with its lattice-tree metric and an archimedean rank-one symmetric space. The drawing does not identify the tree with a two-dimensional hyperbolic manifold, and no particular theorem or work of Weil is named.}
\GSfield{Reconstruction}{Only the explicit analogy is retained. The bibliographic allusion is unresolved, but the mathematical tree construction is identified independently and is not classified F solely because that allusion lacks a title.}
\GSfield{Fate}{\citet[Introduction, pp.~130--131]{mumford1972degenerating} explicitly develops the same local-field/upper-half-plane analogy before this lecture. The notes do not provide enough information to identify the intended Weil citation.}
\GSfield{Present status}{As of 9 November 2026: an established analogy, with one unidentified bibliographic allusion and no additional definite problem.}
\GSfield{Assessment}{literal statement: correct; publication match: partial; verification: read the result; development: none; confidence: medium}
\end{GSitem}

\GStalk{274}{1972-09-25}{D. A. Kazhdan}{P000454-P000455}{Local-field ramification as preparation for Tate's p-divisible groups}
The heading identifies Tate's paper in the Russian translation, \emph{Matematika} 13:2, \S3. The surviving mathematics is a preparatory account of finite Galois extensions of local fields: the ramification filtration, residue and tame characters, additive higher quotients, and solvability. The notebook shifts the usual lower-numbering convention by one. It introduces \(\mathbb C_p\), but does not state Tate's later theorem on its Galois action, so that theorem is not silently supplied as a talk item.

\begin{GSitem}{1972-274.1}{E}{The field C p and the exact Tate reading target}
\GSfield{Kind; source}{definition; \GSpage{P000454}}
\GSfield{The notes}{The heading cites \ru{ст. Тэйта, Математика, 13:2, \S3} and defines \(\mathbb C_p\) as the completion of an algebraic closure of \(\mathbb Q_p\).}
\GSfield{Statement}{Fix an algebraic closure \(\overline{\mathbb Q}_p\), extend the \(p\)-adic absolute value uniquely to it, and let \(\mathbb C_p\) be its completion. The target reference is \S3 of Tate's \emph{p-Divisible Groups}; no specific theorem about \(\mathbb C_p\) beyond this construction is stated on the surviving pages.}
\GSfield{Reconstruction}{The bibliographic identification is exact: the 1969 Russian translation, pp.~3--25, appeared in the specified volume and issue. The English original appeared in the proceedings of the 1966 Driebergen local-fields conference in 1967. The heading's intended destination is kept separate from the finite-extension preliminaries actually recorded.}
\GSfield{Fate}{\citet[\S3]{tate1967pdivisible} is the identified source, in Manin's 1969 Russian translation. The citation predates Kazhdan's 1972 exposition; no new Tate theorem is inferred from the heading.}
\GSfield{Present status}{As of 9 November 2026: the construction and reading reference are identified; the notes do not support an additional Hodge--Tate comparison item.}
\GSfield{Assessment}{literal statement: correct; publication match: exact; verification: read the result; development: none; confidence: high}
\end{GSitem}
\begin{GSitem}{1972-274.2}{E}{Finite local Galois extensions and normalized valuation}
\GSfield{Kind; source}{definition; \GSpage{P000454}}
\GSfield{The notes}{The notes write \(K/\mathbb Q_p\) and \(L/K\) finite, set \(G=\operatorname{Gal}(L/K)\), introduce their integer and residue rings, and choose \(|\pi|=|k|^{-1}\).}
\GSfield{Statement}{Let \(K/\mathbb Q_p\) be finite and let \(L/K\) be finite Galois. Write \(\mathcal O_K,\mathcal O_L\) for integer rings, \(\mathfrak m_K=(\pi),\mathfrak m_L=(\tau)\), and \(k,\ell\) for residue fields. Normalize \(|\pi|_K=|k|^{-1}\), and use its unique extension to \(L\); every element of \(G\) preserves this valuation and \(\mathcal O_L\).}
\GSfield{Reconstruction}{The Galois hypothesis is supplied because it is needed for the surjectivity onto \(\operatorname{Gal}(\ell/k)\) asserted later. Simply naming the automorphism group of an arbitrary finite extension is not enough. The valuation on \(L\) here extends the chosen one on \(K\), rather than being independently normalized by \(|\ell|\).}
\GSfield{Fate}{This is classical local ramification theory, already organized in Serre's 1962 \emph{Corps locaux}; the opened English edition is \citet[Chapter~IV, pp.~61--79]{serre1979local}. The precise formulas and proofs were read in \citet[Lemma~7.57, Theorem~7.58, and Corollary~7.59, pp.~130--131]{milne2020ant}; this modern check is not evidence of a later discovery.}
\GSfield{Present status}{As of 9 November 2026: standard local-field setup with the Galois hypothesis made explicit.}
\GSfield{Assessment}{literal statement: correct with missing hypotheses; publication match: exact; verification: read the result; development: none; confidence: high}
\end{GSitem}
\begin{GSitem}{1972-274.3}{E}{Ramification filtration and its shifted indexing}
\GSfield{Kind; source}{definition; \GSpage{P000454}}
\GSfield{The notes}{The definition is \(G_i=\{g:\ g(x)-x\in\mathfrak m_L^i\text{ for every }x\in\mathcal O_L\}\), with \(G=G_0\supset G_1\supset\cdots\supset\{1\}\); every \(G_i\) is normal.}
\GSfield{Statement}{For the finite Galois extension just specified, define \(G_i^{\rm note}\) by that formula for \(i\geq0\). These are normal subgroups of \(G\), eventually trivial. They satisfy \(G_0^{\rm note}=G\) and \(G_i^{\rm note}=G_{i-1}^{\rm lower}\) for \(i\geq1\), where the conventional lower filtration has inertia at index zero.}
\GSfield{Reconstruction}{The shifted indexing is intentional and consistent across both pages; it is not silently changed. Eventual triviality follows from finiteness of \(G\) and separation of the valuation topology, not from a fixed universal length.}
\GSfield{Fate}{This is classical local ramification theory, already organized in Serre's 1962 \emph{Corps locaux}; the opened English edition is \citet[Chapter~IV, pp.~61--79]{serre1979local}. The precise formulas and proofs were read in \citet[Lemma~7.57, Theorem~7.58, and Corollary~7.59, pp.~130--131]{milne2020ant}; this modern check is not evidence of a later discovery.}
\GSfield{Present status}{As of 9 November 2026: classical filtration with an explicitly recorded one-place indexing shift.}
\GSfield{Assessment}{literal statement: correct; publication match: exact; verification: read the result; development: none; confidence: high}
\end{GSitem}
\begin{GSitem}{1972-274.4}{E}{Residue-field quotient and the exercise on surjectivity}
\GSfield{Kind; source}{theorem; \GSpage{P000454-P000455}}
\GSfield{The notes}{Reduction gives \(\psi:G\to\operatorname{Gal}(\ell/k)\); its kernel is \(G_1\), and proving that the induced injection \(G_0/G_1\to\operatorname{Gal}(\ell/k)\) is an isomorphism is left as an exercise.}
\GSfield{Statement}{For a finite Galois extension of nonarchimedean local fields, reduction yields the exact sequence \(1\to G_1^{\rm note}\to G\to\operatorname{Gal}(\ell/k)\to1\). The quotient is cyclic, generated by the finite-field Frobenius. A residue element must be lifted to \(\mathcal O_L\); kernel membership means fixing every residue element.}
\GSfield{Reconstruction}{The notebook writes a lift \(x\in L\) and suppresses the universal quantifier in one line; the lift must be integral. The exercise is a classical expository task with a definite known answer, not a new research problem; it is therefore indexed as an occurrence but not assigned a new problem identifier.}
\GSfield{Fate}{This is classical local ramification theory, already organized in Serre's 1962 \emph{Corps locaux}; the opened English edition is \citet[Chapter~IV, pp.~61--79]{serre1979local}. The precise formulas and proofs were read in \citet[Lemma~7.57, Theorem~7.58, and Corollary~7.59, pp.~130--131]{milne2020ant}; this modern check is not evidence of a later discovery. The surjectivity is explicit in Milne's Theorem~7.58(a).}
\GSfield{Present status}{As of 9 November 2026: the exercise is solved by the standard maximal unramified subextension argument.}
\GSfield{Assessment}{literal statement: correct with missing hypotheses; publication match: exact; verification: read the result; development: annotated; confidence: high}
\end{GSitem}
\begin{GSitem}{1972-274.5}{E}{Tame character on the first inertia quotient}
\GSfield{Kind; source}{theorem; \GSpage{P000455}}
\GSfield{The notes}{The map \(g\mapsto\overline{g(\tau)/\tau}\) is said to identify \(G_1/G_2\) with a subgroup of \(\ell^\times\), independently of the uniformizer.}
\GSfield{Statement}{For \(g\in G_1^{\rm note}\), let \(\chi(g)=\overline{g(\tau)/\tau}\in\ell^\times\). This is a homomorphism, its kernel is \(G_2^{\rm note}\), and it is independent of the uniformizer \(\tau\). Thus tame inertia \(G_1^{\rm note}/G_2^{\rm note}\) is cyclic of order prime to the residue characteristic.}
\GSfield{Reconstruction}{Independence uses inertia: if \(\tau'=u\tau\) with \(u\in\mathcal O_L^\times\), then \(g(u)/u\) reduces to one. The map is first defined on \(G_1\), and only then induces an injection on its quotient.}
\GSfield{Fate}{This is classical local ramification theory, already organized in Serre's 1962 \emph{Corps locaux}; the opened English edition is \citet[Chapter~IV, pp.~61--79]{serre1979local}. The precise formulas and proofs were read in \citet[Lemma~7.57, Theorem~7.58, and Corollary~7.59, pp.~130--131]{milne2020ant}; this modern check is not evidence of a later discovery. See Milne's Corollary~7.59 and its first character computation.}
\GSfield{Present status}{As of 9 November 2026: the standard tame-inertia character.}
\GSfield{Assessment}{literal statement: correct; publication match: exact; verification: read the result; development: none; confidence: high}
\end{GSitem}
\begin{GSitem}{1972-274.6}{E}{Additive characters on higher ramification quotients}
\GSfield{Kind; source}{theorem; \GSpage{P000455}}
\GSfield{The notes}{For \(i\geq2\), the notes use \(\overline{(g(\tau)/\tau-1)/\tau^{i-1}}\in\ell\) and state that its kernel is \(G_{i+1}\), although the displayed domain is already \(G_i/G_{i+1}\).}
\GSfield{Statement}{For \(i\geq2\), the map \(c_i:G_i^{\rm note}\to(\ell,+)\), \(c_i(g)=\overline{(g(\tau)-\tau)/\tau^i}\), is a homomorphism with kernel \(G_{i+1}^{\rm note}\). Consequently \(G_i^{\rm note}/G_{i+1}^{\rm note}\) embeds in the additive residue group and is an elementary abelian \(p\)-group.}
\GSfield{Reconstruction}{The character's domain is repaired before taking the quotient; an injective map defined on \(G_i/G_{i+1}\) has trivial kernel, not \(G_{i+1}\). Unlike the tame character, these coordinate-valued higher characters can change with the uniformizer, although the quotient statement is intrinsic.}
\GSfield{Fate}{This is classical local ramification theory, already organized in Serre's 1962 \emph{Corps locaux}; the opened English edition is \citet[Chapter~IV, pp.~61--79]{serre1979local}. The precise formulas and proofs were read in \citet[Lemma~7.57, Theorem~7.58, and Corollary~7.59, pp.~130--131]{milne2020ant}; this modern check is not evidence of a later discovery. The same calculation, with the usual one-place shift, is the second part of Milne's Corollary~7.59.}
\GSfield{Present status}{As of 9 November 2026: the higher quotient theorem is standard; the literal domain/kernel combination is inconsistent.}
\GSfield{Assessment}{literal statement: false as written; publication match: partial; verification: read the result; development: annotated; confidence: high}
\end{GSitem}
\begin{GSitem}{1972-274.7}{E}{Solvability of finite local Galois groups}
\GSfield{Kind; source}{theorem; \GSpage{P000454-P000455}}
\GSfield{The notes}{Kazhdan announces that the filtration will show \(G\) is solvable; every successive quotient is declared abelian.}
\GSfield{Statement}{For every finite Galois extension \(L/K\) of nonarchimedean local fields with finite residue field, \(\operatorname{Gal}(L/K)\) is solvable: the residue quotient is cyclic, the tame quotient is cyclic, and the remaining finite filtration has elementary abelian \(p\)-quotients.}
\GSfield{Reconstruction}{The finite residue field and Galois assumptions are essential to this precise formulation; the statement is not a theorem about arbitrary discretely valued fields or arbitrary global number fields.}
\GSfield{Fate}{This is classical local ramification theory, already organized in Serre's 1962 \emph{Corps locaux}; the opened English edition is \citet[Chapter~IV, pp.~61--79]{serre1979local}. The precise formulas and proofs were read in \citet[Lemma~7.57, Theorem~7.58, and Corollary~7.59, pp.~130--131]{milne2020ant}; this modern check is not evidence of a later discovery. Milne's Corollary~7.59 states the solvability conclusion explicitly.}
\GSfield{Present status}{As of 9 November 2026: classical and completely proved; no outstanding solvability question is stated.}
\GSfield{Assessment}{literal statement: correct with missing hypotheses; publication match: exact; verification: read the result; development: none; confidence: high}
\end{GSitem}
\begin{GSitem}{1972-274.8}{E}{Unramified, tame, and totally ramified extensions}
\GSfield{Kind; source}{definition; \GSpage{P000455}}
\GSfield{The notes}{The page calls \(G_1=1\) unramified, \(G_2=1\) \ru{слаборазветвл.}, translated weakly ramified, and \(G_0=G_1\) purely ramified.}
\GSfield{Statement}{In the notebook's shifted numbering, a finite Galois local extension is unramified iff \(G_1^{\rm note}=1\), tamely ramified iff \(G_2^{\rm note}=1\), and totally ramified iff \(G_0^{\rm note}=G_1^{\rm note}\). Modern weak ramification instead means \(G_2^{\rm lower}=1\), equivalently \(G_3^{\rm note}=1\).}
\GSfield{Reconstruction}{The old descriptive phrase must not be conflated with the now distinct technical term weakly ramified. The translation is literal, so this is primarily a terminology annotation rather than evidence that the original writer miscomputed the groups.}
\GSfield{Fate}{This is classical local ramification theory, already organized in Serre's 1962 \emph{Corps locaux}; the opened English edition is \citet[Chapter~IV, pp.~61--79]{serre1979local}. The precise formulas and proofs were read in \citet[Lemma~7.57, Theorem~7.58, and Corollary~7.59, pp.~130--131]{milne2020ant}; this modern check is not evidence of a later discovery. Milne's paragraph after Corollary~7.59 identifies tame and unramified extensions in the standard numbering. The modern distinction from weak ramification is explicit in \citet[\S1, p.~1]{johnston2012weak}, first posted in April 2012; that source is used for terminology only.}
\GSfield{Present status}{As of 9 November 2026: the three intended notions are identified, with tame replacing weak in modern terminology.}
\GSfield{Assessment}{literal statement: correct with missing hypotheses; publication match: partial; verification: read the result; development: annotated; confidence: high}
\end{GSitem}

\GStalk{275}{1972-10-02}{A. M. Gabrielov}{P000456-P000458}{Local systems, flat connections, and polarized Hodge flags}
The continuation rebuilds the sheaf and bundle foundations carefully, gives the flat connection in a local constant frame, then returns to polarized flag domains. It repeats the classical line of record 272 with fuller definitions. The original scan reverses the parity of the polarization, a substantive mathematical error. The final two bare headings, Manin's second lecture and Kazhdan again on Tate, contain no additional mathematical assertions and are recorded as headings in this summary only.

\begin{GSitem}{1972-275.1}{E}{Constant sheaf and its étalé space}
\GSfield{Kind; source}{definition; \GSpage{P000456}}
\GSfield{The notes}{The constant sheaf \(\mathbb Z\) has constant-number sections over a connected domain; it is associated to the presheaf \(U\mapsto\mathbb Z\). Its sheaf space is written \(S\times\mathbb Z\).}
\GSfield{Statement}{On a complex manifold \(S\), the constant sheaf \(\underline{\mathbb Z}\) is the sheaf of locally constant integer-valued functions, equivalently the sheafification of the constant presheaf. Its \(\acute{e}\)tal\'{e} space is \(S\times\mathbb Z\), with \(\mathbb Z\) discrete. On a connected nonempty open set its sections are \(\mathbb Z\); on a disconnected open set they need not be globally constant.}
\GSfield{Reconstruction}{The sheafification and connected-open qualification are part of the mathematical distinction being explained in the notes, not a new theorem. The empty set convention for the constant presheaf is immaterial after sheafification.}
\GSfield{Fate}{The exposition predates this talk: \citet[\S2, pp.~232--235]{griffiths1970periods} summarizes the period-domain formalism in an invited lecture delivered on 29 August 1968, received on 18 September 1969, and published in March 1970. This elementary local-system background underlies the bundle formulation used there.}
\GSfield{Present status}{As of 9 November 2026: standard sheaf theory and the intended foundation for the Hodge discussion.}
\GSfield{Assessment}{literal statement: correct; publication match: partial; verification: read the result; development: none; confidence: high}
\end{GSitem}
\begin{GSitem}{1972-275.2}{E}{Integral local systems from locally constant transition functions}
\GSfield{Kind; source}{construction; \GSpage{P000456}}
\GSfield{The notes}{Locally \(H_{\mathbb Z}|_U\cong U\times\mathbb Z^n\), \(\Gamma(U,H_{\mathbb Z})\cong\mathbb Z^n\), and transition functions take values in \(\mathrm{GL}_n(\mathbb Z)\); connectedness of \(S\) makes the rank constant.}
\GSfield{Statement}{A rank-\(r\) integral local system is obtained by gluing constant sheaves on a trivializing cover with locally constant \(\mathrm{GL}_r(\mathbb Z)\)-valued transition functions satisfying the cocycle condition. Its rank is locally constant, hence constant on connected \(S\). The section formula \(\Gamma(U,H_{\mathbb Z})\cong\mathbb Z^r\) holds for connected trivializing \(U\).}
\GSfield{Reconstruction}{The displayed product is the local \(\acute{e}\)tal\'{e}-space model; \(U\) must be connected for the asserted section group. The cocycle condition is implicit in the gluing language and is stated explicitly here.}
\GSfield{Fate}{The exposition predates this talk: \citet[\S2, pp.~232--235]{griffiths1970periods} summarizes the period-domain formalism in an invited lecture delivered on 29 August 1968, received on 18 September 1969, and published in March 1970. This revisits 1972-272.1 with a transition-function description.}
\GSfield{Present status}{As of 9 November 2026: standard gluing description, not a newly posed classification problem.}
\GSfield{Assessment}{literal statement: correct with missing hypotheses; publication match: partial; verification: read the result; development: none; confidence: high}
\end{GSitem}
\begin{GSitem}{1972-275.3}{E}{Associated real, complex, and holomorphic bundles}
\GSfield{Kind; source}{construction; \GSpage{P000456}}
\GSfield{The notes}{Using the same transition matrices, the notes construct \(H_{\mathbb R}\), \(H_{\mathbb C}\), and the sheaf \(H_{\mathcal O}\) of holomorphic sections, with \(H_{\mathcal O}=H_{\mathbb Z}\otimes_{\mathbb Z}\mathcal O_S\).}
\GSfield{Statement}{Extend the transition matrices of an integral local system along \(\mathbb Z\subset\mathbb R\subset\mathbb C\). This produces real and complex flat vector bundles; the latter's locally constant transition functions are holomorphic and define a holomorphic bundle whose sheaf of sections is \(\mathcal H=H_{\mathbb Z}\otimes\mathcal O_S\).}
\GSfield{Reconstruction}{The local system of flat sections, the total vector bundle, and the sheaf of all holomorphic sections are distinguished explicitly; the notebook uses one symbol in more than one of these senses.}
\GSfield{Fate}{The exposition predates this talk: \citet[\S2, pp.~232--235]{griffiths1970periods} summarizes the period-domain formalism in an invited lecture delivered on 29 August 1968, received on 18 September 1969, and published in March 1970.}
\GSfield{Present status}{As of 9 November 2026: the standard passage from integral monodromy to a flat holomorphic bundle.}
\GSfield{Assessment}{literal statement: correct; publication match: exact; verification: read the result; development: none; confidence: high}
\end{GSitem}
\begin{GSitem}{1972-275.4}{E}{Flat connection in a locally constant frame}
\GSfield{Kind; source}{construction; \GSpage{P000456-P000457}}
\GSfield{The notes}{The connection is \(\mathbb C\)-linear, satisfies \(\nabla(\varphi f)=\varphi\nabla f+d\varphi\otimes f\), and in a constant frame \(u_i\) is \(\nabla\sum_i\varphi_i u_i=\sum_i d\varphi_i\otimes u_i\). Its curvature is zero.}
\GSfield{Statement}{For a holomorphic bundle defined by locally constant transition matrices, the local operator \(d\) in each constant frame glues to a globally defined holomorphic connection. It is complex-linear, satisfies the displayed Leibniz rule, and its extension to bundle-valued forms has \(\nabla^2=0\). Its horizontal sheaf is exactly the initial complex local system.}
\GSfield{Reconstruction}{The connection is not \(\mathcal O_S\)-linear; complex-linearity and the Leibniz rule are the correct axioms. Zero curvature is both explicitly stated and justified by the local formula, so it is not merely an omitted hypothesis here.}
\GSfield{Fate}{The exposition predates this talk: \citet[\S2, pp.~232--235]{griffiths1970periods} summarizes the period-domain formalism in an invited lecture delivered on 29 August 1968, received on 18 September 1969, and published in March 1970. This expands the flat/Gauss--Manin construction in 1972-272.2.}
\GSfield{Present status}{As of 9 November 2026: a complete standard local construction with no unresolved step.}
\GSfield{Assessment}{literal statement: correct; publication match: exact; verification: read the result; development: none; confidence: high}
\end{GSitem}
\begin{GSitem}{1972-275.5}{E}{Horizontal polarization and transversality of the Hodge filtration}
\GSfield{Kind; source}{definition; \GSpage{P000457}}
\GSfield{The notes}{The holomorphic filtration is again required to satisfy Griffiths' condition, and a locally constant bilinear form \(\psi\) is required to produce a Hodge structure at every point.}
\GSfield{Statement}{For the flat holomorphic bundle \((\mathcal H,\nabla)\), require decreasing holomorphic subbundles \(F^p\mathcal H\) with \(\nabla F^p\subset\Omega_S^1\otimes F^{p-1}\), and a flat bilinear form \(Q\). Together with fibrewise opposition, weight parity, and positivity, these give a polarized variation of pure Hodge structure.}
\GSfield{Reconstruction}{The form's local constancy means \(\nabla Q=0\), and the filtration need not be flat. The later fibrewise relations, including the parity correction in the next item, are incorporated only as explicitly identified components of the definition.}
\GSfield{Fate}{The exposition predates this talk: \citet[\S2, pp.~232--235]{griffiths1970periods} summarizes the period-domain formalism in an invited lecture delivered on 29 August 1968, received on 18 September 1969, and published in March 1970. The transversality equation is (2.1). This is a repeat of 1972-272.3, not an independent research problem.}
\GSfield{Present status}{As of 9 November 2026: standard polarized variation; the notes' repetition is preserved in the occurrence index.}
\GSfield{Assessment}{literal statement: correct; publication match: exact; verification: read the result; development: none; confidence: high}
\end{GSitem}
\begin{GSitem}{1972-275.6}{E}{Parity of the polarization is reversed in the notebook}
\GSfield{Kind; source}{claim; \GSpage{P000458}}
\GSfield{The notes}{The form is said to be symmetric for odd \(n\) and skew-symmetric for even \(n\). Both Russian and English transcriptions follow the handwriting.}
\GSfield{Statement}{For a pure polarized Hodge structure of weight \(n\), the bilinear polarization satisfies \(Q(y,x)=(-1)^n Q(x,y)\): it is symmetric for even \(n\) and alternating for odd \(n\). With the notebook's Weil-operator convention, an overall real sign can be chosen so that \(Q(x,C\bar y)\) is positive Hermitian.}
\GSfield{Reconstruction}{The parity reversal is a faithful mathematical error, not a translation correction. Changing an overall sign of \(Q\) cannot change symmetric into alternating; it therefore cannot explain the reversed statement.}
\GSfield{Fate}{\citet[\S1, equations (1.7)--(1.15)]{griffiths1968periods1} gives the correct parity and bilinear relations before the talk; \citet[\S2]{griffiths1970periods} summarizes the same framework.}
\GSfield{Present status}{As of 9 November 2026: the literal parity claim is false, and the standard correction is unambiguous.}
\GSfield{Assessment}{literal statement: false as written; publication match: partial; verification: read the result; development: annotated; confidence: high}
\end{GSitem}
\begin{GSitem}{1972-275.7}{E}{Projective flag variety and the corrected isotropic compact dual}
\GSfield{Kind; source}{construction; \GSpage{P000457-P000458}}
\GSfield{The notes}{For dimensions \(k_0\geq k_1\geq\cdots\geq k_n\), the flag space is projective algebraic. The following page now writes \(\psi(F^p,F^{n-p+1})=0\) with the correct index.}
\GSfield{Statement}{In a fixed complex vector space, flags of dimensions \(k_p\) form a projective partial flag variety. Imposing the closed bilinear relations \(Q(F^p,F^{n-p+1})=0\), with dimensions compatible with a polarized weight-\(n\) Hodge structure, gives the compact dual; a connected component is chosen when necessary.}
\GSfield{Reconstruction}{The dimensions are constrained by the intended Hodge numbers and the nondegenerate form; arbitrary numerical flags need not admit polarized Hodge points. P000458 corrects the earlier index internally, but does not explicitly announce an erratum.}
\GSfield{Fate}{The exposition predates this talk: \citet[\S2, pp.~232--235]{griffiths1970periods} summarizes the period-domain formalism in an invited lecture delivered on 29 August 1968, received on 18 September 1969, and published in March 1970. The flag variety and the bilinear subvariety are separated here to make the algebraic and positivity conditions transparent.}
\GSfield{Present status}{As of 9 November 2026: standard projective compact dual; the correct formula links back to the emendation of 1972-272.4.}
\GSfield{Assessment}{literal statement: correct with missing hypotheses; publication match: exact; verification: read the result; development: none; confidence: high}
\end{GSitem}
\begin{GSitem}{1972-275.8}{E}{Opposition, Weil operator, and positivity define the period domain}
\GSfield{Kind; source}{construction; \GSpage{P000458}}
\GSfield{The notes}{The page imposes \(F^p\oplus\overline{F^{n-p+1}}=H_{\mathbb C}\), defines the Weil operator by Hodge type, requires \(\psi(x,C\bar y)\) positive Hermitian, and calls \(D\) open in \(\overline D\).}
\GSfield{Statement}{After correcting the parity of \(Q\), the opposition identities recover the Hodge decomposition; let \(C|_{H^{p,n-p}}=i^{2p-n}\). The flags for which \(Q(x,C\bar y)\) is positive definite form the period domain, open in the isotropic compact dual, determined by the Hodge numbers and the fixed real polarization form.}
\GSfield{Reconstruction}{This combines the page's related fibrewise defining relations into one item, while retaining parity and the compact-dual construction as separate claims because each has its own correction or hypothesis.}
\GSfield{Fate}{The exposition predates this talk: \citet[\S2, pp.~232--235]{griffiths1970periods} summarizes the period-domain formalism in an invited lecture delivered on 29 August 1968, received on 18 September 1969, and published in March 1970.}
\GSfield{Present status}{As of 9 November 2026: classical period-domain construction; this repeats 1972-272.5--6 with no new question.}
\GSfield{Assessment}{literal statement: correct; publication match: exact; verification: read the result; development: none; confidence: high}
\end{GSitem}
\begin{GSitem}{1972-275.9}{E}{Announced homogeneity and the identity-component qualification}
\GSfield{Kind; source}{claim; \GSpage{P000458}}
\GSfield{The notes}{For next time, the complex isometry group is to be shown transitive on \(\overline D\), and the identity component of the real group transitive on \(D\).}
\GSfield{Statement}{The relevant complex isometry group acts transitively on each appropriate isotropic flag orbit. The identity component \(G_{\mathbb R}^0\) acts transitively on a chosen connected component of the polarized period domain. If \(D\) denotes all components at once, transitivity of the identity component need not hold.}
\GSfield{Reconstruction}{A connected-component choice is supplied. The notes announce an elementary homogeneous-space proof for a subsequent meeting; they do not formulate homogeneity as an open problem.}
\GSfield{Fate}{\citet[\S2, especially equation (2.3)]{griffiths1970periods} already uses the homogeneous presentation. Its underlying exposition circulated in the 1968 invited lecture and was published in 1970.}
\GSfield{Present status}{As of 9 November 2026: established homogeneity with the component convention stated explicitly.}
\GSfield{Assessment}{literal statement: correct with missing hypotheses; publication match: exact; verification: read the result; development: annotated; confidence: high}
\end{GSitem}

\GStalk{276}{1972-10-23}{V. M. Buchstaber}{P000459-P000463}{Formal groups, logarithms, and the Lazard ring}
An extended exposition starts with one-dimensional formal laws over arbitrary commutative rings, strict coordinate changes and invariant differentials, then constructs the universal coefficient ring and states Lazard's polynomial-ring theorem. Its proposed proof passes to indecomposables and polynomial two-cocycles, then to comodules and the divided-power dual. The cubic example contains a genuine arithmetic error, and the initial claim that associativity forces commutativity is explicitly corrected by the speaker's later example in record 285. These notes continue the formal-group theme of occurrence 1971-268.11; their content does not establish a new 1972 theorem.

\begin{GSitem}{1972-276.1}{E}{One-dimensional formal laws and the formal inverse}
\GSfield{Kind; source}{definition; \GSpage{P000459}}
\GSfield{The notes}{Over a commutative ring \(A\) with \(1\), the notes write \(F(x,y)=x+y+\sum_{i+k\ge2}a_{ik}x^iy^k\), require associativity, and obtain an inverse from the formal implicit-function argument.}
\GSfield{Statement}{An associative one-dimensional formal law over \(A\) is a series \(F\in A[[x,y]]\) with \(F(x,0)=x\), \(F(0,y)=y\), and \(F(F(x,y),z)=F(x,F(y,z))\). There is a unique \(\iota(x)\in xA[[x]]\) satisfying \(F(x,\iota(x))=F(\iota(x),x)=0\); recursively \(\iota(x)=-x+O(x^2)\). Commutativity is an additional issue, recorded separately.}
\GSfield{Reconstruction}{The unit identities are made explicit; they also follow from associativity and the displayed linear term. The term ``implicit-function theorem'' is formal coefficient recursion, requiring no analytic convergence.}
\GSfield{Fate}{Classical definition and inverse construction: \citet[\S I]{lazard1955formal}.}
\GSfield{Present status}{As of 9 November 2026: an established algebraic definition.}
\GSfield{Assessment}{literal statement: correct; publication match: exact; verification: read the result; development: none; confidence: high}
\end{GSitem}
\begin{GSitem}{1972-276.2}{E}{Associativity does not force commutativity over every ring}
\GSfield{Kind; source}{claim; \GSpage{P000459}}
\GSfield{The notes}{The parenthesis reads \ru{коммутативность (она сл. из ассоц.)}: commutativity follows from associativity.}
\GSfield{Statement}{The unrestricted assertion is false. It is true over a reduced coefficient ring, and also over a torsion-free coefficient ring. Over \(\mathbb F_p[a]/(a^2)\), \(F(x,y)=x+y+axy^p\) is associative and not commutative; the notebook itself supplies this counterexample in 1972-285.6.}
\GSfield{Reconstruction}{The extra hypotheses are editorial repairs, not present in this parenthesis. The counterexample is a later occurrence in the same notebook, not a conjectural interpretation. The earlier 1971-270.7 had already distinguished the characteristic-zero commutativity statement.}
\GSfield{Fate}{Already distinguished in \citet[Theorem I; \S III, Remark (3.14a)]{lazard1955formal}.}
\GSfield{Present status}{As of 9 November 2026: the literal assertion is refuted; the corrected versions are established.}
\GSfield{Assessment}{literal statement: false as written; publication match: exact; verification: read the result; development: annotated; confidence: high}
\end{GSitem}
\begin{GSitem}{1972-276.3}{E}{Strict coordinate changes versus classification of laws}
\GSfield{Kind; source}{definition; \GSpage{P000459}}
\GSfield{The notes}{The change \(z=\varphi(x)=x+a_1x^2+\cdots\) defines isomorphism. The speaker expressly says that classifying laws is different from classifying their isomorphism classes.}
\GSfield{Statement}{A strict isomorphism \(F\to G\) is \(\varphi\in xA[[x]]\) with \(\varphi'(0)=1\) and \(\varphi(F(x,y))=G(\varphi(x),\varphi(y))\). A general isomorphism instead allows \(\varphi'(0)\in A^\times\). The functor of laws and its quotient by coordinate changes are different classification problems.}
\GSfield{Reconstruction}{The adjective ``strict'' is supplied. The inverse expansion should be \(\varphi^{-1}(z)=z-a_1z^2+\cdots\); the notebook mixes the argument \(x\) with the output variable \(z\).}
\GSfield{Fate}{Classical coordinate-change action: \citet[\S I, (1.4)--(1.8)]{lazard1955formal}; earlier occurrence 1971-270.5 distinguishes full and strict changes.}
\GSfield{Present status}{As of 9 November 2026: standard, with the law/moduli distinction preserved.}
\GSfield{Assessment}{literal statement: correct; publication match: exact; verification: read the result; development: none; confidence: high}
\end{GSitem}
\begin{GSitem}{1972-276.4}{E}{Logarithm of a formal law}
\GSfield{Kind; source}{definition; \GSpage{P000459}}
\GSfield{The notes}{A series \(g(x)=x+\cdots\) is called a logarithm when \(F(x,y)=g^{-1}(g(x)+g(y))\).}
\GSfield{Statement}{A normalized logarithm is a strict homomorphism \(g:F\to\widehat{\mathbb G}_a\): \(g(F(x,y))=g(x)+g(y)\), \(g(0)=0\), \(g'(0)=1\). Its formal inverse exists because its linear coefficient is a unit.}
\GSfield{Fate}{Pre-lecture formulation: \citet[\S1]{honda1968formal}.}
\GSfield{Present status}{As of 9 November 2026: standard; integral or positive-characteristic existence is not automatic.}
\GSfield{Assessment}{literal statement: correct; publication match: exact; verification: read the result; development: none; confidence: high}
\end{GSitem}
\begin{GSitem}{1972-276.5}{E}{Canonical invariant differential}
\GSfield{Kind; source}{construction; \GSpage{P000460}}
\GSfield{The notes}{For the shift \(x\mapsto F(x,t)\), invariance is \(\psi(F(x,t))\partial_xF(x,t)=\psi(x)\). The canonical differential is \(\omega_F=dx/\partial_yF(x,0)\).}
\GSfield{Statement}{For a commutative one-dimensional formal law over a commutative ring \(A\), put \(b(x)=\partial_yF(x,y)|_{y=0}\), so \(b(0)=1\). Then \(\omega_F=b(x)^{-1}dx\) is translation invariant. Every invariant differential is \(c\omega_F\), \(c\in A\); thus normalization at the identity gives uniqueness.}
\GSfield{Reconstruction}{The inverse is the formal inverse of the unit series \(b\); the parameter is held fixed during differentiation. No analytic differential form or integration is assumed.}
\GSfield{Fate}{Pre-lecture invariant-differential construction: \citet[\S1]{honda1968formal}; earlier occurrences 1971-268.12 and 1971-270.8 cover invariant differentials and the logarithm.}
\GSfield{Present status}{As of 9 November 2026: established over arbitrary coefficient rings for commutative laws.}
\GSfield{Assessment}{literal statement: correct; publication match: exact; verification: read the result; development: none; confidence: high}
\end{GSitem}
\begin{GSitem}{1972-276.6}{E}{Existence and uniqueness of the normalized rational logarithm}
\GSfield{Kind; source}{theorem; \GSpage{P000459-P000460}}
\GSfield{The notes}{The ``Lazard lemma'' says that if \(A\) is a \(\mathbb Q\)-module, every formal group has a logarithm. The proof integrates \(\omega_F\), then evaluates at \(x=0\) to identify the integration constant.}
\GSfield{Statement}{If \(A\) is a commutative unital \(\mathbb Q\)-algebra, then \(g(x)=\int_0^x b(t)^{-1}\,dt\) is the unique normalized logarithm of a commutative one-dimensional formal law. Hence any two such laws over \(A\) are uniquely strictly isomorphic by the composite of their logarithms and inverse logarithms.}
\GSfield{Reconstruction}{``\(\mathbb Q\)-module'' is interpreted as a compatible \(\mathbb Q\)-algebra structure on the given ring. The printed pointer ``Mat. 13:6'' resolves to Honda, not to a Russian publication of Lazard.}
\GSfield{Fate}{\citet[\S1]{honda1968formal}: original 1968, Russian translation 1969.}
\GSfield{Present status}{As of 9 November 2026: proved; denominators explain the restriction on the base.}
\GSfield{Assessment}{literal statement: correct with missing hypotheses; publication match: exact; verification: read the result; development: none; confidence: high}
\end{GSitem}
\begin{GSitem}{1972-276.7}{E}{The universal rational law}
\GSfield{Kind; source}{construction; \GSpage{P000460}}
\GSfield{The notes}{With \(g(x)=x+\sum_{i\ge1}g_ix^{i+1}\), the ring \(\mathbb Q[g_1,g_2,\ldots]\) carries a universal law, and every law is obtained by a coefficient homomorphism.}
\GSfield{Statement}{The functor sending a commutative \(\mathbb Q\)-algebra \(A\) to its commutative one-dimensional formal laws is represented by \(\mathbb Q[g_1,g_2,\ldots]\), using \(F=g^{-1}(g(x)+g(y))\). The specialization sends \(g_i\) to the normalized logarithm coefficients. This represents laws in a chosen coordinate, not isomorphism classes.}
\GSfield{Fate}{Pre-lecture rational universal construction: \citet[\S1]{buchstabermishchenkonovikov1971formal}.}
\GSfield{Present status}{As of 9 November 2026: established and consistent with rationalization of the Lazard ring.}
\GSfield{Assessment}{literal statement: correct; publication match: exact; verification: read the result; development: none; confidence: high}
\end{GSitem}
\begin{GSitem}{1972-276.8}{E}{Presentation of the integral universal coefficient ring}
\GSfield{Kind; source}{construction; \GSpage{P000461}}
\GSfield{The notes}{The notes form \(\Lambda=\mathbb Z[a_{ij}]/J\), where \(J\) is generated by associativity coefficients and \(a_{ij}-a_{ji}\), with \(\deg a_{ij}=i+j-1\).}
\GSfield{Statement}{Let \(i,j\ge1\), put \(F=x+y+\sum a_{ij}x^iy^j\), and quotient the polynomial ring by all coefficients of its associator and by all symmetry relations. For every commutative unital \(A\), ring maps \(\Lambda\to A\) correspond bijectively to commutative one-dimensional formal laws over \(A\).}
\GSfield{Reconstruction}{The positive indices exclude pure higher powers and incorporate the identity laws. The grading uses half the usual topological degree.}
\GSfield{Fate}{Universal property: \citet[Theorem~1.2]{buchstabermishchenkonovikov1971formal}.}
\GSfield{Present status}{As of 9 November 2026: established; representability does not itself prove polynomiality.}
\GSfield{Assessment}{literal statement: correct; publication match: exact; verification: read the result; development: none; confidence: high}
\end{GSitem}
\begin{GSitem}{1972-276.9}{E}{Lazard polynomiality and rationalization}
\GSfield{Kind; source}{theorem; \GSpage{P000461}}
\GSfield{The notes}{\(\Lambda=\mathbb Z[a_i]\), one generator of every degree \(i\ge1\), and \(\Lambda\otimes\mathbb Q=\mathbb Q[g_i]\).}
\GSfield{Statement}{The coefficient ring of the universal commutative one-dimensional formal law is noncanonically a graded polynomial ring \(\mathbb Z[t_1,t_2,\ldots]\), \(\deg t_i=i\). Under rationalization logarithm coefficients give polynomial coordinates. A choice of integral generators is not canonical.}
\GSfield{Fate}{\citet[Theorem II]{lazard1955formal}; circulated December 1954, published 1955.}
\GSfield{Present status}{As of 9 November 2026: established; previous occurrence 1971-268.11 concerns the same universal-ring theorem.}
\GSfield{Assessment}{literal statement: correct; publication match: exact; verification: read the result; development: none; confidence: high}
\end{GSitem}
\begin{GSitem}{1972-276.10}{E}{Indecomposables linearize associativity}
\GSfield{Kind; source}{method; \GSpage{P000461}}
\GSfield{The notes}{The quotient by decomposables makes products of positive-degree coefficients zero and turns the associator into a linear equation for \(h=F-x-y\). A marginal gloss wrongly describes decomposables using arbitrary nonunits.}
\GSfield{Statement}{For the augmented graded ring \(\Lambda\), use \(I=\bigoplus_{n>0}\Lambda_n\). Indecomposables are \(I/I^2\), or the augmented ring \(\Lambda/I^2\) if degree zero is retained. Modulo \(I^2\), associativity is \(h(x,y)+h(x+y,z)=h(y,z)+h(x,y+z)\). Products of positive-degree elements, not products of arbitrary nonunits, generate \(I^2\).}
\GSfield{Reconstruction}{The margin is repaired using the unambiguous main-text phrase ``products of two of nonzero degree''. Taking nonunits literally would include integers such as \(2\) and alter degree zero.}
\GSfield{Fate}{Classical linearization: \citet[\S~II, Lemma~2]{lazard1955formal}.}
\GSfield{Present status}{As of 9 November 2026: valid after the augmentation-ideal correction.}
\GSfield{Assessment}{literal statement: false as written; publication match: exact; verification: read the result; development: annotated; confidence: high}
\end{GSitem}
\begin{GSitem}{1972-276.11}{E}{Classify homogeneous additive two-cocycles}
\GSfield{Kind; source}{problem; \GSpage{P000462}}
\GSfield{The notes}{\ru{Задача: найти все решения} of \(h(x,y)+h(x+y,z)=h(y,z)+h(x,y+z)\), homogeneous polynomials over a ring \(K\). The surrounding formal-group problem requires symmetry; the page also temporarily studies nonsymmetric cycles.}
\GSfield{Statement}{For each \(n\ge2\), classify normalized homogeneous \(h\in K[x,y]\) of degree \(n\) satisfying the displayed cocycle equation, and distinguish the symmetric submodule. In the symmetric case every solution is uniquely \(aC_n\), where \(C_n=((x+y)^n-x^n-y^n)/d_n\), \(d_n=p\) if \(n\) is a power of a prime \(p\), and \(d_n=1\) otherwise. This formula holds over every commutative ring via integral coefficients.}
\GSfield{Reconstruction}{The exact symmetric answer is supplied from literature; the notebook states the task and develops its cohomological calculation over subsequent meetings. Nonsymmetric cycles must not be discarded before the distinction is made.}
\GSfield{Fate}{Symmetric classification: \citet[\S~II, Lemma~3]{lazard1955formal}.}
\GSfield{Present status}{As of 9 November 2026: the symmetric problem was already solved; the lecture is an exposition, not an orphan open question.}
\GSfield{Assessment}{literal statement: correct; publication match: exact; verification: read the result; development: none; confidence: high; problem \hyperref[prob:GP-1972-201]{GP-1972-201}}
\end{GSitem}
\begin{GSitem}{1972-276.12}{E}{Primitive polynomial Hopf algebra and comodules}
\GSfield{Kind; source}{construction; \GSpage{P000462}}
\GSfield{The notes}{\(K[x]\) has \(\Delta x=x\otimes1+1\otimes x\); a comodule is described as \(N\to A\otimes N\), and \(K\) is the trivial comodule.}
\GSfield{Statement}{For the graded coalgebra \(C=K[x]\) with primitive generator, a left comodule is a \(K\)-module \(N\) with coaction \(\rho:N\to C\otimes_KN\) satisfying coassociativity and the counit identity. The trivial coaction on \(K\) is \(1\mapsto1\otimes1\). Internal polynomial degree and cohomological degree produce the two gradings.}
\GSfield{Reconstruction}{The base-module structure and both coaction axioms are made explicit. ``Abelian group'' alone is insufficient for a general coefficient ring.}
\GSfield{Fate}{Pre-existing homological apparatus; compare \citet[Chs.VIII,X]{cartaneilenberg1956homological}.}
\GSfield{Present status}{As of 9 November 2026: standard with the specified category.}
\GSfield{Assessment}{literal statement: correct with missing hypotheses; publication match: exact; verification: publisher metadata checked; development: none; confidence: high}
\end{GSitem}
\begin{GSitem}{1972-276.13}{E}{The two-cocycle exact sequence requires coalgebra cohomology}
\GSfield{Kind; source}{claim; \GSpage{P000462}}
\GSfield{The notes}{The page writes \(0\to\{c((x+y)^n-x^n-y^n)\}\to L\to\operatorname{Ext}^{2,n}_{K[x]}(K,K)\to0\).}
\GSfield{Statement}{For the normalized cobar complex of the primitive coalgebra \(K[x]\), \(Z^{2,n}=\{h:\delta h=0\}\), \(B^{2,n}=\{c((x+y)^n-x^n-y^n):c\in K\}\), and \(0\to B^{2,n}\to Z^{2,n}\to H^{2,n}\to0\). This is comodule cohomology (or dual divided-power algebra cohomology under suitable graded duality), not ordinary module Ext over the polynomial algebra.}
\GSfield{Reconstruction}{The category is an essential emendation. Over a field, ordinary \(\operatorname{Ext}_{K[x]}^2(K,K)=0\), while the cobar group in characteristic \(p\) need not vanish. The leftmost term is the image subgroup, not an injective copy of every scalar \(c\).}
\GSfield{Fate}{Cocycle/boundary calculation matches \citet[\S~II--III]{lazard1955formal}.}
\GSfield{Present status}{As of 9 November 2026: correct as a cobar exact sequence; the unqualified Ext notation is misleading.}
\GSfield{Assessment}{literal statement: correct with missing hypotheses; publication match: partial; verification: read the result; development: annotated; confidence: high}
\end{GSitem}
\begin{GSitem}{1972-276.14}{E}{Quadratic and cubic cocycles: a mistaken characteristic}
\GSfield{Kind; source}{claim; \GSpage{P000462}}
\GSfield{The notes}{For degree 2 the notes give \(cxy\). For degree 3 they give \(C_1x^2y+C_2xy^2\), write \(\widetilde\Lambda=\mathbb Z+\mathbb Z_3\), and claim \(x^2y\) is a cocycle over \(\mathbb Z_3\)-modules.}
\GSfield{Statement}{Degree2 normalized cocycles are \(cxy\). In degree 3, direct expansion gives \(\delta(C_1x^2y+C_2xy^2)=2(C_1-C_2)xyz\), with \(\delta h=h(x,y)+h(x+y,z)-h(y,z)-h(x,y+z)\). Thus \(x^2y\) is a cocycle when \(2=0\), not when \(3=0\). The universal coefficient module without symmetry is \(\mathbb Z^2/\langle2(e_1-e_2)\rangle\simeq\mathbb Z\oplus\mathbb Z/2\); imposing symmetry leaves \(\mathbb Z\).}
\GSfield{Reconstruction}{\oursComputed\ The polynomial expansion and independent evaluations in characteristics 2 and 3 were both checked in the included formal-check script. The two visible subscripts3 in the scan are faithful notebook errors, not transcription corrections.}
\GSfield{Fate}{Corrected symmetric case agrees with \citet[\S~II, Lemma~3]{lazard1955formal}.}
\GSfield{Present status}{As of 9 November 2026: the literal characteristic 3 claim is disproved by the displayed defect.}
\GSfield{Assessment}{literal statement: false as written; publication match: partial; verification: read the result; development: annotated; confidence: high}
\end{GSitem}
\begin{GSitem}{1972-276.15}{E}{Minimal resolutions simplify the cohomology differential}
\GSfield{Kind; source}{method; \GSpage{P000462}}
\GSfield{The notes}{``Choose the resolution correctly: all cycles without boundaries (minimal resolution).''}
\GSfield{Statement}{For a connected graded algebra \(A\) over a field \(k\), a minimal graded free resolution of the augmentation module has differential with entries in \(A_{>0}\). Applying \(\operatorname{Hom}_A(-,k)\) makes that differential zero. Such a minimality argument does not hold unchanged over every coefficient ring.}
\GSfield{Reconstruction}{The field and connected-grading hypotheses specify a valid version of the terse method; the notes have not established minimality over arbitrary \(K\).}
\GSfield{Fate}{Established resolution method; \citet[Chs.III,XII]{maclane1963homology}.}
\GSfield{Present status}{As of 9 November 2026: standard under the stated hypotheses.}
\GSfield{Assessment}{literal statement: correct with missing hypotheses; publication match: partial; verification: publisher metadata checked; development: none; confidence: high}
\end{GSitem}
\begin{GSitem}{1972-276.16}{E}{Divided-power multiplication on the dual}
\GSfield{Kind; source}{construction; \GSpage{P000463}}
\GSfield{The notes}{\(K_*[x]=\operatorname{Hom}_K(K[x],K)\), \((x_*^n,x^m)=\delta_{nm}\), and \(x_*^nx_*^m=\binom{n+m}{n}x_*^{n+m}\).}
\GSfield{Statement}{The graded dual of the primitive polynomial coalgebra has basis \(X_n\) and multiplication \(X_nX_m=\binom{n+m}{n}X_{n+m}\). This is the divided-power algebra. The full dual is its degree completion; the distinction matters when forming tensor products and duals.}
\GSfield{Reconstruction}{The completion convention is supplied because the notebook initially uses full Hom, later continuous duality (P476).}
\GSfield{Fate}{Standard Hopf-dual construction; \citet[Ch.X]{cartaneilenberg1956homological}.}
\GSfield{Present status}{As of 9 November 2026: established; the displayed multiplication is independently determined by the binomial coproduct.}
\GSfield{Assessment}{literal statement: correct with missing hypotheses; publication match: exact; verification: publisher metadata checked; development: none; confidence: high}
\end{GSitem}
\begin{GSitem}{1972-276.17}{E}{Characteristic-p factorization of the divided-power algebra}
\GSfield{Kind; source}{claim; \GSpage{P000463}}
\GSfield{The notes}{The dual is written as \(\prod B_q\), with \(B_q=\mathbb Z_p[y_{p^i}]/(y_{p^i}^{p})\), and the text says it is enough to consider \(\mathbb Z_p\) and \(\mathbb Q\).}
\GSfield{Statement}{Over \(k=\mathbb F_p\), the graded divided-power algebra is \(\bigotimes_{i\ge0}k[Y_i]/(Y_i^p)\), \(\deg Y_i=p^i\); its full dual version is the degree-completed tensor product. Over \(\mathbb Q\) it is isomorphic to a polynomial algebra by \(X_n\mapsto t^n/n!\). Passage from field calculations to an integral conclusion requires the finite-generation and torsion argument in each degree.}
\GSfield{Reconstruction}{\(\mathbb Z_p\) here denotes the field with \(p\) elements, not the p-adic integers. An ordinary direct product of algebras would introduce unrelated idempotents and is not the asserted factorization.}
\GSfield{Fate}{Established homological calculation; compare \citet[Ch.X]{cartaneilenberg1956homological}.}
\GSfield{Present status}{As of 9 November 2026: valid after distinguishing tensor product, completion, and coefficient fields.}
\GSfield{Assessment}{literal statement: correct with missing hypotheses; publication match: partial; verification: publisher metadata checked; development: annotated; confidence: high}
\end{GSitem}

\GStalk{277}{1972-10-30}{I. M. Gelfand}{P000464-P000465}{Linear classification, subspaces and Dynkin diagrams}
Under the general Bernstein heading, the explicit \ru{И.М.Г.} intervention begins on P000464 and continues to the separator before \ru{И.Б.} on P000465. Gelfand organizes matrix and subspace classification through directed graphs, contrasts finite classification with arbitrary pairs of operators, and reports the Dynkin and affine-Dynkin boundary. The notebook connects these diagrams to singularities. The catalogue range contains Buchstaber material at its ends; that material is not attributed to this record. The four-subspace discussion continues the line of 1970-198.1--1970-198.8, without repeating the separate deformation problem GP-1970-23.

\begin{GSitem}{1972-277.1}{E}{Representations of a directed graph without imposed relations}
\GSfield{Kind; source}{definition; \GSpage{P000464}}
\GSfield{The notes}{The graph is called a category; vector spaces stand at vertices and maps on arrows, with \ru{функторность не требуется}. The task is classification.}
\GSfield{Statement}{For a finite quiver \(Q\) and field \(k\), a finite-dimensional representation assigns \(V_i\) to every vertex and \(A_a:V_{s(a)}\to V_{t(a)}\) to every arrow. Isomorphisms are invertible vertex maps intertwining all arrows. This is a functor from the free path category of \(Q\); no relations equating distinct paths are imposed.}
\GSfield{Reconstruction}{The free path category repairs the apparent conflict between calling the graph a category and rejecting functoriality. A graph itself has no specified composites.}
\GSfield{Fate}{This formulation is explicit in \citet[Definition 1, p.~17]{bernsteingelfandponomarev1973coxeter}, received 18 December 1972, and is already attributed there to Gabriel.}
\GSfield{Present status}{As of 9 November 2026: established with the hypotheses and distinctions stated here.}
\GSfield{Assessment}{literal statement: false as written; publication match: exact; verification: read the result; development: annotated; confidence: high}
\end{GSitem}
\begin{GSitem}{1972-277.2}{E}{A single arrow requires rank as well as kernel and cokernel}
\GSfield{Kind; source}{claim; \GSpage{P000464}}
\GSfield{The notes}{The only invariants of a map are said to be \(\dim\ker\) and \(\dim\operatorname{coker}\); the problem is reduced to an embedding.}
\GSfield{Statement}{An arbitrary linear map between unspecified finite-dimensional spaces is classified by the triple \((\dim\ker A,\operatorname{rank}A,\dim\operatorname{coker}A)\). If source and target dimensions are fixed, rank alone is sufficient. The representation decomposes into copies of \(k\to0\), \(k\xrightarrow{1}k\), and \(0\to k\).}
\GSfield{Reconstruction}{An elementary check \ours: the missing rank is essential, since the identity maps on \(k\) and \(k^2\) have equal zero kernels and cokernels but are nonisomorphic as arrow representations.}
\GSfield{Fate}{The elementary matrix normal form is reviewed in \citet[Example 1, p.~750]{nazarova1973quivers}; it was classical before the talk.}
\GSfield{Present status}{As of 9 November 2026: established with the hypotheses and distinctions stated here.}
\GSfield{Assessment}{literal statement: false as written; publication match: exact; verification: read the result; development: annotated; confidence: high}
\end{GSitem}
\begin{GSitem}{1972-277.3}{E}{One loop is the similarity problem for one operator}
\GSfield{Kind; source}{theorem; \GSpage{P000464}}
\GSfield{The notes}{A one-loop graph is labelled one transformation.}
\GSfield{Statement}{A representation of the one-loop quiver is a finite-dimensional \(k[t]\)-module. Over an algebraically closed field it is a direct sum of Jordan blocks \(J_m(\lambda)\), uniquely up to order; over an arbitrary field use elementary divisors for monic irreducibles.}
\GSfield{Reconstruction}{The notebook gives the problem, not the Jordan normal-form theorem; the latter is the precise modern resolution.}
\GSfield{Fate}{The standard similarity interpretation is stated in \citet[Example 2, p.~750]{nazarova1973quivers}.}
\GSfield{Present status}{As of 9 November 2026: established with the hypotheses and distinctions stated here.}
\GSfield{Assessment}{literal statement: correct; publication match: exact; verification: read the result; development: none; confidence: high}
\end{GSitem}
\begin{GSitem}{1972-277.4}{E}{Arbitrary matrix pairs reduce to commuting pairs}
\GSfield{Kind; source}{theorem; \GSpage{P000464}}
\GSfield{The notes}{Two loops are called an unsolvable problem; the notes say that the arbitrary pair reduces to commuting operators, citing Gelfand and Ponomarev.}
\GSfield{Statement}{For finite-dimensional vector spaces over a field, simultaneous similarity of arbitrary tuples of endomorphisms embeds into classification of commuting pairs on a larger finite-dimensional space. Thus the commuting-pair problem contains the free-algebra representation problem. This is a statement about complexity of uniform classification, not algorithmic undecidability of the isomorphism test for two given pairs.}
\GSfield{Reconstruction}{The word unsolvable is interpreted in its matrix-problem sense. The notebook supplies no formal undecidability claim.}
\GSfield{Fate}{\citet[pp.~81--82]{gelfandponomarev1969commuting} give the construction; the note was received 23 June 1969 and published in 1969.}
\GSfield{Present status}{As of 9 November 2026: established with the hypotheses and distinctions stated here.}
\GSfield{Assessment}{literal statement: correct with missing hypotheses; publication match: exact; verification: read the result; development: none; confidence: high}
\end{GSitem}
\begin{GSitem}{1972-277.5}{E}{Classification of one or two subspaces}
\GSfield{Kind; source}{theorem; \GSpage{P000464}}
\GSfield{The notes}{A space with one subspace, and a space with two subspaces, are marked trivial.}
\GSfield{Statement}{For ordered subspaces \(U,W\subset V\) of a finite-dimensional space, the four multiplicities \(\dim(U\cap W)\), \(\dim U-\dim(U\cap W)\), \(\dim W-\dim(U\cap W)\), and \(\dim V-\dim(U+W)\) classify the system under ambient isomorphism. The one-subspace case needs \(\dim U\) and \(\dim V/U\).}
\GSfield{Reconstruction}{The explicit four complementary summands make the shorthand precise.}
\GSfield{Fate}{Already the opening of \citet[section 1, p.~762]{gelfandponomarev1971quadruples} treats the two-subspace case as elementary.}
\GSfield{Present status}{As of 9 November 2026: established with the hypotheses and distinctions stated here.}
\GSfield{Assessment}{literal statement: correct; publication match: exact; verification: read the result; development: none; confidence: high}
\end{GSitem}
\begin{GSitem}{1972-277.6}{E}{Nine multiplicities classify an ordered triple of subspaces}
\GSfield{Kind; source}{theorem; \GSpage{P000464}}
\GSfield{The notes}{Three subspaces can be classified; \ru{Инвариантов 9, а не 8!}}
\GSfield{Statement}{For finite-dimensional \((V;U_1,U_2,U_3)\), there are nine indecomposable subspace systems: eight one-dimensional systems indexed by which \(U_i\) equal the ambient line, and one two-dimensional system consisting of three distinct lines. Every system is a direct sum of these, uniquely up to ordering; the nine multiplicities are complete invariants. This is the injective-arrow subcategory of the inward \(D_4\) quiver, not all twelve indecomposable \(D_4\) representations.}
\GSfield{Reconstruction}{The nine indecomposable types and the distinction from all twelve positive roots are reconstruction. The extra type is \(k^2;k e_1,k e_2,k(e_1+e_2)\).}
\GSfield{Fate}{The triple case predates the talk; see \citet[section 1]{gelfandponomarev1971quadruples}. Its quiver interpretation follows the root classification in \citet[Theorem 3.1]{bernsteingelfandponomarev1973coxeter}.}
\GSfield{Present status}{As of 9 November 2026: established with the hypotheses and distinctions stated here.}
\GSfield{Assessment}{literal statement: correct; publication match: exact; verification: read the result; development: annotated; confidence: high}
\end{GSitem}
\begin{GSitem}{1972-277.7}{E}{Four-subspace classification and continuous families}
\GSfield{Kind; source}{theorem; \GSpage{P000464}}
\GSfield{The notes}{Four subspaces are said to have been classified by Gelfand and Ponomarev, with infinitely many indecomposables.}
\GSfield{Statement}{Ordered quadruples of subspaces in a finite-dimensional vector space over an algebraically closed field of characteristic zero admit the Gelfand--Ponomarev indecomposable classification. Infinite families occur already in \(k^2\): four ordered lines \(k e_1,k e_2,k(e_1+e_2),k(e_1+\lambda e_2)\) with \(\lambda\ne0,1\) are pairwise nonisomorphic for distinct \(\lambda\).}
\GSfield{Reconstruction}{The field hypothesis is the one read in the original paper. The displayed family is an illustrative check \ours: an ambient automorphism preserving the first three ordered lines is scalar, and therefore it cannot change the fourth line’s parameter.}
\GSfield{Fate}{\citet[Theorem 1 and section 2]{gelfandponomarev1971quadruples}, received 21 October 1970, published 1971, is an exact prior publication.}
\GSfield{Present status}{As of 9 November 2026: established with the hypotheses and distinctions stated here.}
\GSfield{Assessment}{literal statement: correct; publication match: exact; verification: read the result; development: annotated; confidence: high}
\end{GSitem}
\begin{GSitem}{1972-277.8}{E}{Which quivers have only finitely many indecomposables?}
\GSfield{Kind; source}{question; \GSpage{P000464}}
\GSfield{The notes}{The question is attributed to Gabriel; the answer is Dynkin diagrams, except those with multiple arrows.}
\GSfield{Statement}{For a finite connected quiver over an algebraically closed field, the set of isomorphism classes of finite-dimensional indecomposable representations is finite exactly when the underlying unoriented graph is one of \(A_n,D_n,E_6,E_7,E_8\). Every orientation works. Non-simply-laced Dynkin diagrams are not exceptions to this theorem: they are outside this ordinary-quiver list.}
\GSfield{Reconstruction}{Finite-dimensionality, the coefficient field and the distinction between an ordinary quiver and a valued diagram are supplied.}
\GSfield{Fate}{Gabriel’s theorem was published in 1972; its proof and exact statement are reproduced in \citet[Theorem 3.1]{bernsteingelfandponomarev1973coxeter}.}
\GSfield{Present status}{As of 9 November 2026: established with the hypotheses and distinctions stated here.}
\GSfield{Assessment}{literal statement: correct with missing hypotheses; publication match: exact; verification: read the result; development: none; confidence: high; problem \hyperref[prob:GP-1972-301]{GP-1972-301}}
\end{GSitem}
\begin{GSitem}{1972-277.9}{B}{The affine classification attributed to Roiter or Nazarova}
\GSfield{Kind; source}{claim; \GSpage{P000465}}
\GSfield{The notes}{The theorem is assigned to an uncertain \ru{Ройтов} from Kyiv or \ru{Назаровы}: a class of diagrams admits classification like that of four subspaces.}
\GSfield{Statement}{The identifiable content is the finite/tame boundary for ordinary quivers: connected quivers of infinite tame type have underlying extended Dynkin graphs (including the loop and Kronecker limiting cases under the usual conventions); all other infinite types contain the simultaneous-matrix-pair problem.}
\GSfield{Reconstruction}{Roiter and Nazarova are bibliographically plausible normalized names; the scan does not establish a separate Roiter theorem or joint authorship.}
\GSfield{Fate}{\citet[section 1, theorem, pp.~752--753]{nazarova1973quivers} is the exact match. It was received 19 September 1972 and published in 1973, so the notebook reports an already circulating result.}
\GSfield{Present status}{As of 9 November 2026: established with the hypotheses and distinctions stated here.}
\GSfield{Assessment}{literal statement: not determinable; publication match: exact; verification: read the result; development: none; confidence: medium; problem \hyperref[prob:GP-1972-302]{GP-1972-302}}
\end{GSitem}
\begin{GSitem}{1972-277.10}{E}{Dynkin diagrams also arise in hypersurface singularities}
\GSfield{Kind; source}{remark; \GSpage{P000465}}
\GSfield{The notes}{Gelfand points to a connection with singularities and names Artin, Tyurina and Brieskorn.}
\GSfield{Statement}{The bounded identification supported here is the common simply-laced ADE diagrams: they classify simple complex hypersurface singularities (up to stabilization/right equivalence) and finite-type quivers. The notebook states no functor between these categories and no equivalence of their deformation theories.}
\GSfield{Reconstruction}{Only the common classification diagrams are identified; a categorical correspondence would go beyond the notes.}
\GSfield{Fate}{The singularity classification is in \citet[sections 1--2 and 9]{arnold1972normal}, received 21 June 1972. Records 284 and 294 expand this connection.}
\GSfield{Present status}{As of 9 November 2026: established with the hypotheses and distinctions stated here.}
\GSfield{Assessment}{literal statement: correct; publication match: partial; verification: read the result; development: none; confidence: high}
\end{GSitem}

\GStalk{278}{1972-10-30}{I. N. Bernstein}{P000465-P000466}{Categories of quiver representations and the finite/tame boundary}
The separate \ru{И.Б.} passage begins on P000465, after Gelfand’s introduction. It defines the representation category and direct sums, states the Krull--Schmidt property, gives Gabriel’s ADE classification and the positive-root dimension vectors, and then describes the extended-Dynkin boundary. The mathematical passage ends before Buchstaber’s heading on P000466. The subsequent 1973 BGP paper is an exact expository match for finite type; reflection functors are not actually developed in this surviving fragment, so they are not inserted as a separate notebook item.

\begin{GSitem}{1972-278.1}{E}{The category of representations and direct sums}
\GSfield{Kind; source}{definition; \GSpage{P000465}}
\GSfield{The notes}{\(\operatorname{Cat}(\Gamma)\) has vector spaces over vertices, maps over arrows and commuting vertexwise morphisms. Direct sum is defined, and the note stresses that isomorphic objects need not be identified.}
\GSfield{Statement}{For a fixed finite quiver \(Q\), \(\operatorname{Rep}_k(Q)\) is the category described in 1972-277.1, retaining actual objects and all intertwining morphisms. Direct sums are taken at each vertex and on each arrow. Passing to a set of isomorphism classes loses the category structure needed for functors.}
\GSfield{Fate}{See \citet[Definition 1 and p.~18]{bernsteingelfandponomarev1973coxeter}. This is the already established categorical formulation, not a new theorem of October 1972.}
\GSfield{Present status}{As of 9 November 2026: established with the hypotheses and distinctions stated here.}
\GSfield{Assessment}{literal statement: correct; publication match: exact; verification: read the result; development: none; confidence: high}
\end{GSitem}
\begin{GSitem}{1972-278.2}{E}{Krull--Schmidt decomposition}
\GSfield{Kind; source}{theorem; \GSpage{P000465}}
\GSfield{The notes}{Every object is a sum of indecomposables, uniquely up to isomorphism.}
\GSfield{Statement}{Every finite-dimensional quiver representation is a finite direct sum of nonzero indecomposables. The multiset of their isomorphism classes and multiplicities is unique. The zero object corresponds to the empty sum. Finite total vector-space dimension gives existence and the finite-dimensional endomorphism rings give the Krull--Schmidt property.}
\GSfield{Reconstruction}{Finite-dimensionality and uniqueness up to permutation are made explicit.}
\GSfield{Fate}{Explicitly recalled in \citet[p.~18 and footnote 2]{bernsteingelfandponomarev1973coxeter}.}
\GSfield{Present status}{As of 9 November 2026: established with the hypotheses and distinctions stated here.}
\GSfield{Assessment}{literal statement: correct with missing hypotheses; publication match: exact; verification: read the result; development: none; confidence: high}
\end{GSitem}
\begin{GSitem}{1972-278.3}{E}{Gabriel classification for every orientation}
\GSfield{Kind; source}{theorem; \GSpage{P000465}}
\GSfield{The notes}{The connected graphs \(A_n,D_n,E_6,E_7,E_8\) are drawn; arrows may be placed arbitrarily.}
\GSfield{Statement}{The finite-type statement is exactly 1972-277.8. It concerns finitely many indecomposables across all dimension vectors, not finitely many orbits merely in a fixed representation space.}
\GSfield{Fate}{\citet[Theorem 3.1]{bernsteingelfandponomarev1973coxeter} reprove Gabriel’s 1972 result; the BGP manuscript was received 18 December 1972.}
\GSfield{Present status}{As of 9 November 2026: established with the hypotheses and distinctions stated here.}
\GSfield{Assessment}{literal statement: correct with missing hypotheses; publication match: exact; verification: read the result; development: none; confidence: high; problem \hyperref[prob:GP-1972-301]{GP-1972-301}}
\end{GSitem}
\begin{GSitem}{1972-278.4}{E}{Obstructions to finite type from cycles and branching}
\GSfield{Kind; source}{method; \GSpage{P000465}}
\GSfield{The notes}{The notes exclude cycles because eigenvalues appear, a four-armed star, and a chain with forks at both ends.}
\GSfield{Statement}{A cycle supports infinitely many nonisomorphic representations by varying the monodromy around its underlying unoriented cycle, taking inverse arrow maps where needed and choosing all these maps invertible. The four-subspace star is \(\widetilde D_4\), and the two-fork diagrams are extended \(D\)-type; each has infinite representation type. Extending these representations by zero or along identity edges gives obstructions inside larger graphs. This is only part of the necessity proof.}
\GSfield{Reconstruction}{The notebook sketches excluded shapes; the extension operations and the warning that the list is only partial are editorial.}
\GSfield{Fate}{The full necessity argument is Tits’s quadratic-form proof in \citet[proof of Theorem 3.1, pp.~28--29]{bernsteingelfandponomarev1973coxeter}.}
\GSfield{Present status}{As of 9 November 2026: established with the hypotheses and distinctions stated here.}
\GSfield{Assessment}{literal statement: correct with missing hypotheses; publication match: partial; verification: read the result; development: none; confidence: high}
\end{GSitem}
\begin{GSitem}{1972-278.5}{E}{Positive roots are the indecomposable dimension vectors}
\GSfield{Kind; source}{theorem; \GSpage{P000465}}
\GSfield{The notes}{The number of indecomposables is independent of arrow directions. With \(\gamma=\sum n_i\alpha_i\), the coefficients \(n_i\) are dimensions and each root gives an indecomposable.}
\GSfield{Statement}{For a Dynkin quiver, dimension vector induces a bijection between indecomposable isomorphism classes and positive roots of the associated simply-laced root system. Thus an indecomposable is uniquely determined by its dimension vector, and the number is \(|\Phi^+|\), independent of orientation.}
\GSfield{Reconstruction}{The restriction to finite Dynkin type is essential: dimension vectors do not classify arbitrary quiver representations.}
\GSfield{Fate}{\citet[Theorem 3.1(2)]{bernsteingelfandponomarev1973coxeter}, citing Gabriel 1972, states this precise bijection.}
\GSfield{Present status}{As of 9 November 2026: established with the hypotheses and distinctions stated here.}
\GSfield{Assessment}{literal statement: correct with missing hypotheses; publication match: exact; verification: read the result; development: none; confidence: high}
\end{GSitem}
\begin{GSitem}{1972-278.6}{E}{The notebook counts simple roots instead of positive roots}
\GSfield{Kind; source}{claim; \GSpage{P000465}}
\GSfield{The notes}{After defining positive roots, the next line says \ru{кол-во = кол-ву простых корней}.}
\GSfield{Statement}{The correct count is the number of positive roots: \(n(n+1)/2\) for \(A_n\), \(n(n-1)\) for \(D_n\), and \(36,63,120\) for \(E_6,E_7,E_8\). The number of simple roots is only the number of vertices.}
\GSfield{Reconstruction}{This is an emendation of a faithfully transcribed notebook error, not a translation error.}
\GSfield{Fate}{The correction follows from \citet[Theorem 3.1(2)]{bernsteingelfandponomarev1973coxeter}; the six rows immediately below the handwritten sentence already contradict its literal reading for \(A_3\).}
\GSfield{Present status}{As of 9 November 2026: established with the hypotheses and distinctions stated here.}
\GSfield{Assessment}{literal statement: false as written; publication match: exact; verification: read the result; development: annotated; confidence: high}
\end{GSitem}
\begin{GSitem}{1972-278.7}{E}{The six interval representations of A3}
\GSfield{Kind; source}{construction; \GSpage{P000465}}
\GSfield{The notes}{The drawn \(A_3\) quiver has dimension rows \[(1,0,0),\ (0,1,0),\ (0,0,1),\ (1,1,0),\ (0,1,1),\ (1,1,1).\]}
\GSfield{Statement}{For \(1\to2\to3\), the six indecomposables are the interval representations \([i,j]\), \(1\le i\le j\le3\): one-dimensional spaces at vertices in the interval, zero elsewhere, and identity internal arrows.}
\GSfield{Fate}{This is the \(A_3\) instance of \citet[Theorem 3.1]{bernsteingelfandponomarev1973coxeter}.}
\GSfield{Present status}{As of 9 November 2026: established with the hypotheses and distinctions stated here.}
\GSfield{Assessment}{literal statement: correct; publication match: exact; verification: read the result; development: none; confidence: high}
\end{GSitem}
\begin{GSitem}{1972-278.8}{B}{What the notebook calls an unsolvable category}
\GSfield{Kind; source}{definition; \GSpage{P000466}}
\GSfield{The notes}{The category of two arbitrary operators is denoted \(C\). An embedding \(C\to\operatorname{Cat}(\Gamma)\) taking nonisomorphic objects to nonisomorphic ones is used to define unsolvable.}
\GSfield{Statement}{The intended modern notion is wildness: a representation embedding from finite-dimensional \(k\langle x,y\rangle\)-modules into the category, given by an exact k-linear functor with the standard bimodule or algebraic control, preserving indecomposability and reflecting isomorphism. A bare injection on isomorphism classes does not characterize wildness and is too weak to distinguish tame families.}
\GSfield{Reconstruction}{The exact representation-embedding requirements are modern repairs. No decision-theoretic impossibility is inferred from the notebook word.}
\GSfield{Fate}{The matrix-pair containment criterion is the one used in \citet[Introduction and section 1]{nazarova1973quivers}.}
\GSfield{Present status}{As of 9 November 2026: established with the hypotheses and distinctions stated here.}
\GSfield{Assessment}{literal statement: correct with missing hypotheses; publication match: partial; verification: read the result; development: annotated; confidence: high}
\end{GSitem}
\begin{GSitem}{1972-278.9}{B}{Extended Dynkin quivers have tame infinite type}
\GSfield{Kind; source}{theorem; \GSpage{P000466}}
\GSfield{The notes}{The solvable graphs are said to be extended Dynkin diagrams, with cycle and affine \(D,E\) sketches; their relation with reflection groups is mentioned.}
\GSfield{Statement}{For finite connected ordinary quivers over an algebraically closed field, tame infinite type corresponds to the extended simply-laced Dynkin graphs, with the standard loop/cycle and Kronecker conventions. Finite Dynkin graphs remain finite type, hence are also classifiable. For fixed dimension, tame type means finitely many one-parameter families apart from finitely many exceptions.}
\GSfield{Reconstruction}{The finite-type graphs must be added if solvable includes finite classification. The notebook does not define one-parameter families or deal separately with limiting affine conventions.}
\GSfield{Fate}{Exact matching classification: \citet[section 1, theorem]{nazarova1973quivers}, submitted 19 September 1972 and published 1973 (English translation 1974).}
\GSfield{Present status}{As of 9 November 2026: established with the hypotheses and distinctions stated here.}
\GSfield{Assessment}{literal statement: correct with missing hypotheses; publication match: exact; verification: read the result; development: none; confidence: high; problem \hyperref[prob:GP-1972-302]{GP-1972-302}}
\end{GSitem}

\GStalk{279}{1972-10-30}{V. M. Buchstaber}{P000466-P000467}{Lazard cocycles and the normalized cobar complex}
After the quiver discussion, Buchstaber resumes the proof of Lazard's theorem. Symmetric normalized cocycles, a resolution, and its induced cochain differential are written explicitly. The main editorial issue is the notebook's movement between coalgebra comodules, polynomial-algebra modules, and their duals: the intended calculation is the cobar calculation, while ordinary module Ext over a polynomial algebra would give a different answer.

\begin{GSitem}{1972-279.1}{E}{Lazard theorem restated through indecomposables}
\GSfield{Kind; source}{theorem; \GSpage{P000466}}
\GSfield{The notes}{\(\Lambda=\mathbb Z[a_1,a_2,\ldots]\), \(\deg a_n=n\ge1\); the proposed proof passes to \(\bar\Lambda=\Lambda/(\text{decomposables})\).}
\GSfield{Statement}{The universal commutative formal-law ring is polynomial; its positive indecomposables form one free rank-one group in every positive degree. Linearized relations describe those indecomposables, but a generation and torsion argument is still needed to recover the polynomial ring.}
\GSfield{Reconstruction}{The phrase ``universal ring for polynomials'' is read degree by degree as a representing coefficient module with square-zero multiplication in positive degree.}
\GSfield{Fate}{Restatement of \citet[Theorem II]{lazard1955formal}.}
\GSfield{Present status}{As of 9 November 2026: established; no separate publication or new result is implied by this repeat.}
\GSfield{Assessment}{literal statement: correct; publication match: exact; verification: read the result; development: none; confidence: high}
\end{GSitem}
\begin{GSitem}{1972-279.2}{E}{An equivalent symmetric cocycle identity}
\GSfield{Kind; source}{claim; \GSpage{P000466}}
\GSfield{The notes}{\(h(x,y)=h(y,x)\), \(h(x,0)=0\), and \(h(y,z)+h(x,y+z)=h(x,z)+h(y,x+z)\).}
\GSfield{Statement}{For a symmetric normalized polynomial \(h\), the identity displayed in the notes is equivalent to \(h(x,y)+h(x+y,z)=h(y,z)+h(x,y+z)\). To recover the latter, interchange \(x\) and \(z\) in the former and use symmetry.}
\GSfield{Reconstruction}{The equation is not a mistranscription of the usual cocycle identity; it is equivalent in the explicitly stated symmetric context.}
\GSfield{Fate}{Same cocycle class as \citet[\S~II, Lemma~3]{lazard1955formal}.}
\GSfield{Present status}{As of 9 November 2026: correct with the symmetry hypothesis retained.}
\GSfield{Assessment}{literal statement: correct; publication match: exact; verification: read the result; development: none; confidence: high; problem \hyperref[prob:GP-1972-201]{GP-1972-201}}
\end{GSitem}
\begin{GSitem}{1972-279.3}{E}{Resolution of the primitive coalgebra}
\GSfield{Kind; source}{construction; \GSpage{P000466}}
\GSfield{The notes}{The sequence \(A[x]\to A[x]\otimes\bar A[y]\to A[x]\otimes\bar A[y]\otimes\bar A[z]\to\cdots\) is called a bar resolution; \(\bar A[y]\) has zero constant term.}
\GSfield{Statement}{For the coaugmented coalgebra \(C=A[x]\), use the normalized cofree comodule resolution, whose reduced factors are \(\bar C=\ker\varepsilon\). The associated cobar cochain groups are \(A\), \(\bar C\), \(\bar C^{\otimes2}\), and so on; the differential inserts the reduced coproduct with alternating signs.}
\GSfield{Reconstruction}{``Cobar'' and ``comodule'' specify the variance omitted by the notebook. The polynomial coalgebra is free as an \(A\)-module, so the relative normalized construction has the needed underlying-module exactness.}
\GSfield{Fate}{Classical bar/cobar apparatus; \citet[Ch.X]{cartaneilenberg1956homological}.}
\GSfield{Present status}{As of 9 November 2026: standard once the category is fixed.}
\GSfield{Assessment}{literal statement: correct with missing hypotheses; publication match: partial; verification: publisher metadata checked; development: none; confidence: high}
\end{GSitem}
\begin{GSitem}{1972-279.4}{E}{Ordinary module Hom cannot substitute for comodule Hom}
\GSfield{Kind; source}{claim; \GSpage{P000466-P000467}}
\GSfield{The notes}{The notes label the functor \(\operatorname{Hom}_{A[x]}(\mathcal A,A)\), call the resolution terms \(A[x]\)-modules, and obtain polynomial cochains.}
\GSfield{Statement}{The intended polynomial cochains arise from \(\operatorname{Hom}_{C\text{-comod}}(A,-)\), or from continuous Hom over the dual divided-power algebra. They do not arise from the displayed ordinary polynomial-module Ext computation. For a field \(k\), the two-term free resolution \(0\to k[x]\xrightarrow{x}k[x]\to k\to0\) proves \(\operatorname{Ext}^2_{k[x]}(k,k)=0\).}
\GSfield{Reconstruction}{\ours\ The displayed two-term free resolution independently proves the vanishing of ordinary module Ext in degree two. This is a mathematical notation repair confirmed by the explicit dual algebra introduced on P476. It does not assert that the speaker intended the incorrect ordinary-module calculation.}
\GSfield{Fate}{The resolution distinction is standard; \citet[Chs.~VIII, X]{cartaneilenberg1956homological}.}
\GSfield{Present status}{As of 9 November 2026: the category omission must remain visible in any faithful edition.}
\GSfield{Assessment}{literal statement: false as written; publication match: partial; verification: publisher metadata checked; development: annotated; confidence: high}
\end{GSitem}
\begin{GSitem}{1972-279.5}{E}{Normalized coboundaries and the four-term differential}
\GSfield{Kind; source}{construction; \GSpage{P000466-P000467}}
\GSfield{The notes}{The boundary of \(x^n\) is \((x+y)^n-x^n-y^n\). A later shorthand writes only two terms for the differential of \(h\), and notes that reduction removes constant terms.}
\GSfield{Statement}{The normalized cobar differential is \(d f(x,y)=f(x+y)-f(x)-f(y)\), and \(d h(x,y,z)=h(x+y,z)-h(y,z)-h(x,y+z)+h(x,y)\). Thus degree-two cycles are precisely additive cocycles. The two displayed substitution terms require their reduced-coproduct corrections.}
\GSfield{Reconstruction}{The four-term formula is supplied explicitly to disambiguate the abbreviated P467 formula; P476 later records all four terms correctly.}
\GSfield{Fate}{Classical cocycle differential: \citet[\S~II, Lemma~2]{lazard1955formal}.}
\GSfield{Present status}{As of 9 November 2026: established; the normalized and unreduced versions should not be mixed.}
\GSfield{Assessment}{literal statement: correct with missing hypotheses; publication match: partial; verification: read the result; development: annotated; confidence: high; problem \hyperref[prob:GP-1972-201]{GP-1972-201}}
\end{GSitem}

\GStalk{280}{1972-11-13}{L. N. Vasserstein}{P000468-P000470}{Simplicial algebraic K theory, stabilization, and Steinberg groups}
Vasserstein builds the quasiunit and matrix functors, the permutation-triangular simplicial space, its homotopy groups and block-sum structure, the cone swindle, and the low-degree identifications with elementary and Steinberg groups. A unital-ring hypothesis is explicitly introduced on P000469 and must be retained; the talk's initial allowance of nonunital rings does not erase that restriction. The short Foundations announcement was received the next day. The final tolerance question is an exact assertion in the later 1976 paper, but is false for the literal simplicial definition: three stable elementary matrices give a minimal-rank obstruction. The correction is developed as a published-claim audit, not relabelled as an undisseminated open idea.

\begin{GSitem}{1972-280.1}{E}{Formal unitization of an associative ring}
\GSfield{Kind; source}{construction; \GSpage{P000468}}
\GSfield{The notes}{For an associative ring \(A\), not necessarily with identity, \(iA\) formally adjoins an identity, with \(A\) a two-sided ideal and \(iA/A=\mathbb Z\); \(i\) is a functor.}
\GSfield{Statement}{For any associative ring \(A\), define \(iA=\mathbb Z\oplus A\) additively, with \((m,a)(n,b)=(mn,mb+na+ab)\). Its identity is \((1,0)\), the map \(A\to iA\), \(a\mapsto(0,a)\), identifies an ideal, and \(iA/A\cong\mathbb Z\). Every ring homomorphism induces the identity on \(\mathbb Z\) and the given map on \(A\).}
\GSfield{Reconstruction}{The multiplication, implicit in formal adjoining, is written explicitly. Even if \(A\) already has an identity, this formal unitization is not being identified with \(A\) itself.}
\GSfield{Fate}{The exact framework appears in \citet{vaserstein1973foundations}, received on 14 November 1972, the day after this talk; the detailed treatment \citet{vaserstein1976foundations} was received on 26 February 1975 and published in 1976. See \citet[\S1, pp.~100--101]{vaserstein1976foundations}. The elementary unitization construction itself is antecedent background, so its occurrence is E.}
\GSfield{Present status}{As of 9 November 2026: standard functorial unitization used to organize nonunital K theory.}
\GSfield{Assessment}{literal statement: correct; publication match: exact; verification: read the result; development: none; confidence: high}
\end{GSitem}
\begin{GSitem}{1972-280.2}{E}{Quasiunits under the circle product}
\GSfield{Kind; source}{definition; \GSpage{P000468}}
\GSfield{The notes}{The law \(a\circ b=a+b+ab\) is associative, and \(GA\subset A\) is its group of invertible elements with identity zero. The map \(a\mapsto1+a\) is described as a homomorphism into ordinary units.}
\GSfield{Statement}{The operation \(\circ\) makes \(A\) a monoid with identity zero. Its invertible elements are precisely \(GA=\{a\in A:1+a\in(iA)^\times\}\), and \(a\mapsto1+a\) identifies this group with the units of \(iA\) congruent to one modulo \(A\). Ring homomorphisms preserve quasiunits and give a group-valued functor \(G\).}
\GSfield{Reconstruction}{The map from all of \(A\) lands in the multiplicative monoid \(1+A\), not necessarily its unit group; only its restriction to \(GA\) lands in units. For example, \(-1\) in a nonzero unital ring is not quasiinvertible.}
\GSfield{Fate}{The exact framework appears in \citet{vaserstein1973foundations}, received on 14 November 1972, the day after this talk; the detailed treatment \citet{vaserstein1976foundations} was received on 26 February 1975 and published in 1976. The precise quasiunit convention is \citet[\S2, p.~101]{vaserstein1976foundations}.}
\GSfield{Present status}{As of 9 November 2026: the quasiunit construction is standard; the literal unrestricted units statement needs correction.}
\GSfield{Assessment}{literal statement: false as written; publication match: partial; verification: read the result; development: annotated; confidence: high}
\end{GSitem}
\begin{GSitem}{1972-280.3}{E}{Finite matrix stabilization and the stable general linear group}
\GSfield{Kind; source}{construction; \GSpage{P000468}}
\GSfield{The notes}{The notes introduce \(M_kA\), the finite-support matrix ring \(MA\), and its quasiunit group \(GMA\).}
\GSfield{Statement}{Let \(M_kA\) be the ring of \(k\times k\) matrices and \(MA=\bigcup_k M_kA\) under upper-left embeddings with zero complementary block. For unital \(A\), the map \(a\mapsto I+a\) identifies \(GMA\) with the stable group \(\mathrm{GL}(A)=\varinjlim_k\mathrm{GL}_k(A)\), whose stabilization adjoins identity blocks.}
\GSfield{Reconstruction}{The zero-block and identity-block stabilization conventions agree only after the change from circle coordinates to ordinary multiplicative coordinates. The finite-support ring \(MA\) is not itself unital when \(A\ne0\).}
\GSfield{Fate}{The exact framework appears in \citet{vaserstein1973foundations}, received on 14 November 1972, the day after this talk; the detailed treatment \citet{vaserstein1976foundations} was received on 26 February 1975 and published in 1976. See \citet[\S3, p.~101]{vaserstein1976foundations}. These standard matrix conventions are also the background of 1971-256.1.}
\GSfield{Present status}{As of 9 November 2026: standard stable matrix group, with both coordinate conventions distinguished.}
\GSfield{Assessment}{literal statement: correct; publication match: exact; verification: read the result; development: none; confidence: high}
\end{GSitem}
\begin{GSitem}{1972-280.4}{E}{Row-and-column-finite cone and suspension rings}
\GSfield{Kind; source}{construction; \GSpage{P000468}}
\GSfield{The notes}{The ring \(\widetilde\Omega A\) consists of infinite matrices with finitely many nonzero entries in every row and column; \(MA\) is its finite-support two-sided ideal, and \(\Omega A=\widetilde\Omega A/MA\).}
\GSfield{Statement}{With indices in \(\mathbb N\), row-and-column-finite matrices over any associative ring form a ring under ordinary matrix multiplication: each entry is a finite sum and row and column finiteness are preserved. Finite-support matrices form a two-sided ideal. The quotient is the notebook's suspension ring. All three constructions are functorial in \(A\).}
\GSfield{Reconstruction}{The tilde distinguishes the cone from its quotient throughout this chapter; later sources use different Omega decorations. No boundedness condition on the number of entries per row is imposed by these notes.}
\GSfield{Fate}{The exact framework appears in \citet{vaserstein1973foundations}, received on 14 November 1972, the day after this talk; the detailed treatment \citet{vaserstein1976foundations} was received on 26 February 1975 and published in 1976. The definitions occur at the start of the 1973 announcement and in \citet[\S3]{vaserstein1976foundations}. Earlier occurrences are 1971-253.5 and 1971-256.2.}
\GSfield{Present status}{As of 9 November 2026: established cone/suspension machinery; the later vanishing theorem is a separate item.}
\GSfield{Assessment}{literal statement: correct; publication match: exact; verification: read the result; development: none; confidence: high}
\end{GSitem}
\begin{GSitem}{1972-280.5}{E}{Strict triangular matrices and their quasiinverses}
\GSfield{Kind; source}{definition; \GSpage{P000468}}
\GSfield{The notes}{The subring \(\Theta A\) is called upper triangular but is defined by \(a_{ij}=0\) for \(i\leq j\); the page says \(G\Theta A=\Theta A\) because its matrices are nilpotent.}
\GSfield{Statement}{Choose either strict upper triangular finite-support matrices, with \(a_{ij}=0\) for \(i\geq j\), or the transposed lower convention. Every such matrix \(a\) is nilpotent, and its circle inverse is the finite sum \(-a+a^2-a^3+\cdots\). Thus the triangular ring is a group under \(\circ\). The union over coordinate permutations is the same with either choice of orientation.}
\GSfield{Reconstruction}{The word upper and the written inequality disagree in the handwriting. The choice is harmless for the subsequent permutation-triangular construction, but the two conventions are not presented as literally identical definitions.}
\GSfield{Fate}{\citet[\S3 and \S11]{vaserstein1976foundations} uses triangular matrices in the same simplicial construction; the nilpotent inverse is the elementary geometric-series identity.}
\GSfield{Present status}{As of 9 November 2026: standard nilpotent quasiunit subgroup with a recorded orientation inconsistency.}
\GSfield{Assessment}{literal statement: correct with missing hypotheses; publication match: partial; verification: read the result; development: annotated; confidence: high}
\end{GSitem}
\begin{GSitem}{1972-280.6}{E}{Reindexing requires injective maps}
\GSfield{Kind; source}{claim; \GSpage{P000468}}
\GSfield{The notes}{A bijection \(\mathbb N\to\mathbb N\) simultaneously permutes rows and columns and is called an inner automorphism. The next sentence claims that even any map \(\mathbb N\to\mathbb N\) gives an endomorphism.}
\GSfield{Statement}{Every injection \(e:\mathbb N\hookrightarrow\mathbb N\) gives a ring embedding \(MA\to MA\) by placing \(a_{ij}\) in position \((e(i),e(j))\) and filling other entries with zero; bijections give automorphisms. Arbitrary noninjective reindexings do not preserve matrix multiplication. For unital \(A\), a permutation is implemented by conjugation in the larger row-and-column-finite matrix ring, but need not be inner in the finitary stable group.}
\GSfield{Reconstruction}{\ours\ A collapsed pair of indices sends two orthogonal diagonal matrix units \(E_{11},E_{22}\) over \(\mathbb Z\) to the same \(E_{kk}\): their original product is zero and the product of their images is not. Thus injectivity is indispensable. The ambient meaning of inner is supplied separately.}
\GSfield{Fate}{\citet[\S4, pp.~101--102]{vaserstein1976foundations} correctly uses embeddings of \(\mathbb N\); it matches the repaired statement, not the arbitrary-map wording.}
\GSfield{Present status}{As of 9 November 2026: the notebook's arbitrary-map assertion is false; the injection construction is elementary and established.}
\GSfield{Assessment}{literal statement: false as written; publication match: partial; verification: read the result; development: annotated; confidence: high}
\end{GSitem}
\begin{GSitem}{1972-280.7}{E}{Permutation-triangular simplices and the based space S A}
\GSfield{Kind; source}{construction; \GSpage{P000468-P000469}}
\GSfield{The notes}{A finite \(X\subset GMA\) is a simplex if a common left translation followed by a coordinate permutation carries it into \(\Theta A\). The distinguished vertex is zero, and the resulting simplicial scheme is \(SA\). Volodin is credited for its use in K theory.}
\GSfield{Statement}{In multiplicative coordinates for unital \(A\), let \(U_\prec\) be the stable unitriangular subgroup for a total coordinate order obtained by a permutation. A finite vertex set \(X\subset\mathrm{GL}(A)\) is a simplex exactly when \(X\subset gU_\prec\) for some \(g\) and one order \(\prec\). The identity is the base vertex. This defines an abstract simplicial complex and functorial maps under ring homomorphisms.}
\GSfield{Reconstruction}{A single translation and a single order must work for the whole set. Replacing this by a separate order for each pair changes the complex. The literal coordinate model is kept distinct from more general module-flag versions. The notebook's zero vertex corresponds to \(I\).}
\GSfield{Fate}{\citet[\S1]{volodin1971extraordinary} had already published a flag-based simplicial approach in 1971, from a manuscript received on 27 July 1970. Its general module-flag formulation should not be confused term by term with every coordinate convention of this talk. The exact framework appears in \citet{vaserstein1973foundations}, received on 14 November 1972, the day after this talk; the detailed treatment \citet{vaserstein1976foundations} was received on 26 February 1975 and published in 1976. The left-translation definition is explicit in the 1973 announcement; \citet[\S11, Remark~(b), pp.~110--111]{vaserstein1976foundations} uses the right-translation version and identifies the two via inversion. See also 1971-253.1 and 1971-256.1.}
\GSfield{Present status}{As of 9 November 2026: an identified early simplicial K-theory model; its false tolerance assertion is separately audited below.}
\GSfield{Assessment}{literal statement: correct; publication match: exact; verification: read the result; development: none; confidence: high}
\end{GSitem}
\begin{GSitem}{1972-280.8}{E}{Positive K groups from homotopy groups under the stated unit hypothesis}
\GSfield{Kind; source}{definition; \GSpage{P000469}}
\GSfield{The notes}{The page explicitly begins \ru{Пусть в кольце A есть единица} before defining \(K_nA=\pi_{n-1}SA\), \(n\geq1\). It notes the usual group structure for \(n\geq2\) and automatic abelianity for \(n\geq3\).}
\GSfield{Statement}{For an associative unital ring \(A\), define the positive groups in this model by \(K_n(A)=\pi_{n-1}(|SA|,I)\) for \(n\geq1\), with \(K_1\) initially the based set of components until its group law is supplied. This unital Volodin-type model agrees with algebraic K theory after the comparison theorem. The displayed formula alone does not define \(K_0\) or negative groups.}
\GSfield{Reconstruction}{The unital hypothesis is genuinely on the scan and is not an editorial repair. The notation \(\pi\) is interpreted via realization, equivalently the combinatorial homotopy model used in the later paper. Nonunital groups require the relative/unitization construction and must not be substituted into this formula without further hypotheses.}
\GSfield{Fate}{\citet[\S1]{volodin1971extraordinary} had already published a flag-based simplicial approach in 1971, from a manuscript received on 27 July 1970. Its general module-flag formulation should not be confused term by term with every coordinate convention of this talk. The exact framework appears in \citet{vaserstein1973foundations}, received on 14 November 1972, the day after this talk; the detailed treatment \citet{vaserstein1976foundations} was received on 26 February 1975 and published in 1976. \citet[Corollary~15.1 and Chapter~V]{vaserstein1976foundations} gives the unital identification and comparison. The general comparison is already tracked by GP-1971-21, but this page poses no new comparison question.}
\GSfield{Present status}{As of 9 November 2026: an established model for unital positive algebraic K theory; no claim for arbitrary nonunital rings is inferred.}
\GSfield{Assessment}{literal statement: correct; publication match: exact; verification: read the result; development: none; confidence: high}
\end{GSitem}
\begin{GSitem}{1972-280.9}{A}{Whitehead stabilization homotopy and five block matrices}
\GSfield{Kind; source}{method; \GSpage{P000469}}
\GSfield{The notes}{Reindexing a sphere map is said not to change its homotopy class. The proof moves a finite matrix block using five circle-coordinate matrices: \(\left(\begin{smallmatrix}a&0\\0&0\end{smallmatrix}\right)\), \(\left(\begin{smallmatrix}a&a\\0&0\end{smallmatrix}\right)\), \(\left(\begin{smallmatrix}0&a\\-I&0\end{smallmatrix}\right)\), \(\left(\begin{smallmatrix}0&0\\-I&a\end{smallmatrix}\right)\), \(\left(\begin{smallmatrix}0&0\\0&a\end{smallmatrix}\right)\).}
\GSfield{Statement}{For a unital ring and a based simplicial representative of a homotopy class in \(SA\), every injective coordinate reindexing induces the same homotopy class. After reducing to a finite block and a block swap, consecutive maps in the displayed five-step construction are contiguous, giving a homotopy from a matrix block to its shifted copy. All entries \(-I\) belong to the ring because the unit hypothesis is in force.}
\GSfield{Reconstruction}{Finite-block reduction is made for a finite simplicial representative; it is not an assertion that an arbitrary continuous map has a literal image consisting of finitely many vertices. This is a specific simplicial implementation of the older Whitehead stabilization idea.}
\GSfield{Fate}{The exact framework appears in \citet{vaserstein1973foundations}, received on 14 November 1972, the day after this talk; the detailed treatment \citet{vaserstein1976foundations} was received on 26 February 1975 and published in 1976. The exact five-block method and its contiguity checks occur in \citet[Lemma~13.1 and Theorem~13.1, pp.~112--114]{vaserstein1976foundations}; class A concerns this implementation, not priority for Whitehead's original lemma.}
\GSfield{Present status}{As of 9 November 2026: the detailed stabilization construction is published and valid with the stated unit convention.}
\GSfield{Assessment}{literal statement: correct; publication match: exact; verification: read the result; development: none; confidence: high}
\end{GSitem}
\begin{GSitem}{1972-280.10}{E}{Elementary simplicial homotopy is contiguity}
\GSfield{Kind; source}{definition; \GSpage{P000469}}
\GSfield{The notes}{Maps \(s_1,s_2\) are elementarily homotopic when \(s_1X\cup s_2X\) is a simplex for every source simplex \(X\). The line naming the domains and codomains reuses \(S'\) inconsistently.}
\GSfield{Statement}{Two simplicial maps \(s_1,s_2:S\to T\) are contiguous if, for each simplex \(X\) of \(S\), the union \(s_1(X)\cup s_2(X)\) is a simplex of \(T\). Their realizations are homotopic; a finite chain of such contiguities gives the combinatorial homotopies used here.}
\GSfield{Reconstruction}{The source and target letters are normalized; the substantive simplex condition is unchanged. Pairwise adjacency of images alone is insufficient for contiguity unless the target is flag, and the target in this talk is not flag in general.}
\GSfield{Fate}{\citet[\S\S7--10, especially \S9, Remark~(b)]{vaserstein1976foundations} formalizes this standard simplicial notion before applying it to the five-block argument.}
\GSfield{Present status}{As of 9 November 2026: standard contiguity, with the required whole-simplex condition preserved.}
\GSfield{Assessment}{literal statement: correct with missing hypotheses; publication match: exact; verification: read the result; development: annotated; confidence: high}
\end{GSitem}
\begin{GSitem}{1972-280.11}{E}{Interleaved block sum and its homotopy invariance}
\GSfield{Kind; source}{construction; \GSpage{P000469-P000470}}
\GSfield{The notes}{The injections \(f(i)=2i\), \(g(i)=2i+1\) give \(a\oplus b=a^f+b^g\). The raw operation is described as noncommutative, nonassociative, and without a neutral element; it induces \(SA\times SA\to SA\).}
\GSfield{Statement}{Disjoint coordinate embeddings define a block-sum map on finite matrices and on the simplicial spaces. Different choices of embeddings, parenthesization, and order are related on homotopy groups by stabilization homotopies. Consequently the operation on classes has a neutral element and is independent of those choices, even though the chosen coordinate formula is not strictly associative or strictly unital.}
\GSfield{Reconstruction}{The product is the simplicial product appropriate to the model, and the claim about operations on all \(K_i\), \(i\geq0\), exceeds the displayed definition when \(i=0\); here the conclusion is restricted to \(i\geq1\). A separately defined \(K_0\) also has direct sum, but is not constructed on these pages.}
\GSfield{Fate}{The exact framework appears in \citet{vaserstein1973foundations}, received on 14 November 1972, the day after this talk; the detailed treatment \citet{vaserstein1976foundations} was received on 26 February 1975 and published in 1976. See \citet[\S4 and Corollary~13.1]{vaserstein1976foundations}. The basic block-sum idea is older; the source supplies the precise coordinate and homotopy bookkeeping.}
\GSfield{Present status}{As of 9 November 2026: established block-sum operation in positive K theory; the isolated zero index is annotated.}
\GSfield{Assessment}{literal statement: correct with missing hypotheses; publication match: exact; verification: read the result; development: annotated; confidence: high}
\end{GSitem}
\begin{GSitem}{1972-280.12}{E}{Interchange gives abelian group structures}
\GSfield{Kind; source}{theorem; \GSpage{P000470}}
\GSfield{The notes}{Block sum and the previously introduced homotopy operations are mutually distributive, hence coincide and are commutative. The page describes the resulting sum as an H-space structure.}
\GSfield{Statement}{Under the unital hypothesis, block sum induces abelian group laws on all \(K_n(A)\), \(n\geq1\), agreeing with ordinary homotopy-group composition where that already exists. The key algebraic principle is the interchange law for two binary operations with the same two-sided identity: the operations then agree and are associative and commutative. A homotopy-unit H-space formulation expresses the same block-sum construction on realization.}
\GSfield{Reconstruction}{For \(K_1=\pi_0\), inverses and a group operation must be supplied rather than inferred solely from the topology of components; the ensuing elementary quotient provides them. The word distributive is interpreted as the interchange identity, not the ordinary ring distributive law.}
\GSfield{Fate}{\citet[Lemma~10.2, Corollary~13.1, and Theorem~13.2]{vaserstein1976foundations} proves exactly these assertions. The interchange principle and the already-known K-group structures are expository; no new general H-space theorem is attributed to the talk.}
\GSfield{Present status}{As of 9 November 2026: standard abelian positive K groups, with the low-degree group-law issue accounted for.}
\GSfield{Assessment}{literal statement: correct; publication match: exact; verification: read the result; development: none; confidence: high}
\end{GSitem}
\begin{GSitem}{1972-280.13}{A}{Infinite block sum and vanishing of cone K theory}
\GSfield{Kind; source}{method; \GSpage{P000470}}
\GSfield{The notes}{On \(\widetilde\Omega A\) even an infinite number of matrices can be summed in blocks; the notes conclude that its K functors are trivial.}
\GSfield{Statement}{For an associative unital ring \(A\), the row-and-column-finite cone ring admits an infinite-sum endomorphism compatible with two complementary coordinate embeddings. On positive K groups its identity satisfies \(x+\rho_*(x)=\rho_*(x)\); hence \(K_n(\widetilde\Omega A)=0\) for \(n\geq1\). Infinite sums mean a specified block-diagonal placement with finite support in each row and column, not arbitrary entrywise infinite sums.}
\GSfield{Reconstruction}{The unital hypothesis introduced on P000469 is retained. Extending the vanishing to \(K_0\) and negative degrees requires the separately developed theory; none is defined by the homotopy formula in this talk. This is the algebraic Eilenberg swindle in the specified matrix model.}
\GSfield{Fate}{The exact framework appears in \citet{vaserstein1973foundations}, received on 14 November 1972, the day after this talk; the detailed treatment \citet{vaserstein1976foundations} was received on 26 February 1975 and published in 1976. The first theorem of the 1973 announcement states cone vanishing; \citet[Lemmas~18.1--18.2 and Corollary~18.1]{vaserstein1976foundations} gives the infinite-sum proof. Earlier occurrences are 1971-253.5 and 1971-256.3; A records the subsequent implementation, not a new origin for the swindle.}
\GSfield{Present status}{As of 9 November 2026: published cone-vanishing theorem with an exact infinite-sum construction; no unresolved cone problem here.}
\GSfield{Assessment}{literal statement: correct with missing hypotheses; publication match: exact; verification: read the result; development: none; confidence: high}
\end{GSitem}
\begin{GSitem}{1972-280.14}{E}{K one as components modulo elementary matrices}
\GSfield{Kind; source}{theorem; \GSpage{P000470}}
\GSfield{The notes}{\(K_1A=\pi_0SA\); \(EA\) is the identity component, described as products of elementary or pseudotriangular matrices.}
\GSfield{Statement}{For unital \(A\), the component of the identity vertex in \(SA\) consists of the stable elementary subgroup \(E(A)\subset\mathrm{GL}(A)\). Thus \(\pi_0SA\cong\mathrm{GL}(A)/E(A)=K_1(A)\). A path consists of finitely many simplex edges, and triangular factors decompose into elementary transvections.}
\GSfield{Reconstruction}{Component is understood in the realization or in edge-path connectivity, rather than in a topology on the discrete set of matrices. The ordinary multiplicative elementary matrix is \(I+aE_{ij}\), corresponding to \(aE_{ij}\) in circle coordinates.}
\GSfield{Fate}{The classical quotient predates the talk; \citet[III, Whitehead's Lemma~1.3.3]{weibel2013kbook} records it and its 1950 antecedent. The exact component computation is \citet[Theorems~14.1--14.2]{vaserstein1976foundations}.}
\GSfield{Present status}{As of 9 November 2026: standard K-one identification; no new quotient conjecture.}
\GSfield{Assessment}{literal statement: correct; publication match: exact; verification: read the result; development: none; confidence: high}
\end{GSitem}
\begin{GSitem}{1972-280.15}{A}{Elementary perfectness and the idempotent-ring extension}
\GSfield{Kind; source}{theorem; \GSpage{P000470}}
\GSfield{The notes}{The notes explicitly allow either \(A\ni1\) or \(AA=A\), then assert \(EA=[GMA,GMA]=[EA,EA]\).}
\GSfield{Statement}{If \(A\) is an associative ring satisfying \(A^2=A\), where \(A^2\) is the additive subgroup generated by products, then the stable elementary group satisfies \(E(A)=[GMA,GMA]=[E(A),E(A)]\). In particular it is perfect, and the stable quasiunit abelianization is the elementary quotient. A unital ring satisfies the idempotence condition automatically.}
\GSfield{Reconstruction}{The idempotent condition means finite sums of products, not that every element is a single product. It applies here to the commutator assertion; no blanket extension of every subsequent universal-central-extension claim to all nonunital idempotent rings is inferred.}
\GSfield{Fate}{The exact framework appears in \citet{vaserstein1973foundations}, received on 14 November 1972, the day after this talk; the detailed treatment \citet{vaserstein1976foundations} was received on 26 February 1975 and published in 1976. \citet[Theorem~5.1, pp.~103--104]{vaserstein1976foundations}, applied with \(B=A\), proves the nonunital idempotent extension. The unital Whitehead theorem is older; class A concerns the stated extension and its publication in the speaker's treatment.}
\GSfield{Present status}{As of 9 November 2026: published under the precise idempotent-ring hypothesis.}
\GSfield{Assessment}{literal statement: correct; publication match: exact; verification: read the result; development: none; confidence: high}
\end{GSitem}
\begin{GSitem}{1972-280.16}{E}{Triangular generators and their elementary relations}
\GSfield{Kind; source}{claim; \GSpage{P000470}}
\GSfield{The notes}{\(\Theta A\) is said to have generators \(a^{ij}\), with \(a\) in position \((i,j)\), the addition relation, the adjacent-root commutator relation, and commuting relations for some other pairs.}
\GSfield{Statement}{For a fixed coordinate order, the unitriangular group is generated by \(t_{ij}(a)=I+aE_{ij}\), \(i\prec j\), with \(t_{ij}(a)t_{ij}(b)=t_{ij}(a+b)\), \([t_{ij}(a),t_{jk}(b)]=t_{ik}(ab)\) for \(i\prec j\prec k\), and \([t_{ij}(a),t_{kl}(b)]=1\) when \(j\ne k\) and \(i\ne l\). In the stable finite-support setting these are the triangular presentation relations.}
\GSfield{Reconstruction}{The allowed indices and the omitted commuting condition are reconstructed explicitly. The presentation of a fixed triangular group is distinguished from the all-root Steinberg presentation and from a presentation of \(E(A)\), which generally has additional K-two relations.}
\GSfield{Fate}{\citet[\S6 and the triangular map in the proof of Theorem~15.2]{vaserstein1976foundations} uses these relations to pass from triangular groups to the Steinberg group. The general relations are also stated in \citet[III, \S5.1]{weibel2013kbook}.}
\GSfield{Present status}{As of 9 November 2026: standard triangular elementary relations; the notebook's phrase defining relations is correct only with its group and index set specified.}
\GSfield{Assessment}{literal statement: correct with missing hypotheses; publication match: partial; verification: read the result; development: none; confidence: high}
\end{GSitem}
\begin{GSitem}{1972-280.17}{E}{Stable Steinberg group with the missing index restrictions restored}
\GSfield{Kind; source}{definition; \GSpage{P000470}}
\GSfield{The notes}{The group \(\operatorname{St}A\) is generated by all the displayed elementary symbols and Steinberg relations; the commuting cases are abbreviated as some commutators being zero.}
\GSfield{Statement}{For an associative ring \(A\), define \(\operatorname{St}(A)\) on generators \(x_{ij}(a)\), \(i\ne j\), with \(x_{ij}(a)x_{ij}(b)=x_{ij}(a+b)\), \([x_{ij}(a),x_{jk}(b)]=x_{ik}(ab)\) for distinct \(i,j,k\), and \([x_{ij}(a),x_{kl}(b)]=1\) if \(i\ne l\) and \(j\ne k\). It maps onto \(E(A)\) by \(x_{ij}(a)\mapsto I+aE_{ij}\).}
\GSfield{Reconstruction}{Zero in the notes is the identity of the circle-law convention; the abstract group relations use one. The full index restrictions are essential. Opposite-root pairs are not declared commuting.}
\GSfield{Fate}{This is the classical Steinberg construction, treated before the talk in \citet{milnor1972introduction} (publisher date 21 January 1972). The precise presentation was read in \citet[III, Definition~5.1]{weibel2013kbook} and \citet[\S6]{vaserstein1976foundations}; the Milnor publisher page alone is used for chronology, not as a claim to have read its theorem text.}
\GSfield{Present status}{As of 9 November 2026: established stable Steinberg presentation.}
\GSfield{Assessment}{literal statement: correct with missing hypotheses; publication match: exact; verification: read the result; development: none; confidence: high}
\end{GSitem}
\begin{GSitem}{1972-280.18}{E}{Steinberg group as paths and the K two kernel}
\GSfield{Kind; source}{theorem; \GSpage{P000470}}
\GSfield{The notes}{Loops represent relations among elementary matrices; the exact sequence \(0\to K_2A\to\operatorname{St}A\to EA\to0\) is displayed. Steinberg-group elements are described as homotopy classes of paths starting at zero with endpoints fixed under homotopy.}
\GSfield{Statement}{For a unital ring, the path-group model of \(\operatorname{St}(A)\) maps a based path class to its endpoint in \(E(A)\), with multiplication using translated concatenation. Its kernel is the loop group \(\pi_1(SA,I)=K_2(A)\); thus \(1\to K_2(A)\to\operatorname{St}(A)\to E(A)\to1\) is exact.}
\GSfield{Reconstruction}{The group law on paths is not bare concatenation of paths with unrelated endpoints; translation makes the endpoints composable. The zeroes in the group exact sequence are conventional shorthand and are rewritten as ones.}
\GSfield{Fate}{\citet[Theorem~15.2 and Corollary~15.2]{vaserstein1976foundations} proves the exact path/Steinberg identification and attributes the theorem to Volodin. The underlying simplicial theory had circulated by the July 1970 submission of \citet{volodin1971extraordinary}, published in 1971, so the identification is expository here.}
\GSfield{Present status}{As of 9 November 2026: the standard positive-degree low-dimensional identification; no new universal-relation problem is posed on these pages.}
\GSfield{Assessment}{literal statement: correct; publication match: exact; verification: read the result; development: none; confidence: high}
\end{GSitem}
\begin{GSitem}{1972-280.19}{E}{Universal central extension and the second homology of E A}
\GSfield{Kind; source}{theorem; \GSpage{P000470}}
\GSfield{The notes}{The notes write \(K_2A=H_2EA\) and describe \(\operatorname{St}A\) as the universal central extension of \(EA\).}
\GSfield{Statement}{For an associative unital ring \(A\), the stable Steinberg group \(\operatorname{St}(A)\to E(A)\) is the universal central extension of the perfect elementary group. Therefore \(K_2(A)\cong H_2(E(A);\mathbb Z)\). Stability and the unit convention are part of the assertion; arbitrary small finite matrix ranks require separate results.}
\GSfield{Reconstruction}{The integral trivial coefficients in group homology are supplied. The active unital hypothesis from P000469 is retained, and the immediately preceding idempotent-ring alternative is not automatically propagated to this stronger claim.}
\GSfield{Fate}{The classical Kervaire--Steinberg theorem is stated and proved in \citet[III, Theorem~5.5]{weibel2013kbook}; its preceding discussion points to \citet[\S5]{milnor1972introduction}, published before this talk. The latter's publication was checked from the publisher, while the mathematical theorem was read in Weibel.}
\GSfield{Present status}{As of 9 November 2026: established stable unital theorem, with scope distinct from finite-rank or unrestricted nonunital variants.}
\GSfield{Assessment}{literal statement: correct with missing hypotheses; publication match: exact; verification: read the result; development: none; confidence: high}
\end{GSitem}
\begin{GSitem}{1972-280.20}{A}{Pairwise simplices need not form a simplex: the tolerance assertion}
\GSfield{Kind; source}{problem; \GSpage{P000470}}
\GSfield{The notes}{The explicit \ru{Задача} asks whether three vertices with all three edges present form a simplex, and generalizes this to every finite pairwise adjacent set, calling it a tolerance-space property.}
\GSfield{Statement}{For the complex defined in 1972-280.7, the proposed assertion is: every finite \(X\subset\mathrm{GL}(A)\) whose two-element subsets are simplices is itself a simplex. It is false for every nonzero unital ring. The vertices \(I+E_{12},I+E_{23},I+E_{31}\) are pairwise adjacent but do not lie in any common coset of a permutation-unitriangular subgroup, even after stabilization.}
\GSfield{Reconstruction}{The question is interpreted with the single common translation and order of P000468, not as a flag completion imposed by definition. The complete counterexample, the acyclic-support criterion, and minimal-rank check are proved in the corrective appendix and marked \ours; finite computations are marked \oursComputed. No novelty or priority claim is made for the counterexample.}
\GSfield{Fate}{The exact framework appears in \citet{vaserstein1973foundations}, received on 14 November 1972, the day after this talk; the detailed treatment \citet{vaserstein1976foundations} was received on 26 February 1975 and published in 1976. Crucially, \citet[\S11, Remark~(c), p.~111]{vaserstein1976foundations} asserts exactly this property after defining the complex and identifying left/right conventions in Remark~(b). Thus A records exact publication of the assertion, not its validity. It is not an orphan D item. The same remark says the assertion is not used in that paper, so this correction alone does not refute the paper's K-theory comparisons.}
\GSfield{Present status}{As of 9 November 2026: refuted in the literal model by the fully checked three-vertex example; the published claim is corrected, and no open problem is manufactured from it.}
\GSfield{Assessment}{literal statement: false as written; publication match: exact; verification: read the result; development: developed in this chapter; confidence: high; problem \hyperref[prob:GP-1972-101]{GP-1972-101}}
\end{GSitem}

\GStalk{281}{1972-11-13}{D. A. Kazhdan}{P000470-P000471}{Derived categories by localization of the homotopy category}
Kazhdan defines the homotopy category of complexes, quasi-isomorphisms, the equivalent pre- and post-composition cancellation criteria, and localization to form the derived category. Hartshorne's \emph{Residues and Duality} is explicitly cited. The talk begins below Vasserstein's tolerance problem on P000470 and ends before Buchstaber's heading on P000471; the intervening material is a short foundational exposition rather than an unresolved localization proposal.

\begin{GSitem}{1972-281.1}{E}{Homotopy category of complexes in an abelian category}
\GSfield{Kind; source}{definition; \GSpage{P000470}}
\GSfield{The notes}{For an abelian category \(C\), the notes introduce \(K_C\) as the category of complexes and then explicitly specify homotopy classes of complex morphisms as its arrows.}
\GSfield{Statement}{Let \(C\) be an abelian category, with an appropriate smallness convention. The category \(K(C)\) has cochain complexes as objects and chain maps modulo chain homotopy as morphisms. Composition is induced by composition of chain maps. Homology or cohomology is well defined on these morphism classes.}
\GSfield{Reconstruction}{The full sentence defines the homotopy category, not the unquotiented category of complexes. Cochain indexing is chosen for clarity; the notebook's upward arrows can equally be reindexed homologically.}
\GSfield{Fate}{The notes explicitly refer to \citet[Chapter~I, pp.~19--84]{hartshorne1966residues}, published in 1966 from the Harvard seminar of 1963--1964. The publisher record and contents were checked; the exact modern statements and proofs were read separately in \citet[Tags 05R1 and 05RR]{stacks2026foundations}. This is exposition of an established construction, not a 1972 announcement. See also the opening paragraph of Stacks Tag 05RR.}
\GSfield{Present status}{As of 9 November 2026: standard homotopy-category definition; no ambiguity remains after reading the morphism line.}
\GSfield{Assessment}{literal statement: correct; publication match: exact; verification: read the result; development: none; confidence: high}
\end{GSitem}
\begin{GSitem}{1972-281.2}{E}{Quasi-isomorphisms as the class to be inverted}
\GSfield{Kind; source}{definition; \GSpage{P000470}}
\GSfield{The notes}{The class \(S\subset\operatorname{Hom}K_C\) consists of morphisms inducing isomorphisms on homology.}
\GSfield{Statement}{A morphism in \(K(C)\) is a quasi-isomorphism if its maps \(H^n\) are isomorphisms for every integer \(n\). This condition is invariant under chain homotopy; equivalently the cone of a representing chain map is acyclic. These morphisms form the saturated multiplicative system used in the derived category.}
\GSfield{Reconstruction}{The all-degrees quantifier and cone characterization unpack the abbreviated homology condition. Saturation is a standard structural theorem used for the subsequent localization, not an extra claim that the notes explicitly proved.}
\GSfield{Fate}{The notes explicitly refer to \citet[Chapter~I, pp.~19--84]{hartshorne1966residues}, published in 1966 from the Harvard seminar of 1963--1964. The publisher record and contents were checked; the exact modern statements and proofs were read separately in \citet[Tags 05R1 and 05RR]{stacks2026foundations}. This is exposition of an established construction, not a 1972 announcement. The exact multiplicative-system statement is Stacks Lemma~13.11.2, Tag 05RR.}
\GSfield{Present status}{As of 9 November 2026: standard quasi-isomorphisms and their localization class.}
\GSfield{Assessment}{literal statement: correct; publication match: exact; verification: read the result; development: none; confidence: high}
\end{GSitem}
\begin{GSitem}{1972-281.3}{E}{Equivalent left and right cancellation by quasi-isomorphisms}
\GSfield{Kind; source}{theorem; \GSpage{P000471}}
\GSfield{The notes}{For \(f,g:K\to K'\), the page asserts equivalence between \(fs_1=gs_1\) for a quasi-isomorphism \(s_1:K_1\to K\) and \(s_2f=s_2g\) for a quasi-isomorphism \(s_2:K'\to K_2\).}
\GSfield{Statement}{In the homotopy category of complexes of an abelian category, the following are equivalent for parallel arrows \(f,g:X\to Y\): there exists a quasi-isomorphism \(s:Z\to X\) with \(fs=gs\); there exists a quasi-isomorphism \(t:Y\to W\) with \(tf=tg\); and the images of \(f,g\) in \(D(C)\) agree. All displayed equalities before localization are in \(K(C)\), hence mean chain homotopy for representing maps.}
\GSfield{Reconstruction}{The target arrow's direction is fixed by the scanned diagram. Requiring literal equality of representing chain maps would change the statement; the preceding definition of \(K_C\) is therefore decisive. The equality-in-localization characterization is its standard consequence.}
\GSfield{Fate}{The notes explicitly refer to \citet[Chapter~I, pp.~19--84]{hartshorne1966residues}, published in 1966 from the Harvard seminar of 1963--1964. The publisher record and contents were checked; the exact modern statements and proofs were read separately in \citet[Tags 05R1 and 05RR]{stacks2026foundations}. This is exposition of an established construction, not a 1972 announcement. Stacks Lemma~13.5.5, Tag 05R1, proves both cancellation directions by an acyclic mapping cone and the dual argument.}
\GSfield{Present status}{As of 9 November 2026: a proved calculus-of-fractions lemma; no extra injective or projective resolution hypothesis is needed.}
\GSfield{Assessment}{literal statement: correct; publication match: exact; verification: read the result; development: none; confidence: high}
\end{GSitem}
\begin{GSitem}{1972-281.4}{E}{Derived localization and composition of fractions}
\GSfield{Kind; source}{construction; \GSpage{P000471}}
\GSfield{The notes}{The notes write \(L_C=S^{-1}K_C\), say that the arrows of \(S\) become invertible, and invoke the preceding lemma in defining composition; the reading reference is Hartshorne.}
\GSfield{Statement}{For a small abelian category \(C\), or another setting in which the relevant localization has set-valued Hom classes, define \(D(C)=S^{-1}K(C)\). It is characterized by the universal property that functors from \(K(C)\) inverting all quasi-isomorphisms factor through it. Morphisms may be represented by roofs \(X\xleftarrow{s}Z\xrightarrow{f}Y\), with \(s\in S\); the left/right Ore and cancellation conditions define composition and equality of roofs.}
\GSfield{Reconstruction}{The single displayed cancellation lemma is part of a full calculus of fractions; it alone does not prove every localization axiom. The set-theoretic qualification avoids asserting that every large abelian category has a locally small unbounded derived category. Neither boundedness nor enough injectives is imposed by these particular notes.}
\GSfield{Fate}{The notes explicitly refer to \citet[Chapter~I, pp.~19--84]{hartshorne1966residues}, published in 1966 from the Harvard seminar of 1963--1964. The publisher record and contents were checked; the exact modern statements and proofs were read separately in \citet[Tags 05R1 and 05RR]{stacks2026foundations}. This is exposition of an established construction, not a 1972 announcement. Stacks Definition~13.11.3 gives the localization, and Remark~13.11.4 states the smallness issue explicitly.}
\GSfield{Present status}{As of 9 November 2026: the standard derived category in the stated size setting, with the notebook's shorthand expanded.}
\GSfield{Assessment}{literal statement: correct with missing hypotheses; publication match: exact; verification: read the result; development: none; confidence: high}
\end{GSitem}

\GStalk{282}{1972-11-13}{V. M. Buchstaber}{P000471-P000472}{Yoneda products, bar resolutions, and multivalued formal groups}
The second part of the meeting develops the homological algebra behind the preceding cocycle calculation: graded resolutions, Yoneda composition, normalized bar cochains, and concatenation. The remark connecting noncommutative cohomology to multivalued groups has a traceable continuation in Buchstaber's 1973 announcement and 1976 paper, but the notebook does not state the eventual classification theorem. The quoted Adams bibliographical pointer is defective and is retained as an unresolved exact page reference.

\begin{GSitem}{1972-282.1}{E}{Graded algebra cohomology}
\GSfield{Kind; source}{definition; \GSpage{P000471}}
\GSfield{The notes}{For \(A=\bigoplus_{n\ge0}A_n\), the notes augment to \(A_0\), choose a graded free resolution \(C_\bullet\to A_0\), and put \(H^s(A)=\ker d_{s+1}^*/\operatorname{im}d_s^*\).}
\GSfield{Statement}{For a nonnegatively graded augmented algebra over \(A_0\), define \(\operatorname{Ext}_A^{s,t}(A_0,A_0)\) using a graded projective or free resolution; \(s\) is resolution degree and \(t\) records internal degree. The augmentation is projection to degree zero, and \(A_0\) is made an \(A\)-module through it.}
\GSfield{Fate}{Classical derived-functor definition; the notebook cites \citet{cartaneilenberg1956homological,maclane1963homology}.}
\GSfield{Present status}{As of 9 November 2026: standard; this is algebra-module Ext, as appropriate after dualizing the earlier coalgebra.}
\GSfield{Assessment}{literal statement: correct; publication match: exact; verification: publisher metadata checked; development: none; confidence: high}
\end{GSitem}
\begin{GSitem}{1972-282.2}{E}{Existence of resolutions and independence of the answer}
\GSfield{Kind; source}{theorem; \GSpage{P000471}}
\GSfield{The notes}{The two stated assertions are that a free resolution exists and that \(H^s(A)\) is independent of the chosen resolution.}
\GSfield{Statement}{Every graded module over a graded ring has a graded free resolution. Comparison maps lifting the identity are unique up to chain homotopy; applying graded \(\operatorname{Hom}\) therefore yields canonically isomorphic Ext groups for any two projective resolutions.}
\GSfield{Reconstruction}{This general theorem does not automatically make the particular relative bar resolution free: its terms also require freeness of the reduced algebra over the base, or a suitable projective/relative version.}
\GSfield{Fate}{Established comparison theorem; \citet[Ch.III]{maclane1963homology}.}
\GSfield{Present status}{As of 9 November 2026: proved well before the lecture.}
\GSfield{Assessment}{literal statement: correct; publication match: exact; verification: publisher metadata checked; development: none; confidence: high}
\end{GSitem}
\begin{GSitem}{1972-282.3}{E}{Yoneda multiplication by lifting cocycles}
\GSfield{Kind; source}{construction; \GSpage{P000471-P000472}}
\GSfield{The notes}{A representative \(C_i\to A_0\) is lifted to maps between the resolutions, and composition with a representative \(C_j\to A_0\) yields \(C_{i+j}\to A_0\).}
\GSfield{Statement}{Lift a cocycle of cohomological degree \(i\) to a degree-\(i\) chain map of the resolution, then compose with a degree-\(j\) cocycle. The resulting degree-\(i+j\) class is independent of lifts and representatives and defines the associative Yoneda product, with the order fixed by this composition convention.}
\GSfield{Reconstruction}{The dashed arrows in the scan are chain-map lifts, not canonical maps on the individual free modules. The diagram's existence uses projectivity and exactness.}
\GSfield{Fate}{Classical Yoneda product; \citet[Ch.XII]{maclane1963homology}.}
\GSfield{Present status}{As of 9 November 2026: standard.}
\GSfield{Assessment}{literal statement: correct; publication match: exact; verification: publisher metadata checked; development: none; confidence: high}
\end{GSitem}
\begin{GSitem}{1972-282.4}{E}{Ext need not be graded commutative}
\GSfield{Kind; source}{claim; \GSpage{P000472}}
\GSfield{The notes}{\(H(A)\) is said to be noncommutative ``in a genuine sense'', even though the starting algebra is commutative; a parenthesis contrasts Hopf-algebra examples in topology.}
\GSfield{Statement}{The Yoneda algebra of the augmentation module of a commutative algebra need not obey the graded commutativity rule. For example, if \(A=k\oplus V\) with \(V^2=0\), \(2\le\dim_kV<\infty\), then its normalized bar-cochain differential vanishes and \(\operatorname{Ext}_A(k,k)\cong T(V^*)\). Additional Hopf and symmetry structure is required for standard graded commutativity theorems.}
\GSfield{Reconstruction}{\ours\ The square-zero example follows directly from the next bar differential: every product of two reduced elements vanishes, while concatenation remains free. It gives a precise counterexample; the claim that topology has only commutative rings is read as a comment on the examples under discussion, not as a universal assertion about topology.}
\GSfield{Fate}{Buchstaber's later construction explicitly uses a noncommutative cohomology Hopf algebra: \citet[Introduction]{buchstaber1976twovaluedII}.}
\GSfield{Present status}{As of 9 November 2026: the algebraic warning is correct; the historical aside should not be promoted to a general theorem.}
\GSfield{Assessment}{literal statement: correct; publication match: partial; verification: read the result; development: annotated; confidence: high}
\end{GSitem}
\begin{GSitem}{1972-282.5}{A}{Homological classification of multivalued formal groups}
\GSfield{Kind; source}{method; \GSpage{P000472}}
\GSfield{The notes}{The single sentence says that classification of multivalued groups rests on the preceding noncommutative cohomology ring.}
\GSfield{Statement}{The notes announce a method: use bar-construction cohomology and its noncommutative product to impose associativity and extension conditions for multivalued formal laws. No number of values, complete axioms, coefficient hypotheses, or classification theorem is stated here.}
\GSfield{Reconstruction}{The later two-valued theory supplies a documented continuation, not missing hypotheses to be retroactively inserted into the notebook. In the 1976 paper the ground ring excludes elements of order two; its first-type computation and second-type reduced-base theorem are more specific results than this sentence.}
\GSfield{Fate}{The method is explicitly outlined in \citet[\S~IV, p.~174]{buchstaber1973classification}, received on 3 January 1973; the detailed realization appears in \citet[Introduction and \S1]{buchstaber1976twovaluedII}, received 28 April 1975. The exact all-multivalued classification cannot be claimed from these sources.}
\GSfield{Present status}{As of 9 November 2026: the announced method has a published two-valued realization; the notebook does not formulate a larger open problem.}
\GSfield{Assessment}{literal statement: not determinable; publication match: partial; verification: read the result; development: none; confidence: medium}
\end{GSitem}
\begin{GSitem}{1972-282.6}{E}{Normalized bar resolution and its signs}
\GSfield{Kind; source}{construction; \GSpage{P000472}}
\GSfield{The notes}{With \(\bar A=\bigoplus_{n>0}A_n\), \(C_s=A\otimes_{A_0}\bar A^{\otimes s}\), and the differential merges adjacent entries. The corrected line says the sign exponent is the position when every element is even.}
\GSfield{Statement}{Over a field, or when the reduced pieces are suitably projective over the base, the augmented normalized bar resolution computes Ext. In the all-even internal-degree case, \(d(a_0[a_1|\cdots|a_s])=a_0a_1[a_2|\cdots|a_s]+\sum_{j=1}^{s-1}(-1)^ja_0[a_1|\cdots|a_ja_{j+1}|\cdots|a_s]\). In a suspended graded convention the remaining signs obey the Koszul rule.}
\GSfield{Reconstruction}{The notebook itself repairs the sign qualification. Exactness is supplied by the usual bar contraction; freeness over \(A\) requires the stated base-module hypothesis.}
\GSfield{Fate}{Pre-existing standard resolution, \citet[Ch.X]{cartaneilenberg1956homological}.}
\GSfield{Present status}{As of 9 November 2026: standard under the specified projectivity and sign convention.}
\GSfield{Assessment}{literal statement: correct with missing hypotheses; publication match: exact; verification: publisher metadata checked; development: none; confidence: high}
\end{GSitem}
\begin{GSitem}{1972-282.7}{E}{Dual bar differential and its correct target}
\GSfield{Kind; source}{construction; \GSpage{P000472}}
\GSfield{The notes}{The notes identify \(\operatorname{Hom}_A(C_s,A_0)\) with \((\bar A^*)^{\otimes s}\), insert reduced coproducts, and then write \(\operatorname{im}\bar\Delta\subset\overline{A^*}\).}
\GSfield{Statement}{Under degreewise finite projectivity and graded duality, multiplication dualizes to \(\Delta:A^*\to A^*\widehat\otimes A^*\), and on the augmentation ideal \(\bar\Delta\alpha=\Delta\alpha-\alpha\otimes1-1\otimes\alpha\) takes values in \(\bar A^*\widehat\otimes\bar A^*\). The dual bar differential inserts this reduced coproduct with the bar signs. Without these finiteness conventions one must retain \(\operatorname{Hom}_{A_0}(\bar A^{\otimes s},A_0)\), rather than identify it with a tensor product of duals.}
\GSfield{Reconstruction}{The tensor square missing in the target is visible as an error in both transcription and scan; the preceding diagram already has \(A^*\otimes A^*\). Completion or graded finite pieces resolve the infinite-dimensional duality.}
\GSfield{Fate}{Classical dual bar construction; \citet[Ch.X]{cartaneilenberg1956homological}.}
\GSfield{Present status}{As of 9 November 2026: correct after the duality hypotheses and target correction.}
\GSfield{Assessment}{literal statement: false as written; publication match: partial; verification: publisher metadata checked; development: annotated; confidence: high}
\end{GSitem}
\begin{GSitem}{1972-282.8}{E}{Concatenation represents the Yoneda product}
\GSfield{Kind; source}{construction; \GSpage{P000472}}
\GSfield{The notes}{The product of \([\alpha_1|\cdots|\alpha_s]\) and \([\alpha'_1|\cdots|\alpha'_{s'}]\) is the concatenated word.}
\GSfield{Statement}{On the tensor cobar model, concatenate words. The differential is a derivation for this multiplication with respect to cohomological degree; hence it descends to cohomology and agrees with the Yoneda multiplication, subject to the choice of ordering convention in 1972-282.3.}
\GSfield{Reconstruction}{In general concatenation is not commutative on cochains or on cohomology. This observation explains the method announced immediately above, without supplying its later classification theorem.}
\GSfield{Fate}{Established bar-resolution product; \citet[Ch.XII]{maclane1963homology}.}
\GSfield{Present status}{As of 9 November 2026: standard.}
\GSfield{Assessment}{literal statement: correct; publication match: exact; verification: publisher metadata checked; development: none; confidence: high}
\end{GSitem}

\GStalk{283}{1972-11-20}{A. G. Kostyuchenko}{P000473}{An exponential sine system from an elliptic operator pencil}
The mathematical record is P000473, under \ru{Костюченко. Один пример}; P000472 and P000474 belong to other speakers. The example concerns completeness, biorthogonality, nonbasis behavior, analytic continuation of convergent expansions, Abel summation, interval length and perturbation of frequencies. The signed eigenfunction list and the later one-sided expansion must be distinguished. The basic completeness question is a documented 1969 problem, already answered by Levin in 1971; later work treats minimality, basis properties and perturbations. Several sharper notebook assertions lack their indispensable conditions, and are retained as bounded identification gaps.

\begin{GSitem}{1972-283.1}{E}{The quadratic pencil and its separated eigenfunctions}
\GSfield{Kind; source}{construction; \GSpage{P000473}}
\GSfield{The notes}{\(v^{\prime\prime}+2a\lambda v^{\prime}+b\lambda^2v=0\), \(v(0)=v(\pi)=0\), with \(a^2<b\), \(b>0\); \(v_k=e^{\gamma kx}\sin kx\), \(\gamma=-a/\sqrt{b-a^2}\), \(k=\pm1,\pm2,\ldots\). The associated PDE is \(u_{xx}+2au_{xt}+bu_{tt}=0\).}
\GSfield{Statement}{Let \(a,b\in\mathbb R\), \(d=\sqrt{b-a^2}>0\), and let \(L(\lambda):H^2(0,\pi)\cap H^1_0(0,\pi)\to L^2(0,\pi)\) be the displayed pencil. Its nonzero characteristic values are \(\lambda_k=k/d\), with eigenfunctions \(v_k(x)=e^{-akx/d}\sin kx\), \(k\in\mathbb Z\setminus\{0\}\). The two signs describe the quadratic pencil; one-sign completeness is a different assertion.}
\GSfield{Reconstruction}{The domain and the values \(\lambda_k\) are supplied. Direct substitution \ours after \(v=e^{-a\lambda x}w\) gives \(w^{\prime\prime}+d^2\lambda^2w=0\), independently checking both the sign and the scale.}
\GSfield{Fate}{An exact normalized form and the distinction between the full and half root systems appear in \citet[section 1]{lyubarskiitkachenko1984system}.}
\GSfield{Present status}{As of 9 November 2026: established with the hypotheses and distinctions stated here.}
\GSfield{Assessment}{literal statement: correct with missing hypotheses; publication match: exact; verification: read the result; development: annotated; confidence: high}
\end{GSitem}
\begin{GSitem}{1972-283.2}{E}{Completeness of the one-sided Kostyuchenko system}
\GSfield{Kind; source}{theorem; \GSpage{P000473}}
\GSfield{The notes}{The system is called complete. The final integral reformulation uses only the positive integers.}
\GSfield{Statement}{For every real \(\gamma\), the linear span of \(\{e^{\gamma nx}\sin nx:n\ge1\}\) is dense in \(L^2(0,\pi)\). The same statement holds for \(1\le p<\infty\). Consequently the signed system is complete, but its completeness does not imply minimality.}
\GSfield{Reconstruction}{The half system is the precise object in the later entire-function question. The chronology of the original solution is verified indirectly through primary later research.}
\GSfield{Fate}{\citet[section 1 and Theorem 1]{lyubarskiitkachenko1984system} state the result and identify Kostyuchenko’s 1969 question and Levin’s 1971 function-theoretic proof. The later paper was received 15 February 1983 and published 1984. \citet[Introduction, pp.~509--510]{bilalov2012kostyuchenko} confirms this chronology. Levin’s original 1971 proof was not retrieved in this run.}
\GSfield{Present status}{As of 9 November 2026: the stated completeness problem was solved before the seminar. It is not a surviving open problem.}
\GSfield{Assessment}{literal statement: correct with missing hypotheses; publication match: exact; verification: read the result; development: none; confidence: high; problem \hyperref[prob:GP-1972-303]{GP-1972-303}}
\end{GSitem}
\begin{GSitem}{1972-283.3}{B}{Biorthogonality requires the half system, not the signed list}
\GSfield{Kind; source}{question; \GSpage{P000473}}
\GSfield{The notes}{The coefficients are said to be unknown, but a biorthogonal system exists and is described using an inverse Schwarz--Christoffel integral.}
\GSfield{Statement}{For the one-sided system and real \(\gamma\), completeness and minimality give a unique biorthogonal family \(w_n\in L^2(0,\pi)\), with \(\langle v_m,w_n\rangle=\delta_{mn}\), and uniquely determined formal coefficients \(c_n=\langle f,w_n\rangle\). The full signed system cannot have such a biorthogonal family: the positive subsystem is already complete and there are additional nonzero negative-index functions.}
\GSfield{Reconstruction}{The one-sided correction is forced by completeness, rather than read into the signed display. The brief impossibility check is \ours: any vector biorthogonal to a negative-index element would be orthogonal to the complete positive family, hence zero. The independent check at \(\gamma=0\) gives the explicit duplicate relation \(v_{-n}=-v_n\).}
\GSfield{Fate}{The one-sided minimality is explicitly reviewed in \citet[Introduction, p.~510]{bilalov2012kostyuchenko}, with references to the 1984--1985 work of Lyubarskii and Tkachenko. An exact matching Schwarz--Christoffel formula was not located. This is a partial identification of the asserted coefficient construction, not a claim that the 2012 paper contains that formula.}
\GSfield{Present status}{As of 9 November 2026: the half-system coefficient problem is developed in later literature; the notebook’s particular conformal-map formula remains unverified, and its signed interpretation is false.}
\GSfield{Assessment}{literal statement: false as written; publication match: partial; verification: read the result; development: annotated; confidence: medium; problem \hyperref[prob:GP-1972-303]{GP-1972-303}}
\end{GSitem}
\begin{GSitem}{1972-283.4}{B}{The half system is not a Schauder basis when the tilt is nonzero}
\GSfield{Kind; source}{claim; \GSpage{P000473}}
\GSfield{The notes}{\ru{Это не базис!}: not every vector has a convergent series expansion.}
\GSfield{Statement}{If \(\gamma\in\mathbb R\setminus\{0\}\), \(\{e^{\gamma nx}\sin nx\}_{n\ge1}\) is complete and minimal but is not a Schauder basis of \(L^2(0,\pi)\). For \(\gamma=0\), it is the ordinary orthogonal sine basis. A signed enumeration is not a basis even at \(\gamma=0\), because it is not minimal.}
\GSfield{Reconstruction}{The notebook assumptions allow \(a=0\), hence \(\gamma=0\). The missing hypothesis must be supplied if the assertion is about the half system.}
\GSfield{Fate}{\citet[Theorem, p.~98, and Corollary, p.~99]{kritskov2015minimality} prove nonuniform minimality for a larger class, including this example under the substitution \(\alpha=-i\gamma\). The source cites Kritskov’s 1996 result for the unperturbed system. Thus the mathematical claim is established with the necessary nonzero-tilt qualification.}
\GSfield{Present status}{As of 9 November 2026: established; the zero-tilt half system is an explicit exception, not evidence of an unresolved basis question.}
\GSfield{Assessment}{literal statement: correct with missing hypotheses; publication match: exact; verification: read the result; development: none; confidence: high}
\end{GSitem}
\begin{GSitem}{1972-283.5}{F}{Analyticity of convergent expansions and an unspecified converse}
\GSfield{Kind; source}{claim; \GSpage{P000473}}
\GSfield{The notes}{Every \(L^2\)-convergent expansion has an analytic limit satisfying infinitely many boundary conditions; the converse is asserted.}
\GSfield{Statement}{For real \(\gamma\ne0\), convergence in \(L^2(0,\pi)\) of an ordinary series \[\sum_{n\ge1}c_n e^{\gamma nx}\sin nx\] implies that its limit has a holomorphic representative near every interior point of \((0,\pi)\). The same local conclusion holds for an ordinarily convergent series enumerating the signed system, since its individual summands tend to zero in norm. The notebook does not state the boundary identities or endpoint analyticity class needed for its converse, so that converse cannot be reconstructed uniquely.}
\GSfield{Reconstruction}{Elementary necessity check \ours: convergence implies \(|c_n|\|v_n\|_2\) is bounded, while direct integration gives \(\|v_n\|_2^2=(e^{2\gamma n\pi}-1)/(4\gamma n(1+\gamma^2))\). Thus the series converges normally where \(\gamma\Re z+|\Im z|<\gamma\pi\) if \(\gamma>0\), and where \(\gamma\Re z+|\Im z|<0\) if \(\gamma<0\). These domains contain the open real interval, and normal convergence plus the original norm convergence identifies the limits. For a signed enumeration, the two coefficient bounds hold separately, so both half series converge normally on the intersection of their domains; no separate norm convergence of the halves is needed. The norm estimate was independently rederived by scaling at the endpoint where the exponential is largest; the two signs agree. This proves only the necessary interior analyticity assertion, not the unspecified converse.}
\GSfield{Fate}{The later papers establish nonbasis behavior and analyze analytic function spaces, but the exact necessary-and-sufficient expansion theorem in this wording was not located. \citet[Theorems 5--6]{lyubarskiitkachenko1984system} concern completeness in holomorphic/Smirnov spaces, which is not the asserted converse.}
\GSfield{Present status}{As of 9 November 2026: interior analyticity has the elementary checked proof given here; the full converse and its boundary conditions remain an identification gap, not a defined open conjecture.}
\GSfield{Assessment}{literal statement: not determinable; publication match: thematic only; verification: unresolved lead; development: annotated; confidence: low}
\end{GSitem}
\begin{GSitem}{1972-283.6}{F}{First-order Abel summation and the sign of the exponential}
\GSfield{Kind; source}{claim; \GSpage{P000473}}
\GSfield{The notes}{The notes say first-order Abel summation fails and display \[\sum_{\lambda_k>0}c_ke^{\lambda_k t}v_k(x).\] The limit at \(t\to0\) is again said to be analytic.}
\GSfield{Statement}{The usual decaying Abel factor for \(\lambda_k>0\) is \(e^{-\lambda_k t}\), \(t\downarrow0\); equivalently one may use the written sign with \(t\uparrow0\) from below. The manuscript does not specify this direction, the coefficient functionals or the convergence topology. Consequently no exact general Abel non-summability theorem can be extracted from the displayed formula alone.}
\GSfield{Reconstruction}{The distinction between basis convergence and summability is essential. The sign is retained as a manuscript ambiguity; it is not silently corrected.}
\GSfield{Fate}{The searched Kostyuchenko-system papers identify completeness and basis results, but no exact theorem matching this first-order summation assertion and its sign conventions was verified. Nonbasis behavior by itself does not imply failure of every Abel method.}
\GSfield{Present status}{As of 9 November 2026: the exact summation assertion is unverified because its definition and approach direction are missing.}
\GSfield{Assessment}{literal statement: not determinable; publication match: none found; verification: unresolved lead; development: none; confidence: low}
\end{GSitem}
\begin{GSitem}{1972-283.7}{F}{Nonseparated boundary conditions}
\GSfield{Kind; source}{remark; \GSpage{P000473}}
\GSfield{The notes}{For some nonseparated boundary conditions one can solve the problem.}
\GSfield{Statement}{No particular boundary operators, parameter restrictions, solution space or asserted solvability result are supplied. The phrase can only be indexed as a proposed change from the separated Dirichlet model.}
\GSfield{Fate}{Searches for the exact model and the later Kostyuchenko-system literature did not identify which nonseparated problem was meant. General operator-pencil solvability results are not an exact match.}
\GSfield{Present status}{As of 9 November 2026: the intended boundary-value problem remains unidentified; no open problem is inferred.}
\GSfield{Assessment}{literal statement: not determinable; publication match: none found; verification: unresolved lead; development: none; confidence: low}
\end{GSitem}
\begin{GSitem}{1972-283.8}{F}{Exponential sectorial growth of the resolvent}
\GSfield{Kind; source}{method; \GSpage{P000473}}
\GSfield{The notes}{The explanation of poor behavior is that the resolvent grows exponentially in sectors; a sector sketch is drawn.}
\GSfield{Statement}{For the Dirichlet pencil above, the Green kernel is obtained by conjugating the Dirichlet kernel of \(d^2/dx^2+d^2\lambda^2\) by \(e^{-a\lambda x}\). Away from its poles, this introduces the factor \(e^{-a\lambda(x-y)}\), which explains why a nonselfadjoint exponential tilt can destroy estimates needed for contour summation. The manuscript does not give sector angles, pole separation or an operator norm bound.}
\GSfield{Reconstruction}{The kernel conjugation is an elementary verification \ours from the displayed differential equation; it explains the mechanism without treating the drawing as a quantified estimate.}
\GSfield{Fate}{The pencil is exactly identified in \citet[section 1]{lyubarskiitkachenko1984system}. An exact asymptotic resolvent estimate corresponding to the notebook sketch was not located.}
\GSfield{Present status}{As of 9 November 2026: the elementary conjugation is explicit; the claimed sectorial asymptotic remains insufficiently specified.}
\GSfield{Assessment}{literal statement: correct with missing hypotheses; publication match: thematic only; verification: unresolved lead; development: annotated; confidence: medium}
\end{GSitem}
\begin{GSitem}{1972-283.9}{F}{Shortening the interval and excess completeness}
\GSfield{Kind; source}{claim; \GSpage{P000473}}
\GSfield{The notes}{The system is overdetermined on \([0,\pi-\varepsilon]\).}
\GSfield{Statement}{For the half system with real \(\gamma\ne0\), restriction to a shorter interval is expected here to mean excess completeness, rather than merely density; the notebook does not define whether deletion of one, finitely many, or arbitrarily many terms is intended. Density itself follows by extending an annihilating function by zero to \((0,\pi)\). At \(\gamma=0\), the Fourier-sine comparison supplies the familiar shorter-interval phenomenon.}
\GSfield{Reconstruction}{The zero-extension check \ours proves only density. No exact deletion threshold is attributed to this sentence.}
\GSfield{Fate}{\citet[Theorems 2 and 5]{lyubarskiitkachenko1984system} distinguish completeness on subsets from completeness after deletion in holomorphic spaces. They are relevant but do not license identifying all meanings of the notebook word overdetermined.}
\GSfield{Present status}{As of 9 November 2026: completeness on every shorter interval follows from the known theorem; the intended quantitative excess is unidentified.}
\GSfield{Assessment}{literal statement: not determinable; publication match: partial; verification: unresolved lead; development: annotated; confidence: medium}
\end{GSitem}
\begin{GSitem}{1972-283.10}{F}{Bernstein intervention on extending the completeness interval}
\GSfield{Kind; source}{claim; \GSpage{P000473}; I. N. Bernstein}
\GSfield{The notes}{I. Bernstein says the system remains complete on \([0,\pi+\varepsilon]\), \(0<\varepsilon<\pi\); at \([0,2\pi]\) many orthogonal functions occur.}
\GSfield{Statement}{The printed upper bound is \(2\pi\), but no distinction between the signed system and either half is supplied. The assertion is false under the literal allowed choice \(\gamma=0\): on any \((0,L)\) with \(L>\pi\), choose a nonzero function supported on a small pair of intervals symmetric about \(\pi\), with equal reflected values. Then its integral against every \(\sin nx\) is zero. A nonzero-tilt, specified-index version would need a separate theorem.}
\GSfield{Reconstruction}{The counterexample is an elementary check \ours: \(\sin(n(\pi+s))=-\sin(n(\pi-s))\), so paired integrals cancel. Rechecking with \(n=1\) and \(n=2\) confirms the sign, and the identity proves it for all integers. This refutes the unrestricted statement but does not decide the intended \(\gamma\ne0\) version.}
\GSfield{Fate}{The geometry-dependent completeness results of \citet[Theorems 2--4]{lyubarskiitkachenko1984system} were read, but they do not state this exact real-interval threshold. No verified exact publication for Bernstein’s intervention was found.}
\GSfield{Present status}{As of 9 November 2026: the literal statement is refuted by an allowed parameter value; the intended nonzero-tilt threshold is an identification gap.}
\GSfield{Assessment}{literal statement: false as written; publication match: partial; verification: unresolved lead; development: annotated; confidence: low}
\end{GSitem}
\begin{GSitem}{1972-283.11}{E}{Completeness as an entire-transform uniqueness problem}
\GSfield{Kind; source}{method; \GSpage{P000473}}
\GSfield{The notes}{\(F(\lambda)=\int_0^\pi f(x)e^{\gamma\lambda x}\sin(\lambda x)\,dx\) is entire of exponential type and vanishes at all positive integers. The type and indicator work on a shorter interval; one must use the special form of \(F\).}
\GSfield{Statement}{For \(f\in L^2(0,\pi)\) and real \(\gamma\), the displayed transform is entire and obeys \(|F(\lambda)|\le\sqrt\pi\|f\|_2\exp(\pi\max\{0,\gamma\Re\lambda+|\Im\lambda|\})\). The completeness theorem is equivalent to: if \(F(n)=0\) for all \(n\ge1\), then \(f=0\) almost everywhere. Entire functions of an exponential-type bound alone need not obey this uniqueness conclusion.}
\GSfield{Reconstruction}{The explicit bound is an elementary estimate \ours: \(|\sin(\lambda x)|\le e^{|\Im\lambda|x}\), followed by Cauchy--Schwarz on \((0,\pi)\). The uniqueness conclusion uses the special integral-transform range, not an unrestricted Carlson-type assertion.}
\GSfield{Fate}{\citet[section 1]{lyubarskiitkachenko1984system} identifies the function-theoretic proof as Levin 1971; \citet[Introduction]{bilalov2012kostyuchenko} independently confirms that this was the original problem formulation.}
\GSfield{Present status}{As of 9 November 2026: solved in the stated transform range; no unrestricted uniqueness theorem for all entire functions is asserted.}
\GSfield{Assessment}{literal statement: correct with missing hypotheses; publication match: exact; verification: read the result; development: annotated; confidence: high; problem \hyperref[prob:GP-1972-303]{GP-1972-303}}
\end{GSitem}
\begin{GSitem}{1972-283.12}{B}{Replacing integer frequencies by a sequence}
\GSfield{Kind; source}{question; \GSpage{P000473}}
\GSfield{The notes}{\ru{Неизвестно, что будет, если заменить k на \(\lambda_k\).}}
\GSfield{Statement}{The question asks when completeness/minimality survive for \[\{e^{\gamma\lambda_nx}\sin(\lambda_nx)\}_{n\ge1}.\] No allowed frequency class, separation, asymptotic density or perturbation norm is stated. There is no universal positive answer: repeated frequencies give a finite-dimensional span. A later sufficient-condition theorem must therefore be distinguished from an answer for all unspecified sequences.}
\GSfield{Reconstruction}{No modern all-sequence classification is fabricated, and no residual open conjecture is assigned to a sentence without a specified sequence class.}
\GSfield{Fate}{\citet[Theorems 1--5]{lyubarskii1985system}, received 14 May 1984 and published 1985, explicitly studies this replacement. For frequencies selected from zeros of a Fourier transform with analytic density, Theorem 1 gives completeness and minimality up to finite additions/deletions; Theorems 2--3 impose additional closeness or small-parameter conditions for the unmodified system. This is direct development of the question, but only a partial match to its unspecified scope.}
\GSfield{Present status}{As of 9 November 2026: substantial precise perturbation results exist; the notebook does not define a remaining question with fixed hypotheses that could responsibly be classified C or D.}
\GSfield{Assessment}{literal statement: not determinable; publication match: partial; verification: read the result; development: none; confidence: medium}
\end{GSitem}

\GStalk{284}{1972-11-20}{I. M. Gelfand}{P000474-P000475}{Arnold’s simple singularities, versal deformations and vanishing cycles}
The intervention occupies P000474--P000475, after Kostyuchenko and before Buchstaber’s continuation. Gelfand explains Arnold’s ADE classification, versality through the Jacobian algebra, polynomial discriminants and braid groups, then relates Milnor fibres and intersection forms to Dynkin diagrams. The codimension-six sentence uses Arnold’s convention for critical-point strata, not the dimension of a miniversal base. The closing local-index/asymptotic programme and the planned report on Maslov are preserved as a programme, rather than expanded into an absent theorem.

\begin{GSitem}{1972-284.1}{E}{Function germs and the equivalence of hypersurfaces}
\GSfield{Kind; source}{definition; \GSpage{P000474}}
\GSfield{The notes}{The notes start with \(f:(\mathbb R^n,0)\to(\mathbb R,0)\), \(f(0)=df(0)=0\), then specify that the actual setting is complex. Hypersurfaces are equivalent under changes of variables.}
\GSfield{Statement}{The subsequent formulas concern holomorphic function germs \(f\in\mathcal O_{\mathbb C^n,0}\) with an isolated critical point, under right equivalence \(g=f\circ\phi\), where \(\phi\) is a local biholomorphism fixing the origin. Equivalence of reduced zero hypersurfaces is the weaker contact relation \(g=u(f\circ\phi)\) with \(u\) a unit; it should not be substituted into the Jacobian-algebra versality formula without adjustment.}
\GSfield{Reconstruction}{The distinction between right and contact equivalence is supplied because the notebook alternates the words function and hypersurface.}
\GSfield{Fate}{The function-germ convention is explicit in \citet[section 2]{arnold1972normal}.}
\GSfield{Present status}{As of 9 November 2026: established with the hypotheses and distinctions stated here.}
\GSfield{Assessment}{literal statement: correct with missing hypotheses; publication match: exact; verification: read the result; development: none; confidence: high}
\end{GSitem}
\begin{GSitem}{1972-284.2}{E}{Simple germs and the ADE list}
\GSfield{Kind; source}{theorem; \GSpage{P000474}}
\GSfield{The notes}{A simple hypersurface has only finitely many nearby jet orbits, uniformly in sufficiently high jet order. The \(A,D,E\) diagrams are drawn.}
\GSfield{Statement}{A finitely determined holomorphic germ is right-simple if a sufficiently small neighborhood of a sufficiently high jet meets only finitely many right-equivalence classes. Up to adding nondegenerate quadratic forms, the simple isolated critical-point germs are \(A_k:x^{k+1}\) for \(k\ge1\), \(D_k:x^2y+y^{k-1}\) for \(k\ge4\), and \(E_6:x^3+y^4\), \(E_7:x^3+xy^3\), \(E_8:x^3+y^5\).}
\GSfield{Reconstruction}{The whole list makes the diagrams explicit; the manuscript writes only the A-series normal form at this point.}
\GSfield{Fate}{Exact prior source: \citet[Theorem 2.10 and section 7]{arnold1972normal}, received 21 June 1972, published in the fourth issue of 1972.}
\GSfield{Present status}{As of 9 November 2026: established with the hypotheses and distinctions stated here.}
\GSfield{Assessment}{literal statement: correct with missing hypotheses; publication match: exact; verification: read the result; development: none; confidence: high}
\end{GSitem}
\begin{GSitem}{1972-284.3}{E}{The first nonsimple stratum has codimension six}
\GSfield{Kind; source}{claim; \GSpage{P000474}}
\GSfield{The notes}{\ru{Непростые коразмерности 6}; most orbits are called simple.}
\GSfield{Statement}{In the space of sufficiently high jets with a critical point at the origin and critical value zero, nonsimple germs first occur in codimension six when at least three essential variables are allowed. In two variables the first such locus has codimension seven. This does not say that every nonsimple germ has codimension six, nor that every simple germ has codimension below six: \(E_7\) and \(E_8\) are simple examples of higher codimension.}
\GSfield{Reconstruction}{The ambient stratum and the phrase first occur replace an otherwise ambiguous shorthand. Its codimension convention differs from the Milnor number.}
\GSfield{Fate}{\citet[Introduction and Theorem 7.1]{arnold1972normal} states precisely these thresholds. Thus six is not a transcription error.}
\GSfield{Present status}{As of 9 November 2026: established with the hypotheses and distinctions stated here.}
\GSfield{Assessment}{literal statement: correct with missing hypotheses; publication match: exact; verification: read the result; development: annotated; confidence: high}
\end{GSitem}
\begin{GSitem}{1972-284.4}{E}{Versal deformation and its pullback property}
\GSfield{Kind; source}{definition; \GSpage{P000474}}
\GSfield{The notes}{Every deformation \(G(x,\mu)\) is written \(F(X(x,\mu),\Psi(\mu))\), with \(X(y,0)=y\) and \(\Psi(0)=0\).}
\GSfield{Statement}{A holomorphic unfolding \(F:(\mathbb C^n\times\mathbb C^r,0)\to\mathbb C\), \(F(x,0)=f(x)\), is right-versal if every holomorphic unfolding \(G(y,\mu)\) of \(f\) is induced by a holomorphic parameter map \(\Psi\) fixing zero and a parameter-dependent local biholomorphism \(X\) with \(X(y,0)=y\). The equality is one of function germs in both variables; the moving critical point need not remain the origin for nonzero parameters.}
\GSfield{Fate}{Exact definition: \citet[Definition 8.1]{arnold1972normal}.}
\GSfield{Present status}{As of 9 November 2026: established with the hypotheses and distinctions stated here.}
\GSfield{Assessment}{literal statement: correct with missing hypotheses; publication match: exact; verification: read the result; development: none; confidence: high}
\end{GSitem}
\begin{GSitem}{1972-284.5}{E}{A transverse slice produces a versal family}
\GSfield{Kind; source}{method; \GSpage{P000474}}
\GSfield{The notes}{Take the orbit and a submanifold transverse to it; a sketch illustrates the construction.}
\GSfield{Statement}{For an isolated critical point, choose a sufficiently high jet order so that finite determinacy applies. A finite-dimensional family whose parameter tangent directions surject onto the quotient of function germs by the Jacobian ideal is infinitesimally right-versal and hence right-versal. An arbitrary transverse picture in an unspecified infinite-dimensional function space is not itself a proof.}
\GSfield{Reconstruction}{The infinitesimal criterion and finite-determinacy setting supply the meaning of the geometric sketch.}
\GSfield{Fate}{The operative implication is \citet[Theorem 8.3]{arnold1972normal}.}
\GSfield{Present status}{As of 9 November 2026: established with the hypotheses and distinctions stated here.}
\GSfield{Assessment}{literal statement: correct with missing hypotheses; publication match: exact; verification: read the result; development: none; confidence: high}
\end{GSitem}
\begin{GSitem}{1972-284.6}{E}{The Jacobian algebra and infinitesimal versality}
\GSfield{Kind; source}{theorem; \GSpage{P000474}}
\GSfield{The notes}{Every \(a(x)\) must be \(\sum_i(\partial_i f)h_i+\sum_jc_j\varphi_j\); the quotient by the Jacobian ideal is finite-dimensional, of dimension \(k\).}
\GSfield{Statement}{For a holomorphic germ with an isolated critical point, \[Q_f=\mathcal O_{\mathbb C^n,0}/(\partial_1f,\ldots,\partial_nf)\] has finite dimension \(\mu\). An unfolding is infinitesimally right-versal precisely when the classes of \(\partial_{\lambda_j}F(x,0)\) span \(Q_f\) as a complex vector space. If these classes form a basis, the family is miniversal with \(\mu\) parameters. Finite dimensionality fails in general for nonisolated critical loci.}
\GSfield{Reconstruction}{The notebook word generate means vector-space spanning in the quotient, not generation of the full local ring. The isolated-critical-point assumption is essential.}
\GSfield{Fate}{Exact source: \citet[Definition 8.2, Theorem 8.3 and Corollary 8.4]{arnold1972normal}.}
\GSfield{Present status}{As of 9 November 2026: established with the hypotheses and distinctions stated here.}
\GSfield{Assessment}{literal statement: correct with missing hypotheses; publication match: exact; verification: read the result; development: none; confidence: high}
\end{GSitem}
\begin{GSitem}{1972-284.7}{E}{The Ak Jacobian algebra and miniversal unfolding}
\GSfield{Kind; source}{construction; \GSpage{P000475}}
\GSfield{The notes}{For \(f=z_1^{k+1}+z_2^2+\cdots+z_n^2\), the derivatives are written without their numerical factors; the listed parameters multiply \(1,z_1,\ldots,z_1^{k-1}\).}
\GSfield{Statement}{Over \(\mathbb C\), \(J_f=((k+1)z_1^k,2z_2,\ldots,2z_n)=(z_1^k,z_2,\ldots,z_n)\). Hence \(Q_f\simeq\mathbb C[z_1]/(z_1^k)\), and \(F=f+\sum_{j=0}^{k-1}\lambda_jz_1^j\) is a miniversal right-unfolding. The numerical derivative factors may be omitted from ideal generators, but not from equalities of derivatives.}
\GSfield{Reconstruction}{The ambiguous early subscripts are normalized to 1 using the Jacobian quotient and the later displayed family; the scan itself retains i-like marks.}
\GSfield{Fate}{This is \citet[Corollary 8.4, A-series]{arnold1972normal}.}
\GSfield{Present status}{As of 9 November 2026: established with the hypotheses and distinctions stated here.}
\GSfield{Assessment}{literal statement: false as written; publication match: exact; verification: read the result; development: annotated; confidence: high}
\end{GSitem}
\begin{GSitem}{1972-284.8}{E}{The discriminant and its generic Morse points}
\GSfield{Kind; source}{definition; \GSpage{P000475}}
\GSfield{The notes}{The bifurcation diagram consists of parameters for which singular points remain; on an open set of the diagram the singularities are called Morse.}
\GSfield{Statement}{For a sufficiently small representative of a miniversal unfolding of an isolated critical point, its discriminant is \(\Delta=\{\lambda:\exists x,\ F(x,\lambda)=0,\ d_xF(x,\lambda)=0\}\). A generic smooth point of \(\Delta\) corresponds to one nondegenerate critical point at critical value zero. Thus the generic singular zero fibre has type \(A_1\); off \(\Delta\) the zero fibre is smooth.}
\GSfield{Reconstruction}{Generic smooth stratum replaces the unspecified open set. The word simple here means Morse, which is narrower than the ADE notion of a simple germ.}
\GSfield{Fate}{The definition and ADE discriminants appear in \citet[section 9, opening paragraph]{arnold1972normal}.}
\GSfield{Present status}{As of 9 November 2026: established with the hypotheses and distinctions stated here.}
\GSfield{Assessment}{literal statement: correct with missing hypotheses; publication match: exact; verification: read the result; development: none; confidence: high}
\end{GSitem}
\begin{GSitem}{1972-284.9}{E}{The polynomial discriminant and the swallowtail name}
\GSfield{Kind; source}{construction; \GSpage{P000475}}
\GSfield{The notes}{For \(z_1^{k+1}+z_2^2\), singularity means \(P(z_1)=P^{\prime}(z_1)=0\); the diagram is called a swallowtail.}
\GSfield{Statement}{For \(P(t)=t^{k+1}+\lambda_{k-1}t^{k-1}+\cdots+\lambda_0\), the zero fibre \(P(z_1)+z_2^2=0\) is singular exactly when \(P\) has a multiple root, so \(\Delta\) is its discriminant hypersurface. The ordinary swallowtail is the \(A_3\) discriminant; general \(A_k\) gives the corresponding higher discriminant, also called a generalized swallowtail in Arnold’s terminology.}
\GSfield{Reconstruction}{The general label is not treated as a false formula; its special conventional meaning is clarified.}
\GSfield{Fate}{The qualification generalized and the special \(A_3\) case are explicit in \citet[section 9]{arnold1972normal}.}
\GSfield{Present status}{As of 9 November 2026: established with the hypotheses and distinctions stated here.}
\GSfield{Assessment}{literal statement: correct with missing hypotheses; publication match: exact; verification: read the result; development: none; confidence: high}
\end{GSitem}
\begin{GSitem}{1972-284.10}{E}{The discriminant complement has the corresponding Artin group}
\GSfield{Kind; source}{theorem; \GSpage{P000475}}
\GSfield{The notes}{The fundamental group of the complement of the bifurcation diagram is the braid group of a Coxeter group.}
\GSfield{Statement}{For a simple ADE germ, the complement of its miniversal discriminant is identified with the regular orbit space \(\mathfrak h_{\mathbb C}^{\mathrm{reg}}/W\). Its fundamental group is the Artin group of that ADE Coxeter system: adjacent generators satisfy the braid relation of length three and nonadjacent generators commute. For \(A_k\) this is the ordinary braid group \(B_{k+1}\), not the finite Weyl group itself.}
\GSfield{Reconstruction}{The ADE restriction is necessary. The notebook does not assert this identification for all isolated singularities.}
\GSfield{Fate}{\citet[Proposition 9.3]{arnold1972normal} gives the singularity identification; \citet[theorem, original pp.~57--58; Russian translation pp.~147--148]{brieskorn1971fundamental} gives the presentation. The original publication is 1971; the opened full text is its 1975 Russian translation.}
\GSfield{Present status}{As of 9 November 2026: established with the hypotheses and distinctions stated here.}
\GSfield{Assessment}{literal statement: correct with missing hypotheses; publication match: exact; verification: read the result; development: none; confidence: high}
\end{GSitem}
\begin{GSitem}{1972-284.11}{E}{The Milnor fibre is a bouquet of middle-dimensional spheres}
\GSfield{Kind; source}{theorem; \GSpage{P000475}}
\GSfield{The notes}{A small nonsingular deformed fibre has homology only in degree \(n-1\) and is a bouquet of \(\mu\) spheres.}
\GSfield{Statement}{Let \(f:(\mathbb C^n,0)\to(\mathbb C,0)\), \(n\ge2\), be holomorphic with an isolated critical point. For Milnor radii \(0<|\delta|\ll\varepsilon\), \(F_\delta=f^{-1}(\delta)\cap B_\varepsilon\) has the homotopy type of \(\bigvee_\mu S^{n-1}\), with \(\mu=\dim_{\mathbb C}Q_f\). Thus reduced integral homology vanishes outside degree \(n-1\), where it is \(\mathbb Z^\mu\); ordinary \(H_0\) is also nonzero.}
\GSfield{Reconstruction}{The word reduced repairs the homology sentence. Smallness and transversality are required; an arbitrary regular fibre in an arbitrary ball need not have this type.}
\GSfield{Fate}{\citet[Theorems 6.5 and 7.2, Appendix B, Problem 3]{milnor1968singular} proves the general assertion; the ADE case is \citet[Proposition 9.4]{arnold1972normal}.}
\GSfield{Present status}{As of 9 November 2026: established with the hypotheses and distinctions stated here.}
\GSfield{Assessment}{literal statement: false as written; publication match: exact; verification: read the result; development: annotated; confidence: high}
\end{GSitem}
\begin{GSitem}{1972-284.12}{E}{Intersection form and the sign of vanishing cycles}
\GSfield{Kind; source}{construction; \GSpage{P000475}}
\GSfield{The notes}{For \(n-1\equiv2\pmod4\), vanishing cycles satisfy \(e_i^2=-2\); negative definite, semidefinite and indefinite possibilities are listed.}
\GSfield{Statement}{On \(H_{n-1}(F_\delta;\mathbb Z)\), the middle intersection pairing is symmetric when \(n-1\) is even and skew-symmetric when it is odd. For an oriented vanishing sphere in complex dimension \(n-1\equiv2\pmod4\), the self-intersection is \(-2\). It is \(+2\) when \(n-1\equiv0\pmod4\). Definite, degenerate and indefinite forms all occur among isolated hypersurface singularities; ADE gives the negative Cartan lattice in the stated congruence.}
\GSfield{Reconstruction}{The notebook uses cohomology notation for classes subsequently intersected; homology vanishing cycles are used here. The word always is restricted to the displayed congruence.}
\GSfield{Fate}{The conventions and ADE cycles are explicit in \citet[Proposition 9.4]{arnold1972normal}; signatures beyond ADE are discussed in \citet[section 1]{arnold1973unimodal}.}
\GSfield{Present status}{As of 9 November 2026: established with the hypotheses and distinctions stated here.}
\GSfield{Assessment}{literal statement: correct with missing hypotheses; publication match: exact; verification: read the result; development: none; confidence: high}
\end{GSitem}
\begin{GSitem}{1972-284.13}{E}{Monodromy on the vanishing-cycle lattice}
\GSfield{Kind; source}{construction; \GSpage{P000475}}
\GSfield{The notes}{Monodromy is described only as a representation of the fundamental group, followed by an unfinished target.}
\GSfield{Statement}{The smooth family over a discriminant complement defines \(\rho:\pi_1(B\setminus\Delta,\lambda_0)\to\operatorname{Aut}(H_{n-1}(F_{\lambda_0};\mathbb Z))\), preserving the intersection pairing. For ADE and odd \(n\), Picard--Lefschetz transformations act as the corresponding root reflections, so this representation factors through the Weyl group.}
\GSfield{Reconstruction}{The target and the ADE factorization are supplied from the explicitly identified source; the notebook itself breaks off after the word representation.}
\GSfield{Fate}{\citet[Proposition 9.5]{arnold1972normal} states this precise ADE realization.}
\GSfield{Present status}{As of 9 November 2026: established with the hypotheses and distinctions stated here.}
\GSfield{Assessment}{literal statement: correct with missing hypotheses; publication match: partial; verification: read the result; development: none; confidence: high}
\end{GSitem}
\begin{GSitem}{1972-284.14}{F}{A programme linking local asymptotics and global topology}
\GSfield{Kind; source}{remark; \GSpage{P000475}}
\GSfield{The notes}{Gelfand asks how to remove derivatives from Singer’s calculations and study global topology while sitting at one point; explicit asymptotics for second-order equations and a future Shnirelman report on Maslov are proposed.}
\GSfield{Statement}{This is a programme rather than a defined conjecture: no operator class, coefficient formula, equivalence relation or target estimate is stated. The bounded identifiable themes are local proofs of index formulas and semiclassical/asymptotic methods; the named future talk is not itself evidence of a theorem.}
\GSfield{Fate}{No exact publication can be matched to the undefined request. The later APS intervention in record 296 is a thematic crosslink, not proof that this programme was completed in precisely the intended form.}
\GSfield{Present status}{As of 9 November 2026: the proposed calculation remains unidentified at the necessary level of specificity, so no research problem or C/D item is manufactured.}
\GSfield{Assessment}{literal statement: not determinable; publication match: none found; verification: unresolved lead; development: none; confidence: low}
\end{GSitem}

\GStalk{285}{1972-11-20}{V. M. Buchstaber}{P000476-P000477}{The divided-power Ext calculation and noncommutative formal laws}
The semester's final formal-group lecture explicitly identifies the polynomial cocycle complex with the bar dual of the completed divided-power algebra. It decomposes the characteristic-p algebra, writes its Ext generators, then isolates symmetric cycles. The displayed noncommutative law corrects record 276's unrestricted commutativity claim. Two qualifications remain essential: the exterior/polynomial Ext formula excludes characteristic two, and cocycles classify full laws without higher obstructions only for square-zero infinitesimal coefficients.

\begin{GSitem}{1972-285.1}{E}{Completed divided powers and continuous duality}
\GSfield{Kind; source}{construction; \GSpage{P000476-P000477}}
\GSfield{The notes}{\(A\) is additively \(\Lambda[[X]]\), with \(X_nX_m=\binom{n+m}{n}X_{n+m}\). Its compact dual consists of functionals nonzero on only finitely many basis elements. The next page uses a hatted tensor-product symbol.}
\GSfield{Statement}{Use the degree-completed divided-power algebra and its continuous dual over discrete \(\Lambda\). The dual is \(\Lambda[x]\), paired by \(\langle x^n,X_m\rangle=\delta_{nm}\). In characteristic \(p\), the completion is \(\widehat{\bigotimes}_{i\ge0}\Lambda[Y_i]/(Y_i^p)\) when \(\Lambda\) is an \(\mathbb F_p\)-algebra. Finite internal-degree computations may instead use the graded algebra.}
\GSfield{Reconstruction}{This completion convention clarifies the ordinary product sign on P463. ``Compact'' is the notebook's term for finite-support continuous functionals; it does not mean an arbitrary algebraic dual.}
\GSfield{Fate}{Standard duality used in the later method; \citet[Introduction]{buchstaber1976twovaluedII}.}
\GSfield{Present status}{As of 9 November 2026: established, with the topology specified.}
\GSfield{Assessment}{literal statement: correct with missing hypotheses; publication match: partial; verification: read the result; development: none; confidence: high}
\end{GSitem}
\begin{GSitem}{1972-285.2}{E}{The binomial coproduct}
\GSfield{Kind; source}{construction; \GSpage{P000476}}
\GSfield{The notes}{Pairing multiplication with \(x_n\) gives coefficients \(\binom ni\), and the notes set \(x_n=x^n\), obtaining \(\Delta x^n=(x\otimes1+1\otimes x)^n\).}
\GSfield{Statement}{For the preceding duality, \(\langle\Delta x^n,X_i\otimes X_j\rangle=\binom{i+j}{i}\delta_{n,i+j}\); therefore \(\Delta x^n=\sum_{i=0}^n\binom ni x^i\otimes x^{n-i}\). Its reduced coproduct omits the two endpoint terms.}
\GSfield{Reconstruction}{The stray star on one \(x_j\) in the first expression is a notation slip: all factors belong to the same dual polynomial coalgebra. The subsequent complete formula removes the ambiguity.}
\GSfield{Fate}{Standard binomial duality; \citet[Ch.X]{cartaneilenberg1956homological}.}
\GSfield{Present status}{As of 9 November 2026: correct, with one local notation correction.}
\GSfield{Assessment}{literal statement: correct; publication match: exact; verification: publisher metadata checked; development: none; confidence: high}
\end{GSitem}
\begin{GSitem}{1972-285.3}{E}{Cobar cycles and boundaries recover the Lazard equation}
\GSfield{Kind; source}{claim; \GSpage{P000476}}
\GSfield{The notes}{Expanding the two reduced coproducts produces six terms, cancels the two \(h(x,z)\) terms, and gives the boxed four-term Lazard relation and a short exact sequence.}
\GSfield{Statement}{In internal degree \(n\), \(Z^{2,n}=\ker d_3^*\) consists of normalized homogeneous additive two-cocycles, \(B^{2,n}=\operatorname{im}d_2^*\) consists of \(c((x+y)^n-x^n-y^n)\), and \(0\to B^{2,n}\to Z^{2,n}\to\operatorname{Ext}_A^{2,n}(\Lambda,\Lambda)\to0\) holds in the graded or continuous divided-power category.}
\GSfield{Reconstruction}{Unlike the ambiguous polynomial-module notation on P467, \(A\) now has the correct divided-power multiplication. The single \(d_3\) without a star in the expansion is a notation slip, since the domain and the adjacent formulas are cochains.}
\GSfield{Fate}{Same cocycle calculation as \citet[\S~II, Lemma~3]{lazard1955formal}; the later homological approach is documented in \citet[Introduction]{buchstaber1976twovaluedII}.}
\GSfield{Present status}{As of 9 November 2026: established; the symmetric restriction remains a further step.}
\GSfield{Assessment}{literal statement: correct; publication match: exact; verification: read the result; development: none; confidence: high; problem \hyperref[prob:GP-1972-201]{GP-1972-201}}
\end{GSitem}
\begin{GSitem}{1972-285.4}{E}{Recovering the integral calculation from fields}
\GSfield{Kind; source}{method; \GSpage{P000476}}
\GSfield{The notes}{Because the equation does not involve \(\Lambda\), the notes say it suffices to consider \(\mathbb F_p\)-modules and \(\mathbb Q\).}
\GSfield{Statement}{In any fixed degree the integral coefficient-and-relation module is finitely generated. Its rational rank and its dimensions modulo every prime determine whether it is free of a specified rank, via the structure theorem for finitely generated abelian groups. Base change of kernels or Ext is not automatic; the argument should be made on the finite relation presentation or with the appropriate universal-coefficient sequence.}
\GSfield{Reconstruction}{This supplies the missing integral bridge. Merely observing that an equation has integral coefficients does not justify identifying its solutions over an arbitrary ring with a base change of a field solution space.}
\GSfield{Fate}{Lazard's integral coefficient argument provides the target conclusion: \citet[\S~II, Lemma~3 and Theorem II]{lazard1955formal}.}
\GSfield{Present status}{As of 9 November 2026: a valid proof strategy with the finite-presentation and torsion step supplied.}
\GSfield{Assessment}{literal statement: correct with missing hypotheses; publication match: partial; verification: read the result; development: annotated; confidence: high}
\end{GSitem}
\begin{GSitem}{1972-285.5}{E}{Ext of a truncated polynomial factor: the characteristic-two exception}
\GSfield{Kind; source}{claim; \GSpage{P000477}}
\GSfield{The notes}{For every prime \(p\), the notes write \(H^*(A_i)=\Lambda_p(u_i)\otimes\mathbb F_p[v_i]\), \(|u_i|=(1,p^i)\), \(|v_i|=(2,p^{i+1})\), and list degree-two generators \(u_iu_j\) for \(i<j\) and \(v_k\).}
\GSfield{Statement}{For \(A_i=k[t]/(t^p)\), \(k=\mathbb F_p\), \(|t|=p^i\), odd \(p\) gives \(\operatorname{Ext}_{A_i}(k,k)=\Lambda(u_i)\otimes k[v_i]\) in the displayed bidegrees. At \(p=2\) it is instead \(k[u_i]\), with \(v_i=u_i^2\), not an independent polynomial generator. Thus the listed vector-space basis in cohomological degree two can still be used if this identification is stated.}
\GSfield{Reconstruction}{\ours\ For \(p=2\), the reduced augmentation ideal has square zero, so its one-generator tensor bar cochain algebra has zero differential; all powers of the degree-one class survive. This directly disproves the exterior factor. The odd-prime periodic resolution alternates multiplication by \(t\) and \(t^{p-1}\).}
\GSfield{Fate}{The normalized bar construction supplies the correction; \citet[Ch.X]{cartaneilenberg1956homological}.}
\GSfield{Present status}{As of 9 November 2026: the uniformly stated ring formula is false at two; the corrected ring calculation is standard.}
\GSfield{Assessment}{literal statement: false as written; publication match: partial; verification: publisher metadata checked; development: annotated; confidence: high}
\end{GSitem}
\begin{GSitem}{1972-285.6}{E}{An associative noncommutative one-dimensional formal group}
\GSfield{Kind; source}{claim; \GSpage{P000477}}
\GSfield{The notes}{The emphatic example is \(F=x+y+axy^p\) over \(\mathbb F_p[a]/(a^2)\), and the text says one-dimensional noncommutative groups exist with torsion.}
\GSfield{Statement}{For any prime \(p\), this series is a formal group law. Its associator is zero because \((y+z)^p=y^p+z^p\) and \(a^2=0\); its commutator is \(a(xy^p-yx^p)\ne0\). Therefore commutativity fails over this nonreduced torsion base.}
\GSfield{Reconstruction}{\oursComputed\ The included script checks the associator symbolically modulo \((p,a^2)\) and independently by arithmetic in dual-number rings for \(p=2,3,5\). The general proof is the Frobenius identity just given. This is the speaker's own later correction to 1972-276.2, not an invented counterexample.}
\GSfield{Fate}{The characteristic-two version was already given by \citet[\S~I, after Theorem I]{lazard1955formal}.}
\GSfield{Present status}{As of 9 November 2026: correct; it requires nilpotents as well as torsion, so torsion alone is not asserted to force noncommutativity.}
\GSfield{Assessment}{literal statement: correct; publication match: exact; verification: read the result; development: annotated; confidence: high}
\end{GSitem}
\begin{GSitem}{1972-285.7}{E}{Cocycles and actual laws: the square-zero qualification}
\GSfield{Kind; source}{claim; \GSpage{P000477}}
\GSfield{The notes}{After the example, the notes say ``There are as many of them as there are cycles.''}
\GSfield{Statement}{If \(I\) is an ideal with \(I^2=0\), normalized \(I\)-valued additive two-cocycles correspond exactly to laws \(F=x+y+h\) whose nonlinear coefficients lie in \(I\): higher nonlinear compositions vanish. For unrestricted coefficients a first-order cocycle is only a candidate deformation; higher associativity obstructions must be solved. Modding out by coordinate changes also identifies cocycles differing by a coboundary.}
\GSfield{Reconstruction}{The local square-zero context makes the sentence valid. Extending it to arbitrary full laws or to their isomorphism classes would be false. Lazard gives noncommutative truncated laws over a reduced field that cannot extend to full one-dimensional laws.}
\GSfield{Fate}{The extension distinction is explicit in \citet[\S~II, Theorem III and the following remark]{lazard1955formal}.}
\GSfield{Present status}{As of 9 November 2026: valid as an infinitesimal statement, not an unrestricted counting theorem.}
\GSfield{Assessment}{literal statement: correct with missing hypotheses; publication match: partial; verification: read the result; development: annotated; confidence: high}
\end{GSitem}
\begin{GSitem}{1972-285.8}{E}{Normalized binomial cocycles and one indecomposable per degree}
\GSfield{Kind; source}{theorem; \GSpage{P000477}}
\GSfield{The notes}{The \(v_k\) classes are linked to symmetric groups. The coefficients of \((x+y)^{p^{k+1}}-x^{p^{k+1}}-y^{p^{k+1}}\) are divisible by \(p\), and the bracketed formula divides by \(p\). The last mathematical line says there is one generator in every dimension.}
\GSfield{Statement}{Write \(B_n=(x+y)^n-x^n-y^n\). Its coefficient gcd is \(d_n=p\) when \(n=p^r\), otherwise \(d_n=1\). The integral primitive polynomial \(C_n=B_n/d_n\) generates all symmetric normalized degree-\(n\) additive cocycles over any base. In the Lazard grading this yields one free indecomposable in degree \(n-1\), and the full polynomial-ring theorem follows with the generation argument.}
\GSfield{Reconstruction}{The \(1/p\) normalization is actually present in the scan and in the transcription's diagram; it must not be reported as an omission. Over \(\mathbb F_p\), the unnormalized \(B_{p^r}\) is zero, whereas the reduction of the integral quotient need not be zero.}
\GSfield{Fate}{Exact published theorem: \citet[\S~II, Lemma~3 and Theorem II]{lazard1955formal}.}
\GSfield{Present status}{As of 9 November 2026: established since 1955; the semester closes an exposition of this result.}
\GSfield{Assessment}{literal statement: correct; publication match: exact; verification: read the result; development: none; confidence: high; problem \hyperref[prob:GP-1972-201]{GP-1972-201}}
\end{GSitem}

\GStalk{286}{}{}{P000477-P000479}{Year 1973}
This catalogue entry is a year heading, not an independent lecture. P477 closes the autumn formal-group series; the dated 12 February 1973 meeting begins on P478. No mathematical occurrence is assigned to the year label.


\GStalk{287}{1973-02-12}{V. M. Buchstaber}{P000478-P000479}{Sullivan theory, homotopy structures, and normal invariants}
The opening plan lists topology reports by Buchstaber and Mishchenko and work of Deligne and Langlands; no Langlands statement is supplied. Buchstaber then begins an exposition of Sullivan. The defined ``smoothings'' are homotopy structures, and the normal-data notation omits an essential choice of spherical trivialization. The PL classification and the unresolved smooth homotopy-group calculation must be separated. The remarks refer to work circulating in 1967, not to a theorem first announced in this seminar. The page also announces a plan for topology reports by Buchstaber and Mishchenko and for work of Deligne and Langlands; the latter heading supplies no additional mathematical assertion on this page.

\begin{GSitem}{1972-287.1}{E}{Homotopy smoothings and their equivalence}
\GSfield{Kind; source}{definition; \GSpage{P000478}}
\GSfield{The notes}{A smoothing is a smooth \(N^n\) with a homotopy equivalence \(f:N\to M\); two are equivalent when a diffeomorphism of source manifolds makes the triangle commute up to homotopy.}
\GSfield{Statement}{The object is the marked smooth homotopy structure set \(\mathcal S^{\mathrm{Diff}}(M)\): pairs \((N,f)\) modulo a diffeomorphism \(h:N_1\to N_2\) with \(f_2h\simeq f_1\). For an oriented target the marking determines the orientation of the source. This differs from a smoothing of a fixed topological manifold, where the comparison map is a homeomorphism.}
\GSfield{Reconstruction}{The later notation \(hS\) on P481 and P488 makes the distinction explicit. For the surgery statements below we use closed connected manifolds of dimension at least six and fix a reference smooth or PL structure as appropriate.}
\GSfield{Fate}{Definition and equivalence agree with \citet[pp.81--82]{sullivan1996triangulating}, notes circulated 1967.}
\GSfield{Present status}{As of 9 November 2026: standard; the earlier letter S is retained as historical notation.}
\GSfield{Assessment}{literal statement: correct; publication match: exact; verification: read the result; development: none; confidence: high}
\end{GSitem}
\begin{GSitem}{1972-287.2}{E}{Twenty-eight smooth homotopy seven-spheres}
\GSfield{Kind; source}{theorem; \GSpage{P000478}}
\GSfield{The notes}{The page attributes \(\operatorname{card}S'(S^7)=28\) to Milnor. The prime after S is visually present but not defined.}
\GSfield{Statement}{The group \(\Theta_7\) of oriented smooth homotopy seven-spheres up to oriented diffeomorphism is cyclic of order 28. The corresponding marked structure set over \(S^7\) has28 elements. This is not an assertion that there are28 unoriented diffeomorphism classes after forgetting the marking and orientation.}
\GSfield{Reconstruction}{The interpretation of the unexplained prime is supplied by the precise established theorem. The counting theorem is due to Kervaire and Milnor, rather than a new seminar computation.}
\GSfield{Fate}{\citet{kervairemilnor1963spheres}; the statement \(\Theta_7=\mathbb Z/28\) is also explicit in \citet[\S13.3]{ranicki2002surgery}.}
\GSfield{Present status}{As of 9 November 2026: established since 1963.}
\GSfield{Assessment}{literal statement: correct with missing hypotheses; publication match: exact; verification: read the result; development: none; confidence: high}
\end{GSitem}
\begin{GSitem}{1972-287.3}{E}{Describe the homotopy structure set}
\GSfield{Kind; source}{question; \GSpage{P000478}}
\GSfield{The notes}{``Description of \(S(M^n)\)—?'' The next paragraph says Sullivan solved the analogous PL formulation.}
\GSfield{Statement}{For a fixed closed simply connected smooth \(n\)-manifold \(M\), \(n\ge6\), describe the marked homotopy structure set. Surgery gives \([\Sigma M_+,G/O]\to L_{n+1}(\mathbb Z)\curvearrowright\mathcal S^{\mathrm{Diff}}(M)\to[M,G/O]\to L_n(\mathbb Z)\), with exactness in pointed sets and an action on fibres. An individual numerical classification additionally requires computing the normal-invariant set and the obstruction and action maps.}
\GSfield{Reconstruction}{The high-dimensional closed simply connected scope is an editorial reconstruction justified by the continuation on P488. It is not a universal solution in every dimension or fundamental group. This occurrence continues the normal-data and surgery themes of 1968-171.10, 1968-171.15, and 1968-173.2; the older homotopy-torus problem is narrower and receives no reused problem ID here.}
\GSfield{Fate}{The structural answer predates the lecture: \citet[Theorem~3]{sullivan1996triangulating}; modern precise action form in \citet[Theorem~13.2]{ranicki2002surgery}.}
\GSfield{Present status}{As of 9 November 2026: the specified structural classification is established; the notebook asks no concrete uncomputed manifold case.}
\GSfield{Assessment}{literal statement: correct with missing hypotheses; publication match: exact; verification: read the result; development: none; confidence: high; problem \hyperref[prob:GP-1972-203]{GP-1972-203}}
\end{GSitem}
\begin{GSitem}{1972-287.4}{E}{PL homotopy structures and the standard sphere}
\GSfield{Kind; source}{definition; \GSpage{P000478}}
\GSfield{The notes}{Replacing diffeomorphisms by PL homeomorphisms gives \(PL(M)\), and \(PL(S^n)=\{S^n\}\) is asserted for \(n\ge5\).}
\GSfield{Statement}{For a fixed PL manifold, define \(\mathcal S^{PL}(M)\) by PL source manifolds marked by homotopy equivalences, modulo PL homeomorphisms over the target up to homotopy. A PL homotopy sphere of dimension \(n\ge5\) is PL homeomorphic to the standard sphere, so this marked structure set is a singleton.}
\GSfield{Reconstruction}{The statement concerns combinatorial PL manifold structures. It is not a uniqueness assertion for arbitrary singular triangulations, nor a smooth Poincare theorem.}
\GSfield{Fate}{Classical high-dimensional PL Poincare result used in \citet[pp.69--72,81--82]{sullivan1996triangulating}.}
\GSfield{Present status}{As of 9 November 2026: established in the stated PL high-dimensional setting.}
\GSfield{Assessment}{literal statement: correct; publication match: exact; verification: read the result; development: none; confidence: high}
\end{GSitem}
\begin{GSitem}{1972-287.5}{E}{Thom quotient and the bottom fibre sphere}
\GSfield{Kind; source}{construction; \GSpage{P000478}}
\GSfield{The notes}{\(T_\xi=D(\xi)/S(\xi)\), and a fibre over a chosen point gives \(i:S^m\to T_\xi\) for an \(m\)-plane bundle.}
\GSfield{Statement}{For a real vector bundle \(\xi\) over a finite CW complex, its Thom space is the quotient of its disk bundle by its sphere bundle. The compactified fibre over a basepoint embeds as \(S^m\). Adding a trivial line suspends the Thom space, so stable bundle statements correspond to stable maps of these spaces.}
\GSfield{Reconstruction}{A fibre orientation specifies the sign of the degree of a restriction. The stable rank is allowed to increase; this matters for the next criterion.}
\GSfield{Fate}{\citet[\S2]{atiyah1961thom}.}
\GSfield{Present status}{As of 9 November 2026: standard.}
\GSfield{Assessment}{literal statement: correct; publication match: exact; verification: read the result; development: none; confidence: high}
\end{GSitem}
\begin{GSitem}{1972-287.6}{E}{Thom coreducibility characterizes J-triviality}
\GSfield{Kind; source}{claim; \GSpage{P000478}}
\GSfield{The notes}{\(\xi\) is called J-trivial if a map \(\varphi:T_\xi\to S^m\) has \(\deg(\varphi i)=1\). Connectedness is used to dismiss dependence on the chosen point.}
\GSfield{Statement}{For a stable vector bundle over a connected finite CW complex, \(J(\xi)=0\) exactly when, after sufficient stabilization, its Thom space admits a map to the bottom sphere restricting to degree one on a fibre. Equivalently its associated stable sphere fibration is fibre homotopy trivial. A degree-one restriction also forces orientability; connectedness alone does not make an unoriented sign choice canonical.}
\GSfield{Reconstruction}{The stable qualification is already suggested by ``large dimension'' in the notes. The precise statement is stable coreducibility, not a claim that every unstable sphere bundle is linearly trivial.}
\GSfield{Fate}{\citet[Proposition~2.8]{atiyah1961thom}, using \citet[Corollary~2]{dold1955fibre}.}
\GSfield{Present status}{As of 9 November 2026: established before the lecture.}
\GSfield{Assessment}{literal statement: correct with missing hypotheses; publication match: exact; verification: read the result; development: none; confidence: high}
\end{GSitem}
\begin{GSitem}{1972-287.7}{E}{Kernel of J versus a normal invariant}
\GSfield{Kind; source}{definition; \GSpage{P000478-P000479}}
\GSfield{The notes}{\(K_J(M)\) is defined as stable bundles that are J-trivial. The later sequence treats it as the object represented by the smooth normal-invariant space.}
\GSfield{Statement}{Bare stable bundle classes form \(\ker(\widetilde{KO}(M)\to J(M))\). The standard normal-invariant set instead records pairs \((\xi,\varphi)\), where \(\varphi\) is a specified stable spherical trivialization, modulo the corresponding equivalence. These pairs are classified by \([M,G/O]\), with \(G/O\) the homotopy fibre of \(BO\to BG\). Forgetting \(\varphi\) maps onto the kernel but can identify distinct normal invariants.}
\GSfield{Reconstruction}{This is an essential refinement, not a change of notation alone. P488 draws a trivialization but its prose still suppresses it. The homotopy fibre represents a chosen null-homotopy; the set-theoretic kernel of maps into BO and BG does not by itself retain that data.}
\GSfield{Fate}{The pair is explicit in \citet[p.80]{sullivan1996triangulating}; see also \citet[Proposition~9.21]{ranicki2002surgery}.}
\GSfield{Present status}{As of 9 November 2026: the normal-invariant interpretation is established; the literal bare-bundle functor is insufficient for the asserted classification.}
\GSfield{Assessment}{literal statement: correct with missing hypotheses; publication match: exact; verification: read the result; development: annotated; confidence: high}
\end{GSitem}
\begin{GSitem}{1972-287.8}{E}{Atiyah normal-bundle difference}
\GSfield{Kind; source}{theorem; \GSpage{P000479}}
\GSfield{The notes}{Atiyah's theorem, referenced as Matematika10:5, yields \(S(M)\to K_J(M)\), described only as a difference.}
\GSfield{Statement}{If \(f:N\to M\) is a homotopy equivalence of closed smooth manifolds and \(g\) is a homotopy inverse, then \(\xi=g^*\nu_N+\tau_M\) has trivial stable spherical fibration. Equivalently \(J(\tau_N)=J(f^*\tau_M)\). The associated structure normal invariant further includes the trivialization induced by the homotopy equivalence.}
\GSfield{Reconstruction}{The precise difference formula is supplied by the later P481 and P488 occurrences. The Matematika pointer resolves exactly to the 1966 Russian translation of Atiyah's 1961 paper; Buchstaber was its translator.}
\GSfield{Fate}{\citet[Theorem~3.6]{atiyah1961thom}.}
\GSfield{Present status}{As of 9 November 2026: an established homotopy-invariance theorem for the stable spherical normal data, not for the vector bundle itself.}
\GSfield{Assessment}{literal statement: correct; publication match: exact; verification: read the result; development: none; confidence: high}
\end{GSitem}
\begin{GSitem}{1972-287.9}{E}{The PL surgery classification and the periodic obstruction}
\GSfield{Kind; source}{claim; \GSpage{P000479}}
\GSfield{The notes}{A crossed-out ``always'' is replaced by dimension congruences. The display injects \(PL(M)\) into \(K_{PLJ}(M)\) and lists cokernels \(\mathbb Z\) for \(n\equiv0\) and \(\mathbb Z/2\) for \(n\equiv2\pmod4\).}
\GSfield{Statement}{For a closed simply connected PL manifold of dimension \(n\ge6\), the refined normal invariant gives an injection \(\mathcal S^{PL}(M)\to[M,G/PL]\), whose image is the zero fibre of the surgery obstruction \(\theta:[M,G/PL]\to L_n(\mathbb Z)\). Here \(L_n(\mathbb Z)=\mathbb Z\) for \(n\equiv0\), \(\mathbb Z/2\) for \(n\equiv2\), and zero for odd \(n\). Thus in odd dimensions the normal invariant is a bijection. The exactness is that of pointed sets; a group cokernel is not implied.}
\GSfield{Reconstruction}{Compactness, absence of boundary, simple connectivity, stable trivialization data, and the dimension restriction are supplied. With boundary there are separate relative conventions; with nontrivial fundamental group the obstruction groups involve the group ring. A terminal group in an exact sequence is not automatically its cokernel.}
\GSfield{Fate}{\citet[Theorem~3, p.82]{sullivan1996triangulating}, circulated 1967; compare \citet[\S13.1]{ranicki2002surgery}.}
\GSfield{Present status}{As of 9 November 2026: the qualified PL theorem is established; the notebook display is too compressed to support an unconditional group exact sequence.}
\GSfield{Assessment}{literal statement: correct with missing hypotheses; publication match: exact; verification: read the result; development: annotated; confidence: high}
\end{GSitem}
\begin{GSitem}{1972-287.10}{E}{Sullivan computation of the PL normal-invariant homotopy groups}
\GSfield{Kind; source}{theorem; \GSpage{P000479}}
\GSfield{The notes}{The notes say Sullivan computed the homotopy groups of the representing space for \(K_{PLJ}\), without writing the groups.}
\GSfield{Statement}{For the refined classifying space \(G/PL\), the positive homotopy groups are abstractly \(\mathbb Z\) in degrees divisible by four, \(\mathbb Z/2\) in degrees congruent to two modulo four, and zero in odd degrees. In dimensions at least five these agree with the simply connected surgery groups through the surgery obstruction. The degree-four comparison has a special normalization and should not be inferred from that high-dimensional statement.}
\GSfield{Reconstruction}{The group formula comes from the cited literature. It is not present on this notebook page, which gives only the attribution. Modern G/PL denotes the stable version of Sullivan's F/PL.}
\GSfield{Fate}{\citet[pp.69--72,83--84]{sullivan1996triangulating}; the high-dimensional identification is also in \citet[Remark~13.25]{ranicki2002surgery}.}
\GSfield{Present status}{As of 9 November 2026: established, including the need to distinguish the dimension-four comparison map.}
\GSfield{Assessment}{literal statement: correct with missing hypotheses; publication match: exact; verification: read the result; development: none; confidence: high}
\end{GSitem}
\begin{GSitem}{1972-287.11}{C}{Smooth normal invariants: a reconstructed all-degree computation programme}
\GSfield{Kind; source}{problem; \GSpage{P000479}}
\GSfield{The notes}{Immediately after the PL calculation: ``For \(K_J(M^n)\) it has not been done. A sequence for \(S(M^n)\) exists too.''}
\GSfield{Statement}{Under the intended refinement \(K_J\rightsquigarrow[M,G/O]\), the remaining program is to determine \(\pi_q(G/O)\) in all degrees. For \(q\ge2\), the fibration gives \(\pi_q(O)\xrightarrow{J}\pi_q^S\to\pi_q(G/O)\to\pi_{q-1}(O)\xrightarrow{J}\pi_{q-1}^S\). Finite-degree calculations therefore require stable stems, the J maps, and extension data.}
\GSfield{Reconstruction}{The notes do not name G/O, specify a degree range, or formulate an explicit conjecture. The all-degree program is a cautious reconstruction of the comparison sentence; taken literally for the bare kernel functor, a representing-space problem is not specified. No new finite range is invented.}
\GSfield{Fate}{The reduction and low-degree tables are established in \citet[\S9.2]{ranicki2002surgery}. Finite stable-stem progress continues, for example \citet[Theorem~1.1]{burklundisaksenxu2026synthetic} (preprint 2024, online 2025, volume 2026), computing particular 82nd and 83rd 2-primary stems, not all \(\pi_q(G/O)\).}
\GSfield{Present status}{As of 9 November 2026: many finite ranges are known; no complete all-degree numerical determination was located in the primary literature checked through the status date. The reconstructed all-degree computation programme remains unfinished; no specific finite unresolved group is asserted.}
\GSfield{Assessment}{literal statement: not determinable; publication match: partial; verification: read the result; development: none; confidence: medium; problem \hyperref[prob:GP-1972-202]{GP-1972-202}}
\GSfield{Worth developing}{A specific bounded range or extension question could justify research. The notebook supplies neither; attempting a wholesale new stable-stem computation would not be a sound chapter-sized development. We retain the exact-sequence reduction as annotation and claim no new theorem.}
\end{GSitem}

\GStalk{289}{1973-02-26}{V. M. Buchstaber}{P000481}{The Hauptvermutung and the J-trivial normal difference}
This is a single-page continuation beginning with its own dated heading. It distinguishes homotopy smoothings from homotopy triangulations, states a compressed Hauptvermutung theorem, and checks the stable normal-bundle difference. The scan's superscript \(H^3\) is a substantive error: Sullivan's theorem uses \(H_3\) with integral coefficients, equivalently the 2-torsion in \(H^4\). The final Thom criterion is attributed to Dold and can be located through Atiyah's proof.

\begin{GSitem}{1972-289.1}{E}{Homotopy structures distinguished from actual structures}
\GSfield{Kind; source}{definition; \GSpage{P000481}}
\GSfield{The notes}{\(hS(M^n)\) means homotopy smoothings and \(ht(M^n)\) homotopy triangulations; the latter is said to be describable.}
\GSfield{Statement}{These are the marked smooth and PL structure sets defined by homotopy equivalences, with source isomorphisms over the target up to homotopy. Actual smoothings or PL structures on a fixed topological manifold use homeomorphic comparison maps and require a separate equivalence convention.}
\GSfield{Reconstruction}{The letter h supplies the qualification omitted in the initial P478 label. The notebook does not define every equivalence relation for actual, rather than homotopy, structures here.}
\GSfield{Fate}{Matches \citet[pp.81--82]{sullivan1996triangulating}.}
\GSfield{Present status}{As of 9 November 2026: standard notation with the distinction retained.}
\GSfield{Assessment}{literal statement: correct; publication match: exact; verification: read the result; development: none; confidence: high}
\end{GSitem}
\begin{GSitem}{1972-289.2}{E}{Hauptvermutung: homology index and dimensional hypotheses}
\GSfield{Kind; source}{claim; \GSpage{P000481}}
\GSfield{The notes}{A homeomorphism of manifolds can be replaced by a PL homeomorphism when M is simply connected and has no2-torsion in \(H^3\). The scan clearly places3 above H.}
\GSfield{Statement}{A verified pre-lecture version is: if \(h:L\to M\) is a homeomorphism of compact closed simply connected PL manifolds of dimension at least six and \(H_3(M;\mathbb Z)\) has no2-torsion, then h is homotopic to a PL homeomorphism. Equivalently the integral \(H^4(M;\mathbb Z)\) has no2-torsion. For manifolds with boundary the corresponding simple-connectivity assumptions on boundary components are needed.}
\GSfield{Reconstruction}{Replace superscript \(H^3\) by subscript \(H_3\), or by \(H^4\). These are different groups; the correction follows from Sullivan's original theorem and the universal coefficient theorem. ``Replace'' means homotope the given map, not assert isotopy. The missing compactness and dimension hypotheses must also be supplied.}
\GSfield{Fate}{\citet[Main Theorem]{sullivan1967hauptvermutung}; precise obstruction and corollary in \citet[Theorem H and Corollary~1, pp.93--94]{sullivan1996triangulating}.}
\GSfield{Present status}{As of 9 November 2026: the corrected high-dimensional homotopy Hauptvermutung statement is established; the \(H^3\) sentence is not the published theorem.}
\GSfield{Assessment}{literal statement: false as written; publication match: partial; verification: read the result; development: annotated; confidence: high}
\end{GSitem}
\begin{GSitem}{1972-289.3}{E}{The normal difference is independent of a marked representative}
\GSfield{Kind; source}{construction; \GSpage{P000481}}
\GSfield{The notes}{For \(f:N\to M\), with homotopy inverse g, the notes form \(\tau_M+g^*\nu_N\), where \(\nu_N=-\tau_N\) stably, and ask as an exercise that equivalent smoothings induce isomorphic bundles.}
\GSfield{Statement}{The virtual rank-zero class is \(\xi_f=[\tau_M]-g^*[\tau_N]\in\widetilde{KO}(M)\). If \(f_2h\simeq f_1\) for a diffeomorphism \(h:N_1\to N_2\), then \(g_2\simeq hg_1\), \(h^*\tau_{N_2}\cong\tau_{N_1}\), and \(\xi_{f_2}=\xi_{f_1}\). Homotopy inverses are themselves unique up to homotopy, so the construction is well defined.}
\GSfield{Reconstruction}{\ours\ This is the complete solution of the stated elementary exercise, using pullback homotopy invariance. It proves equality of stable bundle classes; it does not by itself identify the chosen spherical trivializations in a normal invariant.}
\GSfield{Fate}{Compatible with \citet[Theorem~3.6]{atiyah1961thom} and the refined construction of \citet[pp.80--81]{sullivan1996triangulating}.}
\GSfield{Present status}{As of 9 November 2026: the exercise is resolved under the stated stable interpretation.}
\GSfield{Assessment}{literal statement: correct; publication match: exact; verification: read the result; development: annotated; confidence: high}
\end{GSitem}
\begin{GSitem}{1972-289.4}{E}{J-triviality and proper fibrewise equivalence}
\GSfield{Kind; source}{theorem; \GSpage{P000481}}
\GSfield{The notes}{The resulting \(\xi\) is declared J-trivial, drawn as a proper homotopy equivalence \(\xi\to M\times\mathbb R^N\) over the identity of M.}
\GSfield{Statement}{Atiyah's theorem makes the stable sphere bundle of \(\xi_f\) fibre homotopy trivial. In sufficiently large rank it may be represented by a proper fibrewise homotopy equivalence of the associated Euclidean bundles; compactifying the fibres gives a degree-one equivalence of sphere bundles. The proper and fibrewise qualifications are part of the data.}
\GSfield{Reconstruction}{Ordinary homotopy equivalence of the uncompactified total spaces would be too weak: every vector-bundle total space retracts onto its base. The notebook's diagram and the word proper prevent that misreading.}
\GSfield{Fate}{\citet[Theorem~3.6]{atiyah1961thom}; the proper-equivalence formulation is explicit in \citet[p.80]{sullivan1996triangulating}.}
\GSfield{Present status}{As of 9 November 2026: established after stabilization.}
\GSfield{Assessment}{literal statement: correct with missing hypotheses; publication match: exact; verification: read the result; development: none; confidence: high}
\end{GSitem}
\begin{GSitem}{1972-289.5}{E}{Equivalence of the two J-triviality criteria}
\GSfield{Kind; source}{claim; \GSpage{P000481}}
\GSfield{The notes}{A second definition asks for \(T(\xi)\to S^N\) restricting to a homotopy equivalence on the fibre sphere; equivalence is called Dold's theorem and an exercise.}
\GSfield{Statement}{For a stable bundle over a connected finite CW complex, proper fibrewise triviality and the degree-one Thom coreduction criterion are equivalent after stabilization. Dold's fibrewise equivalence criterion supplies the step from a fibrewise map inducing homotopy equivalences to a fibre homotopy equivalence; Atiyah's stable cofibration argument supplies the Thom-space step.}
\GSfield{Reconstruction}{This separates two ingredients hidden in the attribution. The exact cited Dold result is 1955, Corollary~2, identified in Atiyah's Proposition~2.8 proof; it is not the unrelated Dold--Thom infinite-symmetric-product theorem.}
\GSfield{Fate}{\citet[Corollary~2, p.121]{dold1955fibre} and \citet[Proposition~2.8]{atiyah1961thom}.}
\GSfield{Present status}{As of 9 November 2026: proved before the lecture; not an unfinished research exercise.}
\GSfield{Assessment}{literal statement: correct with missing hypotheses; publication match: exact; verification: read the result; development: none; confidence: high}
\end{GSitem}

\GStalk{290}{1973-03-05}{V. I. Lomonosov}{P000482-P000483}{A compact operator, a fixed point, and a common invariant subspace}
The notebook gives a proof for a commuting family containing a nonzero compact operator: a finite cover of a compact image produces a convex-combination map, Schauder gives a fixed point, and compact spectral theory gives the contradiction to transitivity. It then states a locally convex extension and a corollary for multiplication operators on commutative Banach algebras. The page itself replaces a crossed-out attribution to Berezin by Lomonosov. The published note was received on 5 February 1973, before this seminar, and appeared later in 1973; its theorem is stronger than the commuting-family formulation recorded here \citep[pp.~55--56]{lomonosov1973invariant}. Henkin's announcement follows on the lower part of P000483, before Berezin on P000484.

\begin{GSitem}{1972-290.1}{A}{Common invariant subspace for a commuting compact-containing family}
\GSfield{Kind; source}{theorem; \GSpage{P000482}}
\GSfield{The notes}{\ru{сем-во попарно комм.} operators on a Banach space contains a completely continuous \(A\ne0\); the family has a common nontrivial invariant subspace. The notes call this a generalization of Aronszajn--Smith.}
\GSfield{Statement}{Let \(X\) be an infinite-dimensional complex Banach space and \(\mathcal R\subset\mathcal B(X)\) a pairwise commuting family containing a nonzero compact operator. There is a closed linear subspace \(M\), \(0\ne M\ne X\), invariant under every member of \(\mathcal R\). The published theorem gives a subspace invariant under the entire commutant of any nonzero compact operator; pairwise commutativity of that commutant is unnecessary.}
\GSfield{Reconstruction}{Complex scalars, infinite dimension, boundedness and closedness are made explicit; the stronger published theorem is not attributed to the notebook.}
\GSfield{Fate}{The exact proof and a stronger theorem occur in \citet[theorem, p.~55]{lomonosov1973invariant}; received 5 February, published 1973. This develops the compact-operator theme indexed earlier as 1964-50.3.}
\GSfield{Present status}{As of 9 November 2026: Established. Neither the unrestricted Banach-space invariant-subspace problem nor the Hilbert-space problem is a question posed by this entry.}
\GSfield{Assessment}{literal statement: correct with missing hypotheses; publication match: exact; verification: read the result; development: none; confidence: high}
\end{GSitem}
\begin{GSitem}{1972-290.2}{A}{Extension of the argument to locally convex spaces}
\GSfield{Kind; source}{claim; \GSpage{P000482}}
\GSfield{The notes}{\ru{Док-во проходит в случае ЛВП, короткое и простое.} No completeness or compactness convention is specified.}
\GSfield{Statement}{A sufficient precise version replaces \(X\) by an infinite-dimensional complete Hausdorff complex locally convex space and uses a continuous linear compact operator \(K\ne0\), where compact means that some zero-neighbourhood has relatively compact image. A commuting family containing \(K\) has a common nonzero proper closed invariant subspace. One must use the locally convex fixed-point theorem and this neighbourhood notion of compactness, rather than merely mapping bounded sets to relatively compact sets.}
\GSfield{Reconstruction}{The complete Hausdorff formulation and the neighbourhood definition are a sufficient editorial formulation; no claim is made for every possible meaning of compact on an arbitrary locally convex space.}
\GSfield{Fate}{The extension to real or complex locally convex spaces is explicitly stated in \citet[p.~56, paragraph after the proof]{lomonosov1973invariant}. The publication also leaves the precise topological conventions implicit.}
\GSfield{Present status}{As of 9 November 2026: The compact-operator fixed-point argument has this precise sufficient extension. The unqualified notebook sentence requires these conventions.}
\GSfield{Assessment}{literal statement: correct with missing hypotheses; publication match: partial; verification: read the result; development: none; confidence: medium}
\end{GSitem}
\begin{GSitem}{1972-290.3}{A}{A compact image separated from zero and the transitivity alternative}
\GSfield{Kind; source}{method; \GSpage{P000482}}
\GSfield{The notes}{Choose a closed unit ball \(S\) about \(x_0\) with \(\inf_{x\in S}\|Ax\|>1/2\), then use \(\overline{AS}\). The displayed recipe reads \(x_0=\|A\|x'_0\), \(\|x'_0\|=1\), \(Ax'_0\ne0\). Orbit closures \(\overline{\mathcal R y}\) are invariant.}
\GSfield{Statement}{For nonzero bounded \(A\), take a unit vector \(u\) with \(Au\ne0\) and \(t>(\|A\|+1/2)/\|Au\|\). Then the closed unit ball about \(tu\) satisfies \(\|Ax\|\ge t\|Au\|-\|A\|>1/2\). Compactness makes \(K=\overline{AS}\) compact and disjoint from zero. For a unital operator algebra, if every nonzero orbit is dense, then each \(y\in K\) can be carried into the open unit ball about \(x_0\).}
\GSfield{Reconstruction}{\ours\ The explicit valid choice of t repairs the notebook formula. Reversing the triangle inequality independently verifies the separation. It is an elementary editorial check, not a new invariant-subspace theorem.}
\GSfield{Fate}{This is the opening reduction of \citet[p.~55]{lomonosov1973invariant}; the publication asserts existence of the ball without the notebook's incorrect scaling recipe.}
\GSfield{Present status}{As of 9 November 2026: The reduction is valid after correcting the scale. For a small-norm nonzero A, the printed recipe places zero in S and cannot work.}
\GSfield{Assessment}{literal statement: false as written; publication match: exact; verification: read the result; development: annotated; confidence: high}
\end{GSitem}
\begin{GSitem}{1972-290.4}{A}{Finite covering, convex weights, and the Schauder fixed point}
\GSfield{Kind; source}{method; \GSpage{P000482-P000483}}
\GSfield{The notes}{Finitely many \(T_i\) cover \(\overline{AS}\) by \(\|T_i y-x_0\|<1\). With the drawn cutoff \(f(t)=\max(1-t,0)\), form \(\Phi(y)=\sum_i\alpha_i(y)T_i y\), and apply Schauder to \(\Phi A:S\to S\).}
\GSfield{Statement}{Set \(\alpha_i(y)=f(\|T_i y-x_0\|)/\sum_j f(\|T_j y-x_0\|)\). The denominator is positive on \(K\); the weights are continuous, nonnegative and sum to one. Only vectors in the open ball have positive weights, so \(\Phi(K)\subset S\). The map \(\Phi A\) is continuous with relatively compact range, hence has a fixed point \(x\in S\), and \(x\ne0\) because \(0\notin S\).}
\GSfield{Reconstruction}{The value f(0)=1 and the positivity of the denominator are supplied explicitly from the graph; the missing parenthesis in one weight is restored.}
\GSfield{Fate}{The construction is reproduced in \citet[pp.~55--56, displayed fixed-point equation]{lomonosov1973invariant}. Its chronology is the same as item 1972-290.1.}
\GSfield{Present status}{As of 9 November 2026: Established proof mechanism. The dependence of the weights on y is frozen at y=Ax only after obtaining the fixed point.}
\GSfield{Assessment}{literal statement: correct; publication match: exact; verification: read the result; development: none; confidence: high}
\end{GSitem}
\begin{GSitem}{1972-290.5}{A}{The spectral contradiction does not promise a finite-dimensional invariant subspace}
\GSfield{Kind; source}{claim; \GSpage{P000483}}
\GSfield{The notes}{The compact operator \(T=\sum_i\alpha_iT_iA\) has eigenvalue 1; the notes infer an invariant subspace and immediately caution that a finite-dimensional one need not exist in general.}
\GSfield{Statement}{Under the temporary assumption of transitivity and for the commuting algebra in the notebook, \(E=\ker(I-T)\) is nonzero and finite-dimensional, and commutation makes E invariant for that algebra. Since X is infinite-dimensional, E is proper, contradicting transitivity. The finite-dimensional E belongs to the contradiction argument. The final theorem only asserts some nonzero proper closed invariant subspace. An injective quasinilpotent compact operator, such as the Volterra operator on \(L^2[0,1]\), cannot have a nonzero finite-dimensional invariant subspace over \(\mathbb C\).}
\GSfield{Reconstruction}{\ours\ The Volterra example is an explanatory check, not a claim that the speaker gave it: injectivity excludes eigenvalue zero, and quasinilpotence excludes every other eigenvalue; a nonzero finite-dimensional complex invariant subspace would supply an eigenvector. The distinction between the commuting notebook proof and the stronger published proof is explicit.}
\GSfield{Fate}{Compare the contradiction at the end of \citet[p.~56]{lomonosov1973invariant}; the published noncommutative-commutant proof instead uses E to obtain an eigenvector of A.}
\GSfield{Present status}{As of 9 November 2026: The qualification is essential. Compactness of A does not imply a finite-dimensional common invariant subspace.}
\GSfield{Assessment}{literal statement: correct; publication match: exact; verification: read the result; development: annotated; confidence: high}
\end{GSitem}
\begin{GSitem}{1972-290.6}{A}{A nontrivial closed ideal from compact multiplication}
\GSfield{Kind; source}{theorem; \GSpage{P000483}}
\GSfield{The notes}{A commutative Banach algebra containing an element multiplication by which is completely continuous has a nontrivial ideal.}
\GSfield{Statement}{Let B be an infinite-dimensional commutative complex Banach algebra. If some left multiplication \(L_a:b\mapsto ab\) is nonzero and compact, then the commuting family \(\{L_b:b\in B\}\) has a nonzero proper closed invariant subspace I. Commutativity makes I a closed two-sided ideal. The zero multiplier must be excluded from this direct application, and dimension one is an immediate exception to the unqualified statement.}
\GSfield{Reconstruction}{Infinite dimension and a nonzero multiplier are sufficient hypotheses for the stated argument, not a claim that the notebook gave them.}
\GSfield{Fate}{This is the direct multiplication-representation consequence of \citet[p.~55, theorem]{lomonosov1973invariant}. The separate algebraic corollary is recorded in the seminar rather than stated in that two-page article.}
\GSfield{Present status}{As of 9 November 2026: Established under the stated hypotheses; the original unqualified formulation omits necessary exclusions.}
\GSfield{Assessment}{literal statement: correct with missing hypotheses; publication match: partial; verification: read the result; development: none; confidence: high}
\end{GSitem}

\GStalk{292}{1973-03-05}{G. M. Henkin}{P000483}{Leray kernels, extension estimates, and convex embeddings}
A short announcement connects boundary integral formulas, compact localization errors, realization of strictly pseudoconvex domains inside convex domains, bounded holomorphic extension, and the maximal ideal space of the continuous holomorphic algebra. The extension theorem had already appeared in 1972. The convex embedding assertion anticipates the Henkin--Chirka and Fornæss publications; it is substantially stronger than the words ``follows from Oka'' suggest.

\begin{GSitem}{1972-292.1}{E}{The Leray integral representation on a convex domain}
\GSfield{Kind; source}{construction; \GSpage{P000483}}
\GSfield{The notes}{The displayed boundary integral has denominator \(\langle\partial\rho(\zeta),\zeta-z\rangle^n\) and coefficient \((n-1)!/(2\pi i)^n\). It is described for bounded holomorphic functions.}
\GSfield{Statement}{A sufficient classical setting is a bounded smoothly strictly convex domain \(D=\{\rho<0\}\subset\mathbb C^n\), with \(d\rho\ne0\) on the boundary and \(f\in\mathcal O(D)\cap C(\overline D)\). With the boundary orientation and the standard alternating form \(\omega'(w)=\sum_j(-1)^{j-1}w_j\,dw_1\wedge\cdots\widehat{dw_j}\cdots\wedge dw_n\), the notebook formula reproduces f inside D. An arbitrary bounded holomorphic f needs a boundary-limit or limiting-domain interpretation.}
\GSfield{Reconstruction}{Smooth strict convexity and continuous boundary values are sufficient hypotheses supplied for an unambiguous classical integral.}
\GSfield{Fate}{This is established integral-representation machinery predating the announcement; compare \citet[\S1, Theorem 1.1 and Lemma 2.1]{henkin1969integral}. Strict pseudoconvexity requires a suitably constructed supporting kernel, rather than the unmodified convex-domain denominator.}
\GSfield{Present status}{As of 9 November 2026: Established; no new existence problem is posed by this formula.}
\GSfield{Assessment}{literal statement: correct with missing hypotheses; publication match: exact; verification: read the result; development: none; confidence: high}
\end{GSitem}
\begin{GSitem}{1972-292.2}{E}{A compact localization error with a half-Hölder gain}
\GSfield{Kind; source}{claim; \GSpage{P000483}}
\GSfield{The notes}{The operator \(f\mapsto f g-f_g\) is said to be completely continuous and to satisfy a Hölder condition of exponent \(1/2\), but the displayed definition of \(f_g\) does not show its dependence on f.}
\GSfield{Statement}{A precise sufficient interpretation uses a smooth bounded strictly pseudoconvex domain, \(f\in A(D)=\mathcal O(D)\cap C(\overline D)\), and a fixed smooth cutoff g. Apply a bounded \(\overline\partial\)-solution operator T with a half-Hölder estimate to \(f\overline\partial g\), and put \(f_g=fg-T(f\overline\partial g)\). Then \(f_g\) is holomorphic and \(f\mapsto fg-f_g\) maps the unit ball of A(D) into a bounded subset of \(C^{0,1/2}(\overline D)\), hence is compact into \(C(\overline D)\).}
\GSfield{Reconstruction}{The definition through a solution operator is editorial. The notes do not specify the cutoff regularity, source norm, or whether \(f_g\) denotes a projection of \(fg\).}
\GSfield{Fate}{The localization and uniform-estimate method is already present in \citet[Theorem 1.2]{henkin1969integral} and \citet[\S1]{henkin1970integral}. The precise abbreviated operator in the notebook is only partially matched.}
\GSfield{Present status}{As of 9 November 2026: The reconstructed compactness statement is established; identifying the exact boundary operator from the abbreviation alone is not possible.}
\GSfield{Assessment}{literal statement: correct with missing hypotheses; publication match: partial; verification: read the result; development: none; confidence: medium}
\end{GSitem}
\begin{GSitem}{1972-292.3}{A}{A strictly pseudoconvex domain as a convex-domain section}
\GSfield{Kind; source}{theorem; \GSpage{P000483}}
\GSfield{The notes}{A domain with strictly plurisubharmonic defining function is a closed submanifold of a strictly convex domain and meets its boundary transversely; the parenthesis reads \ru{следует из теор. Ока}.}
\GSfield{Statement}{Let \(D\Subset\mathbb C^n\) have smooth strictly pseudoconvex boundary. There are a Stein neighbourhood U of \(\overline D\), chosen with \(\overline D\) holomorphically convex in U, a proper holomorphic embedding \(F:U\to\mathbb C^N\), and a bounded smoothly strongly convex domain \(\Omega\) such that \(F(D)=F(U)\cap\Omega\) is a closed complex submanifold of \(\Omega\), and \(F(U)\) meets \(\partial\Omega\) transversely. The target dimension may exceed n.}
\GSfield{Reconstruction}{Embedding, enlarged ambient dimension, analytic submanifold, regular boundary, and transversality are made explicit. The notebook does not contain an Oka-theoretic proof.}
\GSfield{Fate}{The precise result appears in \citet[Proposition 14, Russian p.~111]{henkinchirka1975boundary} (1975; English 1976, p.~668); the modern formulation is \citet[Theorem 1]{forstneric2026embedding}. Fornæss obtained the independent 1976 realization. Class A records Henkin's coauthorship of the subsequent account.}
\GSfield{Present status}{As of 9 November 2026: Established convex embedding theorem; a generic closed subset of a convex domain would not suffice.}
\GSfield{Assessment}{literal statement: correct with missing hypotheses; publication match: exact; verification: read the result; development: none; confidence: high}
\end{GSitem}
\begin{GSitem}{1972-292.4}{E}{Bounded extension from a transverse analytic submanifold}
\GSfield{Kind; source}{theorem; \GSpage{P000483}}
\GSfield{The notes}{In point 2, \ru{выпис. продолж. с подмн-зия с оценками!} in the class of bounded functions.}
\GSfield{Statement}{Let D be a bounded smooth strictly pseudoconvex domain and let M be the intersection of D with a closed complex submanifold defined in a neighbourhood of \(\overline D\), transverse to \(\partial D\). There is a bounded linear right inverse \(E:H^\infty(M)\to H^\infty(D)\) of restriction, with \(\|Ef\|_\infty\le C\|f\|_\infty\). It can be chosen to take \(A(M)\) into \(A(D)\).}
\GSfield{Fate}{This is the main theorem of \citet[pp.~536--537]{henkin1972continuation}, received 8 December 1971 and published in the 1972 volume. Thus the imperative in the notebook records an announced construction or exposition, not an unaddressed 1973 open problem.}
\GSfield{Present status}{As of 9 November 2026: Established under the smoothness and transversality hypotheses; these must remain attached to the estimate.}
\GSfield{Assessment}{literal statement: correct with missing hypotheses; publication match: exact; verification: read the result; development: none; confidence: high}
\end{GSitem}
\begin{GSitem}{1972-292.5}{E}{The converse geometric implication}
\GSfield{Kind; source}{claim; \GSpage{P000483}}
\GSfield{The notes}{The converse is said to hold: such a submanifold is a bounded domain of holomorphy.}
\GSfield{Statement}{A closed complex submanifold of a Stein convex domain is Stein. If it extends smoothly across the ambient strongly convex boundary and is transverse to it, the restricted defining function is strictly plurisubharmonic and defines a strongly pseudoconvex domain in the submanifold. Calling it a domain in \(\mathbb C^n\) additionally presupposes the given parametrization or embedding of D.}
\GSfield{Reconstruction}{The word analytic and the intrinsic meaning of domain of holomorphy are indispensable; an arbitrary smooth or closed subset is insufficient.}
\GSfield{Fate}{This is the classical Stein/submanifold direction used in the setting of \citet{henkin1972continuation}; it does not require the later difficult convex embedding theorem in the other direction.}
\GSfield{Present status}{As of 9 November 2026: Established geometric implication, distinct from the forward embedding theorem.}
\GSfield{Assessment}{literal statement: correct with missing hypotheses; publication match: partial; verification: read the result; development: none; confidence: high}
\end{GSitem}
\begin{GSitem}{1972-292.6}{E}{The maximal ideal space of A(D)}
\GSfield{Kind; source}{theorem; \GSpage{P000483}}
\GSfield{The notes}{For functions analytic inside and continuous on the closure, the maximal ideal space equals the closure of the domain.}
\GSfield{Statement}{For a bounded smoothly strictly pseudoconvex domain \(D\subset\mathbb C^n\), every nonzero multiplicative continuous functional on the unital uniform algebra \(A(D)\) is evaluation at one point of \(\overline D\); the resulting identification of maximal ideal space with \(\overline D\) is a homeomorphism.}
\GSfield{Fate}{Already explicitly recorded in the introduction to \citet[p.~219]{ovrelid1971generators}, using uniform approximation by functions holomorphic near the closure. It is not the unrestricted corona problem for \(H^\infty(D)\).}
\GSfield{Present status}{As of 9 November 2026: Established for the continuous holomorphic algebra under the stated domain hypotheses.}
\GSfield{Assessment}{literal statement: correct with missing hypotheses; publication match: exact; verification: read the result; development: none; confidence: high}
\end{GSitem}

\GStalk{291}{1973-03-05}{F. A. Berezin}{P000484-P000487}{Quantization, coherent-state symbols, and Kähler geometry}
The talk passes from a correspondence-principle definition of quantization to the Fock model, covariant and contravariant symbols, trace inequalities, and constructions from Bergman kernels. The compact symmetric and disk examples close the record. The symbol calculus is already in the 1972 paper; the quantization programme was announced in Doklady in 1973 (received 14 December 1972), expanded in 1974, and applied to complex symmetric spaces in 1975. The notebook contains several consequential sign, conjugation, and measure-normalization errors. Its formal and convergent analytic meanings must be separated.

\begin{GSitem}{1972-291.1}{E}{The symplectic Poisson bracket}
\GSfield{Kind; source}{definition; \GSpage{P000484}}
\GSfield{The notes}{For a symplectic manifold M with matrix \(\omega_{ik}\) and inverse \(\omega^{ik}\), \([f,g]=\sum\omega^{ik}\partial_i f\partial_k g\).}
\GSfield{Statement}{On a smooth manifold equipped with a closed nondegenerate two-form \(\omega\), its inverse bivector defines \(\{f,g\}=\omega^{-1}(df,dg)\), with the sign fixed by that inverse convention. Closedness is the condition making the Jacobi identity hold. The notebook's complex-coordinate form suppresses the factor i needed for a real Kähler form.}
\GSfield{Fate}{Classical symplectic geometry, restated as the starting data in \citet[\S1]{berezin1974quantization}.}
\GSfield{Present status}{As of 9 November 2026: Established definition; a nondegenerate but nonclosed two-form is not sufficient.}
\GSfield{Assessment}{literal statement: correct with missing hypotheses; publication match: exact; verification: read the result; development: none; confidence: high}
\end{GSitem}
\begin{GSitem}{1972-291.2}{A}{Associative quantization and the correspondence principle}
\GSfield{Kind; source}{definition; \GSpage{P000484}}
\GSfield{The notes}{A family of associative rings \(A_h\) of functions has \(f*_0g=fg\) and \(\partial_h(f*_hg-g*_hf)|_{h=0}=\{f,g\}\).}
\GSfield{Statement}{A formal interpretation specifies an associative \(\mathbb C[[h]]\)-bilinear product \(f*g=fg+hC_1(f,g)+\cdots\) on a chosen function algebra, with \(C_1(f,g)-C_1(g,f)=\{f,g\}\). An analytic interpretation instead specifies admissible real h, symbol domains and a topology for the limit. Factors of i can be absorbed into h; a real Hermitian convention usually has \(ih\{f,g\}\). These are different existence problems.}
\GSfield{Reconstruction}{Formal series are supplied as one exact interpretation, not substituted silently for Berezin’s analytic family.}
\GSfield{Fate}{The quantization definition and examples were announced in \citet{berezin1973bounded} and developed in \citet[\S1]{berezin1974quantization}.}
\GSfield{Present status}{As of 9 November 2026: Established framework. The displayed first-order condition alone neither ensures associativity at higher orders nor selects a unique quantization.}
\GSfield{Assessment}{literal statement: correct with missing hypotheses; publication match: exact; verification: read the result; development: none; confidence: high}
\end{GSitem}
\begin{GSitem}{1972-291.3}{A}{The optional constant unit}
\GSfield{Kind; source}{definition; \GSpage{P000484}}
\GSfield{The notes}{Requirement 3: every \(A_h\) has a unit corresponding to \(f_0\equiv1\).}
\GSfield{Statement}{The constant symbol 1 is required to satisfy \(1*_h f=f*_h1=f\) on the specified symbol algebra for every admissible h. Having some nonconstant unit is weaker than this normalization.}
\GSfield{Fate}{Explicit quantization axiom in \citet[\S1]{berezin1974quantization}; the later question on P000486 concerns precisely this extra requirement.}
\GSfield{Present status}{As of 9 November 2026: Well-defined normalization; see problem GP-1972-401 below.}
\GSfield{Assessment}{literal statement: correct; publication match: exact; verification: read the result; development: none; confidence: high}
\end{GSitem}
\begin{GSitem}{1972-291.4}{A}{Trace and symplectic volume}
\GSfield{Kind; source}{definition; \GSpage{P000484}}
\GSfield{The notes}{Requirement 4: a trace on a distinguished set, with \(\operatorname{Sp}(f*g)=\operatorname{Sp}(g*f)\) and \(\operatorname{Sp}f=c\int f\,d\mu\).}
\GSfield{Statement}{For \(\dim_{\mathbb R}M=2n\), the natural volume is \(\omega^n/n!\). A trace is a linear functional on a specified ideal of integrable or trace-class symbols, cyclic whenever both products lie in its domain. The scalar c may depend on h. It is not asserted that every bounded operator or arbitrary symbol has finite trace.}
\GSfield{Fate}{The trace axiom is formulated in \citet[\S1]{berezin1974quantization}; the actual operator-symbol trace identities already appear in \citet[\S1, formulas (1.13)--(1.16)]{berezin1972symbols}.}
\GSfield{Present status}{As of 9 November 2026: Established in the models with the indicated trace ideal; convergence is part of the assertion.}
\GSfield{Assessment}{literal statement: correct with missing hypotheses; publication match: exact; verification: read the result; development: none; confidence: high}
\end{GSitem}
\begin{GSitem}{1972-291.5}{A}{Naturality relative to a chosen category}
\GSfield{Kind; source}{remark; \GSpage{P000484}}
\GSfield{The notes}{Quantization should be functorial for canonical diffeomorphisms, with the qualification ``some fixed category''. Two Kähler classes are named.}
\GSfield{Statement}{For a specified group G of symplectic symmetries, one may require pullback to act by algebra automorphisms of each \(A_h\), compatibly with composition. This covariance requirement is precise once G, the symbol spaces, and equivalence of quantizations are fixed. The notebook does not claim a functor on all symplectic manifolds and all canonical transformations.}
\GSfield{Fate}{This relative naturality is built into \citet[\S1 and \S3]{berezin1974quantization}. Earlier prequantization items 1969-178.3--4 and 1971-229.4 concern different constructions.}
\GSfield{Present status}{As of 9 November 2026: A meaningful specified-category axiom; unrestricted functorial quantization is not established by this entry.}
\GSfield{Assessment}{literal statement: correct with missing hypotheses; publication match: partial; verification: read the result; development: none; confidence: high}
\end{GSitem}
\begin{GSitem}{1972-291.6}{E}{The Gaussian Fock product and its domain}
\GSfield{Kind; source}{construction; \GSpage{P000484}}
\GSfield{The notes}{An integral of \(f_1(z,\bar v)f_2(v,\bar z)e^{-|z-v|^2/h}\) defines the product on entire real-analytic functions.}
\GSfield{Statement}{For h>0 use \(d\nu_h(v)=(\pi h)^{-n}d^{2n}v\). On polynomial symbols the formula \((f*g)(z,\bar z)=\int f(z,\bar v)g(v,\bar z)e^{-|z-v|^2/h}d\nu_h(v)\) converges and is associative; it equals \(\sum_\alpha h^{|\alpha|}(\partial_{\bar z}^\alpha f)(\partial_z^\alpha g)/\alpha!\). Larger analytic classes require growth estimates. Entire real analyticity alone does not ensure integrability against a Gaussian.}
\GSfield{Reconstruction}{Polynomial symbols give a sufficient closed algebra. The measure convention removes suppressed powers of pi; unrestricted entire functions are not admitted.}
\GSfield{Fate}{The Fock symbol calculus predates the seminar and is treated in \citet{berezin1972symbols}; see also \citet[\S4]{berezin1974quantization}.}
\GSfield{Present status}{As of 9 November 2026: Established on specified symbol classes, with the displayed bidifferential formula also meaningful formally.}
\GSfield{Assessment}{literal statement: false as written; publication match: exact; verification: read the result; development: none; confidence: high}
\end{GSitem}
\begin{GSitem}{1972-291.7}{E}{The Fock representation must be complex linear}
\GSfield{Kind; source}{construction; \GSpage{P000484}}
\GSfield{The notes}{The representation formula writes \(f(z,v)\overline{\varphi(v)}\), while its next anti-Wick formula has \(\varphi(v)\).}
\GSfield{Statement}{On \(\mathcal F_h=\{\varphi\in\mathcal O(\mathbb C^n):\int|\varphi|^2e^{-|z|^2/h}d\nu_h<\infty\}\), the covariant-symbol kernel acts as \((A_f\varphi)(z)=\int f(z,\bar v)\varphi(v)e^{z\bar v/h-|v|^2/h}d\nu_h(v)\). Domain conditions are needed for unbounded polynomial operators. Conjugating \(\varphi\) would instead make the map antilinear.}
\GSfield{Reconstruction}{The bar is moved from the input vector to the second kernel variable as a mathematical emendation.}
\GSfield{Fate}{The corrected kernel is the reproducing-kernel realization in \citet{berezin1972symbols,berezin1974quantization}. The scan confirms that the misplaced bar is in the manuscript.}
\GSfield{Present status}{As of 9 November 2026: The corrected formula is standard; the literal displayed map is not the claimed linear representation.}
\GSfield{Assessment}{literal statement: false as written; publication match: exact; verification: read the result; development: annotated; confidence: high}
\end{GSitem}
\begin{GSitem}{1972-291.8}{E}{Creation, annihilation, and the two orderings}
\GSfield{Kind; source}{construction; \GSpage{P000484}}
\GSfield{The notes}{The notes give \(a=\partial_z\), \(a^*=z\), normal and antinormal orderings, and symbolic monomials with opposite placements of bars.}
\GSfield{Statement}{On the polynomial core of \(\mathcal F_h\), put \(a=h\partial_z\) and \(a^*=z\); these are Hilbert adjoints after closure and \([a_j,a_k^*]=h\delta_{jk}\). The normally ordered \((a^*)^m a^n\) has covariant symbol \(z^m\bar z^n\). A contravariant monomial is quantized in the opposite order \(a^n(a^*)^m\). The derivative without h is the h=1 convention.}
\GSfield{Reconstruction}{The stated h-dependent Gaussian determines the adjoint normalization; the notation a* cannot simultaneously mean its Hilbert adjoint with the unscaled derivative.}
\GSfield{Fate}{These are the Wick/anti-Wick realizations developed before the talk in \citet{berezin1972symbols}.}
\GSfield{Present status}{As of 9 November 2026: Established after fixing h and the symbol orientation.}
\GSfield{Assessment}{literal statement: false as written; publication match: exact; verification: read the result; development: none; confidence: high}
\end{GSitem}
\begin{GSitem}{1972-291.9}{E}{Coherent states and covariant symbols}
\GSfield{Kind; source}{definition; \GSpage{P000485}}
\GSfield{The notes}{\(e_{\bar v}(z)=e^{z\bar v/h}\); the quotient of matrix coefficients equals \(\operatorname{Sp}(AP_v)\).}
\GSfield{Statement}{Write \(e_v(z)=e^{z\bar v/h}\), \(k_v=e_v/\|e_v\|\), and \(P_v=|k_v\rangle\langle k_v|\). For a bounded operator A its covariant symbol is \(\sigma(A)(v)=\langle Ak_v,k_v\rangle=\operatorname{Tr}(AP_v)\). The same vector must occur in both slots and the denominator; the notebook alternates v and \(\bar v\) in the labels.}
\GSfield{Fate}{This is the definition in \citet[\S1, formula (1.2)]{berezin1972symbols}.}
\GSfield{Present status}{As of 9 November 2026: Established. Off-diagonal polarization of the symbol is an additional analytic step.}
\GSfield{Assessment}{literal statement: correct with missing hypotheses; publication match: exact; verification: read the result; development: none; confidence: high}
\end{GSitem}
\begin{GSitem}{1972-291.10}{E}{Boundedness duality for symbols}
\GSfield{Kind; source}{theorem; \GSpage{P000485}}
\GSfield{The notes}{Bounded A implies bounded Wick symbol; bounded anti-Wick symbol implies bounded A.}
\GSfield{Statement}{For normalized coherent states, \(\|\sigma(A)\|_\infty\le\|A\|\). For a normalized resolution \(\int P_vd\nu_h(v)=I\), the weak integral \(T_a=\int a(v)P_vd\nu_h(v)\) satisfies \(\|T_a\|\le\|a\|_\infty\). The two arrows concern different symbol maps and are not inverse norm equivalences.}
\GSfield{Fate}{These estimates are in \citet[\S1, formula (1.14) and the compression construction]{berezin1972symbols}.}
\GSfield{Present status}{As of 9 November 2026: Established; a bounded covariant symbol alone is not the stated boundedness criterion.}
\GSfield{Assessment}{literal statement: correct; publication match: exact; verification: read the result; development: none; confidence: high}
\end{GSitem}
\begin{GSitem}{1972-291.11}{E}{Compactness and decay in the Fock model}
\GSfield{Kind; source}{theorem; \GSpage{P000485}}
\GSfield{The notes}{Compact A implies f tends to zero; a contravariant symbol tending to zero gives a compact operator.}
\GSfield{Statement}{In the Fock model, normalized \(k_v\) tends weakly to zero as \(|v|\to\infty\), so compact A has \(\sigma(A)(v)\to0\). If \(a\in C_0(\mathbb C^n)\), then \(T_a\) is compact: approximate a uniformly by compactly supported symbols, whose coherent-projector integrals are trace class. The assertion needs this geometry and cannot be carried unchanged to an arbitrary frame.}
\GSfield{Reconstruction}{\ours\ The two-sentence compactness verification is editorial; no equivalence with arbitrary vanishing covariant symbols is asserted.}
\GSfield{Fate}{This is the symbol-calculus compactness mechanism associated with \citet{berezin1972symbols}; the stated sufficient Fock hypotheses explain the shorthand.}
\GSfield{Present status}{As of 9 November 2026: Established in this setting; ``all theorems hold later on the page requires additional assumptions for this one.}
\GSfield{Assessment}{literal statement: correct with missing hypotheses; publication match: partial; verification: read the result; development: annotated; confidence: high}
\end{GSitem}
\begin{GSitem}{1972-291.12}{E}{Trace class and the two trace integrals}
\GSfield{Kind; source}{theorem; \GSpage{P000485}}
\GSfield{The notes}{``Nuclearity (analog.)'' followed by equal integrals of Wick and anti-Wick symbols.}
\GSfield{Statement}{If A is trace class, its covariant symbol is integrable against the normalized frame measure and \(\operatorname{Tr}A=\int\sigma(A)d\nu_h\). If \(a\in L^1(d\nu_h)\), then \(T_a\) is trace class, \(\|T_a\|_1\le\|a\|_1\), and \(\operatorname{Tr}T_a=\int a\,d\nu_h\). These sufficient statements do not say that integrability of an arbitrary covariant symbol implies trace class.}
\GSfield{Fate}{Read in \citet[\S1, formulas (1.13)--(1.16)]{berezin1972symbols}.}
\GSfield{Present status}{As of 9 November 2026: Established with trace-class and integrability assumptions made explicit.}
\GSfield{Assessment}{literal statement: correct with missing hypotheses; publication match: exact; verification: read the result; development: none; confidence: high}
\end{GSitem}
\begin{GSitem}{1972-291.13}{E}{Gaussian conversion and the Weyl midpoint}
\GSfield{Kind; source}{construction; \GSpage{P000485}}
\GSfield{The notes}{The Wick symbol is the Gaussian convolution of the anti-Wick symbol; h/2 gives the Weyl symbol.}
\GSfield{Statement}{For a Schwartz contravariant symbol a in the flat model, \(\sigma(T_a)=B_ha\), where \(B_ha(z)=\int a(v)e^{-|z-v|^2/h}d\nu_h(v)=e^{(h/4)\Delta_{\mathbb R^{2n}}}a(z)\). In the consistent Weyl convention the Weyl symbol of the same operator is \(e^{(h/8)\Delta}a\). These constants depend on the chosen real coordinates and h convention.}
\GSfield{Fate}{The conversion formula is explicit in \citet[\S1, formula (1.12)]{berezin1972symbols}; the flat quantizations are compared in \citet[\S4 and \S7]{berezin1974quantization}.}
\GSfield{Present status}{As of 9 November 2026: Established Gaussian/heat relation; inverse heat evolution is not bounded on arbitrary function spaces.}
\GSfield{Assessment}{literal statement: correct with missing hypotheses; publication match: exact; verification: read the result; development: none; confidence: high}
\end{GSitem}
\begin{GSitem}{1972-291.14}{F}{Gelfand’s Poincaré attribution}
\GSfield{Kind; source}{remark; \GSpage{P000485}; I. M. Gelfand}
\GSfield{The notes}{The parenthetical says that these are Poincaré's symbols, without a title, formula, or date.}
\GSfield{Statement}{No uniquely identifiable mathematical attribution follows. It could concern symmetric ordering or an earlier symbolic calculus, but the notebook supplies no discriminating reference.}
\GSfield{Fate}{Searches of the Berezin 1972 and 1974 papers and their bibliographies did not identify the precise Poincaré source intended. The established Weyl conversion in the previous item does not verify this historical attribution.}
\GSfield{Present status}{As of 9 November 2026: Attribution unresolved; no research problem or priority conclusion is inferred.}
\GSfield{Assessment}{literal statement: not determinable; publication match: none found; verification: unresolved lead; development: none; confidence: low}
\end{GSitem}
\begin{GSitem}{1972-291.15}{E}{Continuous Parseval frames and the missing norm weight}
\GSfield{Kind; source}{construction; \GSpage{P000485}}
\GSfield{The notes}{An overcomplete system satisfies \(\|f\|^2=\int|\langle f,e_\alpha\rangle|^2d\alpha\), but the contravariant reconstruction is written \(A=\int\mathring A(\alpha)P_\alpha d\alpha\) with orthogonal projectors.}
\GSfield{Statement}{For a measurable continuous Parseval frame, \(I=\int|e_\alpha\rangle\langle e_\alpha|d\alpha\) weakly. If \(P_\alpha\) is the orthogonal rank-one projector, the appropriate measure is \(d\nu=\|e_\alpha\|^2d\alpha\), omitting zero vectors. Then \(A=\int a(\alpha)P_\alpha d\nu\). In general a contravariant symbol need not exist for every A and need not be unique.}
\GSfield{Fate}{The norm factor is explicitly present in \citet[\S1, formulas (1.1)--(1.3)]{berezin1972symbols}.}
\GSfield{Present status}{As of 9 November 2026: Established frame construction; normalization cannot be dropped when the vectors are not unit vectors.}
\GSfield{Assessment}{literal statement: false as written; publication match: exact; verification: read the result; development: none; confidence: high}
\end{GSitem}
\begin{GSitem}{1972-291.16}{E}{Compression and the Fourier example}
\GSfield{Kind; source}{construction; \GSpage{P000485-P000486}}
\GSfield{The notes}{Embed H into \(L^2(M)\), multiply by the contravariant symbol, and compress by the orthogonal projection. The next page names \(e^{i(p,x)}\) on bounded domains.}
\GSfield{Statement}{For the frame analysis isometry \(V:H\to L^2(M,d\alpha)\), \(T_a=V^*M_aV\); on VH this is \(PM_aP\). For a bounded measurable \(\Omega\subset\mathbb R^n\), the restricted exponentials \((2\pi)^{-n/2}e^{ip\cdot x}\) form a continuous Parseval frame in \(L^2(\Omega)\) by Plancherel. Their normalization determines the projector measure.}
\GSfield{Fate}{This is the multiplication-compression realization in \citet[\S1, formula (1.5)]{berezin1972symbols}; the Fourier example is a standard instance.}
\GSfield{Present status}{As of 9 November 2026: Established; writing PMP as an operator on the full ambient space adds a zero block.}
\GSfield{Assessment}{literal statement: correct; publication match: exact; verification: read the result; development: none; confidence: high}
\end{GSitem}
\begin{GSitem}{1972-291.17}{E}{The convex trace inequalities}
\GSfield{Kind; source}{theorem; \GSpage{P000486}}
\GSfield{The notes}{For a convex function \(\varphi\), \(\int\varphi(A)\le\operatorname{Sp}\varphi(\hat A)\le\int\varphi(\mathring A)\). The lower bound is called Jensen's inequality; the upper one is written as a compression inequality.}
\GSfield{Statement}{For a normalized projector resolution, a self-adjoint A admitting a real contravariant symbol a, and a convex scalar function \(\varphi\), one has \(\int\varphi(\sigma(A))d\nu\le\operatorname{Tr}\varphi(A)\le\int\varphi(a)d\nu\) whenever the terms are finite (or in a consistent extended nonnegative sense). The compression trace is taken on PH; in infinite dimension one cannot include an infinite zero block with \(\varphi(0)\ne0\). Scalar convexity suffices for this trace inequality, not for an operator inequality.}
\GSfield{Fate}{This is \citet[\S2, Theorem 2]{berezin1972symbols}, published before the seminar.}
\GSfield{Present status}{As of 9 November 2026: Established Berezin trace inequalities with the trace domain and frame measure specified.}
\GSfield{Assessment}{literal statement: correct with missing hypotheses; publication match: exact; verification: read the result; development: none; confidence: high}
\end{GSitem}
\begin{GSitem}{1972-291.18}{A}{The Bergman cross-ratio and weighted holomorphic space}
\GSfield{Kind; source}{construction; \GSpage{P000486}}
\GSfield{The notes}{The kernel \(G_h\) has diagonal Bergman factors in the numerator, and the Hilbert-space weight is \(K(z,\bar z)^{+1/h}\).}
\GSfield{Statement}{In a homogeneous bounded domain with admissible h and a consistent analytic power of K, use the weighted space with measure \(K(z,\bar z)^{-1/h}d\mu(z)\), where \(d\mu\) is invariant volume. The normalized overlap in the product is \[\left(\frac{K(z,\bar v)K(v,\bar z)}{K(z,\bar z)K(v,\bar v)}\right)^{1/h},\] with a normalization factor and the invariant integration measure. Cauchy--Schwarz bounds the base of this ratio by 1; its reciprocal has the wrong concentration behaviour.}
\GSfield{Reconstruction}{Both sign changes are mathematical emendations: the scan confirms the positive weight and reciprocal ratio. The unspecified volume density is also supplied.}
\GSfield{Fate}{The corrected ratios and negative weight are in \citet[\S2, especially (2.23), and \S5]{berezin1974quantization}; the short precursor is \citet{berezin1973bounded}.}
\GSfield{Present status}{As of 9 November 2026: Established homogeneous construction. The formulas do not define a quantization on all real-analytic functions of every domain.}
\GSfield{Assessment}{literal statement: false as written; publication match: exact; verification: read the result; development: annotated; confidence: high}
\end{GSitem}
\begin{GSitem}{1972-291.19}{F}{A constant unit outside the homogeneous case}
\GSfield{Kind; source}{question; \GSpage{P000486}}
\GSfield{The notes}{A marginal question asks whether there exists a quantization also satisfying condition 3; an arbitrary scalar \(c_1(h)\) is written in \(c(h)=h^{-n}c_1(h)\).}
\GSfield{Statement}{The supported analytic question is whether the proposed Bergman-kernel integral construction can be given nontrivial symbol algebras and a constant unit on a specified nonhomogeneous domain. The domain class, admissible h, convergence bounds and allowed changes of kernel are not fixed. A separately precise formal variant asks for a unital associative deformation product on a Kähler manifold.}
\GSfield{Reconstruction}{The formal theorem is contextual comparison only. The missing analytic category prevents identification of one precisely formulated question of which that theorem would be a partial solution.}
\GSfield{Fate}{\citet[\S2.8, p.~1133]{berezin1974quantization} explicitly retains an analytic obstruction even to nonzero algebras without his hypothesis A. \citet[\S4, Proposition 4.3]{reshetikhintakhtajan2000deformation} construct a normalized unital formal product on arbitrary Kähler manifolds. That result does not establish convergence of the notebook integral or scalar normalization on every nonhomogeneous domain.}
\GSfield{Present status}{As of 9 November 2026: The formal unital theorem is established, but the intended convergent analytic problem is not identifiable with precise quantifiers. This occurrence is not registered as an open problem.}
\GSfield{Assessment}{literal statement: not determinable; publication match: partial; verification: read the result; development: none; confidence: medium}
\end{GSitem}
\begin{GSitem}{1972-291.20}{A}{Homogeneity fixes the normalization}
\GSfield{Kind; source}{theorem; \GSpage{P000486}}
\GSfield{The notes}{For a homogeneous domain condition 3 holds and fixes \(c_1(h)\); \(e_v(z)=K(z,\bar v)^{1/h}\) gives the coherent system.}
\GSfield{Statement}{For a bounded homogeneous domain and the admissible positive parameters in Berezin's construction, the weighted reproducing kernel is a constant multiple of \(K^{1/h}\). Transitivity makes the diagonal normalization spatially constant. Choosing the reciprocal constant gives the reproducing identity and hence constant symbol 1 as the unit.}
\GSfield{Fate}{This is the homogeneous-domain construction of \citet[\S5]{berezin1974quantization}, announced in \citet{berezin1973bounded}.}
\GSfield{Present status}{As of 9 November 2026: Established for admissible parameters; arbitrary positive h is not part of the statement.}
\GSfield{Assessment}{literal statement: correct with missing hypotheses; publication match: exact; verification: read the result; development: none; confidence: high}
\end{GSitem}
\begin{GSitem}{1972-291.21}{F}{The shorthand uniqueness assertion on symmetric domains}
\GSfield{Kind; source}{claim; \GSpage{P000486}}
\GSfield{The notes}{If the domain is symmetric and its motions act by ring automorphisms, \ru{то только это} (only this).}
\GSfield{Statement}{No fully quantified uniqueness theorem can be recovered from these words: the class of algebras, equivalence, irreducibility and admissible parameters are missing. The later paper formulates exhaustion of irreducible effective w*-quantizations on classical complex symmetric spaces as a conjecture, with the sphere case then proved.}
\GSfield{Fate}{Compare \citet[Introduction, pp.~341--342]{berezin1975symmetric} and the restricted uniqueness discussion in \citet[\S7]{berezin1974quantization}. They prevent treating the notebook's blanket phrase as an established theorem. No exact current problem is extracted from unspecified hypotheses.}
\GSfield{Present status}{As of 9 November 2026: The precise intended uniqueness assertion is not identifiable; covariance alone is not enough to assert it.}
\GSfield{Assessment}{literal statement: not determinable; publication match: partial; verification: read the result; development: none; confidence: medium}
\end{GSitem}
\begin{GSitem}{1972-291.22}{E}{Can the Bergman kernel vanish?}
\GSfield{Kind; source}{question; \GSpage{P000486}}
\GSfield{The notes}{One parenthesis says \(K(z,\bar v)\ne0\); later the general case is called unknown.}
\GSfield{Statement}{For a bounded domain G the off-diagonal Bergman kernel can vanish. Every planar annulus \(r<|z|<1\), \(0<r<1\), has zeros. In contrast the disk and bounded homogeneous complete circular domains have zero-free kernels. Positivity of \(K(z,\bar z)\) on the diagonal is a different assertion.}
\GSfield{Reconstruction}{The negative unrestricted answer is separated from valid zero-freeness in the model homogeneous domains.}
\GSfield{Fate}{The annulus counterexamples are already from 1969 (Rosenthal and Skwarczyński); the precise statements and original references are checked in \citet[\S4, pp.~258--259]{boas2000luqikeng}. The original 1969 full texts were not accessible in this run, so the chronology rests on this author survey.}
\GSfield{Present status}{As of 9 November 2026: The unrestricted question had a negative answer before this seminar. No broader classification problem is silently substituted.}
\GSfield{Assessment}{literal statement: false as written; publication match: exact; verification: indirect; development: none; confidence: high; problem \hyperref[prob:GP-1972-402]{GP-1972-402}}
\end{GSitem}
\begin{GSitem}{1972-291.23}{A}{Global polarization of covariant symbols}
\GSfield{Kind; source}{claim; \GSpage{P000486}}
\GSfield{The notes}{Covariant symbols have analytic continuations to all of \(G\times\bar G\).}
\GSfield{Statement}{For a reproducing-kernel Hilbert space with nowhere-zero kernel \(K_h(z,\bar v)\), the covariant symbol of a bounded operator polarizes as its operator kernel divided by \(K_h\); it is holomorphic in z and antiholomorphic in v on the product. If the denominator has zeros, the quotient may have poles and the blanket global extension no longer follows.}
\GSfield{Fate}{This is the symbol realization under the kernel hypotheses in \citet[\S2]{berezin1974quantization}.}
\GSfield{Present status}{As of 9 November 2026: Established in the zero-free homogeneous setting; related to the preceding kernel-zero question.}
\GSfield{Assessment}{literal statement: correct with missing hypotheses; publication match: exact; verification: read the result; development: none; confidence: high}
\end{GSitem}
\begin{GSitem}{1972-291.24}{A}{Compact Hermitian symmetric quantization}
\GSfield{Kind; source}{construction; \GSpage{P000486-P000487}}
\GSfield{The notes}{A Kähler potential is exponentiated; at \(h=1/n\) a single-valued kernel exists and the Hilbert space is finite-dimensional with computable dimension.}
\GSfield{Statement}{For an integral positive Kähler class on a compact classical Hermitian symmetric manifold, choose its positive holomorphic line bundle L. At integral tensor powers \(L^n\), the quantum space is \(H^0(M,L^n)\), finite-dimensional; local kernels glue as sections with transition factors, not generally as scalar functions on all of \(M\times\bar M\). The potential is local and \(g_{i\bar j}=\partial_i\partial_{\bar j}\Phi\), not the inverse matrix. The relation h=1/n uses a fixed integral normalization.}
\GSfield{Reconstruction}{The line-bundle formulation supplies the indispensable global gluing and integrality conditions.}
\GSfield{Fate}{The compact classical symmetric cases and their parameter/dimension formulas are developed in \citet[compact-space part]{berezin1975symmetric}. The notebook announcement precedes that publication; it does not supply a global scalar potential.}
\GSfield{Present status}{As of 9 November 2026: Established for the specified integral compact symmetric models. Arbitrary Kähler scaling does not automatically permit h=1/n.}
\GSfield{Assessment}{literal statement: false as written; publication match: exact; verification: read the result; development: none; confidence: high}
\end{GSitem}
\begin{GSitem}{1972-291.25}{A}{The disk and its projective symmetry representation}
\GSfield{Kind; source}{construction; \GSpage{P000487}}
\GSfield{The notes}{For the Lobachevsky plane the notebook writes \(K^{1/h}=(1-z\bar z)^{+2/h}\), followed by a projective representation of its automorphism group.}
\GSfield{Statement}{On the unit disk the Bergman kernel is proportional to \((1-z\bar v)^{-2}\). With \(\beta=2/h>1\), the weighted holomorphic norm has density proportional to \((1-|z|^2)^{\beta-2}d^2z\) and reproducing kernel \((1-z\bar v)^{-\beta}\). The automorphism group acts projectively by weighted composition; a genuine representation can be taken on a suitable covering group. The positive exponent belongs to the reciprocal weight, not the reproducing kernel.}
\GSfield{Fate}{The disk specializes the homogeneous-domain construction in \citet[\S5, Theorems 5.1--5.2]{berezin1974quantization}; parameters and symmetry are developed further in \citet{berezin1975symmetric}.}
\GSfield{Present status}{As of 9 November 2026: Established. The displayed sign must be corrected and the Hilbert-space parameter range specified.}
\GSfield{Assessment}{literal statement: false as written; publication match: exact; verification: read the result; development: none; confidence: high}
\end{GSitem}

\GStalk{293}{1973-03-12}{V. M. Buchstaber}{P000488-P000490}{Pontryagin--Thom normal maps and the surgery exact sequence}
The lecture constructs the obstruction map from a trivialized normal bundle to the periodic surgery group by collapse and transversality. It then asks how unique the constructed manifold is and what term precedes the structure set. The modern answer requires normal bordism and the action of the next surgery group, not bare bundle isomorphism. The ``film'' at the end is a homotopy of stable Thom maps whose transverse inverse image is a bordism. This revisits the normal-map and surgery constructions already indexed in 1968.

\begin{GSitem}{1972-293.1}{E}{Four structure notations}
\GSfield{Kind; source}{definition; \GSpage{P000488}}
\GSfield{The notes}{The page lists \(S(M)\), \(hS(M)\), \(t(M)\), and \(ht(M)\), respectively smoothings, homotopy smoothings, triangulations, and homotopy triangulations.}
\GSfield{Statement}{The h-prefixed objects are the marked structure sets of homotopy equivalences, in the smooth and PL categories. The unprefixed terms refer to actual structures compatible with the underlying topological manifold. No new equivalence definition for the unprefixed sets is supplied here; later arguments explicitly concern \(hS(M)\).}
\GSfield{Reconstruction}{We do not equate concordance, isotopy, and abstract diffeomorphism classes of actual structures, since the notebook does not choose between those notions.}
\GSfield{Fate}{The homotopy-structure definitions match \citet[pp.81--82]{sullivan1996triangulating}.}
\GSfield{Present status}{As of 9 November 2026: standard for the h-prefixed objects; the unprefixed notation alone does not specify a classification problem.}
\GSfield{Assessment}{literal statement: correct with missing hypotheses; publication match: partial; verification: read the result; development: none; confidence: high}
\end{GSitem}
\begin{GSitem}{1972-293.2}{E}{The set Phi must retain a specified trivialization}
\GSfield{Kind; source}{definition; \GSpage{P000488}}
\GSfield{The notes}{\(\Phi(M)\) is called the set, then group, of stable smooth J-trivial bundles. A proper equivalence \(\varphi:\xi\to M\times\mathbb R^N\) is drawn, with \(N=\dim\xi\).}
\GSfield{Statement}{The surgery construction needs stable equivalence classes of pairs \((\xi,\varphi)\), not only bundles admitting some such \(\varphi\). This refined \(\Phi(M)\) is \([M,G/O]\), equivalently the normal-bordism set of degree-one smooth normal maps after fixing the reference manifold M.}
\GSfield{Reconstruction}{The drawing contains the required trivialization but the prose does not explicitly include it in the equivalence relation. We retain this distinction, as in 1972-287.7. The rank N in this paragraph is unrelated to the source manifold denoted N below.}
\GSfield{Fate}{\citet[p.80]{sullivan1996triangulating}; \citet[Propositions9.21,9.43]{ranicki2002surgery}.}
\GSfield{Present status}{As of 9 November 2026: standard after the essential pair-data refinement.}
\GSfield{Assessment}{literal statement: correct with missing hypotheses; publication match: exact; verification: read the result; development: annotated; confidence: high}
\end{GSitem}
\begin{GSitem}{1972-293.3}{E}{Normal difference and its zero normalization}
\GSfield{Kind; source}{construction; \GSpage{P000488}}
\GSfield{The notes}{For a homotopy equivalence \(\alpha:N\to M\) with inverse g, the class is \(g^*\nu_N+\tau_M\). The parenthesis mistakenly calls \(\nu_N\) the normal bundle of M. For the identity, the sum is trivial; for a homotopy equivalence, Atiyah--Dold supplies J-triviality.}
\GSfield{Statement}{\(\nu_N\) is the stable normal bundle of the source N. Modulo a trivial summand, \(g^*\nu_N+\tau_M\) represents \(\tau_M-g^*\tau_N\), is zero for the identity structure, and has trivial stable spherical fibration for any homotopy-equivalence marking. Its refinement with the induced trivialization is the normal invariant.}
\GSfield{Reconstruction}{The wrong base M in the explanatory parenthesis is legible in the scan and both transcriptions. Correcting it to N is necessary for the pullback by g even to have the stated domain.}
\GSfield{Fate}{\citet[Theorem~3.6]{atiyah1961thom}.}
\GSfield{Present status}{As of 9 November 2026: established; the local source/target slip is corrected.}
\GSfield{Assessment}{literal statement: false as written; publication match: exact; verification: read the result; development: annotated; confidence: high}
\end{GSitem}
\begin{GSitem}{1972-293.4}{E}{The simply connected surgery obstruction sequence}
\GSfield{Kind; source}{theorem; \GSpage{P000488}}
\GSfield{The notes}{Assuming \(\pi_1(M)=0\), the notes display \(hS(M)\to\Phi(M)\to P_n\), with \(P_n=0,\mathbb Z,\mathbb Z/2\) according to the dimension modulo four.}
\GSfield{Statement}{For a closed simply connected smooth n-manifold with \(n\ge6\), \(\mathcal S^{\mathrm{Diff}}(M)\xrightarrow\eta[M,G/O]\xrightarrow\theta L_n(\mathbb Z)\) is exact in the sense that \(\operatorname{im}\eta=\theta^{-1}(0)\). The group \(L_n(\mathbb Z)\) is zero in odd dimensions, \(\mathbb Z\) in dimension \(4k\), and \(\mathbb Z/2\) in dimension \(4k+2\). The smooth obstruction map need not be surjective or additive.}
\GSfield{Reconstruction}{The missing dimensional/compactness assumptions and the normal-invariant interpretation are supplied. A vanishing obstruction means normal bordism to a homotopy equivalence, not that the particular transverse representative already is one.}
\GSfield{Fate}{Pre-lecture surgery theory: \citet[Theorem~3]{sullivan1996triangulating}; precise modern statement \citet[Theorem~13.2]{ranicki2002surgery}.}
\GSfield{Present status}{As of 9 November 2026: established in the specified range.}
\GSfield{Assessment}{literal statement: correct with missing hypotheses; publication match: exact; verification: read the result; development: none; confidence: high}
\end{GSitem}
\begin{GSitem}{1972-293.5}{E}{Adding the reference normal bundle and passing to Thom spaces}
\GSfield{Kind; source}{construction; \GSpage{P000488}}
\GSfield{The notes}{The proper trivialization is summed with \(\nu_M\), inducing \(T(\xi+\nu_M)\to T(N+\nu_M)\).}
\GSfield{Statement}{For an actual rank-r representative \(\xi\) and its fibrewise trivialization, summing with \(\nu_M\) yields a proper equivalence to \(\varepsilon^r\oplus\nu_M\), hence a Thom homotopy equivalence \(\varphi_0:T(\xi\oplus\nu_M)\simeq\Sigma^rT(\nu_M)\). In the notebook's N+ notation, N means a trivial bundle of that rank.}
\GSfield{Reconstruction}{A homotopy inverse \(\sigma\) may be chosen. To keep dimensions consistent in the next construction, one suspends the Pontryagin--Thom collapse by r; alternatively one works throughout with virtual rank-zero Thom spectra.}
\GSfield{Fate}{\citet[\S2]{atiyah1961thom}; the resulting normal-invariant construction is explicit in \citet[Proposition~9.43]{ranicki2002surgery}.}
\GSfield{Present status}{As of 9 November 2026: standard once rank and suspension conventions are stated.}
\GSfield{Assessment}{literal statement: correct with missing hypotheses; publication match: exact; verification: read the result; development: annotated; confidence: high}
\end{GSitem}
\begin{GSitem}{1972-293.6}{E}{Pontryagin--Thom collapse of the reference embedding}
\GSfield{Kind; source}{construction; \GSpage{P000489}}
\GSfield{The notes}{For \(M^n\subset\mathbb R^L\), a normal tube gives \(\lambda:S^L\to T(\nu_M)\), transverse to the zero section with preimage M.}
\GSfield{Statement}{Embed the closed smooth manifold M in Euclidean space and identify a tubular neighbourhood with the normal disk bundle. Collapsing the complement of its interior to the Thom basepoint gives \(\lambda\); near the zero section it is smooth and transverse, and \(\lambda^{-1}(M)=M\). Its stable homotopy class is independent of the auxiliary embedding and tube in the appropriate normal-invariant formulation.}
\GSfield{Reconstruction}{The zero section is an embedding into the Thom space; it is the collapse map that is transverse to that section, repairing the syntactic ambiguity in the second notebook sentence. This repeats the construction of 1968-171.13.}
\GSfield{Fate}{\citet[Example~9.41 and Theorem~9.42]{ranicki2002surgery}; the underlying construction is classical Pontryagin--Thom theory.}
\GSfield{Present status}{As of 9 November 2026: established.}
\GSfield{Assessment}{literal statement: correct; publication match: exact; verification: read the result; development: none; confidence: high}
\end{GSitem}
\begin{GSitem}{1972-293.7}{E}{Transversality produces a degree-one normal map}
\GSfield{Kind; source}{construction; \GSpage{P000489}}
\GSfield{The notes}{Compose the collapse with a homotopy inverse \(\sigma\), deform it transverse to M, and set \(M_1=(\widetilde{\sigma\lambda})^{-1}(M)\). The boxed identity identifies its normal bundle as a pullback of \(\xi+\nu_M\).}
\GSfield{Statement}{With the suspension convention of 1972-293.5, make \(\sigma\circ\Sigma^r\lambda:S^{L+r}\to T(\xi\oplus\nu_M)\) transverse to the zero section. Its inverse image is a closed smooth n-manifold \(M_1\), and the restricted map \(f:M_1\to M\) has degree one with stable normal identification \(\nu_{M_1}\cong f^*(\xi\oplus\nu_M)\). This is a degree-one normal map representing the given normal invariant.}
\GSfield{Reconstruction}{The degree-one property follows from the chosen fibre degree and Thom class; it is not written explicitly in the notebook but is indispensable for surgery. The equality is a stable bundle identification, not equality of unrelated ambient normal ranks.}
\GSfield{Fate}{\citet[Theorem~9.42 and Proposition~9.43]{ranicki2002surgery}; earlier occurrences 1968-171.10 and 1968-171.13.}
\GSfield{Present status}{As of 9 November 2026: an established construction, not a new existence theorem.}
\GSfield{Assessment}{literal statement: correct with missing hypotheses; publication match: exact; verification: read the result; development: none; confidence: high}
\end{GSitem}
\begin{GSitem}{1972-293.8}{E}{How uniquely is the transverse manifold determined?}
\GSfield{Kind; source}{question; \GSpage{P000489-P000490}}
\GSfield{The notes}{The page asks how uniquely \(M_1\) is determined; the next page says all such manifolds are related by Morse surgeries, attributing the statement to Milnor--Kervaire.}
\GSfield{Statement}{A fixed normal invariant determines a normal-bordism class of degree-one normal maps. Different transverse approximations and compatible Thom homotopies give a normal bordism; a handle decomposition of that bordism gives a finite sequence of surgeries compatible with the normal data. It does not determine the diffeomorphism type of \(M_1\), and equality of bare stable bundle classes is weaker than equality of normal invariants.}
\GSfield{Reconstruction}{The exact theorem matching the notes is the normal-map bordism theorem, followed by the surgery/handle decomposition of a bordism. No unique Kervaire--Milnor theorem locator matching this entire compressed attribution was found, so that attribution remains only partial. The underlying questions were already present in 1968-171.15 and 1968-173.2--3.}
\GSfield{Fate}{\citet[Proposition~9.43; \S4.3]{ranicki2002surgery} supplies the precise standard statement.}
\GSfield{Present status}{As of 9 November 2026: resolved at the level of normal bordism; finer structure-set ambiguity is the next item.}
\GSfield{Assessment}{literal statement: correct with missing hypotheses; publication match: partial; verification: read the result; development: none; confidence: medium; problem \hyperref[prob:GP-1972-204]{GP-1972-204}}
\end{GSitem}
\begin{GSitem}{1972-293.9}{E}{Can surgery produce a homotopy equivalence?}
\GSfield{Kind; source}{question; \GSpage{P000490}}
\GSfield{The notes}{``Is there a manifold there homotopy equivalent to M?'' The group \(P_n\) is called the obstruction to surgery to such an equivalence, and exactness is asserted.}
\GSfield{Statement}{For the simply connected closed degree-one normal map in the high-dimensional range, its normal-bordism class contains a homotopy equivalence exactly when its surgery obstruction vanishes in \(L_n(\mathbb Z)\). In dimension \(4k\), the obstruction is the signature defect divided by eight; in dimension \(4k+2\) it is the Arf invariant of the surgery kernel's quadratic form; in odd dimensions it vanishes.}
\GSfield{Reconstruction}{The obstruction formulas are supplied from the literature to make \(P_n\) substantive. They apply to the normal-map problem and require preliminary surgeries below the middle dimension, not to an arbitrary continuous map or arbitrary manifold.}
\GSfield{Fate}{\citet[Theorem~13.2 and Example~13.3]{ranicki2002surgery}, consolidating the pre-1973 Browder--Novikov--Sullivan--Wall theorem.}
\GSfield{Present status}{As of 9 November 2026: established; this is an exposition of the surgery criterion.}
\GSfield{Assessment}{literal statement: correct with missing hypotheses; publication match: exact; verification: read the result; development: none; confidence: high; problem \hyperref[prob:GP-1972-203]{GP-1972-203}}
\end{GSitem}
\begin{GSitem}{1972-293.10}{E}{The missing left term and h-cobordisms among films}
\GSfield{Kind; source}{problem; \GSpage{P000490}}
\GSfield{The notes}{The sequence begins with a question mark before \(hS(M)\). Two structures yield Thom maps \(\beta_1,\beta_2\); a homotopy F is called a film, and the final question asks whether h-cobordisms occur among these films.}
\GSfield{Statement}{For \(n\ge6\), the missing ambiguity is the action of \(L_{n+1}(\mathbb Z)\) on the smooth structure set. Structures with the same normal invariant lie in one action orbit. A transverse Thom homotopy represents a normal bordism \(W\to M\times I\); it can be changed relative to its ends to an h-cobordism precisely when its relative surgery obstruction is zero. Different films can change that obstruction by the image of \([\Sigma M_+,G/O]\to L_{n+1}(\mathbb Z)\).}
\GSfield{Reconstruction}{This is an action exact sequence, not an ordinary group kernel-cokernel sequence. In the simply connected dimension range an h-cobordism is a product, giving the required diffeomorphism over M up to homotopy. The notebook's endpoint diagram omits normal-map data; stable bundles alone do not justify existence of F. Its Thom targets also contain an unclear \(\widetilde\tau_M\)-like sign where the preceding construction uses \(\nu_M\). We use the stable normal bundle in this reconstruction, while retaining the source reading as an ambiguity rather than silently changing it. The attribution to Novikov's dissertation is compatible with the earlier 1968 normal-map entries but does not uniquely locate a dissertation proposition.}
\GSfield{Fate}{The precise answer is \citet[Proposition~9.48 and Theorem~13.2]{ranicki2002surgery}; the pre-lecture smooth sequence already occurs in \citet[Theorem~3]{sullivan1996triangulating}.}
\GSfield{Present status}{As of 9 November 2026: the structural question is resolved by relative surgery; no concrete uncomputed example is specified.}
\GSfield{Assessment}{literal statement: correct with missing hypotheses; publication match: partial; verification: read the result; development: none; confidence: high; problem \hyperref[prob:GP-1972-204]{GP-1972-204}}
\end{GSitem}

\GStalk{294}{1973-03-19}{V. I. Arnold}{P000491-P000492}{Modality and the intersection forms of simple and unimodal singularities}
Arnold’s actual passage begins on P000491 and ends above Mishchenko’s name on P000492. It recalls finite determinacy and the ADE classification, then links Milnor lattices to definiteness and turns to modality one. The unimodal announcement is identified exactly: Arnold’s paper was received 12 March 1973, one week before this seminar, and published later in 1973. The notebook compresses several different signature statements into false implications; these are checked against Arnold’s announcement, Tyurina’s 1968 paper and Gabrielov’s intersection-matrix calculation.

\begin{GSitem}{1972-294.1}{E}{Finite determinacy and the action on jets}
\GSfield{Kind; source}{definition; \GSpage{P000491}}
\GSfield{The notes}{For \(f:(\mathbb C^n,0)\to(\mathbb C,0)\), the group of \(k\)-jets of diffeomorphisms acts on function jets; a sufficient jet means that later terms do not matter.}
\GSfield{Statement}{A holomorphic germ \(f\) is right \(k\)-determined if every holomorphic germ \(g\) with \(j^kg=j^kf\) satisfies \(g=f\circ\phi\) for a local biholomorphism \(\phi\) fixing zero. An isolated critical point gives finite determinacy for some \(k\); no particular universal jet order is specified here.}
\GSfield{Reconstruction}{Isolatedness, holomorphic coordinate changes and the existence rather than a prescribed value of k are made explicit.}
\GSfield{Fate}{This is the finite-jet framework of \citet[sections 2--3]{arnold1972normal}.}
\GSfield{Present status}{As of 9 November 2026: established with the hypotheses and distinctions stated here.}
\GSfield{Assessment}{literal statement: correct with missing hypotheses; publication match: exact; verification: read the result; development: none; confidence: high}
\end{GSitem}
\begin{GSitem}{1972-294.2}{E}{Modality counts families of nearby orbits}
\GSfield{Kind; source}{definition; \GSpage{P000491-P000492}}
\GSfield{The notes}{\(m\) is the number of moduli; zero modality means only finitely many possibilities after perturbing coefficients.}
\GSfield{Statement}{For a sufficiently high jet of a finitely determined germ, its modality is the least \(m\) such that a neighborhood is covered by finitely many families of right-equivalence orbits depending on at most \(m\) parameters. It is a local orbit-family dimension, not the dimension of the miniversal base \(\mu\), nor simply the number of coefficients in a chosen polynomial.}
\GSfield{Fate}{Exact definition: \citet[section 1]{arnold1973unimodal}, submitted 12 March 1973; the zero-modality case is already defined in \citet[section 2]{arnold1972normal}.}
\GSfield{Present status}{As of 9 November 2026: established with the hypotheses and distinctions stated here.}
\GSfield{Assessment}{literal statement: correct with missing hypotheses; publication match: exact; verification: read the result; development: none; confidence: high}
\end{GSitem}
\begin{GSitem}{1972-294.3}{E}{Zero modality is exactly ADE}
\GSfield{Kind; source}{theorem; \GSpage{P000491}}
\GSfield{The notes}{The list is \(A_k:x^{k+1}\), \(D_k:x^2y+y^{k-1}\), \(E_6:x^3+y^4\), \(E_7:x^3+xy^3\), \(E_8:x^3+y^5\).}
\GSfield{Statement}{The complete list of holomorphic isolated right-simple germs up to stable equivalence is the ADE list of 1972-284.2, with \(k\ge1\) for \(A_k\) and \(k\ge4\) for \(D_k\). Conversely each listed type has modality zero.}
\GSfield{Fate}{Published before this talk in \citet[Theorem 2.10, Corollary 8.6]{arnold1972normal}; the 1973 announcement recalls it as the zero-modal case.}
\GSfield{Present status}{As of 9 November 2026: established with the hypotheses and distinctions stated here.}
\GSfield{Assessment}{literal statement: correct with missing hypotheses; publication match: exact; verification: read the result; development: none; confidence: high}
\end{GSitem}
\begin{GSitem}{1972-294.4}{E}{Stable equivalence and the inappropriate positivity condition}
\GSfield{Kind; source}{definition; \GSpage{P000491}}
\GSfield{The notes}{Stable equivalence means adding nondegenerate quadratic forms; the next sentence says the quadratic form is positive definite.}
\GSfield{Statement}{Two complex holomorphic germs are stably right-equivalent if, after adding nondegenerate complex quadratic forms in auxiliary variables, they become right-equivalent. Every nondegenerate complex quadratic form is equivalent to a sum of squares; positivity is not an intrinsic notion over \(\mathbb C\). For real germs the signature of the quadratic summand is relevant and need not be positive.}
\GSfield{Reconstruction}{The handwritten positivity sentence is faithfully reproduced; its correct replacement depends on the coefficient field.}
\GSfield{Fate}{\citet[section 1]{arnold1973unimodal} uses nondegenerate forms. The real signatures in \citet[Introduction, formula (1.1)]{arnold1972normal} also rule out a blanket positivity requirement.}
\GSfield{Present status}{As of 9 November 2026: established with the hypotheses and distinctions stated here.}
\GSfield{Assessment}{literal statement: false as written; publication match: exact; verification: read the result; development: annotated; confidence: high}
\end{GSitem}
\begin{GSitem}{1972-294.5}{F}{The unresolved sentence about analytic coordinate changes}
\GSfield{Kind; source}{remark; \GSpage{P000492}}
\GSfield{The notes}{\ru{Аналит. ф-я с точн. до анал. зам. перем. однозн. опр. зам. перем.}}
\GSfield{Statement}{The surviving line does not specify what determines what. If read as uniqueness of the analytic normalizing coordinate change, it is false: even a nondegenerate quadratic form has nontrivial linear symmetries. Finite determinacy concerns existence of an equivalence, not uniqueness of a coordinate map.}
\GSfield{Reconstruction}{The tautological diplomatic reading is retained; replacing its final words by Taylor polynomial or another guessed object would go beyond the scan.}
\GSfield{Fate}{The surrounding finite-determinacy and modality definitions identify the subject but not the missing assertion. No exact publication is assigned to this line.}
\GSfield{Present status}{As of 9 November 2026: the intended assertion remains unreadable at the level needed for a mathematical identification.}
\GSfield{Assessment}{literal statement: not determinable; publication match: none found; verification: unresolved lead; development: none; confidence: low}
\end{GSitem}
\begin{GSitem}{1972-294.6}{E}{The local smooth level manifold}
\GSfield{Kind; source}{construction; \GSpage{P000492}}
\GSfield{The notes}{\(f^{-1}(\varepsilon)\cap B_\delta\) is a manifold with boundary of even real dimension.}
\GSfield{Statement}{For a holomorphic germ with isolated critical point and radii \(0<|\varepsilon|\ll\delta\), transverse intersection with the closed ball gives a compact smooth real \(2n-2\)-manifold with boundary, whose interior is a complex \((n-1)\)-manifold. Its diffeomorphism type is independent of admissible sufficiently small choices.}
\GSfield{Reconstruction}{The required order of smallness and boundary transversality are supplied.}
\GSfield{Fate}{The admissible radius/value construction is explicit in \citet[section 1.1]{tyurina1968topological}; see also \citet[sections 4--6]{milnor1968singular}.}
\GSfield{Present status}{As of 9 November 2026: established with the hypotheses and distinctions stated here.}
\GSfield{Assessment}{literal statement: correct with missing hypotheses; publication match: exact; verification: read the result; development: none; confidence: high}
\end{GSitem}
\begin{GSitem}{1972-294.7}{E}{Tyurina’s definiteness criterion}
\GSfield{Kind; source}{theorem; \GSpage{P000492}}
\GSfield{The notes}{The middle-dimensional integral intersection form is negative definite exactly in the listed ADE cases when \(n\equiv3\pmod4\); otherwise add variables. The attribution reads \ru{теор. Тюрина}.}
\GSfield{Statement}{After stabilization to \(n\equiv3\pmod4\), the Milnor intersection form of an isolated complex hypersurface critical point is negative definite precisely for the simple ADE types. In complex surface dimension two this is the equivalence between negative-definite smoothing intersection form and a rational double point. Other congruences require the appropriate suspension and sign convention.}
\GSfield{Reconstruction}{The attribution is to Galina N. Tyurina. Russian Тюрина here is the genitive form and should not be normalized as the surname of a different male author.}
\GSfield{Fate}{The original surface result is \citet[section 3, Lemma 3, together with Theorem 2]{tyurina1968topological}, received 13 May 1967, published 1968. \citet[section 4, Remarks 3--4]{arnold1973unimodal} explicitly attributes the elliptic/ADE classification to G. N. Tyurina and states it in stable form.}
\GSfield{Present status}{As of 9 November 2026: established with the hypotheses and distinctions stated here.}
\GSfield{Assessment}{literal statement: correct with missing hypotheses; publication match: exact; verification: read the result; development: none; confidence: high}
\end{GSitem}
\begin{GSitem}{1972-294.8}{E}{Milnor number, bouquet and local algebra}
\GSfield{Kind; source}{theorem; \GSpage{P000492}}
\GSfield{The notes}{The level manifold is a bouquet of \(\mu\) copies of \(S^{n-1}\), and \(\mu\) is the dimension of the local ring.}
\GSfield{Statement}{The precise statement is 1972-284.11: the Milnor number is \[\mu=\dim_{\mathbb C}\mathcal O/(\partial f),\] not the dimension of the whole local ring \(\mathcal O\), which is infinite. This finite number is the rank of reduced middle homology and the number of spheres in the bouquet.}
\GSfield{Reconstruction}{Local ring is read as the finite local algebra of the gradient mapping, consistently with record 284.}
\GSfield{Fate}{\citet[Theorems 6.5 and 7.2, Appendix B, Problem 3]{milnor1968singular}; the terminology local ring of the gradient map is used by \citet[section 1]{arnold1973unimodal}.}
\GSfield{Present status}{As of 9 November 2026: established with the hypotheses and distinctions stated here.}
\GSfield{Assessment}{literal statement: correct with missing hypotheses; publication match: exact; verification: read the result; development: none; confidence: high}
\end{GSitem}
\begin{GSitem}{1972-294.9}{E}{The E6 intersection diagram}
\GSfield{Kind; source}{construction; \GSpage{P000492}}
\GSfield{The notes}{An \(E_6\) graph is labelled \(-2\) at a vertex and \(1\) at an edge; the quadratic form is said to be given by the Dynkin diagram.}
\GSfield{Statement}{In a suitable distinguished basis of the Milnor lattice of \(E_6\), stabilized to complex fibre dimension \(2\bmod4\), the Gram matrix has diagonal \(-2\), entry \(1\) for adjacent vertices, and entry \(0\) otherwise. It is the negative \(E_6\) Cartan matrix. The diagram depends on a basis; it is not an unlabelled invariant for arbitrary singularities.}
\GSfield{Reconstruction}{An independent exact matrix check \oursComputed gives the positive E6 Cartan leading principal minors \(2,3,4,5,6,3\), confirming negative definiteness after changing sign; see the accompanying code and output. The identification of the singularity’s lattice is the cited published result.}
\GSfield{Fate}{\citet[Proposition 9.4]{arnold1972normal} gives the ADE rule. An explicit basis transformation for \(E_6\) is \citet[section 6, Example 2]{gabrielov1973intersection}, received 29 January 1973 and published later that year.}
\GSfield{Present status}{As of 9 November 2026: established with the hypotheses and distinctions stated here.}
\GSfield{Assessment}{literal statement: correct with missing hypotheses; publication match: exact; verification: read the result; development: annotated; confidence: high}
\end{GSitem}
\begin{GSitem}{1972-294.10}{A}{The unimodal classification announcement}
\GSfield{Kind; source}{theorem; \GSpage{P000492}}
\GSfield{The notes}{For \(m=1\) there is also a long list; the first displayed example is \(x_1^3+x_2^3+x_3^3\).}
\GSfield{Statement}{Up to stable right equivalence, the isolated holomorphic germs of modality one consist of the three-index \(T_{pqr}\) series, including its three parabolic cases, and fourteen exceptional one-parameter families. The latter are named \(E_{12},E_{13},E_{14},Z_{11},Z_{12},Z_{13},Q_{10},Q_{11},Q_{12},W_{12},W_{13},S_{11},S_{12},U_{12}\). The notebook does not list these fourteen normal forms.}
\GSfield{Reconstruction}{The family names and organization are reconstructed from the exact publication; the unrecorded normal-form table is not passed off as notebook content.}
\GSfield{Fate}{Exact announcement: \citet[sections 2--3]{arnold1973unimodal}, received 12 March 1973, published 1973. \citet[Introduction, section 1.4]{arnold1974normal} gives the organized T-series/fourteen-family formulation; received 28 September 1973, published 1974.}
\GSfield{Present status}{As of 9 November 2026: the stated unimodal classification is established. Its published scope is isolated holomorphic right-equivalence classes up to stabilization.}
\GSfield{Assessment}{literal statement: correct with missing hypotheses; publication match: exact; verification: read the result; development: none; confidence: high}
\end{GSitem}
\begin{GSitem}{1972-294.11}{A}{The cubic example is simple elliptic with a two-dimensional radical}
\GSfield{Kind; source}{construction; \GSpage{P000492}}
\GSfield{The notes}{The first modality-one example is \(x_1^3+x_2^3+x_3^3\).}
\GSfield{Statement}{The cubic \(f=x^3+y^3+z^3\) has \(\mu=8\), modality one and intersection-form inertia \((\mu_+,\mu_0,\mu_-)=(0,2,6)\). It belongs to the family \(x^3+y^3+z^3+a xyz\), with \(a^3\ne-27\) for an isolated critical point. Its Milnor lattice splits as a negative \(E_6\) lattice plus a rank-two zero form.}
\GSfield{Reconstruction}{The modulus and its excluded discriminant value are supplied. The notebook example alone already shows why a blanket nullity-one assertion is unsafe; the exact corank is read in the original calculation.}
\GSfield{Fate}{\citet[section 6, Example 4 and Proposition 5]{gabrielov1973intersection} gives the lattice and corank two. The paper records Milnor’s letter to Arnold of 31 October 1972 about the parabolic monodromy groups; Gabrielov’s manuscript was received 29 January 1973.}
\GSfield{Present status}{As of 9 November 2026: established with the hypotheses and distinctions stated here.}
\GSfield{Assessment}{literal statement: correct with missing hypotheses; publication match: exact; verification: read the result; development: annotated; confidence: high}
\end{GSitem}
\begin{GSitem}{1972-294.12}{A}{The signature implications need two separate corrections}
\GSfield{Kind; source}{claim; \GSpage{P000492}}
\GSfield{The notes}{With \(\mu=\mu_++\mu_0+\mu_-\), the notebook appears to assert \(\mu_+\le1\Rightarrow m=1\), and explicitly writes \(\mu_+\le1\Rightarrow\mu_0=1\).}
\GSfield{Statement}{Neither displayed implication is true as an unrestricted assertion. ADE forms have \(\mu_+=0\), \(\mu_0=0\), modality zero. The cubic in 1972-294.11 has \(\mu_+=0\), \(\mu_0=2\), modality one. The matching valid result in Arnold’s announcement is that the cases with \(\mu_+=1\) are the nonparabolic T-series; they have \(\mu_0=1\). This does not assert that every unimodal singularity has \(\mu_+\le1\).}
\GSfield{Reconstruction}{The first inequality is marked doubtful in the diplomatic transcription, so its exact sign remains an ambiguity. The second \(\le1\) is clear in the scan and is a faithful mathematical error. No silent replacement of both inequalities by equalities is made.}
\GSfield{Fate}{\citet[section 3, Corollary Zh]{arnold1973unimodal} gives the precise positive-index-one statement; \citet[Proposition 5]{gabrielov1973intersection} supplies the decisive cubic counterexample.}
\GSfield{Present status}{As of 9 November 2026: the literal implications are false; the correctly restricted positive-index-one classification is established.}
\GSfield{Assessment}{literal statement: false as written; publication match: partial; verification: read the result; development: annotated; confidence: high}
\end{GSitem}
\begin{GSitem}{1972-294.13}{A}{The hyperbolic T-series and its missing restrictions}
\GSfield{Kind; source}{construction; \GSpage{P000492}}
\GSfield{The notes}{Under \(\mu_+=1\), \(axyz+x^k+y^\ell+z^m\) is called the three-index series \(T_{k\ell m}\).}
\GSfield{Statement}{For integers \(p,q,r\ge2\) with \(1/p+1/q+1/r<1\) and \(a\ne0\), the germ \(x^p+y^q+z^r+a xyz\) is of hyperbolic type \(T_{pqr}\), with \(\mu=p+q+r-1\), modality one and \((\mu_+,\mu_0)=(1,1)\) after the stated stabilization. Equality \(1/p+1/q+1/r=1\) gives the parabolic types \((3,3,3),(2,4,4),(2,3,6)\), with parameter values restricted to keep the critical point isolated, and with a different signature.}
\GSfield{Reconstruction}{Exponent restrictions, the nonzero coefficient and the hyperbolic/parabolic distinction are supplied. The three integers label a series; they are not three continuous moduli.}
\GSfield{Fate}{The T-series and inertia-one characterization are \citet[section 3]{arnold1973unimodal}; the reciprocal-sum split is explicit in \citet[section 1.4]{arnold1974normal}.}
\GSfield{Present status}{As of 9 November 2026: established with the hypotheses and distinctions stated here.}
\GSfield{Assessment}{literal statement: correct with missing hypotheses; publication match: exact; verification: read the result; development: none; confidence: high}
\end{GSitem}

\GStalk{295}{1973-03-19}{A. S. Mishchenko}{P000492-P000494}{Sullivan’s localization, completion and proof of the Adams conjecture}
The report starts under Mishchenko’s name in the lower part of P000492 and ends with the comparison of analytic and étale homotopy on P000494. It explains the Adams conjecture through completed Grassmannians, spherical fibrations and Galois symmetry. Its source is identifiable as Sullivan’s June 1970 MIT notes, whose title is almost quoted; the expanded Annals paper appeared in 1974. Thus the material was already circulating, and Quillen’s independent proof had appeared in 1971. The report continues the Adams-operation and étale-homotopy discussions 1971-259.2 and 1971-262.1--.3. The broad catalogue range also contains Arnold, APS and the p-adic report, which are not assigned to Mishchenko.

\begin{GSitem}{1972-295.1}{E}{The finite reduced J-group}
\GSfield{Kind; source}{definition; \GSpage{P000492}}
\GSfield{The notes}{\(K(X)/J\text{-equivalence}=J(X)\) is called a finite group.}
\GSfield{Statement}{For a connected finite CW-complex \(X\), quotient the reduced real K-group \(\widetilde{KO}(X)\) by the subgroup of stably fibre-homotopy-trivial virtual bundles. The resulting reduced J-group is finite. The corresponding complex version uses underlying real sphere bundles. If unreduced K-theory and actual ranks are retained, the quotient also has a \(\mathbb Z\) rank summand and is not finite.}
\GSfield{Reconstruction}{Finite CW type, reduction and stable equivalence are necessary to make the displayed finiteness assertion correct.}
\GSfield{Fate}{\citet[section 1, p.~181]{adams1963groups} explicitly separates \(J(X)=\mathbb Z\oplus\widetilde J(X)\) and quotes Atiyah’s earlier finiteness theorem; received 29 May 1963 and published 1963.}
\GSfield{Present status}{As of 9 November 2026: established with the hypotheses and distinctions stated here.}
\GSfield{Assessment}{literal statement: correct with missing hypotheses; publication match: exact; verification: read the result; development: annotated; confidence: high}
\end{GSitem}
\begin{GSitem}{1972-295.2}{E}{Adams operations determined on line bundles}
\GSfield{Kind; source}{definition; \GSpage{P000492}}
\GSfield{The notes}{\(\psi^k(\xi)=\xi^k\) for \(\dim\xi=1\), together with additivity, determines the operations.}
\GSfield{Statement}{For complex topological K-theory on finite CW-complexes, \(\psi^k\) is the natural ring endomorphism satisfying \(\psi^k[L]=[L^{\otimes k}]\) on line bundles. Naturality, the splitting principle and additivity determine it on all classes. It preserves virtual rank and obeys \(\psi^k\psi^\ell=\psi^{k\ell}\). Its value on a higher-rank bundle is generally a virtual bundle, not a specified bundle of the same actual rank.}
\GSfield{Reconstruction}{Naturality and the splitting principle are supplied; additivity on a single fixed K-group without these conditions would not give uniqueness.}
\GSfield{Fate}{\citet[pp.~162--163]{sullivan1970geometric} states the operations and Newton-polynomial construction; \citet[p.~7]{sullivan1974genetics} gives the line-bundle characterization. Both explain the pre-existing operations.}
\GSfield{Present status}{As of 9 November 2026: established with the hypotheses and distinctions stated here.}
\GSfield{Assessment}{literal statement: correct with missing hypotheses; publication match: exact; verification: read the result; development: none; confidence: high}
\end{GSitem}
\begin{GSitem}{1972-295.3}{E}{Fibrewise degree one and J-equivalence}
\GSfield{Kind; source}{definition; \GSpage{P000492}}
\GSfield{The notes}{J-equivalence is fibrewise homotopy equivalence; a proper bundle map is required to have degree 1 on every fibre.}
\GSfield{Statement}{Two real vector bundles of equal rank are stably J-equivalent when their unit-sphere bundles become fibre-homotopy equivalent after adding trivial bundles. A continuous proper fibre-preserving map, of degree \(+1\) on the one-point compactification of each oriented vector-space fibre, supplies such a stable equivalence: those compactifications are the sphere bundles after one additional trivial summand. A linear isomorphism is not required.}
\GSfield{Reconstruction}{Stabilization, equal dimensions and the relation between unit spheres and fibrewise one-point compactification distinguish the two versions of the notebook definition.}
\GSfield{Fate}{\citet[section 1]{adams1963groups} defines the quotient through sphere bundles. The fibre-degree construction and its localization variant are used in \citet[section 2, Lemma 2.3]{quillen1971adams}.}
\GSfield{Present status}{As of 9 November 2026: established with the hypotheses and distinctions stated here.}
\GSfield{Assessment}{literal statement: correct with missing hypotheses; publication match: exact; verification: read the result; development: none; confidence: high}
\end{GSitem}
\begin{GSitem}{1972-295.4}{E}{Recovering the J-quotient from K-theory and all Adams operations}
\GSfield{Kind; source}{theorem; \GSpage{P000492}}
\GSfield{The notes}{The group \(J(X)\) is determined by \(K(X)\) as a module over the operations.}
\GSfield{Statement}{For finite CW-complexes, Adams’s algebraic lower and conjectural upper J-quotients of real K-theory agree; the Adams conjecture identifies this common quotient with the topological J-quotient. The algebraic construction uses the subgroup generated by suitable \(k^{e(k)}(\psi^k-1)y\) as \(k\) varies, with sufficiently large exponents. It is not the image of one prescribed operation for every fixed \(k\).}
\GSfield{Reconstruction}{The precise algebraic relation is supplied from Adams. Rank reduction and real versus complex K-theory must be fixed before using the notebook shorthand.}
\GSfield{Fate}{\citet[Theorem 1.1 and section 4, pp.~208--209]{adams1965groupsiii}, received 25 November 1963 and published 1965, proves equality of the algebraic quotients. \citet[Theorem 1.1]{quillen1971adams} supplies the remaining conjectural relation. \citet[pp.~8--9]{sullivan1974genetics} explains the resulting universality.}
\GSfield{Present status}{As of 9 November 2026: established with the hypotheses and distinctions stated here.}
\GSfield{Assessment}{literal statement: correct with missing hypotheses; publication match: exact; verification: read the result; development: none; confidence: high}
\end{GSitem}
\begin{GSitem}{1972-295.5}{E}{The purported reverse Adams formula is false for each arbitrary k}
\GSfield{Kind; source}{claim; \GSpage{P000492}}
\GSfield{The notes}{For every J-trivial \(\xi\), and every \(k\), the notes assert \(\xi=k^N(\psi^k-1)y\), after a uniform \(N\) has been chosen.}
\GSfield{Statement}{The assertion is false already in integral real K-theory. \ours\ Let \(z\) generate \(\widetilde{KO}(S^4)\cong\mathbb Z\), so \(\psi^kz=k^2z\). Let \(M\ge1\) be the finite order of \(J(z)\), choose a prime \(q\nmid M\), and set \(k=q+1\). Then \(Mz\) is J-trivial, but \(Mz=k^N(\psi^k-1)(mz)\) would require \(M=k^N(k^2-1)m\). The right side is divisible by \(q\) and the left side is not. This holds for every integer \(m\) and every \(N\ge0\), and still holds after inverting \(k\), since \(q\nmid k\). Thus the surviving universal quantifier in k is decisive.}
\GSfield{Reconstruction}{The counterexample is our verification of the literal quantifiers. The sphere-operation formula is checked in Adams 1965, p.~207; it also follows by injective complexification and the degree-four Chern character. An independent complex check gives the same verdict: in \(\widetilde{KU}(S^2)=\mathbb Zx\), \(x=[L]-1\), one has \(x^2=0\), \(\psi^4x=4x\), and \(2x\) has trivial realification because \(\widetilde{KO}(S^2)=\mathbb Z/2\). Hence J-triviality of \(2x\) would force \(2=3\cdot4^Nm\), impossible even in \(\mathbb Z[1/4]\). Both arguments were independently checked; neither purports to refute a differently localized or differently quantified statement.}
\GSfield{Fate}{The exact prior theorem is the varying-k subgroup statement of \citet[Theorem 1.1 and section 4]{adams1965groupsiii}, not the displayed single-image equality. Sullivan’s completed version instead says a J-trivial completed class is \(\eta-\sigma\eta\) for some Galois unit \(\sigma\); see \citet[p.~9, corollary]{sullivan1974genetics}.}
\GSfield{Present status}{As of 9 November 2026: the displayed universal formula is refuted; Adams’s differently quantified algebraic theorem is established.}
\GSfield{Assessment}{literal statement: false as written; publication match: partial; verification: read the result; development: annotated; confidence: high}
\end{GSitem}
\begin{GSitem}{1972-295.6}{E}{The Adams conjecture, already proved before this report}
\GSfield{Kind; source}{conjecture; \GSpage{P000492}}
\GSfield{The notes}{The converse is called the Adams conjecture: for every finite complex, some \(N\) makes \(k^N(\psi^k-1)\xi\) J-trivial for all \(k\).}
\GSfield{Statement}{For a finite CW-complex \(X\), \(x\in KO(X)\), and an integer \(k\ge2\), there is \(N\ge0\) with \(k^N J(\psi^k x-x)=0\). The analogous complex statement uses underlying real spherical fibrations. A common exponent depending only on \(X\) can be taken: the J-image is finite, and each displayed difference has order supported only on primes dividing \(k\); any \(N\) at least the largest prime-adic valuation of its group exponent kills all such differences.}
\GSfield{Reconstruction}{Virtual K-classes, stable J-equivalence and nonnegative exponent are supplied. Uniformity in k follows from the finite J-group and does not repair the false reverse formula. The uniform-exponent deduction is a finite-group check \ours: if the exponent of J(X) is \(\prod_p p^{a_p}\), take \(N=\max_p a_p\); only p-components with \(p\mid k\) can remain. This was checked prime by prime, including k divisible by more than one prime.}
\GSfield{Fate}{Sullivan’s proof circulated in \citet[Corollary 5.9 and pp.~165--166]{sullivan1970geometric}, June 1970. The independent proof \citet[Theorem 1.1]{quillen1971adams} was received 3 July 1970 and published 1971. The expanded Sullivan treatment was published in 1974, so it should not be dated as a new post-1973 solution.}
\GSfield{Present status}{As of 9 November 2026: proved; the notebook uses the historical name conjecture for a theorem already available in 1970--1971.}
\GSfield{Assessment}{literal statement: correct with missing hypotheses; publication match: exact; verification: read the result; development: annotated; confidence: high; problem \hyperref[prob:GP-1972-304]{GP-1972-304}}
\end{GSitem}
\begin{GSitem}{1972-295.7}{E}{Naturality gives classifying maps for the operations}
\GSfield{Kind; source}{construction; \GSpage{P000493}}
\GSfield{The notes}{\(\psi^k(f^*\eta)=f^*\psi^k\eta\); thus the operation is specified on the universal bundle over \(BU(n)\). The equation \(\psi^k(1)=1\) permits compatible stabilization.}
\GSfield{Statement}{On reduced complex K-theory the natural operation is represented by a map \(BU\to BU\). Its restriction to \(BU(n)\) classifies the virtual class \(\psi^k(\gamma_n)-n\). Compatibility with adding a trivial line follows from additivity and \(\psi^k[1]=[1]\). This yields a stable classifying map, not automatically a lift \(BU(n)\to BU(n)\) for each \(n\).}
\GSfield{Reconstruction}{The rank-zero shift is needed because BU represents reduced K-theory; the unrestricted group is represented by Z times BU.}
\GSfield{Fate}{The representing-space construction and the stable versus unstable distinction are explicit in \citet[pp.~161--164, Proposition 5.8]{sullivan1970geometric}. This repeats the established operation represented in 1971-259.2.}
\GSfield{Present status}{As of 9 November 2026: established with the hypotheses and distinctions stated here.}
\GSfield{Assessment}{literal statement: correct with missing hypotheses; publication match: exact; verification: read the result; development: none; confidence: high}
\end{GSitem}
\begin{GSitem}{1972-295.8}{E}{Filtered equivalences preserve the intrinsic spherical fibration}
\GSfield{Kind; source}{claim; \GSpage{P000493}}
\GSfield{The notes}{If an operation \(A:BU\to BU\) preserves dimension, the notes say \(A\xi\sim_J\xi\); they also say that such operations do not occur.}
\GSfield{Statement}{The valid sufficient condition is a compatible family of homotopy self-equivalences \(A_n:BU(n)\to BU(n)\), preserving the stabilization filtration. The universal spherical fibration is intrinsic to this filtration, so its stable fibre-homotopy type is fixed. Merely preserving rank, or merely having compatible self-maps, does not suffice: the operation sending a bundle to the trivial bundle of the same rank satisfies both weak descriptions and need not preserve its J-class. Identity and conjugation already refute a literal assertion that no such operations exist.}
\GSfield{Reconstruction}{Automorphism rather than arbitrary map is the missing hypothesis. The parenthetical word orientations is not adopted as a definition: preservation of filtration and orientation of a sphere bundle are different data.}
\GSfield{Fate}{\citet[p.~71, theorem and corollary]{sullivan1974genetics} explicitly requires filtered automorphisms. The earlier \citet[section 4, inertia lemma; Corollary 5.9]{sullivan1970geometric} uses filtered homotopy equivalences. Proposition 5.8 explains the actual obstruction to extending all integral Adams operations unstably.}
\GSfield{Present status}{As of 9 November 2026: established with the hypotheses and distinctions stated here.}
\GSfield{Assessment}{literal statement: false as written; publication match: partial; verification: read the result; development: annotated; confidence: high}
\end{GSitem}
\begin{GSitem}{1972-295.9}{E}{The universal sphere bundle in the BU filtration}
\GSfield{Kind; source}{construction; \GSpage{P000493-P000494}}
\GSfield{The notes}{The inclusion \(BU(n)\to BU(n+1)\) has spheres as fibres; \(U(n)\) acts on \(S^{2n-1}\).}
\GSfield{Statement}{The unit-sphere bundle of the universal complex n-plane bundle has total space homotopy equivalent to \(BU(n-1)\), giving \(S^{2n-1}\to BU(n-1)\to BU(n)\). Equivalently the next stabilization has fibre \(S^{2n+1}\). These two indices should not be conflated. The unitary action is the standard action on the unit sphere in \(\mathbb C^n\).}
\GSfield{Reconstruction}{The scan leaves the first exponent sign uncertain; mathematical context selects \(2n-1\) for \(U(n)\), while the separate stabilization fibre is correctly \(2n+1\).}
\GSfield{Fate}{\citet[pp.~68--69]{sullivan1974genetics} uses precisely this identification; \citet[p.~166]{sullivan1970geometric} explains the analogous orthogonal filtration and the intrinsic-fibration argument.}
\GSfield{Present status}{As of 9 November 2026: established with the hypotheses and distinctions stated here.}
\GSfield{Assessment}{literal statement: correct with missing hypotheses; publication match: exact; verification: read the result; development: none; confidence: high}
\end{GSitem}
\begin{GSitem}{1972-295.10}{E}{The stable spherical classifying space and the J-map}
\GSfield{Kind; source}{definition; \GSpage{P000493}}
\GSfield{The notes}{\(B_F\) is labelled the classifying space of bundles with a homotopy sphere as fibre; J is a functor.}
\GSfield{Statement}{Let \(G_m\) be the topological monoid of self-homotopy equivalences of \(S^{m-1}\). Its classifying space parametrizes spherical fibrations up to fibre-homotopy equivalence. Stabilizing by fibrewise join gives \(BG\), and associating a sphere bundle to a vector bundle gives the natural map \(BO\to BG\), or \(BU\to BG\) after realification, representing J. The notebook’s \(B_F\) is identified with this stable classifying space.}
\GSfield{Reconstruction}{The stable model and monoid, not a Lie group of linear transformations, are specified. A homotopy-sphere fibre means a space with sphere homotopy type.}
\GSfield{Fate}{\citet[section 1, p.~69]{quillen1971adams} specifies \(\operatorname{Sph}(X)=\mathop{\mathrm{colim}}_m[X,BG_m]\); \citet[section 4]{sullivan1974genetics} develops the localized and completed analogues.}
\GSfield{Present status}{As of 9 November 2026: established with the hypotheses and distinctions stated here.}
\GSfield{Assessment}{literal statement: correct with missing hypotheses; publication match: exact; verification: read the result; development: none; confidence: high}
\end{GSitem}
\begin{GSitem}{1972-295.11}{E}{Completed spheres and their self-equivalence monoids}
\GSfield{Kind; source}{construction; \GSpage{P000494}}
\GSfield{The notes}{Sullivan replaces \(S^{2n-1}\) by \(\widehat S^{2n-1}\) and introduces a group \(G_n\), preserving fibrewise homotopy type.}
\GSfield{Statement}{Finite or prime completion can be applied fibrewise to spherical fibrations. Fibrations with a completed-sphere fibre are classified by the classifying space of its homotopy self-equivalence monoid. The completed universal fibration over \(\widehat{BU(n)}\) pulls back to the fibrewise completion of the ordinary one. Completed and ordinary fibrations must be distinguished; after stabilization the oriented completed theory recovers the ordinary stable oriented spherical theory.}
\GSfield{Reconstruction}{The manuscript does not define its \(G_n\). The homotopy self-equivalence monoid is a precise model supporting the displayed classifying-space argument, not a recovered literal definition.}
\GSfield{Fate}{\citet[section 4, especially Corollary 4, pp.~65--66; p.~69]{sullivan1974genetics} gives this comparison. It was already part of \citet[section 4 and Corollary 5.9]{sullivan1970geometric}.}
\GSfield{Present status}{As of 9 November 2026: established with the hypotheses and distinctions stated here.}
\GSfield{Assessment}{literal statement: correct with missing hypotheses; publication match: partial; verification: read the result; development: none; confidence: high}
\end{GSitem}
\begin{GSitem}{1972-295.12}{E}{Sullivan completion and its universal mapping property}
\GSfield{Kind; source}{theorem; \GSpage{P000494}}
\GSfield{The notes}{For simply connected complexes with finitely generated homotopy groups there is a functorial \(X\to\widehat X\), through which maps to spaces with finite homotopy groups factor.}
\GSfield{Statement}{For a simply connected CW-complex whose homotopy groups are finitely generated, Sullivan finite completion gives a natural homotopy type \(\widehat X\). For a finite Postnikov target \(F\) with finite homotopy groups, precomposition gives a bijection \([\widehat X,F]\to[X,F]\); the more general formulation uses Sullivan’s compact homotopy functors. Uniqueness is up to homotopy, not uniqueness of a strict map between chosen point-set models.}
\GSfield{Reconstruction}{The homotopy category and the finite Postnikov test objects make the universal claim precise. No naive inverse limit in the category of all ordinary spaces is asserted.}
\GSfield{Fate}{\citet[Theorem 3.1 and corollary, p.~44]{sullivan1974genetics} verifies the claim for good homotopy groups, including all groups in the notebook’s simply connected finite-type hypothesis. The construction already appears in \citet[section 3]{sullivan1970geometric}.}
\GSfield{Present status}{As of 9 November 2026: established with the hypotheses and distinctions stated here.}
\GSfield{Assessment}{literal statement: correct with missing hypotheses; publication match: exact; verification: read the result; development: none; confidence: high}
\end{GSitem}
\begin{GSitem}{1972-295.13}{E}{Completion of homotopy groups and the product over primes}
\GSfield{Kind; source}{theorem; \GSpage{P000494}}
\GSfield{The notes}{\(\pi_i(\widehat X)=\widehat{\pi_i(X)}\); the groups are profinite and compact, and \(\widehat{\mathbb Z}\) is described through integral p-adic numbers.}
\GSfield{Statement}{Under the preceding simply connected finite-generation assumptions, \(\pi_i(\widehat X)\cong\widehat{\pi_i(X)}\) for \(i\ge2\), with its profinite topology. Finite completion of \(\mathbb Z\) is \(\widehat{\mathbb Z}=\varprojlim_m\mathbb Z/m\cong\prod_p\mathbb Z_p\). Completion at one specified prime instead gives \(\mathbb Z_p\).}
\GSfield{Reconstruction}{All primes versus one prime is an explicit distinction. The notebook says described by p-adic integers and does not literally write the incorrect equality \(\widehat{\mathbb Z}=\mathbb Z_p\).}
\GSfield{Fate}{\citet[Theorem 3.1 and its proof, pp.~44--46]{sullivan1974genetics} proves the homotopy-group assertion; the prime factors are explained in its section 3 and in \citet[section 3]{sullivan1970geometric}.}
\GSfield{Present status}{As of 9 November 2026: established with the hypotheses and distinctions stated here.}
\GSfield{Assessment}{literal statement: correct with missing hypotheses; publication match: exact; verification: read the result; development: none; confidence: high}
\end{GSitem}
\begin{GSitem}{1972-295.14}{E}{Finite stable spherical homotopy groups allow extension of J}
\GSfield{Kind; source}{theorem; \GSpage{P000494}}
\GSfield{The notes}{All homotopy groups of the spherical classifying space are finite because they are stable homotopy groups of spheres.}
\GSfield{Statement}{For the stable unoriented spherical classifying space \(BG\), \(\pi_1(BG)=\mathbb Z/2\) and \(\pi_i(BG)=\pi_{i-1}^{S}\) for \(i\ge2\); these groups are finite. Its oriented cover \(BSG\) is simply connected with finite homotopy groups. Consequently J extends over finite completion of BU, as used in Sullivan’s proof. The assertion does not include the degree-zero stable stem \(\pi_0^S=\mathbb Z\), nor does it identify the fundamental group with that stem.}
\GSfield{Reconstruction}{The shifted positive stems and the exceptional fundamental group are specified; this is why the rank-free stable spherical target is the relevant one.}
\GSfield{Fate}{\citet[p.~67, extension diagram, and Corollary 4 on p.~65]{sullivan1974genetics}; the same factorization is displayed in \citet[pp.~164--166]{sullivan1970geometric}.}
\GSfield{Present status}{As of 9 November 2026: established with the hypotheses and distinctions stated here.}
\GSfield{Assessment}{literal statement: correct with missing hypotheses; publication match: exact; verification: read the result; development: none; confidence: high}
\end{GSitem}
\begin{GSitem}{1972-295.15}{E}{The Adams operations that become unstable Galois symmetries}
\GSfield{Kind; source}{theorem; \GSpage{P000493-P000494}}
\GSfield{The notes}{The operations extend to completed Grassmannians and preserve dimension.}
\GSfield{Statement}{For a prime \(p\), an Adams operation \(\psi^k\) with \(p\nmid k\) is the stable shadow of a compatible Galois homotopy symmetry of the p-completed Grassmannians. On full profinite K-theory Sullivan uses the automorphism equal to \(\psi^k\) on p-components with \(p\nmid k\) and to the identity on components with \(p\mid k\). The ordinary \(\psi^k\) for \(|k|>1\) is not itself a profinite automorphism; the omitted prime restriction matters.}
\GSfield{Reconstruction}{The isomorphic part of the operation, rather than the entire completed endomorphism, is supplied. The notebook suppresses precisely this arithmetic restriction.}
\GSfield{Fate}{\citet[Theorem 5.7, Proposition 5.8 and Corollary 5.9, pp.~163--166]{sullivan1970geometric} gives both the symmetry and an obstruction to unrestricted unstable extension. See also \citet[pp.~8--9 and 66--69]{sullivan1974genetics}.}
\GSfield{Present status}{As of 9 November 2026: established with the hypotheses and distinctions stated here.}
\GSfield{Assessment}{literal statement: false as written; publication match: partial; verification: read the result; development: annotated; confidence: high}
\end{GSitem}
\begin{GSitem}{1972-295.16}{E}{Grassmannians over the integers and algebraic Galois action}
\GSfield{Kind; source}{construction; \GSpage{P000494}}
\GSfield{The notes}{Grassmannians have finitely many algebraic charts over \(\mathbb Z\), so \(\operatorname{Gal}(\mathbb C/\mathbb Q)\) acts, discontinuously in the ordinary topology but algebraically in the Zariski and étale settings.}
\GSfield{Statement}{The Grassmannian \(\operatorname{Gr}(n,N)\) is a scheme over \(\mathbb Z\). A field automorphism of \(\mathbb C\) fixing \(\mathbb Q\) acts semilinearly on its complex points and on its algebraic covering data. This need not be continuous in the analytic topology, but it induces an automorphism of the geometric étale homotopy type and hence its profinite realization. These actions are compatible with the Grassmannian inclusions.}
\GSfield{Reconstruction}{Semilinear action on a scheme and on its covering category replaces a misleading assumption of analytic continuity.}
\GSfield{Fate}{\citet[section 5, especially Theorem 5.7]{sullivan1970geometric}; the argument from rational defining equations to completed homotopy action is explicit in \citet[pp.~68--69]{sullivan1974genetics}. It repeats the Galois-action theme 1971-262.2.}
\GSfield{Present status}{As of 9 November 2026: established with the hypotheses and distinctions stated here.}
\GSfield{Assessment}{literal statement: correct with missing hypotheses; publication match: exact; verification: read the result; development: none; confidence: high}
\end{GSitem}
\begin{GSitem}{1972-295.17}{E}{Cyclotomic abelianization of the arithmetic action}
\GSfield{Kind; source}{theorem; \GSpage{P000494}}
\GSfield{The notes}{\(\operatorname{Gal}(\mathbb C/\mathbb Q)\to\operatorname{Gal}(\mathbb Q_{ab}/\mathbb Q)=\widehat{\mathbb Z}^{\times}\).}
\GSfield{Statement}{Restricting to \(\overline{\mathbb Q}\) and acting on all roots of unity gives the cyclotomic character into \(\widehat{\mathbb Z}^{\times}\). By the Kronecker--Weber theorem it identifies the abelianization of the absolute Galois group of \(\mathbb Q\). Sullivan proves that its action on stable profinite K-theory factors through this character. Such factorization is not asserted for every finite Grassmannian’s whole unstable homotopy type.}
\GSfield{Reconstruction}{The absolute algebraic Galois group is distinguished from the much larger automorphism group of C, and stable versus unstable abelianization is kept separate.}
\GSfield{Fate}{\citet[Theorem 5.7 and following discussion, pp.~163--175]{sullivan1970geometric} states the abelianization on K-theory and explicitly separates the harder unstable question. The manuscript only needs the stable factorization.}
\GSfield{Present status}{As of 9 November 2026: established with the hypotheses and distinctions stated here.}
\GSfield{Assessment}{literal statement: correct with missing hypotheses; publication match: exact; verification: read the result; development: none; confidence: high}
\end{GSitem}
\begin{GSitem}{1972-295.18}{E}{Analytic and étale comparison has finite coefficients}
\GSfield{Kind; source}{theorem; \GSpage{P000494}}
\GSfield{The notes}{\(H_i(V_c,A)\approx H_i(V_{et},A)\) is written without conditions on A.}
\GSfield{Statement}{For a complex algebraic variety of finite type and a finite locally constant coefficient system, analytic and étale cohomology agree under comparison. In the smooth Grassmannian case used here, the profinite analytic homotopy type agrees with the profinite étale homotopy type. With finite constant abelian coefficients the corresponding homology statement follows as well. An integral or rational comparison without completion is not supplied by this theorem.}
\GSfield{Reconstruction}{The notebook uses a homological subscript; the standard cohomological comparison and its finite-coefficient homological consequence are distinguished. Its missing coefficient restriction is essential.}
\GSfield{Fate}{\citet[section 12, Comparison theorems]{artinmazur1969etale} is the original reference; this batch checked its publisher metadata. The precise finite-coefficient comparison needed here is read in \citet[p.~68, formula (*) and following paragraph]{sullivan1974genetics} and in the prior \citet[section 5]{sullivan1970geometric}.}
\GSfield{Present status}{As of 9 November 2026: established with the hypotheses and distinctions stated here.}
\GSfield{Assessment}{literal statement: correct with missing hypotheses; publication match: exact; verification: read the result; development: none; confidence: high}
\end{GSitem}
\begin{GSitem}{1972-295.19}{E}{Nerves recover a profinite homotopy type}
\GSfield{Kind; source}{construction; \GSpage{P000494}}
\GSfield{The notes}{\(\widehat V_c\sim\varprojlim N(V_{et})\), with N called a nerve.}
\GSfield{Statement}{Étale hypercoverings supply a cofiltered system of simplicial homotopy types, the Artin--Mazur étale homotopy type. Completing the system and taking a Sullivan homotopy realization recovers the profinite analytic type in the comparison setting. The displayed inverse limit is shorthand in a pro-homotopy or compact-functor category; taking a point-set inverse limit of arbitrary ordinary covering nerves is not the theorem.}
\GSfield{Reconstruction}{Hypercoverings, profinite completion and the category of the limit are the reconstructed data. This repeats 1971-262.1--.3 rather than creating a new open problem.}
\GSfield{Fate}{\citet[sections 8--9 and 12]{artinmazur1969etale} names the original construction; its role in Sullivan’s realization is read in \citet[section 3, p.~34; p.~68 and footnote]{sullivan1974genetics}. The construction was already available in \citet[section 5]{sullivan1970geometric}.}
\GSfield{Present status}{As of 9 November 2026: established with the hypotheses and distinctions stated here.}
\GSfield{Assessment}{literal statement: correct with missing hypotheses; publication match: exact; verification: read the result; development: none; confidence: high}
\end{GSitem}

\GStalk{296}{1973-03-26}{I. M. Gelfand}{P000495}{Atiyah--Patodi--Singer spectral asymmetry}
Gelfand announces the eta series, its regular value at zero, and the shifted-circle example, then suggests more general character-weighted spectral series. The 1973 APS announcement was received eleven days before this seminar. The later proofs require a chronology distinction: APS established the geometric cases and odd-dimensional general theorem, while Gilkey completed general even-dimensional regularity in 1981.

\begin{GSitem}{1972-296.1}{E}{Regularity of an elliptic zeta function at zero}
\GSfield{Kind; source}{theorem; \GSpage{P000495}}
\GSfield{The notes}{For a self-adjoint elliptic A, \(\zeta_A(s)=\operatorname{Tr}A^{-s}=\sum\lambda_i^{-s}\) has no pole at zero (Seeley).}
\GSfield{Statement}{A sufficient unambiguous setting is a positive invertible classical elliptic pseudodifferential operator A of order m>0 on a vector bundle over a closed smooth d-manifold. The spectral sum converges for \(\operatorname{Re}s>d/m\), extends meromorphically, and is regular at zero. For nonnegative A omit the zero eigenvalues; for general nonpositive A a complex-power spectral-cut convention is needed.}
\GSfield{Fate}{Known before the seminar: \citet{seeley1968powers}; the exact zeta statement and the null-space modification are recalled in \citet[\S2, pp.~73--74]{atiyahpatodisinger1976spectral3}.}
\GSfield{Present status}{As of 9 November 2026: Established. Self-adjointness alone does not define the written complex powers without handling negative and zero eigenvalues.}
\GSfield{Assessment}{literal statement: correct with missing hypotheses; publication match: exact; verification: read the result; development: none; confidence: high}
\end{GSitem}
\begin{GSitem}{1972-296.2}{B}{The eta function of a self-adjoint elliptic operator}
\GSfield{Kind; source}{definition; \GSpage{P000495}}
\GSfield{The notes}{\(\eta_A(s)=\sum_{\lambda\ne0}\operatorname{sign}(\lambda)|\lambda|^{-s}\), emphasized for a first-order operator.}
\GSfield{Statement}{For a self-adjoint classical elliptic operator A of positive order m on a Hermitian bundle over a closed smooth d-manifold, take eigenvalues with multiplicity and omit the kernel. The series converges absolutely for \(\operatorname{Re}s>d/m\) and has meromorphic continuation. In that half-plane it equals \(\operatorname{Tr}(A|A|^{-s-1})\).}
\GSfield{Fate}{The new invariant was announced in \citet{atiyahpatodisinger1973spectral}; its geometric role is developed in \citet[Introduction and \S4]{atiyahpatodisinger1975spectral1}. Class B records development by the reported authors, not by the seminar speaker.}
\GSfield{Present status}{As of 9 November 2026: Established eta-function definition; the boundary index theorem is its published fate, not an extra theorem recorded on this page.}
\GSfield{Assessment}{literal statement: correct with missing hypotheses; publication match: exact; verification: read the result; development: none; confidence: high}
\end{GSitem}
\begin{GSitem}{1972-296.3}{B}{The eta function has no pole at zero: announcement and proof scope}
\GSfield{Kind; source}{theorem; \GSpage{P000495}}
\GSfield{The notes}{\ru{Трудная теор.}: zero is not a pole, stated without a dimension or operator subclass restriction.}
\GSfield{Statement}{For every self-adjoint classical elliptic pseudodifferential operator of positive order on a closed smooth manifold, the meromorphic eta function is regular at s=0. This is a statement about the integrated spectral function; its local residue density need not vanish pointwise.}
\GSfield{Reconstruction}{The closed-manifold and classical-operator hypotheses are supplied. The later proof gap is not projected back into the notebook as a question explicitly asked there.}
\GSfield{Fate}{APS announced the unrestricted assertion in 1973 but explicitly reported the remaining even-dimensional gap in \citet[Introduction, pp.~71--72; Theorem 4.5]{atiyahpatodisinger1976spectral3}. The general result is \citet[Theorem 0.1, pp.~290--291]{gilkey1981residue}. The geometric Dirac-type case already occurs in \citet[Theorem 4.2]{atiyahpatodisinger1975spectral1}.}
\GSfield{Present status}{As of 9 November 2026: Established in full stated generality by the later completion; it would be historically inaccurate to attribute that full proof to the 1973 announcement.}
\GSfield{Assessment}{literal statement: correct with missing hypotheses; publication match: exact; verification: read the result; development: none; confidence: high}
\end{GSitem}
\begin{GSitem}{1972-296.4}{B}{The shifted circle eta invariant}
\GSfield{Kind; source}{claim; \GSpage{P000495}}
\GSfield{The notes}{For \(A=i\,d/dx+t\) on periodic functions on \([0,2\pi]\), eigenvalues are \(n+t\), and \(\eta_A(0)=1-2t\).}
\GSfield{Statement}{For real t, the spectrum is \(\mathbb Z+t\). If \(a=\{t\}\in(0,1)\), then \(\eta_A(s)=\zeta(s,a)-\zeta(s,1-a)\), so \(\eta_A(0)=1-2a\). If t is integral, omitting the zero eigenvalue gives \(\eta_A(0)=0\). Thus the notebook formula is correct only on the branch 0<t<1.}
\GSfield{Reconstruction}{\ours\ The Hurwitz-zeta derivation and the integral-t correction are supplied explicitly; shifting t by 1 independently checks the required periodicity.}
\GSfield{Fate}{This is the elementary model of spectral asymmetry in the APS programme; the evaluation uses \(\zeta(0,a)=1/2-a\), checked in \citet[\S25.11]{nist2026dlmf}.}
\GSfield{Present status}{As of 9 November 2026: Established branch formula with integer jumps; the unrestricted affine formula cannot be correct.}
\GSfield{Assessment}{literal statement: false as written; publication match: partial; verification: read the result; development: annotated; confidence: high}
\end{GSitem}
\begin{GSitem}{1972-296.5}{B}{Spectral asymmetry is a regularized difference}
\GSfield{Kind; source}{remark; \GSpage{P000495}}
\GSfield{The notes}{\(\eta(0)\) is described by \(\sum_{\lambda>0}1-\sum_{\lambda<0}1\).}
\GSfield{Statement}{The displayed difference means evaluation at s=0 after meromorphic continuation of \(\sum\operatorname{sign}(\lambda)|\lambda|^{-s}\), not subtraction of two infinite dimensions. The reduced invariant is often \((\eta_A(0)+\dim\ker A)/2\), considered modulo integers, but that extra normalization is not given in the notebook.}
\GSfield{Fate}{The regularized interpretation is central to \citet[Introduction]{atiyahpatodisinger1975spectral1}; dependence on spectral flow connects thematically to the earlier GP-1971-04 thread without identifying the two questions.}
\GSfield{Present status}{As of 9 November 2026: Established regularized invariant, not an ordinary signed count.}
\GSfield{Assessment}{literal statement: correct with missing hypotheses; publication match: exact; verification: read the result; development: none; confidence: high}
\end{GSitem}
\begin{GSitem}{1972-296.6}{F}{Character-weighted spectral L-series}
\GSfield{Kind; source}{question; \GSpage{P000495}}
\GSfield{The notes}{The general problem proposes summing multiplicative characters of eigenvalues and remarks that complex numbers have more complicated characters.}
\GSfield{Statement}{The notebook does not specify the operator class, spectral cut, allowed characters, convergence factor, evaluation point, or invariant to be constructed. Its real sign-character example is identifiable, but a unique complex-spectrum theorem or definite open question is not.}
\GSfield{Fate}{The APS papers supply the eta example, not a result with these unspecified complex-character quantifiers. Searching their announcements, later parts and the spectral-asymmetry literature did not identify an exact completion of this sentence.}
\GSfield{Present status}{As of 9 November 2026: Programme recognizable; a definite problem suitable for the problem index is not recoverable.}
\GSfield{Assessment}{literal statement: not determinable; publication match: thematic only; verification: unresolved lead; development: none; confidence: low}
\end{GSitem}

\GStalk{298}{1973-03-26}{S. I. Gelfand}{P000495-P000497}{p-adic distributions, integration and the zeta function}
The talk introduces finitely additive p-adic distributions on residue classes, contrasts the unbounded invariant distribution with a bounded floor-function construction, and gives a moderate-growth integration criterion. It then sketches interpolation of negative odd zeta values. The missing Euler factor, weight branch and possible pole make the final unqualified continuity assertion false. The source places this continuation before Gabrielov on P000497, although the catalogue numbers 297 and 298 suggest the reverse order.

\begin{GSitem}{1972-298.1}{E}{Distributions on residue classes}
\GSfield{Kind; source}{definition; \GSpage{P000496}}
\GSfield{The notes}{For \(A_m=p^m\mathbb Z_p\), the notes require \(\mu(a+A_m)=\sum_{j=0}^{p-1}\mu(a+jp^m+A_{m+1})\).}
\GSfield{Statement}{A K-valued distribution on \(\mathbb Z_p\), for a complete non-Archimedean extension K of \(\mathbb Q_p\), is a finitely additive function on compact open sets. It is determined by its values on residue classes, subject to the displayed refinement identity and independence of the representative a. A bounded distribution is a measure in the Banach-space sense, with \(\sup_{m,a}|\mu(a+p^m\mathbb Z_p)|<\infty\).}
\GSfield{Fate}{The distinction is made explicitly in \citet[\S8.4]{manin1973periods} and \citet[\S1]{manin1976nonarchimedean}; the underlying distribution formalism predates these lectures.}
\GSfield{Present status}{As of 9 November 2026: Standard definition. Finite additivity alone does not imply a bounded functional on all continuous functions.}
\GSfield{Assessment}{literal statement: correct with missing hypotheses; publication match: exact; verification: read the result; development: none; confidence: high}
\end{GSitem}
\begin{GSitem}{1972-298.2}{E}{The normalized invariant distribution is unbounded}
\GSfield{Kind; source}{construction; \GSpage{P000496}}
\GSfield{The notes}{Translation invariance gives \(\mu(a+A_m)=p^{-m}\); the measure of a small set is large in the p-adic norm.}
\GSfield{Statement}{Among translation-invariant \(\mathbb Q_p\)-valued distributions with total mass 1, refinement forces \(\mu(a+p^m\mathbb Z_p)=p^{-m}\). Its norm on a level-m class is \(p^m\), so it is unbounded and does not define a continuous integration functional on all of \(C(\mathbb Z_p,\mathbb Q_p)\). Without total-mass normalization any constant multiple is possible.}
\GSfield{Fate}{This is the elementary invariant-distribution example recalled in \citet[\S8.4]{manin1973periods}.}
\GSfield{Present status}{As of 9 November 2026: Established. The missing total-mass condition is necessary for the phrase only possibility.}
\GSfield{Assessment}{literal statement: correct with missing hypotheses; publication match: exact; verification: read the result; development: none; confidence: high}
\end{GSitem}
\begin{GSitem}{1972-298.3}{E}{The bounded floor-function distribution}
\GSfield{Kind; source}{construction; \GSpage{P000496}}
\GSfield{The notes}{For positive integer c prime to p, \(\mu_c(a+A_m)=\lfloor ca/p^m\rfloor-(c-1)/2\), with \(0\leq a<p^m\).}
\GSfield{Statement}{Using the canonical representative \(0\leq a<p^m\), this formula defines a bounded finitely additive \(\mathbb Q_p\)-valued distribution. For odd p its values are p-integral; for p=2 the condition c odd again makes every value integral. The representative convention is essential: the unrestricted floor expression is not periodic in a.}
\GSfield{Reconstruction}{\ours\ Refinement reduces to \(\sum_{j=0}^{p-1}\lfloor u+cj/p\rfloor=\lfloor pu\rfloor+(c-1)(p-1)/2\), with \(u=ca/p^{m+1}\). Permuting the residues cj modulo p proves the identity. An independent exact-integer check over several primes and c values is supplied.}
\GSfield{Fate}{\citet[\S8.4, Example 2]{manin1973periods} gives this construction and attributes it to Mazur. It belongs to the pre-existing measure approach to p-adic zeta functions, not a new problem posed here.}
\GSfield{Present status}{As of 9 November 2026: Established bounded distribution; the historical label in the notes suppresses the Mazur regularization step.}
\GSfield{Assessment}{literal statement: correct with missing hypotheses; publication match: exact; verification: read the result; development: annotated; confidence: high}
\end{GSitem}
\begin{GSitem}{1972-298.4}{E}{The two-valued distribution for c equal to two}
\GSfield{Kind; source}{construction; \GSpage{P000496}}
\GSfield{The notes}{\(\mu_2\) is \(+1/2\) above the midpoint and \(-1/2\) below it.}
\GSfield{Statement}{For an odd prime p and \(0\leq a<p^m\), \(\mu_2(a+p^m\mathbb Z_p)=-1/2\) if \(2a<p^m\) and \(+1/2\) if \(2a>p^m\). Equality is impossible. This example does not use c=2 at p=2, where c would not be prime to p.}
\GSfield{Fate}{This is the c=2 specialization of \citet[\S8.4]{manin1973periods}.}
\GSfield{Present status}{As of 9 November 2026: Established, with the odd-prime restriction restored.}
\GSfield{Assessment}{literal statement: correct with missing hypotheses; publication match: exact; verification: read the result; development: none; confidence: high}
\end{GSitem}
\begin{GSitem}{1972-298.5}{E}{Riemann sums and the bounded case}
\GSfield{Kind; source}{method; \GSpage{P000496}}
\GSfield{The notes}{For a representative set \(R_m\), form \(S(R_m)=\sum_{a\in R_m}f(a)\mu(a+A_m)\), and seek a representative-independent limit.}
\GSfield{Statement}{If \(\mu\) is bounded and f is continuous on compact \(\mathbb Z_p\), these sums converge in K, independently of the representatives. Uniform continuity and the ultrametric triangle inequality bound any refinement error by \(\|\mu\|\) times the oscillation of f on residue classes. Locally constant f requires only finite additivity: the sums eventually stabilize.}
\GSfield{Fate}{This is standard non-Archimedean integration, formulated in \citet[\S1]{manin1976nonarchimedean}.}
\GSfield{Present status}{As of 9 November 2026: Established. The distinction between locally constant and continuous integrands matters for unbounded distributions.}
\GSfield{Assessment}{literal statement: correct with missing hypotheses; publication match: exact; verification: read the result; development: none; confidence: high}
\end{GSitem}
\begin{GSitem}{1972-298.6}{B}{Moderate growth suffices for Lipschitz integrands}
\GSfield{Kind; source}{theorem; \GSpage{P000496}}
\GSfield{The notes}{\(p^{-m}\max_a|\mu(a+A_m)|_p\to0\); a Lipschitz condition on f then gives the integral.}
\GSfield{Statement}{Let K be complete and put \(M_m=\max_a|\mu(a+p^m\mathbb Z_p)|\). If \(p^{-m}M_m\to0\) and \(|f(x)-f(y)|\leq C|x-y|_p\), the Riemann sums converge independently of representatives. A consecutive-level error is at most \(Cp^{-m}M_{m+1}\), which tends to zero; the ultrametric inequality gives the Cauchy property.}
\GSfield{Fate}{This is the criterion in \citet[\S8.6]{manin1973periods}, published after the seminar report of Manin's lectures. The class records development by the reported author.}
\GSfield{Present status}{As of 9 November 2026: Established sufficient condition. It excludes the normalized invariant distribution, whose scaled growth is constantly one.}
\GSfield{Assessment}{literal statement: correct with missing hypotheses; publication match: exact; verification: read the result; development: none; confidence: high}
\end{GSitem}
\begin{GSitem}{1972-298.7}{E}{Negative odd zeta values are rational}
\GSfield{Kind; source}{theorem; \GSpage{P000497}}
\GSfield{The notes}{\(\zeta(1-2k)\) is rational for every positive integer k and is expressed through Bernoulli numbers.}
\GSfield{Statement}{With \(t/(e^t-1)=\sum_{j\geq0}B_jt^j/j!\), \(\zeta(1-2k)=-B_{2k}/(2k)\) for k positive. This is a value of the meromorphic continuation, not of the original divergent Dirichlet series at that argument.}
\GSfield{Fate}{Classical before the seminar; see \citet[\S25.11]{nist2026dlmf} for the Hurwitz-zeta specialization and \citet[\S0.1]{manin1976nonarchimedean} for its interpolation role.}
\GSfield{Present status}{As of 9 November 2026: Established. These values connect the next two arithmetic talks to the Bernoulli discussion on P000500.}
\GSfield{Assessment}{literal statement: correct with missing hypotheses; publication match: exact; verification: read the result; development: none; confidence: high}
\end{GSitem}
\begin{GSitem}{1972-298.8}{E}{The uncorrected zeta sequence is not a continuous function of weight}
\GSfield{Kind; source}{claim; \GSpage{P000497}}
\GSfield{The notes}{For arbitrary positive integers \(k_i\to a\) p-adically, the notes assert \(\zeta(1-2k_i)\to\zeta(a)\) and a continuous function on all \(\mathbb Z_p\).}
\GSfield{Statement}{The Kubota--Leopoldt interpolation uses \((1-p^{n-1})\zeta(1-n)\), with n restricted to a fixed residue class modulo p-1 for odd p, and with the appropriate tame-character branch. The zeta branch has a pole at s=1, so a universal continuous function on all weights is not obtained. The raw values in the notebook do not have the stated continuity property even on a nonpolar branch.}
\GSfield{Reconstruction}{\ours\ For p=5 take \(k_i=1+2\cdot5^i\). Then \(k_i\to1\) and \(2k_i\equiv2\pmod4\). Correct interpolation gives \(\zeta(1-2k_i)\to(1-5)\zeta(-1)=1/3\), since \(5^{2k_i-1}\to0\). The constant sequence k=1 instead gives \(-1/12\). The two limits differ. Exact Bernoulli arithmetic checks this discrepancy and the corrected congruences at small levels.}
\GSfield{Fate}{The theorem of \citet{kubotaleopoldt1964} predates the talk; \citet[\S0.1 and \S0.5]{manin1976nonarchimedean} explicitly explains the Euler factor and regularization. Original 1964 archival metadata was opened; the exact interpolation formula was read in Manin's exposition.}
\GSfield{Present status}{As of 9 November 2026: The notebook statement is false; the properly normalized branch theorem is established. No unresolved interpolation problem is inferred.}
\GSfield{Assessment}{literal statement: false as written; publication match: exact; verification: indirect; development: annotated; confidence: high}
\end{GSitem}
\begin{GSitem}{1972-298.9}{E}{The Mellin transform and its regularizing factors}
\GSfield{Kind; source}{claim; \GSpage{P000496-P000497}}
\GSfield{The notes}{The floor measure is first said to be obtained from a zeta function by a Mellin transform; the next page says the zeta function is the Mellin transform of the measure.}
\GSfield{Statement}{For a bounded measure on \(\mathbb Z_p^\times\), its Mellin transform is the function of a continuous character \(\chi\) given by \(\chi\mapsto\int\chi(a)\,d\mu(a)\). For the regularized Bernoulli measure, this gives a regularized p-adic L-function; the Euler factor and the auxiliary-c regularization factor must be retained or explicitly divided out where allowed. It is not literally the unmodified complex zeta function.}
\GSfield{Fate}{The direction of the transform and the regularization are spelled out in \citet[\S0.4--0.5]{manin1976nonarchimedean}.}
\GSfield{Present status}{As of 9 November 2026: Established measure construction and Mellin transform, with the normalization caveat made explicit.}
\GSfield{Assessment}{literal statement: false as written; publication match: exact; verification: read the result; development: none; confidence: high}
\end{GSitem}

\GStalk{297}{1973-03-26}{A. M. Gabrielov}{P000497}{Dynkin diagrams of isolated function singularities: opening fragment}
Gabrielov’s heading occurs only near the bottom of P000497, after the p-adic material of record 298. The surviving fragment introduces the local Milnor fibre and its bouquet of vanishing spheres. No Dynkin graph, intersection-matrix algorithm, sum formula or additional claim is recorded in this assigned fragment. Gabrielov’s 1973 paper has the matching opening setup, but its later original theorems cannot be inferred from the title alone.

\begin{GSitem}{1972-297.1}{E}{Choosing a Milnor fibre of an isolated critical point}
\GSfield{Kind; source}{construction; \GSpage{P000497}}
\GSfield{The notes}{For a complex analytic f near 0, with an isolated singularity at 0, choose a small ball and then a small \(\delta\); consider \(f(x)=\delta\), \(|x|<r_0\).}
\GSfield{Statement}{Let \(f:(\mathbb C^n,0)\to(\mathbb C,0)\) be holomorphic with an isolated critical point and \(f(0)=0\). Choose a sufficiently small Milnor radius \(r_0>0\), then \(0<|\delta|\ll r_0\) so the level is smooth and transverse to the bounding sphere. Its intersection with the open ball is a complex \((n-1)\)-manifold, the interior of the compact Milnor fibre.}
\GSfield{Reconstruction}{The order of choices, nonzero level and isolated critical-point meaning are made explicit. This repeats 1972-294.6 and 1972-284.11.}
\GSfield{Fate}{The construction is classical in \citet[sections 4--6]{milnor1968singular}. It is repeated in \citet[section 1]{gabrielov1973intersection}, received 29 January 1973 and published later in 1973. The elementary setup, not Gabrielov’s later matrix theorem, is the exact notebook match.}
\GSfield{Present status}{As of 9 November 2026: established with the hypotheses and distinctions stated here.}
\GSfield{Assessment}{literal statement: correct with missing hypotheses; publication match: exact; verification: read the result; development: none; confidence: high}
\end{GSitem}
\begin{GSitem}{1972-297.2}{E}{The bouquet and the reference to vanishing cycles}
\GSfield{Kind; source}{theorem; \GSpage{P000497}}
\GSfield{The notes}{The level is homotopy equivalent to a bouquet of \((n-1)\)-spheres; the final words mention vanishing cycles.}
\GSfield{Statement}{The Milnor fibre is homotopy equivalent to a bouquet of \(\mu\) copies of \(S^{n-1}\), where \(\mu=\dim_{\mathbb C}\mathcal O_{\mathbb C^n,0}/(\partial f)\). After a morsification and a distinguished system of paths to its critical values, the vanishing cycles give an integral basis of reduced middle homology. No canonical distinguished basis, Dynkin diagram or intersection signs are fixed by the fragment alone.}
\GSfield{Reconstruction}{The Milnor number and auxiliary choices for a vanishing-cycle basis are supplied. The uncertain sphere-dimension letter is read \(n\), consistently with the complex dimension \(n-1\).}
\GSfield{Fate}{\citet[Theorems 6.5 and 7.2]{milnor1968singular} gives the bouquet theorem; \citet[section 1]{gabrielov1973intersection} describes the distinguished vanishing-cycle choices. Both results were available before the intervention (the latter manuscript received 29 January 1973).}
\GSfield{Present status}{As of 9 November 2026: the recorded bouquet statement is established; the title alone does not identify an additional research claim.}
\GSfield{Assessment}{literal statement: correct with missing hypotheses; publication match: exact; verification: read the result; development: none; confidence: high}
\end{GSitem}

\GStalk{299}{1973-04-09}{S. I. Gelfand}{P000498-P000499}{Manin’s modular periods and ordinary p-adic L-functions}
Explicitly an account of Yu. I. Manin's lectures, this record resumes integration and then builds measures from Hecke-eigenform periods. It culminates in a primitive-character interpolation formula. The source omits the Hecke weights, writes a wrong automorphy factor, leaves a residue class ambiguous, and combines an odd real-period component with even characters. The published 1973 paper identifies the construction, but the notebook's final prefactor cannot be preserved without choosing period, Gauss-sum and measure conventions.

\begin{GSitem}{1972-299.1}{E}{Vector-valued distributions and clopen sets}
\GSfield{Kind; source}{definition; \GSpage{P000498}}
\GSfield{The notes}{\(\mu(U)\in W\) for finite-dimensional \(W/\mathbb Q_p\); \(\mu(U_1\cup U_2)=\mu(U_1)+\mu(U_2)\); \(A_n=p^n\mathbb Z_p\) is called a basis.}
\GSfield{Statement}{For a finite-dimensional p-adic vector space W, a distribution is additive on disjoint compact open subsets of \(\mathbb Z_p\). The cosets \(a+p^n\mathbb Z_p\), not just the subgroups at zero, form its clopen basis. For all of \(\mathbb Q_p\), allow integer n and centers in \(\mathbb Q_p\). Overlapping sets require the inclusion--exclusion term.}
\GSfield{Fate}{This repeats the standard framework of record 298; compare \citet[\S8.4]{manin1973periods}.}
\GSfield{Present status}{As of 9 November 2026: Established after adding disjointness and translating the neighbourhood basis.}
\GSfield{Assessment}{literal statement: false as written; publication match: exact; verification: read the result; development: none; confidence: high}
\end{GSitem}
\begin{GSitem}{1972-299.2}{E}{The missing measure in the Riemann sum}
\GSfield{Kind; source}{method; \GSpage{P000498}}
\GSfield{The notes}{\(S(R_m)=\sum f(a)(a+A_m)\); boundedness and local constancy are offered as sufficient conditions.}
\GSfield{Statement}{The scalar-vector product is \(S(R_m)=\sum f(a)\mu(a+A_m)\in W\). If f is locally constant the sums stabilize for every finitely additive distribution, without a boundedness assumption. If the distribution is bounded, every continuous f can be integrated by uniform approximation.}
\GSfield{Fate}{This is the integration formalism of \citet[\S8.5--8.6]{manin1973periods}; the notes give a sufficient but unnecessarily restrictive condition.}
\GSfield{Present status}{As of 9 November 2026: Established; the absent mu is a visible omission in the manuscript itself.}
\GSfield{Assessment}{literal statement: false as written; publication match: exact; verification: read the result; development: none; confidence: high}
\end{GSitem}
\begin{GSitem}{1972-299.3}{E}{Repeated floor-function measure}
\GSfield{Kind; source}{construction; \GSpage{P000498}}
\GSfield{The notes}{The formula \(\mu_c(a+A_n)=\lfloor ca/p^n\rfloor-(c-1)/2\) is repeated and finite additivity assigned as a useful exercise.}
\GSfield{Statement}{For positive integer c prime to p and canonical \(0\le a<p^n\), the construction and boundedness are exactly those of 1972-298.3. The same floor-sum permutation proves compatibility with refinement.}
\GSfield{Fate}{This repeats \citet[\S8.4, Example 2]{manin1973periods}. It is indexed as a new occurrence, not a new theorem or a separate research problem.}
\GSfield{Present status}{As of 9 November 2026: Established; linked to the previous occurrence and its exact computation.}
\GSfield{Assessment}{literal statement: correct with missing hypotheses; publication match: exact; verification: read the result; development: none; confidence: high}
\end{GSitem}
\begin{GSitem}{1972-299.4}{E}{Analytic characters and the prime two}
\GSfield{Kind; source}{construction; \GSpage{P000498}}
\GSfield{The notes}{\(\chi_t(a)=e^{t\log a}\), with \(a\equiv1\pmod p\) and \(t\in\mathbb Z_p\).}
\GSfield{Statement}{For odd p this formula defines a character on \(1+p\mathbb Z_p\); at p=2 the displayed exponential formula is safely valid on \(1+4\mathbb Z_2\). A character of \(\mathbb Z_p^\times\) also requires its finite torsion component. The analytic family indexed by t in \(\mathbb Z_p\) does not exhaust finite-order wild characters with values in extensions of \(\mathbb Q_p\).}
\GSfield{Fate}{The full character parametrization is in \citet[\S8.1--8.3]{manin1973periods}.}
\GSfield{Present status}{As of 9 November 2026: Established local analytic family, with the p=2 domain and finite character components restored.}
\GSfield{Assessment}{literal statement: correct with missing hypotheses; publication match: exact; verification: read the result; development: none; confidence: high}
\end{GSitem}
\begin{GSitem}{1972-299.5}{E}{Holomorphic cusp forms and the automorphy factor}
\GSfield{Kind; source}{definition; \GSpage{P000498}}
\GSfield{The notes}{\(\varphi(z)=\sum_{n\ge1}\lambda_ne^{2\pi inz}\), but \(\varphi(gz)=(az+c)^{-w}\varphi(z)\).}
\GSfield{Statement}{A holomorphic level-one cusp form f of integral weight w satisfies \(f((az+b)/(cz+d))=(cz+d)^w f(z)\) for every matrix in \(\mathrm{SL}_2(\mathbb Z)\), is holomorphic on the upper half-plane, and has zero constant Fourier coefficient at the cusp. Nonzero level-one forms have even weight; the first cusp-form weight is 12. The notes' negative exponent would belong to a differently defined slash action, and az+c is not its denominator.}
\GSfield{Fate}{The convention is fixed in \citet[\S1.1]{manin1973periods}; the modular-form theory itself long predates the lecture.}
\GSfield{Present status}{As of 9 November 2026: Established definition; both the matrix entries and the exponent require correction.}
\GSfield{Assessment}{literal statement: false as written; publication match: exact; verification: read the result; development: none; confidence: high}
\end{GSitem}
\begin{GSitem}{1972-299.6}{E}{Hecke operators and the eigenvalue normalization}
\GSfield{Kind; source}{construction; \GSpage{P000499}}
\GSfield{The notes}{\(T_m f=\sum_{ad=m,\,0\le b<d}f((az+b)/d)\), followed by \(T_m f=\lambda_mf\) when \(\lambda_1=1\).}
\GSfield{Statement}{For weight w and the Fourier normalization \(a_1=1\), use
     \[T_m f(z)=m^{w-1}\sum_{\substack{ad=m\\a,d>0}}d^{-w}\sum_{b=0}^{d-1}f((az+b)/d).\]
     A simultaneous normalized eigenform has eigenvalue \(a_m\). In particular \(T_pf=p^{w-1}f(pz)+p^{-1}\sum_bf((z+b)/p)\). The unweighted sum written in the notebook is generally not this modular operator.}
\GSfield{Fate}{Compare the normalization in \citet[\S3, equation (11)]{manin1973periods}; this is standard Hecke theory.}
\GSfield{Present status}{As of 9 November 2026: Established with the weights restored. The normalization determines the quadratic used below.}
\GSfield{Assessment}{literal statement: false as written; publication match: exact; verification: read the result; development: none; confidence: high}
\end{GSitem}
\begin{GSitem}{1972-299.7}{E}{Primitive multiplicative characters}
\GSfield{Kind; source}{definition; \GSpage{P000499}}
\GSfield{The notes}{For conductor \(p^r\), primitivity is described by nontriviality on \(p^{r-1}\mathbb Z/p^r\mathbb Z\).}
\GSfield{Statement}{A Dirichlet character modulo \(p^r\) is a character of the unit group, extended by zero to nonunits. For r at least 2, conductor exactly \(p^r\) means nontriviality on \(1+p^{r-1}\mathbb Z/p^r\mathbb Z\). For r=1 it means a nontrivial character modulo p. The additive subgroup printed in the notes is not contained in the multiplicative unit group.}
\GSfield{Fate}{Primitivity is used precisely to cancel a lower-level sum in \citet[\S10.2]{manin1973periods}.}
\GSfield{Present status}{As of 9 November 2026: Established definition, with the missing 1+ and the conductor-one exception separated.}
\GSfield{Assessment}{literal statement: false as written; publication match: exact; verification: read the result; development: none; confidence: high}
\end{GSitem}
\begin{GSitem}{1972-299.8}{E}{Twisting the cusp form changes its level}
\GSfield{Kind; source}{construction; \GSpage{P000499}}
\GSfield{The notes}{\(f_\chi(z)=\sum a_n\chi(n)e^{2\pi inz}\) is also modular.}
\GSfield{Statement}{For a level-one cusp form of weight w and a primitive Dirichlet character of conductor q, the twist is a cusp form of weight w on \(\Gamma_0(q^2)\) with nebentypus \(\chi^2\) (its minimal level may be described more sharply). It is not generally a form for the original full modular group.}
\GSfield{Fate}{The twist is the input of \citet[\S10]{manin1973periods}; the passage to an appropriate congruence group is a standard earlier construction.}
\GSfield{Present status}{As of 9 November 2026: Established; the unnamed level and character must be part of the statement.}
\GSfield{Assessment}{literal statement: correct with missing hypotheses; publication match: exact; verification: read the result; development: none; confidence: high}
\end{GSitem}
\begin{GSitem}{1972-299.9}{E}{The Mellin integral and its constant}
\GSfield{Kind; source}{theorem; \GSpage{P000499}}
\GSfield{The notes}{\(L(f,\chi,s)=C\int_0^{i\infty}f_\chi(z)z^{s-1}dz=\sum a_n\chi(n)n^{-s}\); the constant at s=1 is overwritten.}
\GSfield{Statement}{For a holomorphic cusp form, initially in a right half-plane,
     \[\int_0^\infty f_\chi(iy)y^{s-1}dy=(2\pi)^{-s}\Gamma(s)L(f,\chi,s).\]
     The cusp conditions give the usual continuation of the completed L-function. At s=1, \(L(f,\chi,1)=-2\pi i\int_0^{i\infty}f_\chi(z)dz\). If the contour form uses \(z^{s-1}\), the chosen branch contributes an additional \(i^s\).}
\GSfield{Fate}{The s=1 relation is used in \citet[\S10.2]{manin1973periods}; the Mellin identity is classical.}
\GSfield{Present status}{As of 9 November 2026: Established. A constant independent of s cannot describe the general Mellin identity.}
\GSfield{Assessment}{literal statement: correct with missing hypotheses; publication match: exact; verification: read the result; development: none; confidence: high}
\end{GSitem}
\begin{GSitem}{1972-299.10}{B}{Algebraic normalized periods}
\GSfield{Kind; source}{theorem; \GSpage{P000499}}
\GSfield{The notes}{After dividing \(\operatorname{Re}\int_0^x f(z)dz\) by a period omega, the values at rational x are algebraic integers in a fixed number field, in many cases.}
\GSfield{Statement}{For a normalized level-one Hecke eigenform with its real Fourier embedding, suitable nonzero even and odd periods normalize the corresponding components of \(I(x)=\int_x^{i\infty}f(z)dz\), x rational, into a fixed number field. After a common scaling the values lie in its ring of integers. The choice of parity and period is part of the theorem; no general assertion of transcendence of every possible period is needed.}
\GSfield{Fate}{This is the algebraicity and integrality input of \citet[Theorem 1.2 and Proposition 9.2(c)]{manin1973periods}. The notebook reports the author's lectures before the paper.}
\GSfield{Present status}{As of 9 November 2026: Established in this level-one setting. A chosen p-adic place of the coefficient field is needed for the later measure.}
\GSfield{Assessment}{literal statement: correct with missing hypotheses; publication match: exact; verification: read the result; development: none; confidence: high}
\end{GSitem}
\begin{GSitem}{1972-299.11}{E}{Parity of the real period conflicts with the even character}
\GSfield{Kind; source}{claim; \GSpage{P000499}}
\GSfield{The notes}{The same passage assumes \(\chi(-1)=1\) but uses \(\ell(x)=\omega^{-1}\operatorname{Re}\int_0^x f(z)dz\).}
\GSfield{Statement}{For real Fourier coefficients, let \(J(x)=\int_0^x f(z)dz\) along a path in the upper half-plane between rational cusps. Then \(J(-x)=-\overline{J(x)}\). Thus its real part is odd and its imaginary part is even. The written real component pairs to zero with every even character on a symmetric residue system. An even-character interpolation must use the even period component, with its matching normalization.}
\GSfield{Reconstruction}{\ours\ Change variables \(z\mapsto-\bar z\) in the path integral and use \(f(-\bar z)=\overline{f(z)}\); the differential supplies the minus sign. Pairing a and -a independently checks the forced cancellation. The scan visibly has Re and an even character.}
\GSfield{Fate}{The consistent even/odd separation is in \citet[Proposition 9.2(b), equation (35), and Theorem 10.2]{manin1973periods}.}
\GSfield{Present status}{As of 9 November 2026: A genuine parity error in the literal formula. Choosing the correct component repairs the established construction.}
\GSfield{Assessment}{literal statement: false as written; publication match: exact; verification: read the result; development: annotated; confidence: high}
\end{GSitem}
\begin{GSitem}{1972-299.12}{E}{The ordinary unit root}
\GSfield{Kind; source}{construction; \GSpage{P000499}}
\GSfield{The notes}{\(T^2-\lambda_pT+p^{w-1}=0\), roots \(\pi,\bar\pi\), with \(\lambda_p\not\equiv0\pmod p\).}
\GSfield{Statement}{Fix a p-adic embedding of the eigenvalue field, with valuation ring \(\mathcal O_K\), and assume \(a_p\in\mathcal O_K^\times\). For w at least 2 the polynomial has a unique unit root alpha by Hensel lifting of its simple nonzero residue root; the other root is \(\beta=p^{w-1}/\alpha\) and has valuation w-1 when \(v(p)=1\). The bar does not mean complex conjugation in this p-adic argument.}
\GSfield{Fate}{This is the ordinary specialization of the root condition in \citet[\S9.5]{manin1973periods}.}
\GSfield{Present status}{As of 9 November 2026: Established. Ordinarity is a condition at the selected place, not an unqualified complex inequality.}
\GSfield{Assessment}{literal statement: correct with missing hypotheses; publication match: exact; verification: read the result; development: none; confidence: high}
\end{GSitem}
\begin{GSitem}{1972-299.13}{B}{A corrected ordinary measure from periodic Hecke periods}
\GSfield{Kind; source}{construction; \GSpage{P000499}}
\GSfield{The notes}{A limit of sums of \(\ell(b/p^m)\), divided by \(\pi^m+\bar\pi^m\), is given; the congruence after b is ambiguously 0 or a.}
\GSfield{Statement}{A precise published replacement uses a periodic integral-valued normalized period P with
     \[\sum_{j=0}^{p-1}P((x+j)/p)=a_pP(x)-p^{w-2}P(px).\]
     For the unit root alpha, define for m at least 1
     \[\mu_\alpha(a+p^m\mathbb Z_p)=\alpha^{-m}P(a/p^m)-p^{w-2}\alpha^{-m-1}P(a/p^{m-1}).\]
     This is bounded and compatible with refinement; define its total mass by summing level-one values. The explicit two-term formula avoids an unverified identification of the notebook's averaging normalization.}
\GSfield{Reconstruction}{The two-term replacement, period choice and weight conversion are explicit editorial reconstruction. \ours\ Summing children produces the coefficient \(\alpha^{-m-1}(a_p-p^{w-1}/\alpha)=\alpha^{-m}\) on \(P(a/p^m)\); periodicity gives the other term. Unit alpha and integral P give boundedness. A separate derivation and exact rational coefficient check confirm the identity.}
\GSfield{Fate}{The construction is \citet[Lemma 9.4 and equation (39)]{manin1973periods}, with his weight w+2 renamed w here. The lemma is credited there to Nasybullin as improving the original construction.}
\GSfield{Present status}{As of 9 November 2026: The ordinary measure is established. The ambiguous residue must depend on a, but the exact displayed averaging convention was not matched.}
\GSfield{Assessment}{literal statement: correct with missing hypotheses; publication match: partial; verification: read the result; development: annotated; confidence: high}
\end{GSitem}
\begin{GSitem}{1972-299.14}{B}{Interpolation of primitive-character L-values}
\GSfield{Kind; source}{question; \GSpage{P000499}}
\GSfield{The notes}{The circled question asks for \(\int\chi\,d\mu\) as a multiple of \(L(f,\chi,1)\); the final line supplies a power of pi, \(p^r\), omega and a Gauss sum.}
\GSfield{Statement}{Let \(\chi\) be a nontrivial primitive even character of conductor \(q=p^r\), with its values embedded in a finite p-adic extension. Use the preceding measure for the even component
     \(P(x)=(I(x)+I(-x))/(2\Omega_+)\), \(I(x)=\int_x^{i\infty}f(z)dz\), normalized algebraically. Then
     \[\int_{\mathbb Z_p^\times}\chi^{-1}\,d\mu_\alpha=\alpha^{-r}\sum_{a\bmod q}\chi^{-1}(a)P(a/q).\]
     With \(\tau(\chi^{-1})=\sum_a\chi^{-1}(a)e^{2\pi ia/q}\), the sum is the algebraic value represented over the complex numbers by \(i\tau(\chi^{-1})L(f,\chi,1)/(2\pi\Omega_+)\). The complex expression is transported p-adically through its algebraic normalization, not by embedding arbitrary complex periods.}
\GSfield{Reconstruction}{\ours\ Primitivity cancels the lower-level term in the two-term measure. The finite Gauss identity \(\sum_a\chi^{-1}(a)e^{2\pi ina/q}=\tau(\chi^{-1})\chi(n)\), followed by the Mellin integral, gives the stated prefactor. Independent derivations check both cancellation and parity; no new interpolation theorem is claimed.}
\GSfield{Fate}{The precise interpolation theorem is \citet[Theorem 10.2]{manin1973periods}. The displayed notebook constant is not certified because its period parity, root exponent and Gauss-sum convention are not consistently specified.}
\GSfield{Present status}{As of 9 November 2026: Resolved by the 1973 construction under the precise ordinary and period hypotheses; this is the definite question tracked in the problem index.}
\GSfield{Assessment}{literal statement: correct with missing hypotheses; publication match: partial; verification: read the result; development: annotated; confidence: high; problem \hyperref[prob:GP-1972-403]{GP-1972-403}}
\end{GSitem}

\GStalk{300}{1973-04-09}{I. M. Gelfand}{P000500}{Power sums, Bernoulli coefficients and Hurwitz’s lemniscatic analogue}
Gelfand derives the power-sum coefficients from a triangular system of Taylor identities, then tests it on an exponential to obtain a generating function. A brief proposal about arbitrary fields is not precise enough to identify an additional problem. The page ends with a correctable reference to Hurwitz's 1899 paper on lemniscatic coefficients, linking the arithmetic examples to an elliptic analogue.

\begin{GSitem}{1972-300.1}{E}{Power sums and Bernoulli polynomials}
\GSfield{Kind; source}{theorem; \GSpage{P000500}}
\GSfield{The notes}{\(1^k+\cdots+n^k=P_{k+1}(n)\), called a Bernoulli polynomial.}
\GSfield{Statement}{For each integer k at least zero, the power sum is a rational polynomial of degree k+1:
     \[\sum_{j=1}^n j^k=\frac{B_{k+1}(n+1)-B_{k+1}(1)}{k+1}.\]
     Here \(te^{xt}/(e^t-1)=\sum_{r\ge0}B_r(x)t^r/r!\). The sum polynomial is a difference of Bernoulli-polynomial values, not literally the standard polynomial \(B_{k+1}(n)\).}
\GSfield{Fate}{Classical; the generating conventions are checked in \citet[\S24.2]{nist2026dlmf}.}
\GSfield{Present status}{As of 9 November 2026: Established. The k=0 endpoint is included by using B(1), avoiding the B1(0)/B1(1) convention trap.}
\GSfield{Assessment}{literal statement: correct with missing hypotheses; publication match: exact; verification: read the result; development: none; confidence: high}
\end{GSitem}
\begin{GSitem}{1972-300.2}{E}{Telescoping Taylor identities}
\GSfield{Kind; source}{method; \GSpage{P000500}}
\GSfield{The notes}{For \(F'=f\), sum \(F(j+1)-F(j)=f(j)+f'(j)/2!+\cdots\).}
\GSfield{Statement}{For a polynomial f of degree d, Taylor expansion is finite and exact, giving
     \[F(n+1)-F(1)=\sum_{r=0}^{d}\frac{1}{(r+1)!}\sum_{j=1}^n f^{(r)}(j).\]
     For nonpolynomial analytic f, convergence on each interval must be justified. For merely smooth f, a finite Euler--Maclaurin formula needs its remainder; the infinite expression is not automatically an equality.}
\GSfield{Fate}{The finite identity is elementary; \citet[\S2.10]{nist2026dlmf} gives the rigorous summation formula with remainder.}
\GSfield{Present status}{As of 9 November 2026: Established in the polynomial setting; regularity alone does not justify an infinite derivative expansion.}
\GSfield{Assessment}{literal statement: correct with missing hypotheses; publication match: exact; verification: read the result; development: none; confidence: high}
\end{GSitem}
\begin{GSitem}{1972-300.3}{E}{The determinant is finite after polynomial truncation}
\GSfield{Kind; source}{construction; \GSpage{P000500}}
\GSfield{The notes}{The sums \(S_n,S'_n,\ldots\) are treated as unknowns in a triangular system; an open-ended determinant solves for \(S_n\).}
\GSfield{Statement}{For degree-d f, retain the d+1 unknowns \(\sum_j f^{(r)}(j)\), 0 at most r at most d. The coefficient matrix is upper triangular with diagonal 1. Cramer's rule gives exactly the displayed finite determinant with its first column replaced by endpoint differences. An infinite determinant requires an additional operator or convergence definition absent from the notes.}
\GSfield{Reconstruction}{\oursComputed\ Exact rational linear algebra and direct power sums were independently compared for degrees 0 through 12 and several endpoints; all identities agree. This finite verification supports the implementation, while triangularity proves the formula for every degree.}
\GSfield{Fate}{This is a finite-dimensional derivation of the Bernoulli summation coefficients, consistent with \citet[\S24.2]{nist2026dlmf}.}
\GSfield{Present status}{As of 9 November 2026: Established finite algebra; the source does not define a general infinite determinant.}
\GSfield{Assessment}{literal statement: correct with missing hypotheses; publication match: exact; verification: read the result; development: annotated; confidence: high}
\end{GSitem}
\begin{GSitem}{1972-300.4}{E}{The coefficients called B in the notebook include factorials}
\GSfield{Kind; source}{definition; \GSpage{P000500}}
\GSfield{The notes}{Expanding along the first column gives coefficients \(B_0,B_1,B_2,\ldots\) multiplying successive endpoint derivatives.}
\GSfield{Statement}{If the endpoint differences are left unscaled as on the page, the coefficient of the r-th entry is \(b_r=B_r/r!\), where \(B_r\) is the standard Bernoulli number with \(B_1=-1/2\). Thus \(b_0=1,b_1=-1/2,b_2=1/12,b_4=-1/720\). Using standard B2=1/6 directly in that determinant expansion would be wrong.}
\GSfield{Fate}{The convention follows from \citet[\S24.2]{nist2026dlmf} and the generating identity in the next item.}
\GSfield{Present status}{As of 9 November 2026: Established after distinguishing normalized coefficients from standard Bernoulli numbers.}
\GSfield{Assessment}{literal statement: correct with missing hypotheses; publication match: exact; verification: read the result; development: none; confidence: high}
\end{GSitem}
\begin{GSitem}{1972-300.5}{E}{The exponential test and the generating function}
\GSfield{Kind; source}{construction; \GSpage{P000500}}
\GSfield{The notes}{Set \(f(x)=e^{tx}\), use \(S_n=(e^{t(n+1)}-e^t)/(e^t-1)\), and expand in t.}
\GSfield{Statement}{The finite geometric sum has a removable singularity at t=0. Comparing the endpoint-derivative expansion gives
     \[\sum_{r\ge0}b_rt^r=\frac{t}{e^t-1}.\]
     This is an identity of formal power series over \(\mathbb Q\), or an analytic Taylor series for \(|t|<2\pi\). Its coefficients yield both the Bernoulli recurrence and the power-sum identities.}
\GSfield{Fate}{Classical generating function; see \citet[\S24.2]{nist2026dlmf}.}
\GSfield{Present status}{As of 9 November 2026: Established. Formal coefficients and analytic convergence are separate valid readings.}
\GSfield{Assessment}{literal statement: correct with missing hypotheses; publication match: exact; verification: read the result; development: none; confidence: high}
\end{GSitem}
\begin{GSitem}{1972-300.6}{F}{Bernoulli numbers in arbitrary fields}
\GSfield{Kind; source}{remark; \GSpage{P000500}}
\GSfield{The notes}{Only the sentence \ru{Числа Бернулли в произвольных полях} is recorded.}
\GSfield{Statement}{No construction, range of indices, characteristic restriction, divided-power convention or desired property is supplied. Over characteristic zero the rational coefficients can be transported through the prime field. In positive characteristic the factorials and Bernoulli denominators obstruct an unrestricted reduction of the same generating series.}
\GSfield{Fate}{The classical rational generating series does not identify which extension the speaker intended. The adjacent Hurwitz reference concerns a specific elliptic analogue and does not by itself resolve this arbitrary-field phrase.}
\GSfield{Present status}{As of 9 November 2026: Programme not identifiable as a precise problem; no new open question is manufactured.}
\GSfield{Assessment}{literal statement: not determinable; publication match: thematic only; verification: unresolved lead; development: none; confidence: low}
\end{GSitem}
\begin{GSitem}{1972-300.7}{E}{Hurwitz’s lemniscatic analogue}
\GSfield{Kind; source}{remark; \GSpage{P000500}}
\GSfield{The notes}{The page first cites Math. Z., 1898, then corrects itself to Hurwitz, Math. Ann., 1899, pp. 196--226, on lemniscatic coefficients.}
\GSfield{Statement}{The identifiable work is Hurwitz's study of the expansion coefficients of lemniscatic elliptic functions and the corresponding Gaussian-lattice analogues of Bernoulli arithmetic. This is a specific characteristic-zero elliptic theory, not a theorem asserting a universal Bernoulli construction over arbitrary fields.}
\GSfield{Fate}{\citet{hurwitz1899lemniscatic} appeared in Mathematische Annalen 51 (1899), pp.~196--226. The publisher's collected-works entry explicitly records these original details; its preview identifies the Gaussian-lattice setting. The full original proof was not read.}
\GSfield{Present status}{As of 9 November 2026: Identified older reference. The initial journal/year on the page is a faithful bibliographic error corrected later on the same page.}
\GSfield{Assessment}{literal statement: correct with missing hypotheses; publication match: exact; verification: publisher metadata checked; development: none; confidence: high}
\end{GSitem}

\section{Developments}\label{dev:1972}
% Corrective development of an A item; no bounded C/D research target qualified.
\subsection{A correction to the published tolerance assertion}
\label{sec:gy1972-tolerance-correction}

The question at the end of Vasserstein's talk, occurrence 1972-280.20 and
problem GP-1972-101, asks whether pairwise adjacency implies that a finite
set is a simplex. This is a precise question, but it is not an undisseminated
C or D item. The same assertion was published in
\citet[\S11, Remark~(c), p.~111]{vaserstein1976foundations}, immediately after
the definition of the simplicial scheme. The short announcement
\citep{vaserstein1973foundations} was received on 14 November 1972, one day
after the talk, and gives the same left-translation construction. The
detailed paper was received on 26 February 1975 and published in 1976.

The result below is included as corrective development of an A item.
The label A records publication provenance, not mathematical correctness.
This is a deliberate methodological extension of the usual C/D-only
development filter: a fully checkable correction of an exact published
assertion should not be suppressed or reclassified D merely to fit that
filter. The marks \ours\ and \oursComputed\ identify work performed for
this companion; they do not assert that the counterexample or criterion
has never been noticed before.

\paragraph{The exact model.}
Let \(R\) be an associative ring with identity. Write
\(\mathrm{GL}(R)=\bigcup_{n\geq1}\mathrm{GL}_n(R)\), with stabilization
\(g\mapsto\operatorname{diag}(g,1)\). For a total order \(\prec\) of the
coordinates obtained by a permutation of \(\mathbb N\), let
\(U_\prec\) be the subgroup of matrices differing from the identity in
finitely many entries and satisfying
\[
 u_{ii}=1,\qquad u_{ij}=0\quad\text{if }i\ne j\text{ and }j\prec i.
\]
A finite nonempty set \(X\subset\mathrm{GL}(R)\) is a simplex of the
left-coset model when
\[
                 X\subset gU_\prec
       \quad\text{for some }g\in\mathrm{GL}(R)\text{ and one order }\prec.
\]
This is the definition on P000468 after replacing the notebook's circle
coordinate \(a\) by \(I+a\). Reversing upper and lower triangular
conventions gives the same collection of simplices after coordinate
permutations. The 1976 paper uses right cosets and explains in
\S11, Remark~(b), that inversion identifies its model with the
left-translation model. In particular, the same order must work for an
entire simplex. The pairwise condition permits different orders for
different pairs and is strictly weaker.

\begin{proposition}[Exact acyclic-support criterion]\ours\ 
\label{prop:gy1972-support-criterion}
Let \(X\subset\mathrm{GL}(R)\) be finite and nonempty, and choose
\(x_0\in X\). Form a finite directed graph on a set of coordinates
containing all nonidentity entries of all \(x_0^{-1}x\), \(x\in X\).
Include the edge \(i\to j\), including a loop if \(i=j\), whenever
\[
                (x_0^{-1}x-I)_{ij}\ne0
                \quad\text{for some }x\in X.
\]
Then \(X\) is a simplex of the left-coset model if and only if this graph
is acyclic. The criterion is unchanged by adjoining identity coordinates.
For the right-coset model, replace \(x_0^{-1}x\) by \(xx_0^{-1}\).
\end{proposition}

\begin{proof}
If \(X\subset gU_\prec\), write \(x_0=gu_0\) and \(x=gu_x\).
Then \(x_0^{-1}x=u_0^{-1}u_x\in U_\prec\). Thus every edge points
strictly forward in the same order, which excludes directed cycles and
loops. Conversely, a finite acyclic directed graph has a topological
ordering. Order the finitely many relevant coordinates accordingly and
extend this to a permutation order of \(\mathbb N\). Absence of loops
says that every \(x_0^{-1}x\) has diagonal entries one, and all other
nonzero entries lie strictly above the diagonal in the chosen order.
Consequently every \(x_0^{-1}x\) belongs to \(U_\prec\), giving
\(X\subset x_0U_\prec\). Adding identity coordinates adds isolated
vertices, so cannot create or remove a directed cycle. The right-coset
proof uses \(xx_0^{-1}\) in exactly the same way.
\end{proof}

This is an exact replacement for the proposed pairwise test. It reduces
simplex recognition to finding a common order of the union of the
relative supports, rather than checking each pair with a separate order.
When equality to zero is decidable in the coefficient ring, the criterion
also gives a finite decision procedure: compute the relative matrices
and perform a topological-sort test on their support graph.

\begin{proposition}[A three-vertex counterexample over every nonzero unital ring]\ours\ 
\label{prop:gy1972-nonflag-triple}
Suppose \(1_R\ne0\), and in \(\mathrm{GL}_3(R)\) put
\[
 a=I+E_{12},\qquad b=I+E_{23},\qquad c=I+E_{31}.
\]
Every two-element subset of \(\{a,b,c\}\) is a simplex, but the entire
set is not a simplex. This holds in both the left-coset and right-coset
conventions and after every stabilization.
\end{proposition}

\begin{proof}
The matrix units carry coefficient \(1_R\), so each is square-zero and
the three inverses are obtained by replacing plus with minus. Direct
multiplication gives
\[
\begin{aligned}
 a^{-1}b&=I-E_{12}+E_{23}-E_{13},\\
 a^{-1}c&=I-E_{12}+E_{31},\\
 b^{-1}c&=I-E_{23}+E_{31}-E_{21}.
\end{aligned}
\]
The off-diagonal supports point forward, respectively, in the orders
\[
             1\prec2\prec3,\qquad
             3\prec1\prec2,\qquad
             2\prec3\prec1.
\]
Thus every pair is a left-coset simplex. If the triple were a simplex,
the same order would triangularize both \(a^{-1}b\) and \(a^{-1}c\).
Their supports would require simultaneously
\(1\prec2\), \(2\prec3\), and \(3\prec1\), an impossibility.
Equivalently, their union support contains a directed triangle, so
Proposition~\ref{prop:gy1972-support-criterion} applies.

For completeness, the same three vertices also work directly in the
right-coset convention, without replacing them by their inverses:
\[
\begin{aligned}
 ba^{-1}&=I-E_{12}+E_{23},\\
 ca^{-1}&=I-E_{12}+E_{31}-E_{32},\\
 cb^{-1}&=I-E_{23}+E_{31}.
\end{aligned}
\]
The same respective three orders triangularize these pairwise
quotients. The first two again have union support containing
\(1\to2\to3\to1\), excluding a common right coset.

All decisive coefficients are \(1_R\) or \(-1_R\), which are nonzero
in every nonzero unital ring, including characteristic two. The proof
uses no field hypothesis and no division by a coefficient. Additional
identity coordinates cannot remove a cycle among coordinates one,
two, and three. This proves all the stated cases.
\end{proof}

\begin{proposition}[Minimal coordinate rank of the obstruction]\ours\ 
\label{prop:gy1972-rank-minimal}
For every unital ring, a finite pairwise adjacent set of vertices
supported in one two-by-two block is a simplex, even when adjacency is
tested after stabilization. Consequently the example in
Proposition~\ref{prop:gy1972-nonflag-triple} has minimal coordinate rank,
as well as the smallest possible number of vertices for failure of the
pairwise criterion.
\end{proposition}

\begin{proof}
Translate by the inverse of one vertex so that the set contains \(I\).
In the left model, every other vertex adjacent to \(I\) and supported
on two coordinates must have one of the forms
\(I+rE_{12}\) or \(I+sE_{21}\). Stabilization does not change this
fact, since only the relative order of those two coordinates matters.
If both orientations occur nontrivially, with \(r\ne0\) and \(s\ne0\),
then
\[
 (I+rE_{12})^{-1}(I+sE_{21})
       =I-rE_{12}+sE_{21}-rsE_{11}.
\]
Both off-diagonal entries are nonzero and point in opposite directions.
No coordinate order can make this matrix unitriangular, irrespective of
whether \(rs\) vanishes. This contradicts pairwise adjacency.
Hence all nonidentity vertices use one orientation and lie in one
unitriangular group, so the whole set is a simplex. The same argument
works with right quotients, or follows by inversion. Rank one is
immediate. Sets of at most two vertices cannot contradict a condition
defined by their two-element subsets, while the displayed triple does.
\end{proof}

\paragraph{Independent verification.}\oursComputed\par
The included Python program \texttt{check\_tolerance.py} uses exact
matrix arithmetic and exhausts all coordinate permutations in dimensions
three, four, and five over \(\mathbb Z\) and
\(\mathbb F_2,\mathbb F_3,\mathbb F_5,\mathbb F_7\), in both coset
conventions. Its thirty parameter cases all have three present edges
and no triple simplex. The complete successful-order lists, integer
quotient matrices, and results are in \texttt{tolerance\_output.json}.
The integer calculation rechecks the signs and cross terms; the
characteristic-two calculation checks that the obstruction does not
rely on distinguishing plus from minus. A separate independently
written audit by the assembling researcher exhausts the four finite
fields in the same three dimensions and both conventions and agrees in
all twenty-four cases. The general-ring and arbitrary-stabilization
conclusions follow from the proof, not extrapolation from these finite
computations. The rank-two argument was also checked independently by
normalizing a base vertex and recomputing the opposite-direction
quotient above.

\paragraph{Scope and historical assessment.}
This resolves the literal notebook question negatively; no residual
mathematical obstruction remains for that question. The 1976 source
expressly says the tolerance assertion is not used in the paper.
Accordingly, this local correction provides no basis for declaring its
comparison theorems or the Volodin construction invalid. Searching the
speaker's Foundations papers and bibliography, the original Volodin
paper, and targeted literature queries identified the exact published
assertion but did not establish the history of any later correction.
Absence of a retrieved erratum is not evidence that none exists, and no
claim is made about why the error persisted.

The strongest general statement supported after correction is the
acyclic-support criterion above, with its full common-order requirement.
Replacing the simplicial complex by its flag completion would create a
different construction and does not rescue the literal assertion.
A further bounded bibliographic run could trace corrections or uses of
Remark~(c); it is not needed to establish the explicit counterexample,
and it is not recorded as a manufactured open mathematical problem.

\section{Problem index}\label{problems:1972}
Earlier-year occurrence identifiers are retained as text; current-year identifiers link to their entries.

\subsection{GP-1972-101: Does pairwise adjacency imply a simplex in the permutation-triangular K-theory complex?}\label{prob:GP-1972-101}
For a nonzero unital ring \(A\), let the vertices be stable invertible matrices, and let a finite set be a simplex when it lies in one left coset of a subgroup of matrices unitriangular for one permutation of the coordinate order. Must every finite pairwise adjacent set be a simplex? The notebook first asks the three-vertex case and then the finite version.
\par\noindent\textit{Occurrences:} \hyperref[item:1972-280.20]{1972-280.20}.
\par\noindent\textit{Status (2026-11):} As of 9 November 2026, the literal assertion is false. It was explicitly asserted in \citet[\S11, Remark~(c)]{vaserstein1976foundations}, not left undisclosed. The appendix gives a complete proof for the triple \(I+E_{12},I+E_{23},I+E_{31}\) over every nonzero unital ring, checks invariance under stabilization and coset conventions, and includes independent exact computation. This is a correction of a published claim, with no assertion of priority for the counterexample.

\subsection{GP-1972-201: Homogeneous additive cocycles and formal-law indecomposables}\label{prob:GP-1972-201}
For each \(n\ge2\) and commutative ring \(K\), determine normalized homogeneous polynomials \(h(x,y)\) satisfying \(h(x,y)+h(x+y,z)=h(y,z)+h(x,y+z)\), and distinguish the symmetric solutions from the unrestricted ones. The formal-law indecomposable calculation concerns the symmetric part.
\par\noindent\textit{Occurrences:} \hyperref[item:1972-276.11]{1972-276.11}, \hyperref[item:1972-279.2]{1972-279.2}, \hyperref[item:1972-279.5]{1972-279.5}, \hyperref[item:1972-285.3]{1972-285.3}, \hyperref[item:1972-285.8]{1972-285.8}.
\par\noindent\textit{Status (2026-11):} As of November 2026: the symmetric classification is established by \citet[\S~II, Lemma~3]{lazard1955formal}, circulated December 1954 and published 1955. It is freely parametrized by the primitive integral binomial cocycle in each degree. The notebook's temporary nonsymmetric cubic example is corrected in 1972-276.14; its extra parameter is two-torsion, not three-torsion. These are expository occurrences, not an unsolved 1972 conjecture.

\subsection{GP-1972-202: Determine the smooth normal-invariant homotopy groups}\label{prob:GP-1972-202}
Under the reconstructed interpretation of the notebook's \(K_J\) as stable bundles with a specified spherical trivialization, determine \(\pi_q(G/O)\) in all degrees, including the extensions in the homotopy exact sequence of \(G/O\to BO\to BG\). No finite range is specified in the notes.
\par\noindent\textit{Occurrences:} \hyperref[item:1972-287.11]{1972-287.11}.
\par\noindent\textit{Status (2026-11):} As of November 2026: partially developed. The exact-sequence reduction and many finite ranges are established; no complete all-degree numerical determination was located in the primary sources checked. \citet[\S~9.2]{ranicki2002surgery} gives the normal-invariant framework. \citet[Theorem~1.1]{burklundisaksenxu2026synthetic}, first circulated 2 August 2024, online 2025, volume 2026, illustrates continued finite stable-stem progress. This does not itself compute all groups \(\pi_q(G/O)\). The source sentence is programmatic, and the identification has medium confidence.

\subsection{GP-1972-203: Homotopy structures and the obstruction to surgery}\label{prob:GP-1972-203}
For a closed simply connected smooth \(n\)-manifold \(M\), \(n\ge6\), describe its marked homotopy structure set and determine when a degree-one normal map is normally bordant to a homotopy equivalence. These dimensional and compactness hypotheses reconstruct the range intended by the surgery continuation.
\par\noindent\textit{Occurrences:} \hyperref[item:1972-287.3]{1972-287.3}, \hyperref[item:1972-293.9]{1972-293.9}.
\par\noindent\textit{Status (2026-11):} As of November 2026: the structural question was solved before these lectures. The surgery exact sequence identifies the image of the normal-invariant map with the zero fibre of the \(L_n(\mathbb Z)\)-valued obstruction. See \citet[Theorem~3]{sullivan1996triangulating}, circulated 1967, and the precise action-sequence formulation in \citet[Theorem~13.2]{ranicki2002surgery}. The source gives no particular manifold whose numerical structure set is left to compute.

\subsection{GP-1972-204: Uniqueness of transverse normal maps and the left surgery term}\label{prob:GP-1972-204}
For a fixed smooth normal invariant over a closed simply connected manifold of dimension at least six, determine the ambiguity of transverse Pontryagin--Thom representatives, and decide when a normal bordism between homotopy-equivalence representatives can be changed relative to its ends to an h-cobordism.
\par\noindent\textit{Occurrences:} \hyperref[item:1972-293.8]{1972-293.8}, \hyperref[item:1972-293.10]{1972-293.10}.
\par\noindent\textit{Status (2026-11):} As of November 2026: resolved structurally by normal bordism and relative surgery. A normal invariant specifies a normal-bordism class; the \(L_{n+1}(\mathbb Z)\) action accounts for the ambiguity of homotopy structures in a fibre, and the relative surgery obstruction tests a chosen bordism. See \citet[Propositions~9.43, 9.48 and Theorem~13.2]{ranicki2002surgery}; the pre-lecture sequence is in \citet[Theorem~3]{sullivan1996triangulating}. The compressed notebook attribution to Milnor--Kervaire or Novikov's dissertation was not matched to a unique proposition.

\subsection{GP-1972-301: Classification of finite-type quivers}\label{prob:GP-1972-301}
For which finite connected quivers are there only finitely many indecomposable finite-dimensional representations up to isomorphism?
\par\noindent\textit{Occurrences:} \hyperref[item:1972-277.8]{1972-277.8}, \hyperref[item:1972-278.3]{1972-278.3}.
\par\noindent\textit{Status (2026-11):} As of 9 November 2026: Solved before the talks by Gabriel 1972; ADE is the exact list. The 1973 BGP proof is an exact subsequent exposition.

\subsection{GP-1972-302: The tame infinite boundary for quivers}\label{prob:GP-1972-302}
Which connected quivers of infinite representation type avoid the arbitrary-matrix-pair classification problem?
\par\noindent\textit{Occurrences:} \hyperref[item:1972-277.9]{1972-277.9}, \hyperref[item:1972-278.9]{1972-278.9}.
\par\noindent\textit{Status (2026-11):} As of 9 November 2026: Solved: extended Dynkin type, with standard affine limiting cases. Nazarova’s exact paper was received 19 September 1972 before these talks and published 1973.

\subsection{GP-1972-303: Completeness and coefficients of the one-sided Kostyuchenko system}\label{prob:GP-1972-303}
For real \(\gamma\), determine completeness, minimality and coefficient functionals for \(\{e^{\gamma nx}\sin nx:n\ge1\}\) in \(L^2(0,\pi)\), and distinguish completeness from the basis property.
\par\noindent\textit{Occurrences:} \hyperref[item:1972-283.2]{1972-283.2}, \hyperref[item:1972-283.3]{1972-283.3}, \hyperref[item:1972-283.11]{1972-283.11}.
\par\noindent\textit{Status (2026-11):} As of 9 November 2026: Completeness was proved by Levin 1971, before the seminar, according to the later primary papers. The system is minimal; for nonzero real \(\gamma\) it is not a Schauder basis, whereas \(\gamma=0\) is the sine basis. The signed system in the notebook is redundant and must not be substituted for this precise problem. The asserted particular Schwarz--Christoffel formula was not verified.

\subsection{GP-1972-304: The Adams conjecture for stable spherical fibrations}\label{prob:GP-1972-304}
For every finite CW-complex X, integer \(k\ge2\), and virtual real bundle \(x\in KO(X)\), does some \(k^N\) annihilate \(J(\psi^k x-x)\)?
\par\noindent\textit{Occurrences:} \hyperref[item:1972-295.6]{1972-295.6}.
\par\noindent\textit{Status (2026-11):} As of 9 November 2026: proved before the seminar: Sullivan’s June 1970 notes, Corollary 5.9, and Quillen’s 1971 Theorem 1.1 (received 3 July 1970). The notebook’s separate reverse formula is false and is not this conjecture.

\subsection{GP-1972-402: Off-diagonal zeros of the Bergman kernel}\label{prob:GP-1972-402}
For a bounded domain G the off-diagonal Bergman kernel can vanish. Every planar annulus \(r<|z|<1\), \(0<r<1\), has zeros. In contrast the disk and bounded homogeneous complete circular domains have zero-free kernels. Positivity of \(K(z,\bar z)\) on the diagonal is a different assertion.
\par\noindent\textit{Occurrences:} \hyperref[item:1972-291.22]{1972-291.22}.
\par\noindent\textit{Status (2026-11):} As of 9 November 2026: The unrestricted question had a negative answer before this seminar. No broader classification problem is silently substituted.

\subsection{GP-1972-403: Ordinary p-adic interpolation of modular L-values}\label{prob:GP-1972-403}
Let \(\chi\) be a nontrivial primitive even character of conductor \(q=p^r\), with its values embedded in a finite p-adic extension. Use the preceding measure for the even component
     \(P(x)=(I(x)+I(-x))/(2\Omega_+)\), \(I(x)=\int_x^{i\infty}f(z)dz\), normalized algebraically. Then
     \[\int_{\mathbb Z_p^\times}\chi^{-1}\,d\mu_\alpha=\alpha^{-r}\sum_{a\bmod q}\chi^{-1}(a)P(a/q).\]
     With \(\tau(\chi^{-1})=\sum_a\chi^{-1}(a)e^{2\pi ia/q}\), the sum is the algebraic value represented over the complex numbers by \(i\tau(\chi^{-1})L(f,\chi,1)/(2\pi\Omega_+)\). The complex expression is transported p-adically through its algebraic normalization, not by embedding arbitrary complex periods.
\par\noindent\textit{Occurrences:} \hyperref[item:1972-299.14]{1972-299.14}.
\par\noindent\textit{Status (2026-11):} As of 9 November 2026: Resolved by the 1973 construction under the precise ordinary and period hypotheses; this is the definite question tracked in the problem index.

\section{Corrections to the transcription and translation}\label{corr:1972}
The evidence distinguishes a transcription correction from a mathematical emendation of wording faithfully preserved in the notebook.
\subsection*{P000449-P000451 (both; provenance): record 272 boundary}
\noindent\textit{Transcription reads.} catalogue range P000449-P000453\par
\noindent\textit{Critical reading.} Gabrielov's notes begin at the dated heading on P000450 and end on P000451; P000449 belongs to the preceding material.\par
\noindent\textit{Evidence.} Original headings and consecutive source pages; catalogue ranges overlap adjacent talks.\par
\subsection*{P000450 (both; faithful error): local system parenthesis}
\noindent\textit{Transcription reads.} local system = bundle with connection\par
\noindent\textit{Critical reading.} local system = bundle with flat (integrable) connection; the integral lattice is additional data\par
\noindent\textit{Evidence.} P000450 scan omits flatness; the connection is locally \(d\). Griffiths 1970, section 2.\par
\subsection*{P000451 (both; faithful error): first displayed orthogonality}
\noindent\textit{Transcription reads.} \(\psi(F^p,F^{n-p})=0\)\par
\noindent\textit{Critical reading.} \(\psi(F^p,F^{n-p+1})=0\)\par
\noindent\textit{Evidence.} The scan genuinely has \(n-p\); P000458 and the polarization relations show the required \(+1\).\par
\subsection*{P000451 (both; faithful error): last assertion about \(G/K\)}
\noindent\textit{Transcription reads.} \(G/K\) has complex structure; the projection is not holomorphic\par
\noindent\textit{Critical reading.} \(G/K\) is Riemannian symmetric in general; a complex structure and nonholomorphicity require further hypotheses, with \(V=K\) an immediate exception.\par
\noindent\textit{Evidence.} Original scan; Griffiths 1970, equation (2.3) and surrounding text.\par
\subsection*{P000452 (both; provenance): Manin heading}
\noindent\textit{Transcription reads.} first lecture of Manin\par
\noindent\textit{Critical reading.} retelling of Manin's first lecture\par
\noindent\textit{Evidence.} The original heading ends (\ru{пересказ}); this qualifies the attribution of the surviving notes.\par
\subsection*{P000453 (both; faithful error): two-summand quotient variant}
\noindent\textit{Transcription reads.} \(L_1/L_2=A/\pi^{\nu_1}A\oplus A/\pi^{\nu_2}A\) without an inclusion qualification\par
\noindent\textit{Critical reading.} Use representatives with \(L_2\subset L_1\) and nonnegative exponents; the distance remains \(\lvert\nu_1-\nu_2\rvert\) independently.\par
\noindent\textit{Evidence.} The initial normalized inclusion is explicit in the scan; the later quotient variant suppresses that convention.\par
\subsection*{P000454 (both; faithful error): finite-extension setup}
\noindent\textit{Transcription reads.} \(L/K\) finite; \(G=\operatorname{Gal}(L/K)\)\par
\noindent\textit{Critical reading.} Assume \(L/K\) finite Galois for the later residue quotient isomorphism.\par
\noindent\textit{Evidence.} Original page omits Galois; the surjectivity exercise on P000455 needs it. Milne 2020, Theorem 7.58.\par
\subsection*{P000454 (both; faithful error): residue-field lift and kernel}
\noindent\textit{Transcription reads.} \(x\in L\); \(\bar x^g=\bar x\) without explicit quantifier\par
\noindent\textit{Critical reading.} Choose \(x\in\mathcal O_L\) lifting \(\bar x\); \(g\) is in the kernel iff every \(\bar x\in\ell\) is fixed.\par
\noindent\textit{Evidence.} Original scan and the definition of reduction modulo the maximal ideal.\par
\subsection*{P000455 (both; faithful error): domain and kernel of higher character}
\noindent\textit{Transcription reads.} \(\psi_i:G_i/G_{i+1}\to\ell\); \(\ker\psi_i=G_{i+1}\)\par
\noindent\textit{Critical reading.} Define \(\psi_i\) on \(G_i\) with that kernel, or state that its induced map from \(G_i/G_{i+1}\) is injective.\par
\noindent\textit{Evidence.} P000455 scan has the same domain/kernel inconsistency; Milne 2020, Corollary 7.59.\par
\subsection*{P000455 (both; ambiguity): ramification terminology}
\noindent\textit{Transcription reads.} weakly ramified means \(G_2=1\)\par
\noindent\textit{Critical reading.} tamely ramified in modern terminology; modern weak ramification would mean \(G_3=1\) in the notebook's indexing.\par
\noindent\textit{Evidence.} The Russian scan has \ru{слаборазветвл.}; the distinction is terminological, not a mistranscription.\par
\subsection*{P000456 (both; faithful error): local section group}
\noindent\textit{Transcription reads.} \(\Gamma(U,H_{\mathbb Z})=\mathbb Z^n\) for an unspecified trivializing \(U\)\par
\noindent\textit{Critical reading.} Assume \(U\) connected, or use locally constant maps \(U\to\mathbb Z^n\).\par
\noindent\textit{Evidence.} The scan's local trivialization does not by itself require connected \(U\); the earlier sentence handles connected domains.\par
\subsection*{P000458 (both; faithful error): polarization parity}
\noindent\textit{Transcription reads.} symmetric for odd \(n\) and skew-symmetric for even \(n\)\par
\noindent\textit{Critical reading.} symmetric for even \(n\) and alternating for odd \(n\)\par
\noindent\textit{Evidence.} Confirmed directly in P000458 scan; Griffiths 1968 I, equations (1.7)--(1.15).\par
\subsection*{P000458 (both; faithful error): real-group transitivity}
\noindent\textit{Transcription reads.} the identity component acts transitively on \(D\)\par
\noindent\textit{Critical reading.} the identity component acts transitively on each chosen connected component of \(D\)\par
\noindent\textit{Evidence.} The scan explicitly says component of unity; disconnected period domains require this qualification.\par
\subsection*{P000458 (both; provenance): two final headings}
\noindent\textit{Transcription reads.} Manin's 2nd lecture; D. Kazhdan, more on Tate\par
\noindent\textit{Critical reading.} Bare continuation headings without surviving mathematical content on this page\par
\noindent\textit{Evidence.} Original scan; no assertion or new problem is recoverable from the headings alone.\par
\subsection*{P000459 (both; faithful error): first definition, parenthetical commutativity}
\noindent\textit{Transcription reads.} commutativity follows from associativity\par
\noindent\textit{Critical reading.} false over arbitrary rings; valid over a reduced ring or a torsion-free ring\par
\noindent\textit{Evidence.} Scan confirms the assertion; Lazard 1955 Theorem I and Remark (3.14a), and the notebook counterexample on P477.\par
\subsection*{P000459 (both; faithful error): inverse coordinate-change expansion}
\noindent\textit{Transcription reads.} \(\varphi^{-1}(x)=z-a_1z^2+\cdots\)\par
\noindent\textit{Critical reading.} \(\varphi^{-1}(z)=z-a_1z^2+\cdots\)\par
\noindent\textit{Evidence.} The displayed inverse uses z throughout its output; the scan preserves the inconsistent argument.\par
\subsection*{P000459 (both; ambiguity): reference after logarithm lemma}
\noindent\textit{Transcription reads.} Mat. 13:6\par
\noindent\textit{Critical reading.} Honda, Formal groups and zeta-functions, Matematika 13(6) (1969), 3--17; original Osaka J. Math. 5 (1968), 199--213\par
\noindent\textit{Evidence.} Math-Net mat530 and the Osaka institutional original resolve the pointer; this is not a Russian publication of Lazard 1955.\par
\subsection*{P000461 (both; faithful error): margin explaining decomposables}
\noindent\textit{Transcription reads.} products of noninvertible elements\par
\noindent\textit{Critical reading.} products of elements of strictly positive degree; the ideal \(I^2\)\par
\noindent\textit{Evidence.} Scan and main text distinguish the margin from the explicit nonzero-degree condition; nonunits would incorrectly include integer products.\par
\subsection*{P000462 (both; ambiguity): Ext exact sequence}
\noindent\textit{Transcription reads.} \(\operatorname{Ext}^{2,n}_{K[x]}(K,K)\)\par
\noindent\textit{Critical reading.} cohomology in the comodule category of the primitive coalgebra, or the corresponding divided-power dual algebra category\par
\noindent\textit{Evidence.} P462 defines comodules; P476 explicitly introduces the dual algebra. Ordinary module Ext over k[x] vanishes in degree two over a field.\par
\subsection*{P000462 (both; faithful error): cubic monomial coefficient characteristic}
\noindent\textit{Transcription reads.} \(x^2y\) is a cycle over \(\mathbb Z_3\)-modules\par
\noindent\textit{Critical reading.} \(x^2y\) is a cycle when \(2=0\), and is not a cycle over \(\mathbb F_3\)\par
\noindent\textit{Evidence.} Scan confirms the second subscript 3. Symbolic expansion and an independent finite-field evaluation give defect 2xyz.\par
\subsection*{P000462 (both; faithful error): cubic universal module}
\noindent\textit{Transcription reads.} \(\mathbb Z+\mathbb Z_3\)\par
\noindent\textit{Critical reading.} \(\mathbb Z\oplus\mathbb Z/2\) for the nonsymmetric degree-three coefficient module\par
\noindent\textit{Evidence.} The scan clearly has subscript 3. The independently checked defect is \(2(C_1-C_2)xyz\). Symmetry removes the torsion summand.\par
\subsection*{P000463 (both; ambiguity): factorization of the dual algebra}
\noindent\textit{Transcription reads.} \(\prod B_q\)\par
\noindent\textit{Critical reading.} graded tensor product, or its degree completion for the full dual\par
\noindent\textit{Evidence.} The binomial divided-power multiplication yields truncated polynomial tensor factors; an ordinary product has extraneous idempotents. P477 writes a completed tensor product.\par
\subsection*{P000464 (both; faithful error): graph/category paragraph}
\noindent\textit{Transcription reads.} Graph is a category; functoriality is not required\par
\noindent\textit{Critical reading.} Representation of the free path category; no relations among distinct paths are imposed.\par
\noindent\textit{Evidence.} Scan confirms both phrases. Composition along paths is necessary for a functor; the intended lack of relations is confirmed by BGP 1973 Definition 1.\par
\subsection*{P000464 (both; faithful error): single-arrow invariants}
\noindent\textit{Transcription reads.} Only \(\dim\ker\) and \(\dim\operatorname{coker}\)\par
\noindent\textit{Critical reading.} Add rank unless the source and target dimensions are fixed.\par
\noindent\textit{Evidence.} Scan confirms the two invariants. Identity maps in dimensions 1 and 2 disprove the unrestricted literal statement.\par
\subsection*{P000465 (both; ambiguity): first theorem attribution}
\noindent\textit{Transcription reads.} Roytov or Nazarovy\par
\noindent\textit{Critical reading.} Probably Roiter or Nazarova; preserve the manuscript uncertainty.\par
\noindent\textit{Evidence.} Scan reads approximately \ru{Ройтов} and \ru{Назаровы}; Nazarova 1973, received 19 September 1972, proves the matching affine classification. The handwritten attribution cannot be repaired with certainty.\par
\subsection*{P000465 (both; faithful error): root count above \(A_3\) table}
\noindent\textit{Transcription reads.} Number equals number of simple roots\par
\noindent\textit{Critical reading.} Number equals number of positive roots.\par
\noindent\textit{Evidence.} Scan clearly reads \ru{простых корней}. Its own \(A_3\) table has six rows while \(A_3\) has only three simple roots; BGP Theorem 3.1 supplies the exact positive-root statement.\par
\subsection*{P000466-P000467 (both; faithful error): resolution and Hom category}
\noindent\textit{Transcription reads.} \(A[x]\)-modules and \(\operatorname{Hom}_{A[x]}(-,A)\)\par
\noindent\textit{Critical reading.} cofree primitive-coalgebra comodules and comodule Hom, or continuous modules over the dual algebra\par
\noindent\textit{Evidence.} The displayed polynomial cobar cochains do not compute ordinary module Ext of A over A[x]; P476 makes the duality explicit.\par
\subsection*{P000466 (both; faithful error): solvable quiver sentence}
\noindent\textit{Transcription reads.} These are all graphs for which the problem is solvable\par
\noindent\textit{Critical reading.} These are the infinite tame cases; finite Dynkin cases must also be included.\par
\noindent\textit{Evidence.} Scan confirms the affine diagrams and unrestricted wording. The preceding page already classifies finite Dynkin cases; modern statement must distinguish finite and tame infinite type.\par
\subsection*{P000467 (both; ambiguity): two-term shorthand for the differential of h}
\noindent\textit{Transcription reads.} only the substitutions \(h(x+y,z)-h(x,y+z)\)\par
\noindent\textit{Critical reading.} \(h(x+y,z)-h(y,z)-h(x,y+z)+h(x,y)\) in the normalized cobar complex\par
\noindent\textit{Evidence.} P467 refers to reduction; P476 explicitly has all four terms. The reduced-coproduct terms must be retained when the shorthand is expanded.\par
\subsection*{P000468 (both; faithful error): arbitrary reindexing}
\noindent\textit{Transcription reads.} any map \(\mathbb N\to\mathbb N\) induces an endomorphism\par
\noindent\textit{Critical reading.} an injective map induces the stated matrix-ring embedding\par
\noindent\textit{Evidence.} Original scan says any map; collapsing two diagonal matrix units contradicts multiplicativity. Vaserstein 1976 section 4 correctly requires embeddings.\par
\subsection*{P000468 (both; ambiguity): inner permutation automorphism}
\noindent\textit{Transcription reads.} inner automorphism\par
\noindent\textit{Critical reading.} implemented by a permutation multiplier (or by a unit of the row-and-column-finite ring for unital \(A\)), not necessarily by a finitary stable unit\par
\noindent\textit{Evidence.} Infinite permutations need not be finitary. The scan does not specify the ambient ring.\par
\subsection*{P000468 (both; faithful error): map into ordinary units}
\noindent\textit{Transcription reads.} \(A\to1+A\) is a homomorphism into invertible elements\par
\noindent\textit{Critical reading.} The monoid map is defined on \(A\); restrict its domain to \(GA\) to obtain a homomorphism into units.\par
\noindent\textit{Evidence.} P000468 scan and the preceding definition of quasiunits; Vaserstein 1976, section 2.\par
\subsection*{P000468 (both; ambiguity): triangular convention}
\noindent\textit{Transcription reads.} upper triangular; \(a_{ij}=0\) for \(i\leq j\)\par
\noindent\textit{Critical reading.} Either strict lower triangular with the written inequality, or strict upper triangular with \(i\geq j\).\par
\noindent\textit{Evidence.} Both the word and inequality are clearly present in the scan; transposition leaves the permutation-triangular construction unchanged.\par
\subsection*{P000469 (both; ambiguity): elementary homotopy source and target}
\noindent\textit{Transcription reads.} \(s_1,s_2:S'\to S'\) with \(X\subset S\)\par
\noindent\textit{Critical reading.} \(s_1,s_2:S\to S'\), and \(s_1(X)\cup s_2(X)\) a simplex in \(S'\) for every source simplex \(X\)\par
\noindent\textit{Evidence.} The source prime is a short doubtful stroke in the scan; the following unprimed \(X\subset S\) and primed target fix the intended roles.\par
\subsection*{P000470-P000471 (both; ambiguity): category of complexes terminology}
\noindent\textit{Transcription reads.} \(K_C\) called category of complexes\par
\noindent\textit{Critical reading.} \(K_C\) is the homotopy category: its morphisms are chain-homotopy classes, as the next line explicitly says.\par
\noindent\textit{Evidence.} P000470 scan; this is a reading clarification, not a contradictory transcription.\par
\subsection*{P000470 (both; provenance): end of record280/start of record281}
\noindent\textit{Transcription reads.} catalogue ranges overlap\par
\noindent\textit{Critical reading.} Vasserstein ends after the tolerance problem; Kazhdan begins at the underlined derived-category heading lower on the same page.\par
\noindent\textit{Evidence.} Original P000470 scan and P000471 continuation.\par
\subsection*{P000470 (both; ambiguity): omitted Steinberg commuting restrictions}
\noindent\textit{Transcription reads.} some commutators are zero\par
\noindent\textit{Critical reading.} \([x_{ij}(a),x_{kl}(b)]=1\) for \(i\ne l\) and \(j\ne k\); adjacent-root relation requires distinct \(i,j,k\)\par
\noindent\textit{Evidence.} The scan abbreviates the restrictions; Vaserstein 1976 section 6 and Weibel III.5.1 give the exact relations.\par
\subsection*{P000470 (both; ambiguity): range of block-sum \(K\) indices}
\noindent\textit{Transcription reads.} all \(K_i\), \(i\geq0\)\par
\noindent\textit{Critical reading.} all \(K_i\), \(i\geq1\), for the displayed homotopy definition; \(K_0\) requires its separate definition\par
\noindent\textit{Evidence.} P000469 explicitly defines only \(K_n=\pi_{n-1}(SA)\) for \(n\geq1\); no \(K_0\) construction is given here.\par
\subsection*{P000470 (both; faithful error): tolerance task}
\noindent\textit{Transcription reads.} if every pair is a simplex, the whole finite set is a simplex\par
\noindent\textit{Critical reading.} False for the defined complex: \(I+E_{12},I+E_{23},I+E_{31}\) are pairwise adjacent but not a simplex.\par
\noindent\textit{Evidence.} P000468 fixes a common translation and coordinate order; complete proof and two independent exact audits. The same false assertion appears in Vaserstein 1976 section 11 Remark(c), p.111.\par
\subsection*{P000471 (both; ambiguity): Adams bibliographical pointer}
\noindent\textit{Transcription reads.} Mat. 3:4, pp. 7--19\par
\noindent\textit{Critical reading.} exact article unresolved; retain the pointer without substituting an unverified title\par
\noindent\textit{Evidence.} The scan confirms 3:4. The accessible Adams translation On the non-existence of elements of Hopf invariant one is Matematika 5(4) (1961), 3--86 and does not match this pointer.\par
\subsection*{P000471 (both; provenance): end of record281}
\noindent\textit{Transcription reads.} catalogue range continues through P000472\par
\noindent\textit{Critical reading.} Kazhdan's derived-category exposition ends before the Buchstaber heading on P000471.\par
\noindent\textit{Evidence.} Original P000471 scan and the explicit change of speaker.\par
\subsection*{P000472 (both; ambiguity): identification of dual bar terms}
\noindent\textit{Transcription reads.} \(\operatorname{Hom}_{A_0}(\bar A^{\otimes s},A_0)=(\bar A^*)^{\otimes s}\) without qualifications\par
\noindent\textit{Critical reading.} use degreewise finite projectivity and graded/completed duals, or retain the Hom expression\par
\noindent\textit{Evidence.} The duality identification is not valid for arbitrary infinite modules; the later continuous-dual convention provides the intended interpretation.\par
\subsection*{P000472 (both; faithful error): target of the reduced coproduct}
\noindent\textit{Transcription reads.} \(\operatorname{im}\bar\Delta\subset\overline{A^*}\)\par
\noindent\textit{Critical reading.} \(\operatorname{im}\bar\Delta\subset\overline{A^*}\widehat\otimes\overline{A^*}\)\par
\noindent\textit{Evidence.} Scan confirms the missing tensor factor; the preceding unreduced coproduct has target a tensor square.\par
\subsection*{P000473 (both; ambiguity): Abel series}
\noindent\textit{Transcription reads.} \(e^{\lambda_k t}\), \(\lambda_k>0\), \(t\to0\)\par
\noindent\textit{Critical reading.} Direction of \(t\to0\) must be stated; the usual decaying Abel factor for \(t\downarrow0\) is \(e^{-\lambda_k t}\).\par
\noindent\textit{Evidence.} The scan confirms the positive exponent. The PDE domain suggests nonnegative time, but the summation paragraph gives no one-sided limit. Preserve as ambiguity, not a certain typographical correction.\par
\subsection*{P000473 (both; faithful error): Bernstein interval extension}
\noindent\textit{Transcription reads.} Complete on \([0,\pi+\varepsilon]\) for \(0<\varepsilon<\pi\)\par
\noindent\textit{Critical reading.} At least exclude \(\gamma=0\) and state whether the half or signed system is intended.\par
\noindent\textit{Evidence.} Both the interval and attribution are visible in the scan. Functions symmetric about \(\pi\) with small paired support are orthogonal to every \(\sin(nx)\), refuting the literal zero-tilt case.\par
\subsection*{P000473 (both; faithful error): nonbasis and analyticity claims}
\noindent\textit{Transcription reads.} \(a^2<b\), \(b>0\); not a basis; every convergent limit analytic\par
\noindent\textit{Critical reading.} For the half system, add \(a\ne0\) (equivalently \(\gamma\ne0\)) to these pathological assertions.\par
\noindent\textit{Evidence.} The scan has no nonzero-\(a\) hypothesis. With \(a=0\) the half system is the ordinary sine basis, whose \(L^2\) sums need not be analytic.\par
\subsection*{P000473 (both; faithful error): signed eigenfunction list and biorthogonality paragraph}
\noindent\textit{Transcription reads.} \(k=\pm1,\pm2,\ldots\); a biorthogonal system exists\par
\noindent\textit{Critical reading.} The biorthogonal assertion can concern a one-sided system only.\par
\noindent\textit{Evidence.} The scan confirms the signed indices and the later positive-index transform. Completeness of the positive subsystem rules out a biorthogonal family for the full signed list; at \(\gamma=0\) the negative entries are explicit duplicates.\par
\subsection*{P000474 (both; ambiguity): nonsimple codimension}
\noindent\textit{Transcription reads.} Nonsimple of codimension 6\par
\noindent\textit{Critical reading.} First nonsimple locus has codimension 6 for \(n\ge3\) under Arnold’s critical-point-stratum convention; codimension 7 for \(n=2\).\par
\noindent\textit{Evidence.} The scan genuinely says 6. Arnold 1972 Introduction and Theorem 7.1 confirm it. Do not replace 6 by a Milnor-number threshold.\par
\subsection*{P000474 (both; ambiguity): simple-hypersurface equivalence and versal formula}
\noindent\textit{Transcription reads.} Equivalence of hypersurfaces under coordinate changes; Jacobian quotient criterion\par
\noindent\textit{Critical reading.} Distinguish right equivalence of functions from contact equivalence of defining equations.\par
\noindent\textit{Evidence.} The later deformation equality is for functions and the quotient uses only partial derivatives, identifying right-versality. Equivalence of zero sets alone also allows multiplication by a unit.\par
\subsection*{P000475 (both; faithful error): derivatives in \(A_k\) example}
\noindent\textit{Transcription reads.} \(\partial f/\partial z_1=z_1^k\); \(\partial f/\partial z_2=z_2\)\par
\noindent\textit{Critical reading.} \(\partial f/\partial z_1=(k+1)z_1^k\); \(\partial f/\partial z_j=2z_j\) for \(j\ge2\).\par
\noindent\textit{Evidence.} The omitted factors are visible in the manuscript; the generated ideal is unchanged over \(\mathbb C\), but the displayed derivative equalities are false.\par
\subsection*{P000475 (both; faithful error): homology of the smooth fibre}
\noindent\textit{Transcription reads.} Homology nonzero only in middle degree \(n-1\)\par
\noindent\textit{Critical reading.} Reduced homology nonzero only in middle degree; ordinary \(H_0\) also occurs.\par
\noindent\textit{Evidence.} The scan has no reduced qualifier. A nonempty connected Milnor fibre has \(H_0=\mathbb Z\); Milnor 1968 Theorems 6.5 and 7.2 give the exact result.\par
\subsection*{P000475 (both; ambiguity): local-algebra basis}
\noindent\textit{Transcription reads.} \(1,z_i,z_i^2,\ldots,z_1^{k-1}\)\par
\noindent\textit{Critical reading.} \(1,z_1,z_1^2,\ldots,z_1^{k-1}\).\par
\noindent\textit{Evidence.} The first subscripts are \(i\)-like in the scan. The Jacobian quotient kills every \(z_j\) with \(j\ge2\), and the next unfolding uses \(z_1\) throughout.\par
\subsection*{P000476 (both; ambiguity): label of the third cochain differential}
\noindent\textit{Transcription reads.} \(d_3\)\par
\noindent\textit{Critical reading.} \(d_3^*\) under the notebook dual-complex convention\par
\noindent\textit{Evidence.} The surrounding arrows are dual cochain differentials; the displayed formula is the dual bar differential.\par
\subsection*{P000476 (both; faithful error): upper coproduct formula, isolated superscript}
\noindent\textit{Transcription reads.} \(x_j^*\) in one tensor factor of the dual coproduct\par
\noindent\textit{Critical reading.} \(x_j\)\par
\noindent\textit{Evidence.} The target is already A* tensor A* and \(x_j\) is a dual-basis element. The repeated coproduct formula immediately below omits this extra star. The scan preserves the stray star.\par
\subsection*{P000477 (both; faithful error): Ext ring of a truncated polynomial factor}
\noindent\textit{Transcription reads.} \(H^*(k[t]/t^p)=\Lambda(u)\otimes k[v]\) for every prime p\par
\noindent\textit{Critical reading.} odd p: the displayed formula; p=2: \(k[u]\), with \(v=u^2\)\par
\noindent\textit{Evidence.} The scan explicitly has exponent p and no odd-prime exclusion. The normalized bar differential for \(k[t]/t^2\) is zero, so u squared is nonzero.\par
\subsection*{P000477 (both; ambiguity): identification of cocycles with formal groups}
\noindent\textit{Transcription reads.} as many groups as cycles\par
\noindent\textit{Critical reading.} additive cocycles correspond directly to laws with square-zero coefficient perturbation; full laws have nonlinear extension conditions\par
\noindent\textit{Evidence.} The equality is correct in the preceding indecomposable/square-zero setting. Lazard's prolongation theorem distinguishes arbitrary linearized data from full nonlinear laws.\par
\subsection*{P000478 (both; ambiguity): count of smooth seven-spheres}
\noindent\textit{Transcription reads.} \(\operatorname{card}S\prime(S^7)=28\)\par
\noindent\textit{Critical reading.} 28 oriented homotopy-sphere classes, or the corresponding marked structure set\par
\noindent\textit{Evidence.} The prime is legible but undefined. Kervaire--Milnor's theorem counts oriented classes; forgetting orientation and marking changes the equivalence relation.\par
\subsection*{P000478-P000479 (both; ambiguity): definition of \(K_J\) and its representability}
\noindent\textit{Transcription reads.} stable J-trivial bundles, with no chosen trivialization\par
\noindent\textit{Critical reading.} pairs consisting of a stable bundle and a specified spherical trivialization, when identifying with maps into G/O\par
\noindent\textit{Evidence.} Sullivan's normal-invariant definition and Ranicki Proposition 9.21 retain the null-homotopy data; the bare kernel forgets it.\par
\subsection*{P000478 (both; ambiguity): word smoothings in the initial definition}
\noindent\textit{Transcription reads.} smoothings defined using a homotopy equivalence\par
\noindent\textit{Critical reading.} homotopy smoothings or marked smooth homotopy structures\par
\noindent\textit{Evidence.} The definition itself is unambiguous; P481 inserts homotopic and P488 distinguishes S from hS. This is terminology, not a false claim.\par
\subsection*{P000479 (both; ambiguity): PL obstruction sequence and cokernel language}
\noindent\textit{Transcription reads.} unqualified exact sequence of structures with terminal obstruction group called a cokernel\par
\noindent\textit{Critical reading.} closed simply connected PL manifolds of dimension at least six; exactness of pointed sets; no automatic group cokernel\par
\noindent\textit{Evidence.} Sullivan Theorem 3 supplies the hypotheses and pointed-set meaning. The smooth continuation has a group action at the left.\par
\subsection*{P000480 (both; provenance): Pre-indexed Dobrushin record between requested records 287 and 289}
\noindent\textit{Transcription reads.} The supplied page range includes the 19 February 1973 discussion, but the requested 28-record list omits record 288.\par
\noindent\textit{Critical reading.} Retain its existing occurrences 1972-288.1 through 1972-288.5 from the supplied previous index; do not create duplicate occurrences or silently assign this page to another speaker.\par
\noindent\textit{Evidence.} The dated manuscript page and the supplied previous index identify the same material. This is an ownership discrepancy in the supplied year partition, not missing new mathematical coverage.\par
\subsection*{P000481 (both; ambiguity): Hauptvermutung replacement statement}
\noindent\textit{Transcription reads.} a homeomorphism can be replaced by a PL homeomorphism under the two listed conditions\par
\noindent\textit{Critical reading.} for compact PL manifolds in dimension at least six, with simply connected manifold and boundary, the homeomorphism is homotopic to a PL homeomorphism under Sullivan's torsion condition\par
\noindent\textit{Evidence.} Sullivan 1967, main theorem, distinguishes stable from unstable conclusions and explicitly states the compactness, dimension, and boundary hypotheses. For the closed case the boundary condition is vacuous.\par
\subsection*{P000481 (both; faithful error): Hauptvermutung torsion condition}
\noindent\textit{Transcription reads.} no 2-torsion in \(H^3\)\par
\noindent\textit{Critical reading.} no 2-torsion in \(H_3(M;\mathbb Z)\), equivalently in the torsion subgroup of \(H^4(M;\mathbb Z)\)\par
\noindent\textit{Evidence.} The scan clearly has superscript 3. Sullivan 1967, main theorem, has \(H_3\); the 1967 seminar notes, pp.93--94, give the \(H^4\) obstruction and \(H_3\) corollary.\par
\subsection*{P000481 (both; ambiguity): proper homotopy equivalence of bundles}
\noindent\textit{Transcription reads.} proper map and a homotopy equivalence of total spaces\par
\noindent\textit{Critical reading.} proper fibrewise homotopy equivalence, stably, with the compatible fibre data\par
\noindent\textit{Evidence.} The diagram is over the identity of M; ordinary total-space equivalence alone is too weak since vector bundles retract to their bases. Dold Corollary 2 and Atiyah Proposition 2.8 give the intended spherical criterion.\par
\subsection*{P000482 (both; faithful error): Scaling the centre of the ball}
\noindent\textit{Transcription reads.} \(x_0=\|A\|x'_0\).\par
\noindent\textit{Critical reading.} Choose \(x_0=tu\) with \(t\|Au\|>\|A\|+1/2\).\par
\noindent\textit{Evidence.} The scan visibly has the norm multiplier. Its value does not ensure that AS avoids zero; the stated inequality follows by the reverse triangle inequality. This is an emendation, not a transcription replacement.\par
\subsection*{P000483 (both; ambiguity): Henkin localization operator}
\noindent\textit{Transcription reads.} \(f\mapsto f g-f_g\) with \(f_g\) defined from g alone.\par
\noindent\textit{Critical reading.} The dependence of \(f_g\) on f must be specified; retain as incomplete notation.\par
\noindent\textit{Evidence.} Both transcriptions and the scan omit the needed dependence. An integral of g alone gives an affine, not the asserted linear, operator.\par
\subsection*{P000483 (en; translation): Henkin point 2: submanifold}
\noindent\textit{Transcription reads.} English: closed subset.\par
\noindent\textit{Critical reading.} Read closed analytic submanifold; preserve the manuscript abbreviation.\par
\noindent\textit{Evidence.} The Russian has \ru{подмн-зие}, repeated in the extension clause, and the transverse analytic-submanifold theorem identifies the intended category.\par
\subsection*{P000483 (both; faithful error): Numerator of the fixed-point weights}
\noindent\textit{Transcription reads.} \(f(\|T_jy-x_0\|\)\par
\noindent\textit{Critical reading.} \(f(\|T_jy-x_0\|)\).\par
\noindent\textit{Evidence.} Scan inspected: the closing parenthesis is absent, while the parallel denominator determines its restoration.\par
\subsection*{P000483-P000484 (both; provenance): Order of the three 5 March records}
\noindent\textit{Transcription reads.} Catalogue order: 290, 291, 292.\par
\noindent\textit{Critical reading.} Manuscript order: 290 (Lomonosov), 292 (Henkin), 291 (Berezin).\par
\noindent\textit{Evidence.} P000483 places the Henkin announcement immediately after the invariant-subspace proof; P000484 begins Berezin. Catalogue windows are not exact talk spans.\par
\subsection*{P000484 (both; faithful error): Annihilation and adjoint}
\noindent\textit{Transcription reads.} \(a=\partial_z,\ a^*=z\) for Gaussian parameter h.\par
\noindent\textit{Critical reading.} \(a=h\partial_z,\ a^*=z\), or restrict this line to h=1.\par
\noindent\textit{Evidence.} Monomial Gaussian norms are proportional to \(h^m m!\), which determine the adjoint. The displayed h was not specialized in the notes.\par
\subsection*{P000484 (both; faithful error): Fock representation}
\noindent\textit{Transcription reads.} \(f(z,v)\overline{\varphi(v)}\)\par
\noindent\textit{Critical reading.} \(f(z,\bar v)\varphi(v)\)\par
\noindent\textit{Evidence.} The scan places the bar over the input vector. The literal operation is antilinear; the reproducing-kernel formula fixes the correction.\par
\subsection*{P000484 (both; faithful error): Normal-symbol monomials}
\noindent\textit{Transcription reads.} \((a^*)^m a^n\leftrightarrow\bar z^m z^n\).\par
\noindent\textit{Critical reading.} \((a^*)^m a^n\leftrightarrow z^m\bar z^n\) in the stated Fock convention.\par
\noindent\textit{Evidence.} The representation sends the creation operator to multiplication by z; direct coherent-state evaluation interchanges the displayed bars.\par
\subsection*{P000485 (both; faithful error): General contravariant integral}
\noindent\textit{Transcription reads.} \(A=\int aP_\alpha d\alpha\) after an unnormalized Parseval identity.\par
\noindent\textit{Critical reading.} \(A=\int aP_\alpha\|e_\alpha\|^2d\alpha\), or absorb the weight in the measure.\par
\noindent\textit{Evidence.} Rank-one orthogonal projectors differ from outer products by the squared norm; the 1972 paper includes it.\par
\subsection*{P000486 (both; faithful error): Bergman overlap ratio}
\noindent\textit{Transcription reads.} Diagonal kernel factors in numerator.\par
\noindent\textit{Critical reading.} Off-diagonal factors in numerator, diagonal factors in denominator.\par
\noindent\textit{Evidence.} The scan confirms the reciprocal ratio; Cauchy--Schwarz and the Gaussian limit determine the mathematical correction.\par
\subsection*{P000486 (both; faithful error): Holomorphic Hilbert-space weight}
\noindent\textit{Transcription reads.} \(K^{+1/h}\).\par
\noindent\textit{Critical reading.} \(K^{-1/h}\) against invariant volume.\par
\noindent\textit{Evidence.} No minus is visible in the scan. The positive sign makes the model boundary weight grow in the wrong direction.\par
\subsection*{P000487 (both; faithful error): Disk kernel exponent}
\noindent\textit{Transcription reads.} \((1-z\bar z)^{+2/h}\).\par
\noindent\textit{Critical reading.} Kernel: \((1-z\bar z)^{-2/h}\); reciprocal weight has the positive exponent.\par
\noindent\textit{Evidence.} The plus is explicit in the scan. Comparison with the reproducing kernel and the corrected weight resolves its role.\par
\subsection*{P000487 (both; faithful error): Kähler potential matrix}
\noindent\textit{Transcription reads.} \(\omega^{ij}=\partial_i\partial_{\bar j}\Phi\).\par
\noindent\textit{Critical reading.} \(g_{i\bar j}=\partial_i\partial_{\bar j}\Phi\); inverse coefficients belong to the Poisson tensor.\par
\noindent\textit{Evidence.} The scan uses upper indices; the earlier page defines upper indices as the inverse. These two conventions conflict.\par
\subsection*{P000488-P000489 (both; ambiguity): composition with a Thom homotopy inverse}
\noindent\textit{Transcription reads.} \(\sigma\lambda:S^L\to T(\xi+\nu_M)\) when \(\xi\) has positive rank r\par
\noindent\textit{Critical reading.} \(\sigma\Sigma^r\lambda:S^{L+r}\to T(\xi+\nu_M)\), or use virtual rank-zero Thom spectra\par
\noindent\textit{Evidence.} The target of the preceding equivalence is \(T(r+\nu_M)=\Sigma^rT(\nu_M)\). The source sphere must be suspended by the same r to obtain an n-dimensional transverse inverse image.\par
\subsection*{P000488 (both; ambiguity): definition and group interpretation of Phi}
\noindent\textit{Transcription reads.} group of stable smooth J-trivial bundles\par
\noindent\textit{Critical reading.} stable pairs \((\xi,\varphi)\), retaining the specified fibre trivialization, for the normal-invariant interpretation\par
\noindent\textit{Evidence.} The page draws the needed map but does not include it in the equivalence relation; Ranicki Propositions 9.21 and 9.43 specify the data.\par
\subsection*{P000488 (both; faithful error): explanation of the normal bundle symbol}
\noindent\textit{Transcription reads.} \(\nu_N\) is the normal bundle of M\par
\noindent\textit{Critical reading.} \(\nu_N\) is the stable normal bundle of N\par
\noindent\textit{Evidence.} Scan and both transcriptions preserve the wrong M; g:M→N pulls back a bundle over N.\par
\subsection*{P000489 (both; ambiguity): transversality sentence after the collapse construction}
\noindent\textit{Transcription reads.} an embedding of M in its Thom complex, which is transverse\par
\noindent\textit{Critical reading.} the collapse map is transverse to the zero section M in the Thom space\par
\noindent\textit{Evidence.} The construction and inverse-image formula determine the intended subject of transversality.\par
\subsection*{P000490 (both; ambiguity): targets of the two endpoint Thom maps}
\noindent\textit{Transcription reads.} \(T(\xi_i+\tau_M)\), with an unclear extra mark above the tangent symbol\par
\noindent\textit{Critical reading.} retain the unclear reading; reconstruction uses \(T(\xi_i+\nu_M)\), equivalently the stable inverse tangent bundle\par
\noindent\textit{Evidence.} The scan has a τ-like sign with an additional mark; the transcriptions flag uncertainty. P488--P489 use \(\nu_M\). It is not established whether the mark already denotes stable inverse tangent data.\par
\subsection*{P000490 (both; ambiguity): uniqueness via Morse surgeries}
\noindent\textit{Transcription reads.} all resulting manifolds related by Morse surgeries, based only on the bundle class\par
\noindent\textit{Critical reading.} the same normal invariant gives a normal-bordism class; compatible surgeries follow from a handle decomposition of that normal bordism\par
\noindent\textit{Evidence.} The relative surgery construction requires the stable trivialization and normal-map data. Bare bundle equality is weaker than equality of normal invariants.\par
\subsection*{P000491 (both; faithful error): quadratic stabilization}
\noindent\textit{Transcription reads.} The quadratic form is positive definite\par
\noindent\textit{Critical reading.} Over \(\mathbb C\) use a nondegenerate quadratic form, equivalent to a sum of squares; over \(\mathbb R\) retain its signature.\par
\noindent\textit{Evidence.} The scan plainly says positive definite while the germ is explicitly complex. Arnold 1973 definition uses nondegenerate quadratic forms.\par
\subsection*{P000492 (en; translation): Tyurina attribution}
\noindent\textit{Transcription reads.} Tyurin’s theorem\par
\noindent\textit{Critical reading.} Tyurina’s theorem (G. N. Tyurina).\par
\noindent\textit{Evidence.} The Russian reads \ru{теор. Тюрина}, which is consistent with the genitive of \ru{Тюрина}. Arnold 1973 section 4 Remark 3 explicitly attributes the theorem to G. N. Tyurina and cites her 1968 paper.\par
\subsection*{P000492 (both; ambiguity): finiteness of the J-quotient}
\noindent\textit{Transcription reads.} \(K(X)/J\)-equivalence is finite\par
\noindent\textit{Critical reading.} Use the reduced stable \(J\)-group of a finite CW-complex; retaining rank adds \(\mathbb Z\).\par
\noindent\textit{Evidence.} Adams 1963 section 1 explicitly distinguishes the finite reduced group from its unreduced version. The notebook omits reduction.\par
\subsection*{P000492 (both; ambiguity): first line about analytic changes}
\noindent\textit{Transcription reads.} An analytic function up to analytic variable change uniquely determines a variable change\par
\noindent\textit{Critical reading.} Unresolved compressed sentence; do not impose uniqueness of a normalizing map.\par
\noindent\textit{Evidence.} The scan repeats abbreviated \ru{зам. перем.} at both ends and does not supply the missing object. Quadratic-form symmetries refute uniqueness if that is the intended literal meaning.\par
\subsection*{P000492 (both; faithful error): nullity claim}
\noindent\textit{Transcription reads.} \(\mu_+\le1\) implies \(\mu_0=1\) (always degenerate)\par
\noindent\textit{Critical reading.} The positive-index-one \(T\) cases have \(\mu_0=1\); simple elliptic cases have \(\mu_+=0\) and \(\mu_0=2\), while ADE has \(\mu_0=0\).\par
\noindent\textit{Evidence.} The second \(\le\) is clear in the scan. Gabrielov 1973 Proposition 5 gives corank 2 for the cubic written immediately above; Arnold 1973 gives the \(\mu_+=1\) theorem.\par
\subsection*{P000492 (both; faithful error): reverse Adams equality}
\noindent\textit{Transcription reads.} Every \(J\)-trivial \(\xi\) equals \(k^N(\psi^k-1)y\) for every \(k\)\par
\noindent\textit{Critical reading.} The topological kernel is described by a subgroup generated using varying \(k\); no arbitrary single \(k\) has this property.\par
\noindent\textit{Evidence.} The scan clearly includes \(\forall k\) and a single term. Adams 1965 section 4 uses a sum over \(k\). In \(KO(S^4)\), choose \(M=\operatorname{ord}J(z)\), a prime \(q\) not dividing \(M\), and \(k=q+1\); \(q\)-divisibility refutes the equality. The independent complex \(S^2\) check with \(k=4\) gives the same verdict.\par
\subsection*{P000492 (both; ambiguity): signature and modality}
\noindent\textit{Transcription reads.} \(\mu_+\le1\) implies \(m=1\)\par
\noindent\textit{Critical reading.} The \(\le\) sign is doubtful; the matching published nonparabolic theorem assumes \(\mu_+=1\). As a literal \(\le1\) assertion it fails for ADE.\par
\noindent\textit{Evidence.} Crop inspection confirms the conclusion \(m=1\) but the first inequality is retraced. Arnold 1973 section 3 Corollary Zh uses \(\mu_+=1\); ADE provides the check.\par
\subsection*{P000493 (both; faithful error): dimension-preserving operation}
\noindent\textit{Transcription reads.} If \(A\) preserves dimension then \(A\xi\) is \(J\)-equivalent to \(\xi\); such \(A\) do not occur\par
\noindent\textit{Critical reading.} The sufficient condition is a compatible filtered homotopy automorphism; rank preservation alone fails. Identity and conjugation exist.\par
\noindent\textit{Evidence.} Sullivan 1974 p. 71 explicitly says filtered automorphism, while the manuscript only draws maps. The trivial-bundle-of-equal-rank operation shows the weaker claim fails.\par
\subsection*{P000493-P000494 (both; faithful error): unstable completed Adams operations}
\noindent\textit{Transcription reads.} The constructed completed \(\psi^k\) extends the Adams operation and preserves dimension\par
\noindent\textit{Critical reading.} Use \(\psi^k\) at primes \(p\) not dividing \(k\), and its isomorphic part (identity at \(p\) dividing \(k\)) for full profinite completion.\par
\noindent\textit{Evidence.} Sullivan 1970 Theorem 5.7 and Proposition 5.8 explicitly make this distinction and show why unrestricted unstable extension fails.\par
\subsection*{P000494 (both; ambiguity): analytic/étale comparison}
\noindent\textit{Transcription reads.} \(H_i(V_c,A)\approx H_i(V_{\mathrm{et}},A)\)\par
\noindent\textit{Critical reading.} Specify finite coefficient systems and the profinite comparison; the usual primary statement is cohomological.\par
\noindent\textit{Evidence.} The manuscript supplies no condition on \(A\). Sullivan 1974 p. 68 formula \((*)\) explicitly restricts to finite coefficients.\par
\subsection*{P000494 (both; ambiguity): inverse limit of nerves}
\noindent\textit{Transcription reads.} \(\widehat V_c\) is the inverse limit of \(N(V_{\mathrm{et}})\)\par
\noindent\textit{Critical reading.} Interpret completed hypercovering nerves as a pro-homotopy type and use a homotopy/compact-functor realization.\par
\noindent\textit{Evidence.} Sullivan 1974 p. 68, including its footnote, distinguishes completion of nerves and the sense of the limit.\par
\subsection*{P000494 (both; ambiguity): unitary sphere exponent}
\noindent\textit{Transcription reads.} \(U(n)\) acts on \(S^{2n-1}\), possibly \(2n+1\)\par
\noindent\textit{Critical reading.} \(S^{2n-1}\) for the standard \(U(n)\)-action; \(S^{2n+1}\) belongs to \(BU(n)\to BU(n+1)\).\par
\noindent\textit{Evidence.} Scan review confirms the exponent sign is retraced. The two distinct sphere dimensions occur in the adjacent classifying-space construction.\par
\subsection*{P000495 (both; faithful error): Circle eta formula}
\noindent\textit{Transcription reads.} \(\eta(0)=1-2t\), unrestricted t.\par
\noindent\textit{Critical reading.} For \(0<t<1\) this holds; otherwise use \(1-2\{t\}\), with value 0 at integers.\par
\noindent\textit{Evidence.} The scan gives no branch restriction. Integer translation preserves the spectrum, whereas the affine expression is not periodic.\par
\subsection*{P000495-P000497 (both; provenance): Order of the three 26 March records}
\noindent\textit{Transcription reads.} Catalogue order: 296, 297, 298.\par
\noindent\textit{Critical reading.} Manuscript order: 296 (spectral asymmetry), 298 (p-adic integration), 297 (Gabrielov opening fragment).\par
\noindent\textit{Evidence.} The arithmetic continuation occupies the upper part of P000497, before the Gabrielov heading. The chapter follows the manuscript while retaining catalogue record identifiers.\par
\subsection*{P000496 (both; faithful error): Invariant distribution uniqueness}
\noindent\textit{Transcription reads.} \(p^{-m}\) is the only possibility.\par
\noindent\textit{Critical reading.} This is the unique translation-invariant distribution of total mass 1.\par
\noindent\textit{Evidence.} The scan states uniqueness without a normalization. Multiplying every value by a constant preserves finite additivity and invariance.\par
\subsection*{P000496-P000497 (both; faithful error): Mellin-transform direction}
\noindent\textit{Transcription reads.} The measure is obtained by the Mellin transform from the zeta function.\par
\noindent\textit{Critical reading.} The Mellin transform sends the regularized measure to a character function; recovering a measure is the inverse problem.\par
\noindent\textit{Evidence.} The scan confirms the reversed sentence on P000496; the continuation on P000497 has the forward direction.\par
\subsection*{P000497 (both; faithful error): Zeta continuity assertion}
\noindent\textit{Transcription reads.} Raw \(\zeta(1-2k_i)\) has a limit depending only on \(\lim k_i\), everywhere.\par
\noindent\textit{Critical reading.} Insert \(1-p^{2k_i-1}\), fix the tame-character branch, and exclude or account for the pole.\par
\noindent\textit{Evidence.} The displayed notebook formula lacks these conditions. The p=5 counterexample and the published interpolation formula distinguish a faithful mathematical error from a transcription fault.\par
\subsection*{P000497 (both; ambiguity): dimension of the vanishing spheres}
\noindent\textit{Transcription reads.} Bouquet of \((n-1)\)-spheres, possibly \((k-1)\)-spheres\par
\noindent\textit{Critical reading.} \((n-1)\)-spheres for a germ on \(\mathbb C^n\).\par
\noindent\textit{Evidence.} The scan leaves the first letter retraced, but the preceding manifold dimension is \(n-1\) and Milnor’s theorem gives that sphere dimension.\par
\subsection*{P000498 (both; faithful error): Analytic character at p=2}
\noindent\textit{Transcription reads.} \(a\equiv1\pmod p\), all primes.\par
\noindent\textit{Critical reading.} For the displayed exponential at p=2 use \(a\in1+4\mathbb Z_2\).\par
\noindent\textit{Evidence.} The exponential convergence radius differs at two.\par
\subsection*{P000498 (both; faithful error): Automorphy factor}
\noindent\textit{Transcription reads.} \((az+c)^{-w}\)\par
\noindent\textit{Critical reading.} \((cz+d)^w\) for the stated transformation convention.\par
\noindent\textit{Evidence.} The scan clearly has az+c and a negative exponent.\par
\subsection*{P000498 (both; faithful error): Finite additivity}
\noindent\textit{Transcription reads.} Additivity for U1 union U2 without qualification.\par
\noindent\textit{Critical reading.} Require disjointness, or subtract the intersection term.\par
\noindent\textit{Evidence.} No disjointness sign appears in the scan.\par
\subsection*{P000498 (both; faithful error): Riemann sum}
\noindent\textit{Transcription reads.} \(f(a)(a+A_m)\)\par
\noindent\textit{Critical reading.} \(f(a)\mu(a+A_m)\)\par
\noindent\textit{Evidence.} The scan itself lacks mu; a scalar times a set is not the claimed W-valued sum.\par
\subsection*{P000499 (both; ambiguity): Final interpolation constant}
\noindent\textit{Transcription reads.} Notebook prefactor with \(\pi^{r-1}\), \(p^r\), omega and \(g(\chi)\).\par
\noindent\textit{Critical reading.} Do not certify this constant without fixing the measure scale, period component and Gauss-sum convention.\par
\noindent\textit{Evidence.} The preceding parity error and ambiguous normalization prevent an exact literal match; the chapter supplies one consistent convention.\par
\subsection*{P000499 (both; faithful error): Hecke operator}
\noindent\textit{Transcription reads.} Unweighted sum of f((az+b)/d).\par
\noindent\textit{Critical reading.} Insert \(m^{w-1}d^{-w}\), with positive a,d and ad=m.\par
\noindent\textit{Evidence.} The source contains no determinant or weight factors; the eigenvalue and root formulas fix the standard normalization.\par
\subsection*{P000499 (both; faithful error): Period component versus parity}
\noindent\textit{Transcription reads.} Re of the path integral, with chi(-1)=1.\par
\noindent\textit{Critical reading.} Use the even imaginary period component, or instead change to the odd-character branch.\par
\noindent\textit{Evidence.} Both Re and the even parity are visible. Complex conjugation of the path proves the mismatch.\par
\subsection*{P000499 (both; faithful error): Primitive character kernel}
\noindent\textit{Transcription reads.} \(p^{r-1}\mathbb Z/p^r\mathbb Z\)\par
\noindent\textit{Critical reading.} \(1+p^{r-1}\mathbb Z/p^r\mathbb Z\) for r at least 2.\par
\noindent\textit{Evidence.} The missing 1+ is decisive: the written set is additive and contains nonunits.\par
\subsection*{P000499 (both; ambiguity): Residue in averaging formula}
\noindent\textit{Transcription reads.} \(b\equiv0\pmod{p^n}\) (possibly a).\par
\noindent\textit{Critical reading.} A measure of \(a+p^n\mathbb Z_p\) must sum \(b\equiv a\pmod{p^n}\); retain 0/a as a palaeographic ambiguity.\par
\noindent\textit{Evidence.} The sign is overwritten. A formula independent of a cannot give the intended varying coset distribution.\par
\subsection*{P000500 (both; ambiguity): Bernoulli coefficient convention}
\noindent\textit{Transcription reads.} \(B_r\) multiplying the raw r-th endpoint derivative.\par
\noindent\textit{Critical reading.} These coefficients are standard \(B_r/r!\).\par
\noindent\textit{Evidence.} The determinant expansion and exponential test fix b2=1/12. This is a notation clarification, not evidence of a missing factorial in the transcription.\par
\subsection*{P000500 (both; faithful error): Hurwitz reference}
\noindent\textit{Transcription reads.} Math. Z., 1898.\par
\noindent\textit{Critical reading.} Mathematische Annalen 51 (1899), 196--226.\par
\noindent\textit{Evidence.} The notebook supplies the correct citation lower down; the publisher collected-works metadata confirms it.\par
